% =============================================================================
% Chapter 3 - Literature Review  (shortened, assembled)
% =============================================================================
% Single-file assembly of the ten shortened sections. Include from the master
% with % =============================================================================
% Chapter 3 - Literature Review  (shortened, assembled)
% =============================================================================
% Single-file assembly of the ten shortened sections. Include from the master
% with % =============================================================================
% Chapter 3 - Literature Review  (shortened, assembled)
% =============================================================================
% Single-file assembly of the ten shortened sections. Include from the master
% with % =============================================================================
% Chapter 3 - Literature Review  (shortened, assembled)
% =============================================================================
% Single-file assembly of the ten shortened sections. Include from the master
% with \include{contents/chapter3_shortened} in place of contents/chapter3.tex.
% Do not include both: they share every label.
%
% Assembled from chapter3_short.tex (3.1-3.2) and chapter3_short_part3..10.tex.
% Those part files are retained unchanged as the editing units; regenerate this
% file from them rather than editing it directly, or edit here and retire them.
%
% Structure that other chapters depend on, and that must not be disturbed:
%   * Chapter 2 cites Sections 3.1.2, 3.1.6, 3.1.8, 3.3.1, 3.3.3, 3.5.7,
%     3.9, 3.9.5 and 3.9.7 BY NUMBER, so those subsection positions are fixed.
%   * The divergence ledger (Table, sec:divergence-ledger) cites twelve
%     subsections by number. scratchpad/ledgercheck.py re-verifies all of them
%     plus every external numeric reference after any renumbering.
%   * Chapters 4, 5, 6 and the introduction reference sec:research-gap,
%     sec:four-limitations, sec:addressing-gap, sec:divergence-ledger,
%     sec:synthesis-and-gap, sec:imbalance-review, sec:imbalance-gap,
%     sec:imbalance-double-correction, sec:cnn-review, sec:cnn-implications,
%     sec:tim-implications, sec:evm-implications, sec:evm-evidence,
%     sec:evm-limitations, sec:evm-motion-hypothesis, sec:corpus,
%     sec:corpus-threats, sec:corpus-implications, sec:label-taxonomy,
%     sec:megc-metrics, sec:sibling-corpora, sec:flow-review, sec:flow-implications,
%     sec:strain-review, sec:strain-implications, sec:strain-divergence,
%     sec:simam-gap, sec:simam-implications, sec:slstt-results,
%     sec:transformer-implications, sec:addressing-gap, tab:megc-cde-results
%     and tab:spontaneous-corpora.
%
% Conventions held throughout: no sentence points at another chapter; every
% float carries a short \caption[...] for the List of Figures / List of Tables;
% figures are fig3_1_dataset.png and fig3_2_fold_composition.png only.
% =============================================================================

\chapter{Literature Review}
\label{ch:literature}

This chapter reviews the literature bearing on each component of the recognition pipeline evaluated in this thesis: the corpus on which every experiment is run, the four techniques placed under ablation, the three input transformations they operate on, and the training regime held constant across all of them.

It is organised as a single continuous argument rather than as ten independent surveys: a fact is established once, in the section that owns it, and cross-referenced thereafter. Each section runs from what the technique addresses, through how it works and what the published evidence shows, to its limitations and the implications for this thesis, closing with the gap and how this work differs. Two claims recur --- that on a corpus of 156 clips several decisions usually presented as modelling choices are in fact forced by properties of the data, and that several results usually attributed to architecture are attributable to the evaluation protocol. \autoref{sec:synthesis-and-gap} consolidates the nine closing statements into a single research gap and collects every declared departure from the reviewed methods.


% =============================================================================
\section{The Evaluation Corpus: CASME II and the Spontaneous Micro-Expression Benchmark Family}
\label{sec:corpus}
% =============================================================================

This section covers the corpus literature only --- the papers that introduce, specify or standardise evaluation on the database used throughout this thesis. It is treated first because on a database of this size several design decisions usually presented as modelling choices are \textit{forced} by properties of the data.

\subsection{Why a dedicated spontaneous corpus was necessary}
\label{sec:why-spontaneous}

Automatic facial expression recognition matured on large corpora of \textit{posed} expressions --- CK+, MUG, MMI, JAFFE and Multi-PIE --- reporting accuracies above 90\,\% on the six basic posed emotions by the early 2010s~\cite{yan2014}. Micro-expression research could not follow that path: as \autoref{sec:phenomenon} establishes, a micro-expression is involuntary and cannot be \textit{performed} on instruction. Yan et al.~\cite{yan2014} review a pre-2013 landscape of only four micro-expression databases, of which only two were spontaneous (\autoref{tab:spontaneous-corpora}). Two problems in it motivated CASME II directly.

The first is \textbf{construct validity}. Yan et al.~\cite{yan2014} and Li et al.~\cite{liX2013} both note that posed ``micro-expressions'', produced deliberately on instruction, differ systematically from spontaneous ones --- Li et al.~\cite{liX2013} further observing that USF-HD adopted a 2/3-second duration criterion longer than the definitional upper bound of \autoref{sec:phenomenon}. A system trained on posed data is therefore not demonstrably a micro-expression recogniser; it may be a recogniser of fast deliberate movement.

The second is \textbf{acquisition quality}. Defined jointly by very short duration and very low intensity, the phenomenon sits close to the noise floor of ordinary video: Yan et al.~\cite{yan2014} argue that CASME's 60~fps coarsely samples a movement lasting a fraction of a second --- while marking the benefit of higher rates as unresolved --- and that flicker under alternating-current lighting contaminated earlier recordings. Elicitation was a third difficulty: high-stakes deception paradigms produced samples heavily contaminated by conversational and other emotion-irrelevant movement~\cite{yan2014}, and the paradigm that proved workable, which CASME II adopts, is to have participants watch high-arousal video episodes while attempting to keep the face neutral~\cite{liX2013}.

\begin{table}[htbp]
	\centering
	\footnotesize
	\setlength{\tabcolsep}{4pt}
	\begin{tabular}{l c c c c c c c}
		\hline
		\textbf{Corpus} & \textbf{Subj.} & \textbf{Micro samples} & \textbf{fps} & \textbf{Facial res.} & \textbf{Spont.} & \textbf{FACS} & \textbf{Classes} \\
		\hline
		\multicolumn{8}{l}{\textit{The pre-CASME II landscape, reproduced from Yan et al.~\cite{yan2014} (Table 1)}} \\
		USF-HD~\cite{shreve2011}                       & ---      & 100 & 30  & ---              & No  & No  & 6 \\
		Polikovsky et al.\textsuperscript{\textdagger} & 10       & 42  & 200 & ---              & No  & No  & 6 \\
		SMIC~\cite{liX2013}                            & 16 valid & 164 & 100 & 190 $\times$ 230 & Yes & No  & 3 \\
		CASME (2013)\textsuperscript{\textdagger}      & 19 valid & 195 & 60  & 150 $\times$ 190 & Yes & Yes & 7 \\
		\hline
		\multicolumn{8}{l}{\textit{The spontaneous corpora that define the current context}} \\
		\textbf{CASME II}~\cite{yan2014} & \textbf{26 valid} & \textbf{247 / 255} & \textbf{200} & \textbf{280 $\times$ 340} & \textbf{Yes} & \textbf{Yes} & \textbf{5 (7 raw)} \\
		SAMM~\cite{davison2018}          & 32 & 159             & 200 & 400 $\times$ 400 & Yes & Yes & 7 \\
		CAS(ME)$^2$~\cite{qu2016}        & 22 & 53 (+250 macro) & 30  & ---              & Yes & Yes & 4 \\
		\hline
		\multicolumn{8}{l}{\textit{Ethnicities / mean age (SD):} SMIC 3 / 26.7 \quad CASME 1 / 22.03 (1.60) \quad \textbf{CASME II 1 / 22.03 (1.60)}} \\
		\multicolumn{8}{l}{\quad SAMM 13 / 33.24 (11.32) \quad CAS(ME)$^2$ 1\textsuperscript{*} / 22.59 (2.20)} \\
		\hline
	\end{tabular}
	\caption[Posed and spontaneous micro-expression corpora]{Posed and spontaneous micro-expression corpora, compiled from Yan et al.~\cite{yan2014} (Table 1), Davison et al.~\cite{davison2018} (Table 1), Li et al.~\cite{liX2013} and Qu et al.~\cite{qu2016}. \textdagger{}~Described in Yan et al.; the primary papers lie outside the review corpus and are not cited separately. \textsuperscript{*}~Inferred from the recruitment description; Qu et al. do not state it. SMIC's demographic figures are those of all 20 recruited participants; the 16 who produced usable clips have a mean age of 28.1 and span two ethnicities.}
	\label{tab:spontaneous-corpora}
\end{table}

\subsection{CASME II: construction, apparatus, and annotation}
\label{sec:casme2-construction}

CASME II~\cite{yan2014} is the corpus on which every experiment in this thesis is performed, and several of its specifics propagate into the data pipeline used here.

\textbf{Participants and apparatus.} Thirty-five participants were recruited, mean age 22.03 years (SD = 1.60), and watched high-valence video episodes pre-rated by an independent panel of twenty raters, under instruction to suppress all facial movement. Recording sustained \textbf{200 fps} at 640~$\times$~480 with no inter-frame compression, under four LED lamps with umbrella reflectors giving steady flicker-free illumination --- a deliberate correction of the artefacts observed in SMIC. The facial region is approximately \textbf{280 $\times$ 340 pixels}, over three times CASME's facial area and roughly double SMIC's.

\textbf{Selection and annotation.} From nearly 3,000 elicited movements a small subset was retained, very subtle movements being discarded because their onset and offset frames could not be reliably located. Each retained sample carries onset and offset frame indices, Facial Action Coding System (FACS) action-unit labels from two independent coders with disagreements arbitrated by discussion, and an emotion label. Yan et al.~\cite{yan2014} argue against forcing spontaneous micro-expressions into the six basic categories, since a single action unit is frequently ambiguous --- AU4 (brow lowerer) may indicate disgust, anger, attention or tension --- and assign the emotion label from a \textit{combination} of the coded AUs, the participant's self-report, and the emotional content of the eliciting episode.

\textbf{Preprocessing in the original release.} The baseline applies a three-step normalisation from which this thesis departs: a frontal neutral model face is chosen; 68 Active Shape Model landmarks are detected on the \textit{first frame only} of each clip and registered to the model face by a Local Weighted Mean transform; that transform is applied unchanged to every frame, and the face is cropped by eye position. Yan et al.~\cite{yan2014} justify landmarking only the first frame on two grounds --- that rigid head motion within so short a window is negligible, and that per-frame ASM landmarks are unstable enough to introduce jitter where the face did not in fact move.

\textbf{A note on sample count.} The published paper reports \textbf{247} micro-expressions from 26 valid participants; the coding file distributed with the released database contains \textbf{255}. This discrepancy, a consequence of post-publication revision of the coding, is rarely acknowledged and is one reason sample counts reported for ``CASME II'' vary between papers. The present work uses the released 255-sample coding.

\subsection{The label taxonomy, and the sparsity problem it creates}
\label{sec:label-taxonomy}

Yan et al.~\cite{yan2014} provide five main emotion categories with AU-based criteria, and are candid about the resulting imbalance: ``some types of facial expressions are difficult to elicit in laboratory situations, thus the samples in different categories distributed unequally.'' The released coding is finer-grained still, exposing seven raw labels; counting directly from the label table used in this thesis (\texttt{Processed\_Data/master\_thesis\_labels.csv}, 255 rows) gives the middle columns of \autoref{tab:casme2-labels}.

\begin{table}[htbp]
	\centering
	\footnotesize
	\begin{tabular}{l l c c l c}
		\hline
		\multicolumn{3}{c}{\textbf{Five categories, as published}} & \multicolumn{2}{c}{\textbf{Released coding (255)}} & \textbf{Grouped} \\
		\textbf{Emotion} & \textbf{AU criterion} & \textbf{N} & \textbf{N} & \textbf{Raw label} & \textbf{class} \\
		\hline
		Happiness  & AU6 or AU12                           & 33  & 32 & happiness  & Positive \\
		Disgust    & AU9, AU10, or AU4+AU7                 & 60  & 63 & disgust    & Negative \\
		Surprise   & AU1+2, AU25, or AU2                   & 25  & 25 & surprise   & Surprise \\
		Repression & AU15 or AU17, alone or combined       & 27  & 27 & repression & Negative \\
		Others     & other emotion-related facial movement & 102 & 99 & others     & \textit{discarded} \\
		---        & ---                                   & --- & 7  & sadness    & Negative \\
		---        & ---                                   & --- & 2  & fear       & Negative \\
		\hline
		\multicolumn{6}{l}{\textbf{Working set:} Negative $63+27+7+2 = $ \textbf{99}; Positive \textbf{32}; Surprise \textbf{25}; \textbf{N = 156}} \\
		\hline
	\end{tabular}
	\caption[CASME II label taxonomy and the three-class working set]{The CASME II label taxonomy at three granularities: the five published categories with their AU criteria and frequencies, after Yan et al.~\cite{yan2014} (Table 3); the seven raw labels of the released coding as parsed for this thesis; and the three-class valence grouping that yields the 156-clip working set used throughout.}
	\label{tab:casme2-labels}
\end{table}

\textbf{Seven-class classification is not viable on this corpus.} Fear is represented by two clips in the entire database and sadness by seven. Under any subject-disjoint protocol --- and \autoref{sec:original-baseline} shows subject-disjointness has been the norm for this benchmark since its introduction --- a two-clip class cannot be meaningfully trained or tested: a fold either contains the class or does not, and per-class metrics degenerate.

The field's response has been to group by affective valence. The convention originates with SMIC, released with exactly three classes --- positive, negative, surprise~\cite{liX2013} --- and was carried into cross-corpus work by the Second Facial Micro-Expression Grand Challenge (MEGC 2019)~\cite{see2019}. MEGC states the scheme twice and not identically: as a challenge-wide rule in prose, with Negative covering Repression, Anger, Contempt, Disgust, Fear and Sadness, and again as per-corpus footnotes under its sample table, where CASME II's Negative narrows to Disgust and Repression alone. \autoref{tab:casme2-labels} follows the challenge-wide rule.

\begin{figure}[htbp]
	\centering
	\includegraphics[width=\textwidth]{fig3_1_dataset.png}
	\caption[Composition of CASME II]{Composition of CASME II. Left: all 255 released clips by raw label, with the excluded ``others'' category in grey. Right: the three-class pool of 156 clips used throughout this thesis, in the ratio 99 : 32 : 25; the grey ``others'' bar is shown again for scale and is not part of the pool.}
	\label{fig:dataset-composition}
\end{figure}

\textbf{A declared divergence from the MEGC convention.} MEGC 2019 reports the CASME II portion of its composite database as \textbf{145} samples and footnotes that ``Negative class of CASME II: Disgust and Repression''~\cite{see2019} --- building its subset from the paper's \textit{five-category} release, in which sadness and fear are absorbed into ``others''. This thesis uses the released seven-label coding and assigns sadness and fear to Negative, following the mapping MEGC gives for \textbf{SAMM} rather than the one it gives for CASME II. The consequence is an \textbf{eleven-clip} difference: N = 156 against N = 145, with a Negative class of 99 rather than 88. Nine of the eleven are the sadness (7) and fear (2) clips the seven-label coding exposes; the remaining two are unaccounted for, since disgust and repression together number 90 in the released coding against the 88 the challenge reports --- a residual noted rather than resolved. The direction is mild: the nine added clips are all in the majority class and drawn from the categories the coders found hardest to assign, so if anything they make the present task marginally harder than the challenge's.

The residual imbalance is severe in either case --- 99 : 32 : 25, roughly 4 : 1.3 : 1 --- and everything in \autoref{sec:corpus-threats} concerning metric choice, sampling and loss follows from it.

\subsection{The original baseline, and what it does and does not establish}
\label{sec:original-baseline}

Yan et al.~\cite{yan2014} provide a baseline using Local Binary Patterns from Three Orthogonal Planes (LBP-TOP)~\cite{zhaoG2007} with a Support Vector Machine. LBP-TOP extends the static LBP operator by extracting codes on the XY, XT and YT planes of the video volume and concatenating the three histograms, so that the XT and YT planes encode dynamic rather than static texture. The baseline divides the face into 5~$\times$~5 blocks, sets $P = 4$ neighbouring points, and sweeps the spatial radii $R_X, R_Y \in \{1,2,3,4\}$ and temporal radius $R_T \in \{2,3,4\}$; $R_T = 1$ is deliberately excluded because at 200~fps the change between adjacent frames is negligible --- an early acknowledgement that a high sampling rate is not automatically an advantage. Five-class accuracies under \textbf{leave-one-subject-out cross-validation} cluster in the high fifties to low sixties, the best configuration ($R_X = R_Y = 1$, $R_T = 4$) reaching \textbf{63.41\,\%}.

Three observations matter here.

\textbf{First, leave-one-subject-out is native to this benchmark.} It is not a stricter protocol adopted later by cautious authors; it is the protocol of the paper that released the corpus. Any evaluation reporting a single random train/test split departs from the benchmark's own convention.

\textbf{Second, the baseline is reported as accuracy alone}, with no per-class or macro-averaged score. On a five-class problem in which ``others'' accounts for 102 of 247 samples, accuracy is a weak instrument: the majority-class-constant classifier scores 41\,\% without modelling anything. In the three-class working set the floor is higher still --- always predicting Negative yields \textbf{0.6346 accuracy} at a macro F1 of only \textbf{0.2588}. A headline accuracy near 63.41\,\% is therefore not, on its own, evidence of a working system.

\textbf{A caution on this figure.} Yan et al.~\cite{yan2014} state in both their abstract and their results prose that leave-one-subject-out was applied to obtain 63.41\,\% --- but the caption of the very table reporting it reads ``leave-one-out cross-validation'', without the subject qualifier. Li et al.~\cite{liX2018}, reproducing the same figure in their comparison table, mark it with an asterisk denoting ``results achieved using leave-one-sample-out cross validation'' --- a materially easier protocol. The sources therefore disagree on the protocol behind the most widely quoted number on this benchmark, and the originating paper disagrees with itself. This thesis follows that paper's prose, but records the discrepancy rather than resolving it, because it affects how the figure should be read in any comparison.

\textbf{Third, the baseline's sensitivity to a purely descriptive hyper-parameter is itself informative.} Accuracy moves by several points across the radius sweep with no change to the data, the split or the classifier. Where effect sizes of that magnitude arise from descriptor settings, differences of similar magnitude between architectures cannot be interpreted without a variance estimate --- a limitation this thesis inherits and declares.

\subsection{Sibling corpora, and why this study remains single-corpus}
\label{sec:sibling-corpora}

Three further databases define the context in which CASME II is used; \autoref{tab:spontaneous-corpora} gives their properties, and each isolates a specific limitation of CASME II.

\textbf{SMIC~\cite{liX2013}}, the first spontaneous micro-expression database, is the origin of the three-class taxonomy. Its high-speed subset of 164 clips is distributed 51 positive, 70 negative, 43 surprise, and its baseline --- LBP-TOP with a temporal interpolation model (TIM) normalising each clip to ten frames, under leave-one-subject-out --- reaches \textbf{48.78\,\%}. Two elements of its design are inherited directly here: the three-class valence grouping, and temporal interpolation to normalise clips of unequal length. It is not used as a corpus because it is neither FACS-coded nor apex-annotated.

\textbf{SAMM~\cite{davison2018}} addresses the demographic narrowness of the CASME family, contributing 13 ethnicities and a mean age of 33.24 years against CASME II's single ethnicity and 22.03. Davison et al.~\cite{davison2018} observe that a corpus drawn from one ethnicity and one narrow age band restricts analysis to ``similar looking participants'', and tailor stimulus selection to each participant in advance rather than relying on post-hoc self-report. The design is stronger than CASME II's on both counts, but the corpus is smaller in usable samples and spreads them across seven categories with the same sparsity difficulty.

\textbf{CAS(ME)$^2$~\cite{qu2016}} provides 250 macro-expression and 53 micro-expression samples from the same participants under identical conditions, which makes it the natural resource for studying macro-to-micro domain transfer. That strategy produced the strongest entry in MEGC 2019 --- Liu et al.'s EMR~\cite{liu2019}, the only reported submission to use macro-expression data, though it drew that data from CK+ rather than from CAS(ME)$^2$~\cite{see2019}. Two properties limit its use here: 30~fps is an order of magnitude below CASME II's and poorly matched to sub-200-ms movements, and only 53 micro-expression samples exist. Its baseline is also evaluated leave-\textit{one-video}-out, which does not enforce subject-disjointness and is therefore not comparable to the protocol used here.

\textbf{Why this thesis uses CASME II alone.} It has the highest temporal resolution of any spontaneous corpus of comparable size, a precondition for magnification in a defined frequency band and for dense frame-to-frame flow; it is FACS-coded with onset, apex and offset annotation, which makes temporal normalisation to a fixed length well-defined rather than heuristic; and it remains the most widely reported single-corpus benchmark, which makes comparison against published results meaningful. The cost --- that no claim here extends to other populations, capture conditions or corpora --- is stated as a declared limitation rather than argued away.

\subsection{How CASME II is evaluated in the field: MEGC 2019 and the metric convention}
\label{sec:megc-metrics}

The most consequential development in how CASME II is \textit{used}, as opposed to how it was built, is the Composite Database Evaluation (CDE) protocol of the Second Facial Micro-Expression Grand Challenge~\cite{see2019}. This thesis adopts the challenge's metric definitions while declining its composite training regime, so the protocol warrants precise description.

CDE merges SMIC, CASME II and SAMM into one composite database under the three-class grouping of \autoref{sec:label-taxonomy}, giving \textbf{442 samples} --- Negative 250, Positive 109, Surprise 83 --- of which CASME II contributes 145. Leave-one-subject-out is run across the composite, which the organisers note ``ensures subject-independent evaluation''. Plain accuracy is rejected explicitly on grounds of class imbalance, and two balanced metrics mandated: \textbf{Unweighted F1 (UF1)}, the macro-averaged F1-score, which \textit{accumulates} true positives, false positives and false negatives per class \textbf{over all $k$ folds}, computes a per-class F1 from those pooled counts and averages without weighting; and \textbf{Unweighted Average Recall (UAR)}, the balanced accuracy, obtained by averaging per-class recall.

\[
	F1_c = \frac{2 \cdot TP_c}{2 \cdot TP_c + FP_c + FN_c}, \qquad \text{UF1} = \frac{1}{C}\sum_c F1_c
\]

UF1 is a \textbf{pooled} quantity: counts accumulate across folds \textit{before} any F1 is computed, which \autoref{sec:pooled-vs-fold} establishes is a different estimator from the average of per-fold macro F1 scores. The distinction is not cosmetic here. Because clips are distributed very unequally across subjects, and many single-subject folds contain only one of the three classes, a per-fold macro F1 is bounded well below 1 by fold composition alone and its average across folds bounded correspondingly (\autoref{fig:fold-composition}). The primary metric of this thesis, \textit{pooled macro F1}, is therefore identical in construction to MEGC's UF1; the per-fold average appearing in the raw experimental logs is not a metric the field uses at all.

\begin{figure}[htbp]
	\centering
	\includegraphics[width=\textwidth]{fig3_2_fold_composition.png}
	\caption[Class composition of the 25 leave-one-subject-out folds]{Class composition of the 25 leave-one-subject-out folds, and the per-fold macro-F1 ceiling it implies.}
	\label{fig:fold-composition}
\end{figure}

\begin{table}[htbp]
	\centering
	\begin{tabular}{l c c}
		\hline
		\textbf{Method} & \textbf{Full composite UF1 / UAR} & \textbf{CASME II subset UF1 / UAR} \\
		\hline
		LBP-TOP~\cite{zhaoG2007}                  & 0.5882 / 0.5785 & 0.7026 / 0.7429 \\
		Bi-WOOF~\cite{liong2018}                  & 0.6296 / 0.6227 & 0.7805 / 0.8026 \\
		OFF-ApexNet~\cite{liong2019a}             & 0.7196 / 0.7096 & 0.8764 / 0.8681 \\
		Quang et al.\textsuperscript{\textdagger} & 0.6520 / 0.6506 & 0.7068 / 0.7018 \\
		Zhou et al.\textsuperscript{\textdagger}  & 0.7322 / 0.7278 & 0.8621 / 0.8560 \\
		Liong et al.~\cite{liong2019b}, STSTNet   & 0.7353 / 0.7605 & 0.8382 / 0.8686 \\
		Liu et al.~\cite{liu2019}, EMR            & \textbf{0.7885 / 0.7824} & 0.8293 / 0.8209 \\
		\hline
	\end{tabular}
	\caption[MEGC 2019 Composite Database Evaluation results]{MEGC 2019 CDE results, reproduced from See et al.~\cite{see2019} (Table IV). \textdagger{} outside the review corpus; reported here only as it appears in that table.}
	\label{tab:megc-cde-results}
\end{table}

\textbf{One structural caveat governs these numbers.} Every row of \autoref{tab:megc-cde-results} comes from a model trained on the \textit{composite} database --- roughly 442 samples across three corpora and 68 subjects --- then scored on the CASME II portion of the pooled predictions. The CASME II column is therefore not a CASME II-only result: it reports how a cross-corpus model performs on CASME II clips, with approximately \textbf{2.8 times the training data} available in any single-corpus fold, where a single-corpus LOSO experiment of the kind conducted here trains each fold on roughly 150 clips. The two regimes are not interchangeable, and single-corpus results should not be compared directly against \autoref{tab:megc-cde-results} without stating the difference.

Two further observations from the challenge bear on this thesis's design. The top three submissions all chose \textbf{optical flow} as their feature representation, a convergence that supports the motion-based input used here; and one submission substituted mid-position frames for annotated apex frames with little loss, suggesting that precise apex localisation may be less critical than assumed and that methods consuming the full onset-to-offset sequence are not obviously disadvantaged.

\subsection{Threats to validity carried by the corpus}
\label{sec:corpus-threats}

The properties in \autoref{tab:corpus-threats} are not criticisms of CASME II, which remains the best available resource for this task. They are constraints any study built on it inherits, and each is answered by a specific design decision --- which is the sense in which, on a corpus of this size, several choices usually presented as modelling decisions are determined by the data. Fold and duration statistics are computed directly from the 156-clip label table used in this thesis.

\begin{table}[htbp]
	\centering
	\scriptsize
	\begin{tabular}{@{}p{0.3cm} p{2.0cm} p{5.3cm} p{5.6cm}@{}}
		\hline
		& \textbf{Threat} & \textbf{What the corpus contains} & \textbf{Design decision it forces} \\
		\hline
		(a) & Scale
		& 156 usable clips across three classes --- small enough that parameter count is itself a risk factor.
		& \textbf{Parameter-free or low-parameter components} wherever a choice exists, motivating SimAM~\cite{yang2021}, which adds zero learnable parameters; shallow rather than deep spatial backbones, extending the model-shrinking argument of Xia et al.~\cite{xia2020b}, which they establish for the composite-database setting rather than a single small corpus. \\
		(b) & Class imbalance
		& 99 : 32 : 25. The majority-class-constant classifier attains 0.6346 accuracy at a macro F1 of 0.2588.
		& Minority oversampling via a balanced sampler with focal loss, inverse-frequency class weighting \textbf{disabled}, since both corrections together induced single-class collapse. Accuracy is never reported alone, and the always-Negative floor is quoted alongside every result. \\
		(c) & Subject-level non-uniformity
		& \textbf{Subject 17 alone contributes 33 clips, 21\,\% of the working set}; subjects 8, 10 and 21 contribute one each, so fold sizes vary by a factor of 33. \textbf{Ten of the 25 folds contain a single class}, eight contain two, only seven all three.
		& \textbf{Full 25-fold leave-one-subject-out}, every subject held out exactly once, no clip untested, subject-disjointness guaranteed by construction. Primary metric is \textbf{pooled macro F1}, identical in construction to MEGC's UF1~\cite{see2019}; per-fold averaged macro F1 is rejected, since fold composition bounds it below the study's own target. \\
		(d) & A subject that disappears
		& Of the 26 participants, subject 18's three clips are all ``others'' and are removed by the three-class filter.
		& The LOSO run has \textbf{25 folds, not 26} --- stated explicitly for reproducibility, since fold counts differ between papers using this corpus. \\
		(e) & Demographic homogeneity
		& All participants Chinese, mean age 22.03 (SD 1.60) --- named by Davison et al.~\cite{davison2018} as a principal motivation for SAMM. The MEGC organisers saw the effect in reverse: adaptation from CK+ helped on SMIC and SAMM but ``not so for CASME II, which contain predominantly Chinese subjects''~\cite{see2019}.
		& Input is a \textbf{motion representation} --- dense optical flow $(u,v)$ plus optical strain --- rather than raw pixels, since motion is invariant to identity and illumination whereas appearance is exactly what a 26-subject corpus would let a network memorise. \\
		(f) & Laboratory ecology
		& Frontal recording, fixed flicker-free lamps, participants told not to move and not speaking; Yan et al.~\cite{yan2014} flag ``natural conversation and interaction'' as future work. Almost none of the head motion, illumination variation, occlusion or speech-related movement a deployed system would face is present.
		& No cross-dataset or in-the-wild claim is made anywhere in this thesis; generalisation is declared as further work. \\
		(g) & Annotation dependence
		& Labels derive from coded AUs, self-report and stimulus content --- a construct produced by trained coders under a rubric, with disagreements arbitrated by discussion rather than reported as residual uncertainty.
		& The coding is taken as given; its construct status is recorded as a declared limitation rather than treated as measurement error. \\
		(h) & Duration heterogeneity
		& Onset-to-offset lengths range from \textbf{31 to 126 frames}, median 66 --- at 200~fps, \textbf{0.155\,s to 0.63\,s}, median 0.33\,s. The upper end exceeds the 1/2-second criterion of \autoref{sec:phenomenon}.
		& Every clip temporally resampled to a fixed \textbf{T = 32} frames, following the TIM precedent of SMIC's baseline~\cite{liX2013}. \\
		(i) & Version ambiguity
		& Three sample counts circulate under the name ``CASME II'' --- 247 (the paper), 255 (the released coding), 145 (the MEGC subset) --- and both five- and seven-class label sets are in use, rarely stated.
		& The released 255-clip coding is used, N is reported explicitly alongside every figure, and the eleven-clip divergence from the MEGC subset is declared in \autoref{sec:label-taxonomy}. \\
		\hline
	\end{tabular}
	\caption[Threats to validity carried by CASME II and the decisions they force]{Threats to validity carried by CASME II, and the design decision each forces.}
	\label{tab:corpus-threats}
\end{table}

\subsection{Implications for this research}
\label{sec:corpus-implications}

Three further corpus properties are not threats but still determine settings. The \textbf{200 fps} capture rate fixes the Eulerian magnification pass-band at \textbf{5--25 Hz}, matched to the recording rate. The \textbf{seven raw labels}, with fear at two clips and sadness at seven, force the three-class valence grouping of \autoref{sec:label-taxonomy}. And the baseline's \textbf{sensitivity to a descriptive hyper-parameter} (\autoref{sec:original-baseline}) sets the standard of evidence: effect sizes below the resolution of a 156-clip corpus are reported as unresolved rather than as findings, and no confidence interval is available, since the single-seed design and unretained per-clip predictions preclude one.

Three consequences then shape the interpretation of the results rather than merely the construction of the pipeline. \textbf{First, the corpus determines the metric, and the metric determines the conclusion:} the fold composition of \autoref{tab:corpus-threats}(c) caps a macro F1 averaged over folds by arithmetic alone, before any classifier is considered, so reporting the pooled quantity is not a stylistic preference but the correction MEGC 2019 already standardised. \textbf{Second, a single held-out split is not a reliable measuring instrument here:} with 25 subjects distributed as unevenly as \autoref{tab:corpus-threats}(c) describes, which subjects fall in the test set can move the result by more than any architectural change under study, so only a complete leave-one-subject-out run is reported and the incomplete evaluations that preceded it are set out with their exclusion reasons rather than suppressed. \textbf{Third, the small-sample regime predicts which components should help:} 156 clips cannot supply the data needed to learn a general spatio-temporal feature extractor from scratch, particularly once the input has been reduced to flow and strain --- that is, once much of that extraction has been performed analytically before the network sees the data --- whereas a component modelling the temporal arc of an expression exploits exactly the structure the corpus's onset/apex/offset annotation guarantees is present and correctly delimited.

\subsection{Limitations of the existing corpus literature, and how this work differs}
\label{sec:corpus-lit-limitations}

The corpus literature establishes three things well and one thing poorly. It establishes the phenomenon's definitional boundaries, an elicitation methodology producing genuinely spontaneous samples, and a capture standard --- 200 fps, flicker-free illumination, FACS coding with onset/apex/offset annotation --- that makes fine-grained motion analysis feasible. What it does not establish is a stable basis for comparing methods. The original baseline reports accuracy only, on a five-class problem whose majority class is 41\,\% of the data, and its own hyper-parameter sweep moves the headline figure by several points with no change to the model. Sample counts, label granularity and class groupings differ across papers using the same database --- 247, 255 and 145 all appear under the name ``CASME II'' --- and are frequently not stated. MEGC 2019's CDE protocol fixed the metric problem convincingly, but within a composite training regime that makes its per-database columns non-comparable to single-corpus work, a distinction subsequent papers routinely elide. Above all, only one study in this group quantifies how much of a reported difference between methods is attributable to the evaluation protocol rather than the methods themselves --- Li et al.~\cite{liX2018} add a single matched row of their own method under both protocols --- and none does so across methods, despite a subject distribution that makes the question unavoidable.

This thesis differs in three respects. It runs the \textbf{complete} 25-fold leave-one-subject-out protocol over \textbf{every} valid cell of a four-component ablation matrix rather than reporting a single proposed configuration, so each component's contribution is measured from matched pairs differing in exactly one factor. It adopts MEGC's pooled macro F1 as its primary metric while explicitly declining the composite training regime, so its numbers describe single-corpus performance and are labelled as such. And it reports one complete protocol rather than the strongest of several, setting out which earlier evaluations were run and why each is excluded rather than suppressing that history. The intended contribution is therefore not a recognition result on CASME II, but a measurement of how much of such a result is attributable to each stage of the pipeline that produces it.

The 200 fps capture rate established here is what lets Eulerian Video Magnification, reviewed next, target a matched temporal band.


% =============================================================================
\section{Motion Magnification: Eulerian Video Magnification as a Pre-Processing Stage}
\label{sec:evm-review}
% =============================================================================

This section reviews Eulerian Video Magnification (EVM) through the four review-corpus papers that apply it: Bai et al.~\cite{bai2021}, Li et al.~\cite{liX2018}, and Li, Huang and Zhao~\cite{liY2018,liY2021}. The original formulations they build on --- Wu et al.'s amplitude-based method and Wadhwa et al.'s phase-based method --- lie outside the corpus and are described only as attributed through those four papers.

\subsection{The problem magnification addresses}
\label{sec:evm-problem}

The corpus design of \autoref{sec:why-spontaneous} responds to the \textit{short-duration} half of the micro-expression definition but not to the \textit{low-intensity} half. CASME II's 200~fps sampling spreads a movement completing in under 200~ms across dozens of frames~\cite{yan2014}, yet that movement may still occupy only a few grey levels of change across a small patch of skin --- and Yan et al.~\cite{yan2014} report discarding samples outright because the movement was too faint for coders to place an onset or offset. Li et al.~\cite{liX2018} state the difficulty directly: ``the intensity levels of facial movements are too low to be distinguishable.'' Their remedy is to amplify the motion itself before any descriptor is computed, so that the \textit{difference between expression categories is enlarged} rather than the classifier being made more sensitive. This is a data-level intervention: the model is untouched and the modification is entirely upstream, which is what makes EVM cleanly separable as an ablation factor in this thesis.

\subsection{How Eulerian magnification works, and the two variants}
\label{sec:evm-how}

\textbf{Amplitude-based magnification (AMM)} is the formulation set out in \autoref{sec:motion-magnification} --- Laplacian pyramid, temporal band-pass, multiplication by $\alpha$, reconstruction --- realising its band-pass step with a \textbf{Butterworth} filter. It is the variant implemented in this thesis. \textbf{Phase-based magnification (PMM)}, of which Bai et al.~\cite{bai2021} give the clearest account in the corpus, decomposes each frame with \textbf{octave complex steerable pyramids} over four orientations and applies the temporal filter to the \textit{local phase} rather than the amplitude. Bai et al.~\cite{bai2021} explain the advantage in signal terms: the steerable pyramid ``has impulse response with spatial support,'' which makes it easier to isolate the intended temporal frequencies while suppressing the remainder. The practical consequences they report are that the phase-based method \textbf{supports larger amplification factors} and is \textbf{notably less sensitive to noise}.

Three parameters are therefore free in any EVM application, and all three recur below: the \textbf{temporal pass-band}, which decides \textit{which} motions are amplified; the \textbf{magnification factor $\alpha$}, which decides \textit{how much}; and the \textbf{pyramid depth}, which decides at which spatial scales amplification occurs.

\textbf{Setting the temporal band.} Bai et al.~\cite{bai2021} report the rule given by the original EVM authors: the pass-band should correspond to the duration of the phenomenon being amplified --- for micro-expressions, 1/5\,s to 1/25\,s, equivalently \textbf{5 Hz to 25 Hz}. They also report that applying this full band produced ``noticeable noise'' in their setting, and that they narrowed it to 15--25 Hz accordingly. This is the only explicit guidance on band selection anywhere in the review corpus, and it is the basis for the band used in this thesis.

\subsection{Evidence that magnification improves recognition}
\label{sec:evm-evidence}

The strongest empirical case in the review corpus is Li et al.~\cite{liX2018}, who sweep the magnification factor systematically rather than fixing it. Clips are magnified at what the paper describes as ten levels while listing nine values, $\alpha \in \{1, 2, 4, 8, 12, 16, 20, 24, 30\}$ --- an inconsistency in the source rather than in this account --- where $\alpha = 1$ denotes no magnification. Each magnified clip is temporally interpolated to ten frames (TIM10), and LBP, HOG and HIGO descriptors are extracted and classified, on both SMIC and CASME II under leave-one-subject-out. Two findings emerge.

\textbf{Magnification helps, consistently}, ``for all three kinds of features on all four testing datasets''. On CASME II the effect is substantial (\autoref{tab:evm-effect}), and the best configuration of their full framework --- magnification, TIM10 and HIGO --- reaches 67.21\,\% on CASME II and 68.29\,\% on SMIC-HS under LOSO, competitive with or better than the contemporaneous state of the art.

\begin{table}[htbp]
	\centering
	\begin{tabular}{l c c c}
		\hline
		\textbf{Descriptor} & \textbf{No magnification} & \textbf{With magnification} & \textbf{$\Delta$} \\
		\hline
		LBP  & 55.87\,\% & 60.73\,\% & \textbf{+4.86} \\
		HOG  & 57.49\,\% & 63.97\,\% & \textbf{+6.48} \\
		HIGO & 57.09\,\% & 67.21\,\% & \textbf{+10.12} \\
		\hline
	\end{tabular}
	\caption[Effect of Eulerian motion magnification on CASME II]{Effect of Eulerian motion magnification on CASME II under leave-one-subject-out, after Li et al.~\cite{liX2018} (Table 6).}
	\label{tab:evm-effect}
\end{table}

\textbf{The relationship between $\alpha$ and accuracy is non-monotonic.} Li et al.~\cite{liX2018} describe the resulting curves as ``rainbow-shaped,'' with best performance ``generally achieved when the motion is magnified in the range of $[8, 16]$,'' and give a mechanistic account of both failure modes: ``magnification at lower levels might not be enough to reveal the ME motion progress; on the other hand, magnification at higher levels degrades the performance because too many artifacts are induced.'' Magnification is therefore not a monotone improvement to be maximised but a parameter with an interior optimum that must be selected.

Bai et al.~\cite{bai2021} provide the second quantitative demonstration using a fundamentally different downstream model --- a pre-trained VGGFace2 (ResNet-50) spatial encoder feeding a bidirectional LSTM. On SMIC they reach \textbf{75.76\,\% with phase-based magnification against a 60.60\,\% baseline} drawn from an earlier VGGFace2 + Bi-LSTM publication rather than from a magnification-off arm of their own, and the phase-based variant outperforms the amplitude-based one. Because the pipeline shares nothing with Li et al.'s hand-crafted descriptors, the two results together establish that the benefit is not an artefact of a particular feature family --- at least when the downstream representation is appearance-based.

\subsection{Magnification in apex-frame pipelines, and the $\alpha$ discrepancy}
\label{sec:evm-apex}

Two further papers use EVM in a different regime, and the contrast is informative. Li, Huang and Zhao~\cite{liY2018} address whether a micro-expression can be recognised from a \textit{single} apex frame, arguing for magnification specifically in that setting: the apex frame carries the most emotional information, but ``the intensity level of apex frame is not distinguished due to subtle change, especially for deep model that describes high-layer abstract information.'' They magnify the spotted apex frame with \textbf{$\alpha = 30$} and fine-tune a VGG-Face descriptor on the result; because a single frame per clip is too little data to train on, they exploit the high frame rate to augment, retaining the apex frame plus its two preceding and two following frames and so expanding the training set fivefold. They restate the trade-off Li et al.~\cite{liX2018} quantified --- ``bigger values of motion amplification level lead to larger scale of motion amplification, but also can cause bigger displacement and artifacts.'' Li, Huang and Zhao~\cite{liY2021} extend this into a joint local--global learning framework with frequency-domain apex detection, carrying the same setting forward --- magnification level ``set as 30 in our framework'' --- with the same $\pm$2-frame augmentation producing what they name the Extended Magnified ME (EMME) database.

\textbf{The $\alpha$ discrepancy is easy to misread as a contradiction.} Li et al.~\cite{liX2018} find the optimum at $\alpha \in [8, 16]$ and observe degradation beyond it; Li, Huang and Zhao~\cite{liY2018,liY2021} adopt $\alpha = 30$, at the extreme end of that same sweep. The regimes differ in a way that plausibly explains the gap. The $[8, 16]$ optimum is measured on \textbf{full magnified sequences} described by spatio-temporal descriptors, where inter-frame artefacts accumulate along the temporal axis and corrupt exactly the dynamic texture the descriptor encodes. The $\alpha = 30$ setting is applied to a \textbf{narrow five-frame window around the apex}, consumed by a deep face encoder as largely static appearance, where temporal artefacts have far less opportunity to express themselves --- and where the encoder needs a strong appearance signal to distinguish categories at all. The implication for any new system is that $\alpha$ cannot be transferred between pipelines: it is a parameter of the \textit{combination} of magnification and downstream representation, not of the corpus.

\subsection{Known limitations of magnification}
\label{sec:evm-limitations}

The review corpus is unusually candid about EVM's failure modes, and four are documented.

\textbf{It amplifies everything in band, not only expressions.} This is the most important limitation and the most easily overlooked. Li et al.~\cite{liX2018} decline to use magnification for micro-expression \textit{spotting} at all, ``because it magnifies unwanted motions (e.g. head movements) at the same time.'' EVM has no notion of which motion is of interest: a rigid head translation whose temporal frequency falls inside the pass-band is amplified exactly as a brow raise is. In a recognition pipeline this is mitigated by prior face alignment and cropping --- Li et al.~\cite{liX2018} magnify only after normalising and cropping to the eye-defined rectangle --- but mitigated, not eliminated.

\textbf{Artefacts grow with $\alpha$.} All four papers report this: Li et al.~\cite{liX2018} as the descending limb of the rainbow curve, Li, Huang and Zhao~\cite{liY2018,liY2021} as a design caution, and Bai et al.~\cite{bai2021} directly, as noise that forced them to narrow their temporal band.

\textbf{There is no principled procedure for choosing $\alpha$.} One value in the corpus is obtained by exhaustive sweep~\cite{liX2018}; the apex-frame papers assert theirs --- Li et al.~\cite{liY2018} state $\alpha = 30$ without support, and Li et al.~\cite{liY2021} carry the value forward by citing them. No paper offers a way to select $\alpha$ from properties of the data.

\textbf{Evaluation is confined to appearance-based downstream representations:} LBP, HOG and HIGO in Li et al.~\cite{liX2018}, VGGFace2 features in Bai et al.~\cite{bai2021}, VGG-Face in both Li, Huang and Zhao papers. Every one of these descriptors is computed on \textbf{magnified pixel intensities}. No study in the review corpus measures what magnification contributes when the representation handed to the classifier is itself an explicit motion field rather than an appearance signal.

\subsection{Implications for this research}
\label{sec:evm-implications}
% sec:evm-motion-hypothesis is referenced from Chapter 5 (chapter5.tex:236). The
% hypothesis it names now opens this subsection instead of having one of its own,
% so the label is kept here and resolves to this subsection.
\label{sec:evm-motion-hypothesis}

\autoref{tab:evm-settings} sets out the settings the review determines. Behind them lies a hypothesis the review corpus does not test, stated here because it frames the interpretation of the measurement. EVM amplifies the \textit{amplitude} of intensity variation in a temporal band, and an appearance descriptor computed on magnified frames sees a directly stronger signal. An optical flow field is a different object: flow estimates \textit{displacement}, and amplifying the underlying intensity change by $\alpha$ does not simply scale the recovered displacement by $\alpha$ --- it changes the conditioning of the estimation problem, potentially making faint motion recoverable where it previously was not, while introducing displacement artefacts of the kind all four reviewed papers warn about. Any residual global scaling of flow magnitude is further attenuated by the per-channel normalisation a learned pipeline applies before training. The plausible outcome is that magnification in front of a flow-and-strain representation retains the part of its benefit that consists of \textit{making otherwise-unrecoverable motion recoverable} while losing the part that consists of \textit{increasing contrast for an appearance descriptor} --- so its measured contribution should be smaller on average than the roughly +5 to +10 percentage points of \autoref{tab:evm-effect}, and concentrated in configurations whose downstream model can exploit finer spatial deformation. EVM is measured across six matched pairs of configurations differing in nothing but the tensor directory, which is what allows that prediction to be tested.

\begin{table}[htbp]
	\centering
	\footnotesize
	\begin{tabular}{@{}p{6.4cm} p{7.4cm}@{}}
		\hline
		\textbf{Finding from the review} & \textbf{Decision taken in this thesis} \\
		\hline
		Pass-band should match the phenomenon's duration, 1/5--1/25\,s $\equiv$ 5--25 Hz~\cite{bai2021}, though Bai et al. found that band noisy and narrowed it to 15--25 Hz & Temporal band-pass set to \textbf{5--25 Hz}, the unnarrowed duration rule; Bai et al.'s narrowing was not followed, and this is a declared divergence \\
		Accuracy against $\alpha$ is rainbow-shaped, optimum in \textbf{$[8, 16]$} for full-sequence pipelines~\cite{liX2018} & \textbf{$\alpha = 10$}, inside the reported optimum; the $\alpha = 30$ of the apex-frame papers is explicitly \textit{not} adopted, since this thesis consumes full sequences \\
		Amplitude-based magnification uses a Laplacian pyramid plus temporal band-pass~\cite{bai2021} & Amplitude-based EVM over a \textbf{4-level Laplacian pyramid}, temporal filtering by FFT band-pass \\
		Magnification amplifies head motion as well as expression~\cite{liX2018} & Magnification applied only after the clip has been trimmed onset-to-offset; the risk is acknowledged as a residual threat rather than claimed to be eliminated \\
		Magnification is a purely upstream, data-level operation & Stage 1 is run \textbf{twice}, producing a magnified and a non-magnified tensor set; the EVM switch selects a directory and leaves the network unchanged, making it a clean ablation factor \\
		Magnification is part of the proposed data pipeline, not of the network & EVM is treated as the \textbf{baseline condition} (\texttt{config\_4}), not as a component to be added; the non-magnified arm (\texttt{config\_1}) is the control that isolates it \\
		\hline
	\end{tabular}
	\caption[Review findings and the magnification settings they determine]{Review findings and the magnification settings they determine.}
	\label{tab:evm-settings}
\end{table}

\textbf{A defect found and repaired, which the review made detectable.} In this project's earliest evaluations, each magnification-enabled configuration produced results \textit{bit-identical} to its magnification-disabled counterpart --- identical accuracies, macro F1 scores and confusion matrices. Given that all four reviewed papers report clear and repeatable effects from magnification, an effect of exactly zero across six independent pairs is not a plausible null result but a symptom: both arms were reading the same tensor directory, so half of the ablation matrix was duplicated and the magnification hypothesis had never in fact been exercised. The defect was repaired before the reported run, and wiring the existing verification helper in as a start-up assertion is a declared item of further work.

\subsection{Limitations of the existing work, and how this study differs}
\label{sec:evm-gap}

What the reviewed literature does not establish is threefold. First, every evaluation in the corpus places magnification in front of an \textbf{appearance-based} representation, so the reported gains conflate two distinct mechanisms --- making faint motion recoverable, and increasing contrast for a texture or gradient descriptor --- and no experiment separates them. Second, the magnification factor is selected empirically in every case and transferred between studies without re-validation, despite the $\alpha$ discrepancy of \autoref{sec:evm-apex}. Third, and most consequentially, magnification is always evaluated as an \textit{addition to a fixed pipeline} --- a two-cell comparison, with and without --- so its interaction with the other components of a system is never measured.

This thesis differs on all three counts. It places magnification in front of an explicit \textbf{motion} representation --- dense optical flow and optical strain --- rather than an appearance representation, isolating the recoverable-motion mechanism from the appearance-contrast one. It selects $\alpha$ from the full-sequence optimum rather than the apex-frame value, with the choice justified rather than inherited. And it evaluates magnification as one factor in a \textbf{full factorial ablation}, so its effect is reported not as a single delta but as six matched-pair measurements whose spread across architectures is itself the finding.

What magnification is placed in front of therefore matters as much as its own parameters, and \autoref{sec:flow-review} turns to that motion representation.


\section{Motion Representation: Dense Optical Flow}
\label{sec:flow-review}

Dense optical flow supplies the horizontal and vertical displacement channels used in this thesis. The relevant literature concerns three choices: the estimator, the frame pairs it processes, and whether flow serves as a recognition feature or an alignment tool. These choices determine what information remains available to the learned model.

\subsection{Why motion rather than appearance}
\label{sec:why-flow}

Xu et al.~\cite{xu2017} motivate facial-dynamics descriptors by the difficulty of interpreting micro-expressions through texture alone. Flow describes apparent movement directly, whereas image intensity also records appearance and illumination. This makes displacement a useful representation of subtle facial activity, although it does not guarantee independence from identity or recording conditions: facial geometry, alignment and brightness changes still affect its estimation.

The benchmark evidence supports this choice. The MEGC 2019 organisers report that their three highest-ranked submissions used optical flow~\cite{see2019}. OFF-ApexNet~\cite{liong2019a} uses horizontal and vertical flow as neural-network inputs; STSTNet~\cite{liong2019b} adds strain; Zhao et al.~\cite{zhaoS2021} retain a sequence of flow fields. Their convergence motivates a motion-based input, but is not a controlled demonstration that it always outperforms appearance-based representations.

\subsection{Estimation assumptions}
\label{sec:brightness-constancy}

OFF-ApexNet~\cite{liong2019a} formulates flow using brightness constancy and small-motion assumptions. A local brightness constraint supplies one equation for two displacement components, so estimation also needs spatial regularisation or a local motion model. The output is a pair of fields, $u(x,y)$ and $v(x,y)$; Zhao et al.~\cite{zhaoS2021} additionally describe their magnitude, $\sqrt{u^2+v^2}$.

These assumptions matter more than repeating the derivation. Illumination changes can resemble movement, while smoothing can suppress a weak local contraction together with noise. Magnification can make subtle motion easier to estimate but can also amplify unwanted movement. Flow therefore transforms the measurement problem; it does not remove uncertainty in the source frames.

\subsection{Choice of estimator}
\label{sec:flow-estimator}

The reviewed methods use three estimators. Zhao et al.~\cite{zhaoS2021} adopt Farneb\"ack through its OpenCV implementation, citing its established availability. It estimates displacement through local polynomial approximations. OFF-ApexNet~\cite{liong2019a} instead uses total-variation regularisation with an L1 data term (TV-L1), citing preservation of motion discontinuities and robustness. Xu et al.~\cite{xu2017} use the method of Sun, Roth and Black. These algorithm attributions are taken from the reviewed papers; the original estimator publications are outside the supplied reference corpus.

STSTNet~\cite{liong2019b} specifies onset--apex optical flow but does not name its estimator. Its similarity to OFF-ApexNet is therefore insufficient to attribute TV-L1 to it. Likewise, the advantages given for TV-L1 should not be read as a measured recognition gain over Farneb\"ack on CASME II. The cited studies do not provide a controlled three-estimator comparison supporting such a ranking. This thesis adopts Farneb\"ack as a fixed preprocessing choice and leaves that comparison open.

\subsection{Onset--apex flow or a temporal sequence}
\label{sec:flow-what-between}

OFF-ApexNet~\cite{liong2019a} represents each clip by flow from onset to apex, where expression intensity is greatest, and resizes the two component maps to $28\times28$. STSTNet~\cite{liong2019b} uses the same frame-pair construction. This is an economical representation of the displacement from neutral to peak expression, suitable for a small training set. It omits the intermediate trajectory and the relaxation after the apex.

Zhao et al.~\cite{zhaoS2021} instead construct an eleven-frame sequence using onset, apex and offset annotations, then calculate ten adjacent-pair flow fields. This retains changes in motion across the expression. The distinction is between one displacement summary and a sequence whose order a temporal model can exploit; it is not simply a difference in input size.

The present study follows the sequence approach but uses uniform sampling of existing frames rather than Zhao et al.'s key-frame construction. This preserves a temporal input for the transformer ablation. It does not, by itself, measure the recognition cost of replacing a sequence with an onset--apex pair: that representation comparison is not an experimental arm.

\subsection{Flow for fine alignment}
\label{sec:flow-alignment}

Xu et al.~\cite{xu2017} also use flow to refine alignment. Their pipeline first locates landmarks and coarsely aligns the face, then performs fine in-sequence alignment using pixel-level movement. The motivation is specific to micro-expressions: a small registration error can be comparable to the facial displacement being recognised.

This use differs from supplying flow to a classifier. An uncorrected displacement field contains residual head movement as well as expression. Magnifying the frames beforehand can increase both. The present pipeline performs no flow-based fine alignment, so corpus preprocessing must not be treated as proof that this source of error has been eliminated.

\subsection{Limitations of the representation}
\label{sec:flow-limitations}

Flow remains sensitive to brightness-constancy violations, alignment errors and estimator settings. Apex-based systems additionally depend on selecting an informative frame. OFF-ApexNet~\cite{liong2019a} attributes weaker performance on SMIC to imprecise apex selection and cites an earlier apex-spotting study reporting an average discrepancy of thirteen frames. That reported average is neither a maximum error nor a measurement made by this thesis.

Using an onset--offset sequence avoids dependence on one selected apex, but still relies on clip boundaries and temporal sampling. Repeated sampled frames can yield zero-flow pairs, and large sampling intervals can make the small-motion approximation less appropriate. These limitations constrain the representation before any network component is compared.

\subsection{Implications for this research}
\label{sec:flow-implications}

\begin{table}[htbp]
    \centering
    \small
    \begin{tabular}{@{}p{0.29\linewidth} p{0.65\linewidth}@{}}
        \hline
        \textbf{Review finding} & \textbf{Consequence for this study} \\
        \hline
        Flow is a recognition input~\cite{liong2019a,see2019}
        & Horizontal and vertical displacement form the first two channels; the classifier receives motion tensors rather than raw frames. \\
        Sequence flow preserves temporal variation~\cite{zhaoS2021}
        & Adjacent-pair fields are retained across the sampled clip, allowing the temporal encoder to operate on an ordered sequence. \\
        Estimator choices differ~\cite{liong2019a,zhaoS2021,xu2017}
        & Farneb\"ack remains fixed; a matched comparison with TV-L1 is an untested extension. \\
        Fine alignment matters~\cite{xu2017}
        & No flow-based fine alignment is performed. It remains a priority for improving preprocessing. \\
        \hline
    \end{tabular}
    \caption[Optical-flow evidence and study decisions]{Optical-flow evidence and the decisions it motivates. Fixed preprocessing choices constrain the interpretation of the architectural ablation.}
    \label{tab:flow-decisions}
\end{table}

The central interpretive point is that optical flow is an \textit{analytical spatio-temporal feature extractor}: it estimates motion from paired frames before the network is trained. A learned spatial stem therefore acts on a representation that already makes displacement explicit. With only 156 clips, learning additional spatial filters may offer limited benefit or duplicate structure already exposed by preprocessing. This is a hypothesis about the stem's marginal value, not proof that learned spatial processing is unnecessary.

A temporal encoder addresses a different task: modelling relations among the successive displacement fields. Pairwise flow estimation does not itself learn those sequence-level relations. Keeping these roles distinct is necessary when attributing performance to the stem and transformer.

\subsection{Research gap and scope}
\label{sec:flow-gap}

The relevant unresolved question is what learned spatial and temporal processing contributes \textit{after} motion has been extracted analytically. The thesis addresses this through matched architectural comparisons under complete leave-one-subject-out (LOSO) evaluation, retaining the same flow construction while switching the learned components. Magnification is varied separately through the paired preprocessing conditions.

The design neither ranks optical-flow estimators nor isolates alignment quality or onset--apex versus sequence representations. Those remain limitations rather than gaps claimed as solved. The resulting architectural conclusions apply to the implemented motion representation and evaluation setting; they do not establish a universal ordering of network families.


\section{Deformation Representation: The Optical Strain Tensor}
\label{sec:strain-review}

Optical strain supplies the third input channel. Its literature establishes a deformation descriptor derived from flow, two ways to use that descriptor, and limitations relevant to interpreting the architectural ablation. Shreve et al.~\cite{shreve2011}, the Liong strain series~\cite{liong2014a,liong2014b,liong2016}, and STSTNet~\cite{liong2019b} provide the principal evidence.

\subsection{What strain contributes}
\label{sec:strain-what-adds}

Flow describes displacement; strain describes its local spatial variation. A spatially uniform translation has zero displacement gradient, whereas neighbouring skin regions moving differently produce deformation. Shreve et al.~\cite{shreve2011} use this distinction to detect facial activity and describe robustness to moderate head translation. It does not imply immunity to arbitrary head movement or registration error.

Their illumination- and make-up-robustness statements are attributed to earlier studies rather than measured anew. They should therefore be read as reported motivation, not as evidence that the present pipeline is robust under those conditions. Strain can suppress a uniform displacement while still inheriting errors from the estimated flow.

\subsection{Tensor components and magnitude}
\label{sec:strain-tensor}

The relevant representation is the symmetric displacement-gradient tensor, with normal components $\varepsilon_{xx}$ and $\varepsilon_{yy}$ and shear components $\varepsilon_{xy}=\varepsilon_{yx}$. Shreve et al.~\cite{shreve2011} approximate spatial derivatives by central differences with offsets of approximately two to three pixels; Liong et al.~\cite{liong2014a} use one-pixel offsets. This discretisation is a modelling choice because it sets the scale at which local variation is measured.

Liong et al.~\cite{liong2014a,liong2014b,liong2016} use the four-component magnitude
\[
    \varepsilon_{\mathrm{F}}
    =\sqrt{\varepsilon_{xx}^{2}+\varepsilon_{yy}^{2}
        +\varepsilon_{xy}^{2}+\varepsilon_{yx}^{2}}
    =\sqrt{\varepsilon_{xx}^{2}+\varepsilon_{yy}^{2}
        +2\varepsilon_{xy}^{2}}.
\]
The scalar describes deformation intensity and discards its sign and direction. Supplying it together with $u$ and $v$ retains the displacement fields from which those spatial relations were calculated.

\subsection{The implemented magnitude and its divergence}
\label{sec:strain-divergence}

The implementation instead computes
\[
    \varepsilon_{\mathrm{os}}
    =\sqrt{\varepsilon_{xx}^{2}+\varepsilon_{yy}^{2}
        +\varepsilon_{xy}^{2}},
    \qquad
    \varepsilon_{xy}=\tfrac{1}{2}
        \left(\frac{\partial u}{\partial y}
             +\frac{\partial v}{\partial x}\right).
\]
It counts shear once rather than twice. NumPy's \texttt{gradient} supplies spatial differences at one-pixel spacing, using central differences in the array interior. The spacing follows Liong et al.~\cite{liong2014a}, but the scalar reduction differs from their convention.

STSTNet~\cite{liong2019b} cannot establish a clear precedent for the three-term form. Its printed magnitude equation repeats $\partial u$ in the mixed derivative where the symmetric tensor calls for $\partial v/\partial x$ alongside $\partial u/\partial y$. It also places the factor $\tfrac{1}{2}$ \emph{outside} the squared derivative sum. Correcting the repeated derivative while retaining that placement gives
\[
    \tfrac{1}{2}
       \left(\frac{\partial u}{\partial y}
            +\frac{\partial v}{\partial x}\right)^2
       =2\varepsilon_{xy}^{2},
\]
which agrees with the four-term convention. Moving the factor inside the square would instead give the implemented three-term form, but that is a further change to the printed expression. The intended implementation cannot be settled from this equation alone; the thesis therefore declares a departure from the unambiguous four-term convention rather than claiming to reproduce STSTNet's magnitude.

The four-term form weights pure shear by $\sqrt{2}$ relative to the implemented form. Independent min--max normalisation removes a uniform scale factor, but does not generally remove this difference when normal and shear contributions vary across pixels and time. Neither convention was compared experimentally here, so their practical equivalence is not assumed.

\subsection{Strain as a feature, a weight, or an input channel}
\label{sec:strain-feature-weight}

\textbf{Optical strain features (OSF)} use a temporally pooled strain map as a descriptor~\cite{liong2014a,liong2014b,liong2016}. Averaging accounts for sequence length: the early method already divides its accumulated map by the number of frame pairs~\cite{liong2014a}. The maps are filtered and normalised before classification. The early work explicitly applies Wiener and Gaussian filters to suppress noise~\cite{liong2014a}.

\textbf{Optical strain weighted features (OSW)} use pooled strain to weight the regional XY-plane histograms of local binary patterns on three orthogonal planes (LBP-TOP), emphasising active facial regions~\cite{liong2014b}. Liong et al.~\cite{liong2016} perform spatial mean pooling followed by temporal mean pooling and concatenate OSF and OSW for a combined descriptor. Strain then serves both as a feature and as a spatial weighting signal.

\textbf{STSTNet} supplies strain directly alongside horizontal and vertical onset--apex flow~\cite{liong2019b}. This is the precedent for the thesis's three-channel input structure, independently of the magnitude ambiguity. The present study retains all three channels across successive frame pairs rather than reducing each clip to one map.

\subsection{Recognition evidence and protocol limits}
\label{sec:strain-evidence}

\begin{table}[htbp]
    \centering
    \small
    \begin{tabular}{@{}lccc@{}}
        \hline
        \textbf{Method} & \multicolumn{2}{c}{\textbf{SMIC accuracy (\%)}} & \textbf{CASME II (\%)} \\
        & Micro-average & Macro-average & Accuracy \\
        \hline
        Baseline & 45.73 & 47.76 & 61.94 \\
        OSF & 41.46 & 46.00 & 51.01 \\
        OSW & 49.39 & 50.71 & 61.94 \\
        OSF + OSW & 52.44 & 58.15 & 63.16 \\
        OFF + OFW & 40.24 & 41.94 & 55.87 \\
        \hline
    \end{tabular}
    \caption[Recognition results for strain-based descriptors]{Recognition results reproduced from Liong et al.~\cite{liong2016}, Table 8. SMIC uses LOSO; CASME II uses leave-one-video-out. The summary combines block settings, so equal entries are not a matched-setting comparison. Macro-average here denotes subject-averaged accuracy, not macro F1.}
    \label{tab:strain-performance}
\end{table}

The CASME II evidence depends on how strain is used and on the comparator. OSF alone falls below the texture baseline. The equal baseline and OSW entries in \autoref{tab:strain-performance} reflect different block settings: under matched settings OSW improves accuracy by 0.81 percentage points. The combined OSF + OSW descriptor improves on the texture baseline by 0.41--1.22 points across block settings~\cite{liong2016}. These accuracy gains are not directly comparable with pooled macro F1, and leave-one-video-out does not enforce subject-disjoint testing.

Crucially, Liong et al.~\cite{liong2016} also test optical-flow features and weights (OFF + OFW) in the same descriptor framework. Their CASME II accuracy is 55.87\,\%, against 63.16\,\% for OSF + OSW, a 7.29-point difference. This is a matched flow-based control, so it would be incorrect to claim that the literature never isolates strain against flow. It concerns hand-crafted descriptors; it does not establish the gain from adding a strain channel to the learned sequence model used here.

\subsection{Limitations}
\label{sec:strain-limitations}

Spatial differentiation can amplify flow-estimation noise. The filtering used by Liong et al.~\cite{liong2014a,liong2016} is therefore relevant to a pipeline without a dedicated strain post-filter. A deterministic transformation of $u$ and $v$ adds no independent observation: its potential value lies in making deformation directly accessible to the classifier. Whether this representation helps must be measured rather than inferred from its construction.

The reviewed recognition methods also simplify time. OSF and OSW pool their sequences; STSTNet uses one onset--apex pair. Shreve et al.~\cite{shreve2011} retain temporal strain activity for spotting, which localises expressions rather than classifying their emotion. The temporal-collapse limitation thus concerns these recognition representations, not every use of strain in the literature.

\subsection{Implications for this research}
\label{sec:strain-implications}

\begin{table}[htbp]
    \centering
    \small
    \begin{tabular}{@{}p{0.29\linewidth} p{0.65\linewidth}@{}}
        \hline
        \textbf{Design choice} & \textbf{Interpretation or unresolved test} \\
        \hline
        Strain accompanies $u$ and $v$
        & Follows STSTNet's channel structure~\cite{liong2019b}; does not establish that its ambiguous magnitude equation was reproduced. \\
        Three-term magnitude
        & Differs from the four-term convention~\cite{liong2014a,liong2014b,liong2016}. Relative shear weighting remains an untested choice. \\
        Per-pair strain sequence
        & Preserves deformation over time instead of pooling it into one descriptor. Each channel is independently min--max normalised across the sequence. \\
        No dedicated strain post-filter
        & Departs from the explicit filtering of Liong et al.~\cite{liong2014a}. Farneb\"ack's internal smoothing is not a separately evaluated strain filter. \\
        Strain held constant
        & Present in all twelve evaluated configurations; its marginal contribution is not measured. \\
        \hline
    \end{tabular}
    \caption[Strain representation and unresolved choices]{The implemented strain representation and the limits it places on interpretation.}
    \label{tab:strain-decisions}
\end{table}

A strain-on/strain-off comparison is the natural next ablation. Neither published gains from a different descriptor nor the performance of the complete system substitutes for that measurement. Comparing magnitude conventions and adding a dedicated post-filter are separate follow-ups.

Strain also extends the analytical-feature argument: the input contains a \textit{hand-specified differential operator}, exposing local deformation without learning that operator from the training clips. A spatial stem may still learn useful combinations of those features, but it is operating after both motion estimation and deformation extraction. This motivates testing its incremental contribution; it does not establish the cause of any observed architectural effect.

\subsection{Research gap and scope}
\label{sec:strain-gap}

The study preserves the temporal evolution of strain for learned sequence processing. Its contribution is an architectural comparison on that representation, not an isolated demonstration of strain's benefit. The independent value of the third channel, the choice of shear weighting and the omitted post-filter remain unresolved. These boundaries preserve the distinction between evidence for a component, a reason to include it, and a result that the experimental design can actually establish.


\section{Temporal Normalisation: The Temporal Interpolation Model}
\label{sec:tim-review}

Flow and strain fix \textit{what} is computed from each frame pair; temporal normalisation fixes \textit{which} frame pairs, and how many. The Temporal Interpolation Model (TIM) is credited in this corpus to Zhou et al. by Li et al.~\cite{liX2013,liX2018} and to Pfister et al. by Lu et al.~\cite{lu2015}, who applied it to micro-expressions outside the review corpus. It is described here only through the reviewed papers that implement and evaluate it.

\subsection{The problem: sequences of unequal length}
\label{sec:tim-problem}

Micro-expressions vary in duration, and the three-class CASME II working set spans roughly a four-fold range of sequence length (\autoref{sec:corpus-threats}h). Three difficulties follow. Any batched network requires a fixed input shape, and the corpus supplies no constant. A descriptor with a temporal radius cannot be applied to a sequence shorter than that radius, so the shortest clips constrain the descriptor for the whole dataset: Li et al.~\cite{liX2018} note that at 25~fps some micro-expressions last four to five frames, restricting LBP-TOP to radius $r = 1$. And duration is itself a nuisance variable --- Lu et al.~\cite{lu2015} normalise to remove ``temporal fluctuations of micro-expressions, which may introduce noise irrelevant to classification''.

\subsection{What TIM is}
\label{sec:tim-what}

Lu et al.~\cite{lu2015} give the corpus's most compact description: ``TIM is a manifold-based interpolation method, which builds a low-dimensional embedding of an image sequence, and then interpolates a curve in the low-dimensional space. The interpolated frames are mapped back to the original high-dimensional space.''

The property that matters downstream is that TIM \textbf{synthesises frames that were never recorded}: Lu et al.~\cite{lu2015} note that Pfister et al. proposed it ``to generate sufficient frames from an image sequence'', that is, as an upsampler rather than only a length-standardiser.

\subsection{TIM in practice: the convergence on ten frames}
\label{sec:tim-ten-frames}

TIM appears throughout the corpus as an unremarked preprocessing step, almost always at ten frames: SMIC's baseline~\cite{liX2013} normalises to ten before LBP-TOP; Li et al.~\cite{liX2018} apply TIM10 after magnification; and Lu et al.~\cite{lu2015} apply TIM before their Delaunay-based coding model, following Pfister et al.'s protocol so as to compare against TIM10 baselines. CAS(ME)$^2$'s baseline~\cite{qu2016} is the exception, normalising to \textbf{120 frames by linear interpolation} and evaluating leave-one-video-out.

One study tests the setting on its own terms. Li et al.~\cite{liX2018} sweep interpolation length from 10 to 80 frames on all three SMIC variants under leave-one-subject-out, holding the descriptor fixed and skipping magnification so only length varies, and report that ``the best performance for all three datasets is achieved with interpolation to 10 frames'', because ``in longer sequences, changes along the time-dimension are diluted''.

\subsection{How many frames? The remaining systematic comparisons}
\label{sec:tim-how-many}

Ben et al.~\cite{ben2021} run a second controlled experiment and reach the opposite conclusion. Holding alignment, descriptor and classifier fixed --- JCFDA alignment, LBP-TOP, and an SVM with an RBF kernel --- they interpolate sequences to 10, 20, 30 and so on up to 110 frames by two algorithms, \textbf{Newton interpolation} and \textbf{TIM}, recording the best recognition rate at each length on MMEW.

Newton peaks at \textbf{33.3\,\%} at 60 frames; TIM peaks at \textbf{38.9\,\%} ``when the micro-expression sequence is respectively interpolated to 30 or 60 frames'', and Ben et al. adopt TIM at 60 accordingly. Accuracy is non-monotonic in length: rates ``increase initially and then decrease towards the end''.

Two caveats limit the transfer. The experiment uses MMEW with LBP-TOP, not CASME II and not a learned model; and it runs largely in the \textit{upsampling} regime, since Ben et al.~\cite{ben2021} interpolate ``to the maximal value of 110 frames'', so the longer settings add frames for most clips rather than discarding them. This thesis operates in the opposite regime (\autoref{sec:tim-implications}), and an optimum measured while synthesising frames need not hold when discarding them.

A third comparison is the only multi-corpus one. Xu et al.~\cite{xu2017} interpolate every sequence to 10, 15 and 20 frames for both their facial dynamics map and LBP-TOP, ``on all four datasets'' --- CASME, CASME II and SMIC's high-speed and visual subsets. They also depart from the convention in a way the others do not, using \textbf{linear} interpolation ``instead of Temporal Interpolation Model because in a high frame-rate scenario, linear interpolation suffices to characterize the motion pattern with less error''.

\subsection{Alternatives to uniform interpolation}
\label{sec:tim-keyframe}

Zhao et al.~\cite{zhaoS2021} do not interpolate uniformly. They build an eleven-frame key-frame sequence from the annotated onset, apex and offset by generating eight transition frames allocated between the onset--apex and apex--offset intervals \textbf{in proportion to each sub-interval's duration}, so sampling density follows the expression's own phase structure and the apex is emphasised. Optical flow between adjacent key-frames then gives a ten-frame flow sequence. OFF-ApexNet~\cite{liong2019a} and STSTNet~\cite{liong2019b} take the limit, keeping the onset--apex pair alone (\autoref{sec:flow-what-between}).

The corpus therefore contains temporal representations of length 2, 10, 11, 15, 20, 30, 60, 110 and 120 frames, chosen variously by convention, by sweep and by architecture. There is no consensus, and only Xu et al.'s~\cite{xu2017} sweep spans more than one corpus.

\subsection{Limitations}
\label{sec:tim-limitations}

\textbf{Ten frames is a convention, not a finding.} The most-used setting traces to one baseline paper's choice, and the three studies that test it do not agree: ten frames for Li et al.~\cite{liX2018}, 30 to 60 for Ben et al.~\cite{ben2021}, and a 10--20 frame range for Xu et al.~\cite{xu2017} (\autoref{sec:tim-how-many}). Interpolation length interacts with the temporal receptive field of whatever consumes it, and in all three that consumer is a hand-crafted descriptor.

\textbf{Frame synthesis and motion analysis are in tension.} TIM generates frames that were not recorded. Where the downstream representation is appearance-based this is benign; where it is \textit{motion}, synthesised frames yield synthesised displacements. One reviewed paper recognises this: Xu et al.~\cite{xu2017} choose linear interpolation over TIM precisely on the grounds that it ``suffices to characterize the motion pattern with less error'' at a high frame rate. But they assert the point rather than measure it --- TIM appears once in their paper, in that sentence --- and no reviewed study compares the two interpolators against a motion representation. The rest of the corpus does not raise the question, because those using TIM extract appearance descriptors, while those extracting flow either use two frames or, as Zhao et al.~\cite{zhaoS2021} do, generate transition frames without discussing the consequence.

\textbf{Downsampling forfeits part of the corpus's advantage.} \autoref{sec:corpus} establishes why CASME II was recorded at 200~fps rather than its predecessor's 60~fps; reducing a clip to a few tens of frames gives part of that back. Only Xu et al.~\cite{xu2017} tie the interpolation choice to the acquisition rate at all, and then only to justify the simpler interpolator. The rest apply TIM10 to 100~fps SMIC data as readily as to 200~fps CASME II data, without weighing what the higher rate bought.

\subsection{Implications for this research}
\label{sec:tim-implications}

\begin{table}[htbp]
	\centering
	\footnotesize
	\begin{tabular}{@{}p{5.6cm} p{8.0cm}@{}}
		\hline
		\textbf{Finding from the review} & \textbf{Decision taken in this thesis} \\
		\hline
		Fixed-length input is required, and unequal duration is nuisance variance~\cite{lu2015} & Every clip is normalised to a fixed temporal length before training \\
		TIM beats Newton interpolation by 5.6 points, with a joint optimum at 30 or 60 frames~\cite{ben2021} & \textbf{L = 33 frames sampled, giving T = 32 flow pairs} --- the lower of the two reported optima rather than the field-conventional 10 \\
		Sequence length has an interior optimum found only by sweep~\cite{ben2021} & T = 32 is fixed throughout and is \textbf{not} swept; a sweep is declared as further work \\
		Onset and offset are annotated in CASME II (\autoref{sec:casme2-construction}) & Sampling is bounded by the annotated onset and offset, so the window spans the expression rather than surrounding neutral frames. This holds on the image-folder path used for CASME II; the loader's \texttt{.avi} branch samples across a whole file and is not exercised here \\
		TIM synthesises frames, and flow between synthesised frames is synthesised motion (\autoref{sec:tim-limitations}) & \textbf{No frames are synthesised} --- see the declaration below \\
		\hline
	\end{tabular}
	\caption[Review findings and the temporal-normalisation decisions they determine]{Review findings and the temporal-normalisation decisions they determine.}
	\label{tab:tim-decisions}
\end{table}

\textbf{A declaration about what this pipeline implements.} The module performing this step is named for the Temporal Interpolation Model, but what it computes is \textbf{uniform index sampling} of existing frames (\autoref{sec:three-channel-tensor}). No embedding is built, no curve is interpolated, no frame is synthesised. That is a materially different operation from the TIM of \autoref{sec:tim-what}, and it is named accurately here rather than in the pipeline's own terms. It is also defensible: because the downstream representation is flow between adjacent sampled frames, interpolation would manufacture displacement fields for motion that never occurred, whereas subsampling discards temporal resolution without fabricating it. What is given up is TIM's ability to \textit{lengthen} short clips.

Three quantitative consequences follow.

\textit{Effective temporal resolution falls by a clip-dependent factor.} Sampling 33 frames from a clip of $N$ frames recorded at 200~fps gives an effective rate of $6600/N$ fps: \textbf{100 fps} at the median ($N = 66$), which is SMIC's recording rate and half what CASME II was built to provide, and roughly \textbf{52 fps} for the longest ($N = 126$), below the original CASME database's 60~fps~\cite{yan2014}. Eighty-one of the 156 clips are downsampled by a factor of two or more.

\textit{Magnification is applied after subsampling, which shifts the realised pass-band.} Magnification runs on the 33 sampled frames at the nominal 200~fps (\autoref{sec:evm-departures}), but those frames span the clip's full duration, so the true inter-frame interval is $N/(200 \times 32)$ seconds rather than $1/200$. At the median a band specified as 5--25 Hz therefore corresponds to roughly 2.5--12.5 Hz physically, and the discrepancy grows with clip length. The realised band is thus neither the 5--25 Hz \autoref{sec:evm-implications} derives from the duration rule of Bai et al.~\cite{bai2021} nor constant across clips. This is a declared divergence; magnifying before subsampling is declared as further work.

\textit{Three clips are shorter than the sampling window.} \textbf{Three of the 156 clips} contain 31 frames against a target of 33, so sampling repeats indices and yields near-zero flow at those positions.

\subsection{Research gap and scope}
\label{sec:tim-gap}

The unresolved question is what temporal normalisation costs a representation made of motion rather than appearance. Three studies sweep interpolation length and give three answers, none on CASME II and none with a learned temporal model; and although Xu et al.~\cite{xu2017} raise both the synthesis-versus-motion tension and the acquisition-rate trade-off, they settle neither by measurement (\autoref{sec:tim-limitations}). This thesis takes the position that follows from its own representation --- sampling only recorded frames --- and quantifies what that costs in effective frame rate (\autoref{sec:tim-implications}), which no reviewed paper reports.

The design does not sweep T, nor compare interpolators, nor separate the effect of sequence length from the effect of the sampling regime. Those remain limitations rather than gaps claimed as solved, and the resulting conclusions apply to the implemented sampling scheme at T = 32.


\section{The Learned Spatial Backbone: Shallow 3D Convolutional Networks}
\label{sec:cnn-review}

This section covers the convolutional backbone, switched on and off as Variable C of this thesis's ablation, through four review-corpus papers.

\subsection{The case for, and against, a learned spatial stem}
\label{sec:cnn-case}

By 2019 the field had converged on optical flow as the input representation (\autoref{sec:why-flow}); the open question was what should consume it. A learned backbone promises what hand-crafted descriptors cannot give~\cite{liong2019b}, but runs into the small-data constraint (\autoref{sec:capacity}, \autoref{sec:corpus-threats}). Xia et al.~\cite{xia2020b} describe a related failure mode from a different cause --- ``the important subtle dynamics are prone to disappearing in the domain shift \ldots especially for deep models'' --- where the mechanism is the domain shift induced by mixing corpora rather than sample count alone, and where they report deeper, higher-resolution models doing better on the individual-database task. Applying their shrinking argument to a single small corpus therefore extends it beyond the setting they establish.

\subsection{What ``3D CNN'' means in this literature}
\label{sec:cnn-meaning}

The term covers three materially different architectures.

\textbf{A convolution over spatial and modality axes, with no temporal axis.} STSTNet~\cite{liong2019b}, the canonical shallow model in this field, calls itself a ``Shallow Triple Stream Three-dimensional CNN''. Its input is a single \textbf{28 $\times$ 28 $\times$ 3} tensor from the onset and apex frames, whose third dimension holds horizontal flow, vertical flow and optical strain --- not time. Three parallel streams, each a single convolutional layer with 3, 5 and 8 kernels respectively, supplement ``the small scale input data \ldots to avoid the problem of underfitting''. \textbf{There is no time dimension in this network}; the clip was collapsed to a single frame pair upstream (\autoref{sec:flow-what-between}).

\textbf{A convolution genuinely extended across time.} STRCN~\cite{xia2020a} extends convolutional connectivity ``across temporal domain'' with recurrent convolutional layers, modelling appearance and geometry separately.

\textbf{A convolution over a short flow sequence.} Zhao et al.~\cite{zhaoS2021} apply a C3D-style network to the ten-frame flow sequence built from their eleven key-frames (\autoref{sec:tim-keyframe}), so the temporal axis is present but short.

\subsection{The central finding: shrink the model \textit{and} the input}
\label{sec:cnn-shrink}

Xia et al.~\cite{xia2020b} treat depth and input resolution as controlled variables. They build four models of increasing depth --- from a single convolutional layer (Model 1) to one convolutional plus three recurrent convolutional layers (Model 4) --- and evaluate each at nine input resolutions, \textbf{20, 40, 60, 80, 100, 150, 200, 250 and 300 pixels}, on optical flow with UAR as the metric. Two results follow.

\textbf{Deeper is worse.} Models 3 and 4 ``have worse recognition performances compared to shallower \ldots models'', even though all four are far shallower than ResNet-101 or DenseNet-169.

\textbf{Higher resolution is worse, with a threshold.} Performance ``will be degraded with larger input resolutions (Model 2, Model 3 and Model 4)'' --- the two-layer model they go on to adopt included --- and for the deeper models ``begins to decrease dramatically when the resolution is bigger than \textbf{100 $\times$ 100}''. Shallow models are nonetheless ``basically robust to the change of input resolution'', improving toward lower resolution above a floor of about 40 $\times$ 40.

STSTNet takes the same argument to its limit and remains competitive on the CASME II subset (\autoref{sec:megc-metrics}): Liong et al.~\cite{liong2019b} report \textbf{0.00167 M parameters across 2 layers at 28 $\times$ 28 $\times$ 3 input}, against OFF-ApexNet's 2.77 M across 5 layers and GoogLeNet's 7 M across 22.

\subsection{Coping with limited and imbalanced training data}
\label{sec:cnn-limited-data}

Two strategies in the corpus have counterparts in this thesis's training configuration. Xia et al.~\cite{xia2020a} pair temporal augmentation with a balanced loss (\autoref{sec:imbalance-balanced-loss}, \autoref{sec:imbalance-double-correction}); a two-stage few-shot scheme~\cite{zhaoS2021} is the one route not taken. Xia et al.~\cite{xia2020b} take a third --- adding capability without capacity --- by developing ``three parameter-free modules \ldots integrate with RCN without increasing any learnable parameters'', so representational power grows while the parameter count, and therefore the overfitting risk, does not.

\subsection{Limitations}
\label{sec:cnn-limitations}

\textbf{The terminology obscures what is being compared.} Papers titled as three-dimensional CNNs include networks with no temporal axis at all, so ``3D-CNN'' results across the corpus are not like with like.

\textbf{Input resolution is almost never controlled.} Xia et al.~\cite{xia2020b} are the exception; elsewhere resolution is an implementation detail --- 28 $\times$ 28 in STSTNet and OFF-ApexNet --- with no indication of sensitivity, despite the substantial sensitivity Xia et al. measure.

\textbf{The evidence base is composite-database.} Xia et al.'s shrinking study~\cite{xia2020b} and STSTNet~\cite{liong2019b} are developed for the MEGC composite setting, which supplies substantially more training data per fold than a single-corpus study (\autoref{sec:megc-metrics}); only the earlier recurrent-convolutional work~\cite{xia2020a} is evaluated per database. Conclusions about affordable capacity therefore understate the constraint here.

\textbf{No study isolates the backbone itself.} The corpus is not short of matched comparisons: it sets descriptor-only methods against learned ones on the same folds, and Xia et al.~\cite{xia2020b} vary individual architectural modules against a common baseline with every other parameter held constant. But those comparisons all \textit{add} to a backbone that stays in place. None removes the learned spatial stem entirely and passes the motion representation straight to a classifier or temporal model, so what the stem itself contributes, set against what it costs, is nowhere measured.

\subsection{Implications for this research}
\label{sec:cnn-implications}

The component ablated as Variable C is a \textbf{three-stream shallow convolutional backbone}: one unshared two-layer stream per input channel --- flow-$u$, flow-$v$, optical strain.

\begin{table}[htbp]
	\centering
	\footnotesize
	\begin{tabular}{@{}p{5.6cm} p{8.0cm}@{}}
		\hline
		\textbf{Finding from the review} & \textbf{Decision taken in this thesis} \\
		\hline
		Shallow beats deep on small micro-expression data~\cite{xia2020b}; STSTNet uses 2 layers and $\sim$1,670 parameters~\cite{liong2019b} & Two convolutional layers per stream, \textbf{14,544 parameters} in the backbone (15,027 including the classifier head) --- shallow by the corpus's standards, though roughly nine times STSTNet's \\
		Triple-stream design supplements small-scale input~\cite{liong2019b} & Three unshared streams, one per input modality. STSTNet's streams differ in \textit{filter count} over the same input; these differ in \textit{input channel} \\
		Parameter-free modules add capability without capacity~\cite{xia2020b} & Motivates the parameter-free attention evaluated as Variable B (\autoref{sec:simam-review}) \\
		Limited and imbalanced data require augmentation plus loss correction~\cite{xia2020a} & Balanced sampling with focal loss; loss class-weighting deliberately disabled after single-class collapse (\autoref{sec:corpus-implications}) \\
		\textbf{Performance degrades above 100 $\times$ 100 input, dramatically for deeper models~\cite{xia2020b}} & \textbf{Not followed.} The backbone is fed at \textbf{224 $\times$ 224}, more than twice that threshold, though the same resolution as Bai et al.~\cite{bai2021} and Li, Huang and Zhao~\cite{liY2021} use --- see below \\
		\hline
	\end{tabular}
	\caption[Review findings and the backbone decisions they determine]{Review findings and the backbone decisions they determine.}
	\label{tab:cnn-decisions}
\end{table}

\textbf{A declaration about what this component is.} Every kernel in this backbone is spatial only: no filter spans two time steps, so despite the \texttt{Conv3d} implementation it is a \textbf{purely spatial} feature extractor (\autoref{sec:conv3d}). The ablation therefore does not test spatio-temporal convolution; it tests a \textbf{learned spatial stem}, and whatever it measures for Variable C is evidence about \textit{this} stem --- spatial-only, at 224 $\times$ 224, on flow and strain input --- and \textbf{not} evidence that 3D convolution over time is unhelpful for micro-expression recognition, which this study never evaluated.

Two observations mitigate that, and one aggravates it. In the field's own terms the choice is defensible: STSTNet, one of the two strongest shallow baselines on the CASME II subset, likewise performs no temporal convolution (\autoref{sec:cnn-meaning}). The stem also competes against a representation that already performs spatio-temporal extraction analytically (\autoref{sec:flow-implications}, \autoref{sec:strain-implications}).

Against both, \textbf{input resolution is the clearest departure from Xia et al.'s guidance}. It is not a departure from the corpus as a whole: Bai et al.~\cite{bai2021} resize frames to 224 $\times$ 224, and without alignment or cropping, as here; Li, Huang and Zhao~\cite{liY2021} resize faces to the same; and Xia et al.'s own earlier work~\cite{xia2020a} runs one stream at 300 $\times$ 245. What 224 $\times$ 224 departs from is the specific threshold of \autoref{sec:cnn-shrink}, which is a composite-database result and which bites hardest on models deeper than this one. Since this stem is shallow, the reviewed evidence predicts a smaller resolution effect than the deep-model figures suggest, so re-running the backbone arm at a lower resolution is worth doing because it is cheap and would settle the question, not because the corpus implies a large gain. The choice is also the expensive one: a 224 $\times$ 224 map carries sixty-four times the spatial elements of STSTNet's 28 $\times$ 28, and the backbone convolves it at every one of the 32 time steps. Whatever this study measures for the stem is therefore a resolution result as much as an architecture result, and the re-run is the first item of declared further work.

\subsection{Limitations of the existing work, and how this study differs}
\label{sec:cnn-gap}

The four gaps of \autoref{sec:cnn-limitations} define what the reviewed literature does not establish.

This thesis differs on the third, which is what its research question requires: it removes the convolutional backbone entirely in four of its twelve configurations, and measures the difference under complete leave-one-subject-out with every other factor held constant. On the first gap it answers the field's terminological ambiguity by declaring precisely what its own component is (\autoref{sec:cnn-implications}); on the second it knowingly departs from the literature's guidance, at the 224 $\times$ 224 resolution flagged above.


\section{Parameter-Free Attention: SimAM}
\label{sec:simam-review}

SimAM re-weights the convolutional feature maps without adding learned parameters. Yang et al.~\cite{yang2021} provide the mechanism and its image-recognition evidence; Xia et al.~\cite{xia2020b} motivate parameter-free modules in composite-database micro-expression recognition. The latter support the design principle, rather than validating SimAM itself on this task.

\subsection{Why parameter-free attention}
\label{sec:simam-why}

Local facial activity occupies a small part of a motion tensor. Attention offers a way to emphasise informative activations, but its parameter and computational costs must be justified on a 156-clip corpus. Yang et al.'s~\cite{yang2021} comparison illustrates the range of parameter costs.

\begin{table}[htbp]
    \centering
    \small
    \begin{tabular}{@{}lc@{}}
        \hline
        \textbf{Module} & \textbf{Additional learnable parameters} \\
        \hline
        SE & $2C^2/r$ \\
        CBAM & $2C^2/r+2k^2$ \\
        GC & $2C^2/r+C$ \\
        ECA & $k$ \\
        SRM & $6C$ \\
        SimAM & $0$ \\
        \hline
    \end{tabular}
    \caption[Parameter costs of attention modules]{Parameter costs reported by Yang et al.~\cite{yang2021}, Table 1. $C$ denotes channels, $r$ the reduction ratio and $k$ the relevant kernel size. The competing modules are attributed through this source.}
    \label{tab:attention-modules}
\end{table}

Zero additional parameters makes SimAM a plausible candidate, not the only defensible one: alternatives such as ECA are also small. Xia et al.~\cite{xia2020b} demonstrate a related strategy of adding parameter-free modules to a recurrent convolutional model. Neither argument establishes that a parameter-free module has zero runtime cost or cannot worsen generalisation.

\subsection{The distinctiveness premise}
\label{sec:simam-premise}

Yang et al.~\cite{yang2021} motivate attention through spatial suppression: neurons distinguishable from their surroundings are assigned greater importance. They express this as the linear separability of one target activation from the other activations in its channel. The resulting weights scale the feature map rather than being added to it.

This is a prior about activation statistics, not a detector of facial action units or minority classes. A distinctive response may encode useful deformation, noise or an artefact. Its suitability for flow-derived features therefore requires evaluation rather than following automatically from the neuroscience motivation.

\subsection{Energy and implemented refinement}
\label{sec:simam-energy}

The paper minimises a regularised separation energy and obtains a closed-form importance score. Its reference implementation can be written compactly as
\[
    d_i=(x_i-\mu)^2,\qquad
    s^2=\frac{1}{M-1}\sum_{i=1}^{M}d_i,\qquad
    E_i^{-1}=\frac{d_i}{4(s^2+\lambda)}+\frac{1}{2},
\]
\[
    \widetilde{x}_i=x_i\,\operatorname{sigmoid}(E_i^{-1}),
    \qquad \mu=\frac{1}{M}\sum_{i=1}^{M}x_i.
\]
Here $M$ is the number of positions in one channel's statistical context and $\lambda$ is a positive regularisation constant. The $+\tfrac12$ follows from the reciprocal-energy expression rather than an arbitrary offset.

An implementation distinction matters: the shared variance printed with Yang et al.'s equation (5) divides by $M$, whereas their Figure 3 code divides by $M-1$~\cite{yang2021}. The equations above follow the reference \emph{code}, as does this thesis. The sigmoid preserves the ordering of the scores while bounding the multiplicative weights; it does not create a hard mask of inactive regions.

\subsection{What three-dimensional weights mean}
\label{sec:simam-3d-weights}

The original module produces a separate weight at each channel--height--width position~\cite{yang2021}. Its three-dimensional weights therefore describe $C\times H\times W$, not a temporal volume. Statistics are computed separately within each channel's spatial plane, with $M=HW$. Extending the context across video frames is an additional design decision, not part of the original image formulation.

\subsection{Evidence and hyper-parameter selection}
\label{sec:simam-evidence}

Yang et al.~\cite{yang2021} evaluate SimAM with several network families on CIFAR and ImageNet, reporting competitive classification performance against learned attention modules. Their timing results also show that unchanged parameter totals do not imply unchanged inference speed. The evidence supports testing SimAM, but does not establish its value on micro-expression motion tensors.

For CIFAR, they sweep $\lambda$ from $10^{-1}$ to $10^{-6}$ and select $10^{-4}$ as a balance between accuracy and variability. They repeat selection for ImageNet and choose $0.1$~\cite{yang2021}. Thus $10^{-4}$ is a dataset-selected operating point, not a universal setting justified by the formula.

\subsection{Limits of transfer to this task}
\label{sec:simam-limitations}

The source classification experiments concern natural images, different training-set sizes and learned feature distributions. In this thesis the input is flow and strain, so distinctiveness need not have the same relationship to useful evidence. The statistics are calculated from each feature map, not estimated by pooling the entire training corpus; the small dataset instead constrains how the surrounding network is learned.

The image experiments also do not decide whether video statistics should be computed per frame or over a clip. Nor does the lack of learned attention weights establish that runtime, memory or overfitting concerns disappear. These are empirical limits on transferring the source result, rather than reasons to assume either success or failure.

\subsection{Implications for this research}
\label{sec:simam-implications}

\begin{table}[htbp]
    \centering
    \small
    \begin{tabular}{@{}p{0.29\linewidth} p{0.65\linewidth}@{}}
        \hline
        \textbf{Choice} & \textbf{Implementation and interpretation} \\
        \hline
        Placement & Applied separately to each convolutional stream before concatenation; no trainable weights or biases are introduced. \\
        Refinement & Uses the reference-code reciprocal energy and sigmoid gate~\cite{yang2021}. \\
        Statistical context & Mean and variance span $(T,H,W)$ within each channel, with divisor $THW-1$, rather than one $(H,W)$ plane. \\
        Regularisation & $\lambda=10^{-4}$ is retained without re-tuning on CASME II. \\
        Architectural dependency & SimAM is enabled only with the CNN stem in this model. Four of sixteen candidate configurations are excluded, leaving twelve evaluated configurations and four SimAM matched pairs. \\
        \hline
    \end{tabular}
    \caption[SimAM implementation choices]{SimAM implementation choices and their relation to the reviewed image formulation. The dependency on the CNN is a rule of this architecture.}
    \label{tab:simam-decisions}
\end{table}

The five-dimensional input includes batch, channel, time, height and width. Pooling statistics across time makes an activation distinctive relative to its whole clip rather than its own frame. This could emphasise a brief contraction, but does not guarantee low weights at onset or high weights at the apex: the score depends on squared deviation, not expression labels. The two context choices are not compared experimentally.

The CNN dependency similarly describes this implementation, not a mathematical prohibition on applying SimAM to other tensors. Its measured effect is conditional on the stem being present. Retention should be judged against both recognition performance and computational cost; parameter count alone cannot settle that decision.

\subsection{Research gap and scope}
\label{sec:simam-gap}

The study evaluates SimAM in a motion-input, small-corpus setting through four matched comparisons under complete leave-one-subject-out evaluation. It tests the implemented module as a whole, rather than separately establishing the benefit of each design choice.

Three follow-ups remain: testing whether the distinctiveness prior suits flow-derived features, re-tuning $\lambda$, and comparing per-frame with whole-clip statistics. Each requires new evaluations despite introducing no learned attention weights. Per-class changes may also be examined, but a rare class benefiting most is a hypothesis of this study, not a prediction established by Yang et al.~\cite{yang2021}.


\section{The Temporal Encoder: Transformers for Micro-Expression Recognition}
\label{sec:transformer-review}

The temporal encoder is motivated by Zhang et al.'s Short and Long Range Relation Based Spatio-Temporal Transformer (SLSTT)~\cite{zhang2022} and its Vision Transformer (ViT) foundation~\cite{dosovitskiy2021}. The component evaluated here differs substantially from published SLSTT, particularly in the axis over which self-attention operates.

\subsection{Why temporal self-attention}
\label{sec:transformer-why}

A micro-expression evolves through onset, peak and relaxation. Retaining a motion sequence permits a model to relate separated stages rather than compressing them immediately into one descriptor. Zhang et al.~\cite{zhang2022} motivate their design through the need to capture both local and global patterns.

This thesis compares a temporal transformer with mean pooling over the same input sequence. Mean pooling discards order, whereas positional information and self-attention allow the encoder to model relations among frame-pair features. This motivates the comparison without establishing in advance which representation will generalise better.

\subsection{The Vision Transformer foundation}
\label{sec:transformer-vit}

ViT splits an image into patches, projects them into token vectors, adds learned positional embeddings and processes them with transformer encoder layers~\cite{dosovitskiy2021}. Attention then relates spatial patches. Dosovitskiy et al. also show that data scale and training regime matter: without the locality and translation-equivariance biases of convolution, image transformers can perform poorly when trained with insufficient data, while large-scale pre-training improves their performance.

That evidence is a caution for a 156-clip study, not a universal requirement that every transformer must be pre-trained. A short sequence of motion features presents a different learning problem from image patches, but remains subject to limited data and overfitting.

\subsection{Published SLSTT}
\label{sec:slstt-architecture}

Zhang et al.~\cite{zhang2022} combine three stages. First, \textbf{onset-referenced optical flow} describes sampled frames relative to the onset rather than to their immediate predecessors. They argue that this gives a more structured trajectory towards and away from the apex. Their experiments use this representation throughout; they do not provide a matched short-term-flow ablation establishing its independent gain.

Second, a \textbf{spatial ViT encoder} processes each flow field as patches. The implementation uses ImageNet-pre-trained ViT-B/16, twelve encoder layers and $384\times384$ input. Third, an \textbf{LSTM temporal aggregator} combines the per-frame features. A mean aggregator supplies the comparison variant. Self-attention is spatial in this architecture; recurrence handles time.

\subsection{Published results and the aggregation comparison}
\label{sec:slstt-results}

\begin{table}[htbp]
    \centering
    \small
    \begin{tabular}{@{}lccc@{}}
        \hline
        \textbf{Method} & \textbf{Composite UF1} & \textbf{CASME II UF1} & \textbf{CASME II UAR} \\
        \hline
        LBP-TOP & 0.588 & 0.703 & 0.743 \\
        Bi-WOOF & 0.630 & 0.781 & 0.803 \\
        OFF-ApexNet~\cite{liong2019a} & 0.720 & 0.876 & 0.868 \\
        STSTNet~\cite{liong2019b} & 0.735 & 0.838 & 0.869 \\
        EMR~\cite{liu2019} & 0.789 & 0.829 & 0.821 \\
        RCN~\cite{xia2020b} & 0.705 & 0.809 & 0.856 \\
        SLSTT-Mean & 0.788 & 0.844 & 0.830 \\
        SLSTT-LSTM & 0.816 & 0.901 & 0.885 \\
        \hline
    \end{tabular}
    \caption[SLSTT and comparison results]{Selected results from Zhang et al.~\cite{zhang2022}, Table 2, reported under their three-class Composite Database Evaluation heading. Columns retain the source's dataset labels. The listed methods are not all arms of one controlled architectural comparison.}
    \label{tab:slstt-results}
\end{table}

The SLSTT variants compare temporal aggregation with the spatial transformer retained in both. CASME II UF1 increases from 0.844 to 0.901, a difference of 0.057, or 5.7 UF1 percentage points~\cite{zhang2022}. This supports evaluating aggregation choices within that framework. It is not an estimate of what an LSTM would add to this thesis, where mean pooling follows \emph{temporal} self-attention rather than independent spatial encoding.

Protocol labels need care. The source describes three-class LOSO evaluations on the databases separately and together under its Composite Database Evaluation heading. That heading alone therefore does not establish that every CASME II column entry uses pooled cross-database training. It also reports a distinct Sole Database Evaluation using original labels, including five CASME II classes~\cite{zhang2022}. Neither setting should be equated automatically with this thesis's 156-clip three-class pool.

\subsection{Limitations of the evidence}
\label{sec:transformer-limitations}

Published SLSTT retains a pre-trained spatial encoder. Its results do not isolate the benefit of a transformer trained from scratch on this study's motion representation. The Mean--LSTM comparison isolates an aggregator, not the presence of a transformer, since both variants contain one.

Label interpretation also limits comparison. Zhang et al.~\cite{zhang2022} note that facial action units do not map uniquely to emotion labels, so visual evidence alone need not determine the annotation. This is a limitation of the recognition task, not a numerical upper bound measured by the present experiments.

\subsection{Implementation and six declared divergences}
\label{sec:transformer-implications}

The evaluated component receives 32 tokens of width 96. They come either from the convolutional stem or from $4\times4$ spatial pooling followed by a \emph{learned linear projection}. The encoder uses four pre-norm layers, eight heads, feed-forward width 256 and dropout 0.1. Fixed sinusoidal positions precede attention; mean pooling follows it. With the encoder disabled, the input tokens are averaged directly.

\begin{table}[htbp]
    \centering
    \small
    \begin{tabular}{@{}p{0.23\linewidth} p{0.34\linewidth} p{0.35\linewidth}@{}}
        \hline
        \textbf{Property} & \textbf{Published SLSTT~\cite{zhang2022}} & \textbf{This thesis} \\
        \hline
        Attention axis & Spatial patches within a frame & Temporal tokens within a clip \\
        Aggregation & LSTM over spatially encoded frames & Mean after temporal self-attention \\
        Initialisation & ImageNet-pre-trained ViT-B/16 & Trained from scratch \\
        Capacity & Twelve layers, width 768; $384\times384$ inputs & Four layers, width 96; 32 tokens \\
        Positions & Learned spatial embeddings & Fixed sinusoidal temporal encoding \\
        Input flow & Relative to onset & Relative to the preceding sampled frame \\
        \hline
    \end{tabular}
    \caption[Divergences from published SLSTT]{Six material divergences from published SLSTT. The component evaluated here is a small temporal transformer encoder, not a reproduction of the published architecture.}
    \label{tab:slstt-divergences}
\end{table}

Six matched pairs evaluate the temporal encoder with the other component choices fixed within each pair. They test the whole encoder, including positional encoding, feed-forward layers and normalisation, against direct mean pooling. They do not isolate self-attention from those accompanying operations or compare it with another learned temporal model.

Three follow-ups follow from the divergences: testing an LSTM aggregator, comparing onset-referenced with consecutive-frame flow, and replacing fixed positions with learned embeddings. The last would add $32\times96=3{,}072$ embedding parameters for this sequence length. All three remain untested; changing the input representation also requires re-training and evaluation before any recognition benefit can be claimed.

\subsection{Research gap and scope}
\label{sec:transformer-gap}

The study asks what a small temporal encoder contributes after motion extraction, using a single corpus without pre-training. This differs from SLSTT's spatial transformer and recurrent aggregation while retaining its motivation to model relations beyond isolated frames.

The contribution is a matched comparison of the implemented encoder and an order-blind baseline within this pipeline. It neither validates published SLSTT nor establishes a general ranking of attention, recurrence and convolution. The six declared divergences and the limited training regime define the scope of the architectural conclusions.


\section{Training Under Severe Class Imbalance}
\label{sec:imbalance-review}

This section covers the training configuration held constant across all twelve ablated configurations, and which \autoref{sec:corpus-implications} identified as forced by the corpus rather than chosen freely: the handling of a 4 : 1.3 : 1 class imbalance.

\textbf{A deliberate exclusion.} Two topics the project's reading list groups here --- Grad-CAM++ visual auditing and adversarial identity disentanglement --- are \textbf{not reviewed}, because neither is part of the implemented system. The loss module records that ``the Cross-Batch Memory (XBM) queue, and the adversarial identity loss --- are intentionally removed'', and no gradient-based visualisation exists anywhere in the codebase. Both are declared as proposed further work.

\subsection{The imbalance is a property of the phenomenon, not of the sampling}
\label{sec:imbalance-property}

\autoref{sec:label-taxonomy} and \autoref{sec:original-baseline} established the class counts and the always-Negative floor; what matters here is why the skew exists. It is not an artefact of careless collection. Yan et al.~\cite{yan2014} are explicit that it follows from what can be elicited in a laboratory: ``some types of facial expressions are difficult to elicit in laboratory situations, thus the samples in different categories distributed unequally, e.g., there are 60 disgust samples but only 7 sadness samples.'' Nor is it specific to CASME II --- the MEGC composite shows the same pattern (\autoref{sec:megc-metrics}). Skew must be coped with; it cannot be collected away.

\subsection{Three places to intervene}
\label{sec:imbalance-three-places}

The reviewed corpus corrects for imbalance at three distinct points, worth separating because this thesis intervenes at all three and the interaction between two of them turned out to matter.

\begin{enumerate}
	\item \textbf{At the data}, by changing how often each class is presented during training --- resampling or augmentation.
	\item \textbf{At the loss}, by changing how much each example contributes to the gradient --- class weighting or focal re-weighting.
	\item \textbf{At the metric}, by changing how performance is scored --- macro-averaged rather than sample-averaged measures.
\end{enumerate}

\subsection{Loss-level correction: focal loss}
\label{sec:imbalance-focal}

Zhao et al.~\cite{zhaoS2021} bring focal loss into micro-expression recognition ``to alleviate the inefficient model training caused by the class imbalance in the MEs Datasets'', since cross-entropy ``could be biased toward particular emotions that constitute a larger portion of the training set''. They adopt the formulation developed for one-stage object detection, where the imbalance is between foreground and background regions:

\[
	\text{FL}(p_t) = -\alpha_t (1 - p_t)^{\gamma}\log p_t
\]

where $p_t$ is the model's probability for the true class, $\gamma$ ``the balance factor for loss'' and $\alpha_t$ ``the weight balance factor for samples'' (\autoref{sec:focal-loss}); under skew, the examples it up-weights are disproportionately the minority classes. Adapting it to multi-class classification, Zhao et al. report that ``$\gamma$ is set as 2 in practice, and $\alpha$ is treated as a hyper-parameter to set by cross validation.''

The division of labour there matters: $\gamma$ re-weights by \textit{difficulty} and is treated as a constant, $\alpha$ by \textit{class frequency} and is treated as dataset-specific. Conflating them is where this project went wrong (\autoref{sec:imbalance-double-correction}).

\subsection{Data-level correction: resampling and augmentation}
\label{sec:imbalance-balanced-loss}

Xia et al.~\cite{xia2020a} address the same problem at the data end and pair the two interventions, using ``temporal data augmentation strategies as well as a balanced loss \ldots jointly'' to overcome ``limited and imbalanced training samples''. The pairing is not left unexamined: their subsection \textit{``The Impact of Data Augmentation and Balanced Loss''} removes each correction in turn from the joint system and reports the balanced loss improving performance ``slightly'', with one leave-one-subject-out cell on SMIC where removing it scores higher. This is the corpus's only controlled look at composing two corrections, and \autoref{sec:imbalance-double-correction} records a far larger effect in the same direction.

\subsection{Metric-level correction}
\label{sec:imbalance-metrics}

\autoref{sec:megc-metrics} defines UF1 and the case for macro-averaged metrics. Liong et al.~\cite{liong2016} add the argument from the subject side rather than the class side, reporting F1, precision and recall because these ``provide a more meaningful perspective than accuracy rates when the datasets used are naturally imbalanced since each subject has a different number of video samples.'' The two arguments compound on CASME II, where classes are skewed \textbf{and} per-subject clip counts range from 1 to 33 (\autoref{tab:corpus-threats}c).

\subsection{Limitations}
\label{sec:imbalance-limitations}

\textbf{The interaction between corrections is examined once, and only at small magnitude.} Xia et al.~\cite{xia2020a} ablate augmentation and balanced loss out of the joint system and report the balanced loss helping only ``slightly'' --- at most about two points across their twelve cells, and in one cell costing accuracy; Zhao et al.~\cite{zhaoS2021} apply focal loss alongside a large per-sample augmentation, but without class-proportional resampling and without ablating the pairing. What no reviewed paper reports is the pairing failing outright.

\textbf{$\gamma$ is inherited rather than tuned.} Zhao et al.~\cite{zhaoS2021} take the value from the object-detection literature ``in practice''. They do compare focal loss against cross-entropy across a sweep of $\alpha$, but that experiment holds $\gamma$ fixed, so the focusing exponent remains unvalidated on this task.

\textbf{Label smoothing is named in the corpus but never examined on this task.} Dosovitskiy et al.~\cite{dosovitskiy2021} list it among the regularisation parameters they tune, but no reviewed micro-expression paper uses or discusses it, or reports what it does under class skew. It is used here and declared as a choice this corpus supports neither way.

\subsection{Implications for this research}
\label{sec:imbalance-double-correction}

\begin{table}[htbp]
	\centering
	\footnotesize
	\begin{tabular}{@{}p{5.6cm} p{8.0cm}@{}}
		\hline
		\textbf{Finding from the review} & \textbf{Decision taken in this thesis} \\
		\hline
		Focal loss addresses class-frequency bias in micro-expression recognition~\cite{zhaoS2021} & Focal loss adopted as the training objective, with \textbf{$\gamma = 2.0$}, matching Zhao et al.'s reported setting \\
		$\alpha$ is a dataset-specific class-weight term~\cite{zhaoS2021} & Implemented as an optional per-class $\alpha$ vector, \textbf{switched off at run time} --- see the declaration below \\
		Correct at the data end as well as the loss end~\cite{xia2020a} & Minority oversampling via a weighted sampler, active in every configuration --- a sampler, where Xia et al. use uniform augmentation plus a balanced loss \\
		Imbalance-aware metrics are mandatory~\cite{see2019,liong2016} & Pooled macro F1 --- identical in construction to MEGC's UF1 --- as the primary metric, with accuracy quoted only alongside the always-Negative floor \\
		\textit{(no micro-expression source)} & Label smoothing of 0.05 applied inside the focal loss; named among the regularisers tuned by Dosovitskiy et al.~\cite{dosovitskiy2021}, but declared as unsupported either way by this corpus \\
		\hline
	\end{tabular}
	\caption[Review findings and the imbalance-handling decisions they determine]{Review findings and the imbalance-handling decisions they determine.}
	\label{tab:imbalance-decisions}
\end{table}

\textbf{The finding this project contributes: correct at exactly one point, and know which.} The configuration is resolved at run time by a single rule --- class weighting in the loss is enabled only when the balanced sampler is \textit{not} active --- and the run log records the decision each fold: \textit{``Class weights in loss disabled (balanced sampler already active).''}

That rule exists because of two failures in opposite directions, both recoverable from this project's own run history. In the earliest code era the loss carried inverse-frequency class weights and no balanced sampler existed; under that single correction the proposed configuration abandoned the \textbf{Positive} class outright, predicting none of its clips --- a failure initially misdiagnosed as over-capacity overfitting. A later era added the sampler while leaving the loss weights unconditional, so both corrections ran at once; the collapse then inverted, and configurations abandoned the \textbf{Negative} majority class instead, one of them predicting a single class for every clip in the corpus. Correcting once at the loss starved the minority; correcting twice over-served it. With the guard in place and exactly one correction active, no configuration in the reported study abandons any class.

This is a contribution rather than an incident, because of how little guidance the corpus offers (\autoref{sec:imbalance-limitations}): no reviewed paper warns that combining corrections can invert the intended effect rather than merely blunt it. On a 156-clip corpus with a 4 : 1 skew, correcting twice is as damaging as correcting once in the wrong place --- the two failures differ only in which class is abandoned.

\textbf{A second declaration.} The severity of the imbalance also constrains what the ablation can conclude. Because every configuration shares the same sampler, loss and $\gamma$, this thesis measures \textbf{no} effect attributable to the imbalance treatment itself: it is a constant, not a factor. The claim supported is that the four architectural components were compared under a training regime that does not abandon minority classes, not that this regime is optimal.

\subsection{Research gap and scope}
\label{sec:imbalance-gap}

The corrections themselves are well established; what the corpus does not settle are the three gaps of \autoref{sec:imbalance-limitations} --- how two corrections compose, what $\gamma$ should be on this task, and what over-correction looks like when it fails.

This thesis contributes to the first and third, documenting two opposite class-collapse failures and the run-time rule that resolved them (\autoref{sec:imbalance-double-correction}): a negative result the corpus does not contain and one directly actionable for anyone reproducing this class of pipeline. On the second it offers no improvement, since $\gamma = 2.0$ and label smoothing of 0.05 are held constant across all twelve configurations. Sweeping either, and the Grad-CAM++ auditing and identity-disentanglement work this project's earlier stages removed, are declared as further work.


\section{Synthesis and Research Gap}
\label{sec:synthesis-and-gap}

This section consolidates the gap statements closing \autoref{sec:corpus}--\autoref{sec:imbalance-review}: what the corpus collectively establishes, the four limitations recurring across every topic, the resulting research gap, and which parts of it this thesis closes. It also collects in one place every declared divergence between the system built here and the literature it draws on.

\subsection{What the corpus establishes}
\label{sec:corpus-establishes}

Read as a whole, the reviewed literature settles a good deal. \textbf{On data:} the corpora are small, demographically narrow, laboratory-bound and irreducibly imbalanced by what can be elicited, and CASME II remains the primary single-corpus benchmark under its native leave-one-subject-out protocol (\autoref{sec:corpus}, \autoref{sec:imbalance-review}). \textbf{On representation:} motion rather than appearance is the appropriate input; a strain field adds a further cue, though on CASME II only weakly and under a protocol that is not subject-disjoint; and both the magnification factor and the interpolated sequence length have interior optima located by sweep (\autoref{sec:evm-review}--\autoref{sec:tim-review}). \textbf{On architecture:} at this data size capability bought without parameters outperforms capability bought with them, and self-attention is the strongest-performing mechanism reported in the corpus (\autoref{sec:cnn-review}--\autoref{sec:transformer-review}). \textbf{On evaluation:} accuracy alone is inadequate under skew and must give way to macro-averaged measures accumulated across folds before averaging (\autoref{sec:megc-metrics}, \autoref{sec:imbalance-metrics}).

\subsection{Four limitations that recur across every topic}
\label{sec:four-limitations}

The per-section gap statements are not nine independent observations. Four structural problems generate most of them.

\textbf{(a) Components are isolated one at a time, inside one proposed design, and never jointly.} Matched-pair measurement is not absent from the corpus. On the input side \autoref{tab:evm-effect} isolates magnification and \autoref{tab:strain-performance} isolates strain; on the architectural side Xia et al. compare a basic recurrent convolutional network against three variants adding a wide expansion, a shortcut connection and an attention unit, holding every other parameter constant across all four (\autoref{sec:cnn-limited-data}), and separately sweep depth against input resolution (\autoref{sec:cnn-shrink}); SLSTT's Mean-versus-LSTM comparison isolates its aggregator (\autoref{sec:slstt-results}).

What no reviewed study does is vary \textit{several} components together, so that their interactions can be read. Every ablation in the corpus is internal to one proposed design and moves one factor at a time from that design's own defaults; across papers, systems are compared whole --- STSTNet against OFF-ApexNet, SLSTT against STSTNet --- so a difference in the reported number cannot be attributed to any one stage. There is therefore no basis for deciding which of the learned stages are load-bearing, nor whether one component's value depends on another's presence.

\textbf{(b) The evaluation protocol is unstable, and under-reported.} Sample counts, label-set size and the choice among leave-one-subject-out, leave-one-video-out and leave-one-sample-out all vary inconsistently, and the field's most-quoted baseline is reported under two different protocols by its two primary sources and under two more within the originating paper itself (\autoref{sec:original-baseline}). Given this corpus's fold composition (\autoref{sec:corpus-threats}), these are not bookkeeping differences: they change what the numbers mean.

\textbf{(c) The strongest results depend on regimes a single-corpus study cannot access.} The best CASME II figures in the corpus come from composite-database training (\autoref{sec:megc-metrics}) and, in SLSTT's case, ImageNet pre-training (\autoref{sec:transformer-limitations}) --- neither a property of the architecture being credited. Almost no evidence exists about what these components do when trained from scratch on one small corpus, which is the situation most practitioners are in.

\textbf{(d) Preprocessing parameters have interior optima, and are transferred without re-validation.} The magnification factor, the interpolated sequence length and the input resolution all rise then fall, all have their optima located only by sweep, and none has a principled selection rule (\autoref{sec:evm-evidence}, \autoref{sec:tim-how-many}, \autoref{sec:cnn-shrink}). Values are carried between pipelines whose downstream models differ entirely: $\alpha = 30$ moves from a single-frame appearance encoder to other settings, TIM10 propagates from one baseline paper to the whole field, and $\lambda = 10^{-4}$ transfers from CIFAR though its own authors re-tuned it for ImageNet (\autoref{sec:simam-evidence}).

\subsection{The research gap}
\label{sec:research-gap}

Taking (a)--(d) together:

\textbf{The field has established a standard pipeline for micro-expression recognition --- magnify, compute motion, normalise duration, extract spatially, model temporally --- but has never measured which of its stages are responsible for the result. Because components are only ever varied one at a time inside a single proposed system, under protocols that vary in unstated ways, and predominantly in training regimes richer than a single small corpus provides, it is not currently possible to say which parts of that pipeline earn their place, which are inherited convention, and which depend on the presence of another.}

This matters practically as well as scientifically. A practitioner assembling this pipeline today has no published basis for allocating effort: for judging whether the most computationally expensive stage is the most valuable, whether a component with a negligible average effect may still matter in a specific failure, or whether a ranking read off an incomplete evaluation of a corpus this size is stable at all.

\subsection{How this thesis addresses the gap --- and what it does not}
\label{sec:addressing-gap}

\textbf{What it does.} This thesis treats the pipeline as a factorial experiment rather than a proposal: four components varied independently across every architecturally valid cell of a 2\textsuperscript{4} matrix, twelve configurations, each evaluated under \textbf{complete 25-fold leave-one-subject-out on CASME II} with every non-varied factor held identical. That gives the joint, multi-component measurement (a) says the corpus lacks; reporting the protocol in full addresses (b) for this study at least; and training from scratch is the regime (c) identifies as unevidenced. Each component's marginal contribution is read from pairs of runs differing in one factor alone, and the distribution of those four effects --- how much of the pipeline's performance each stage accounts for, whether any subtracts, and what each costs --- is the substantive answer to the gap.

Three further contributions follow from the review rather than from the experiment:

\begin{itemize}
	\item \textbf{A caution about partial evaluations.} This project's earlier evaluations, on smaller held-out sets and incomplete fold counts, ranked the twelve configurations differently from the complete run reported here. They are not evidence of protocol sensitivity, since they differ in sample count, training budget and class coverage too, and the earliest carry a data-routing defect --- but the episode illustrates limitation (b) in practice, and only the complete run is reported.
	\item \textbf{A repaired measurement.} That defect left the magnification switch inert, producing bit-identical results across every pair those runs contained. What made it detectable as a defect rather than a null result was the review's evidence that magnification has a clear, repeatable effect (\autoref{sec:evm-implications}).
	\item \textbf{Two opposite class-collapse failures.} A single loss-level correction suppressed the minority class; adding a data-level correction on top of it suppressed the majority class instead. No reviewed paper reports either collapse (\autoref{sec:imbalance-double-correction}).
\end{itemize}

\textbf{What it does not do.} Three of the four limitations are only partially addressed. On (b), this study reports its own protocol completely but cannot repair the field's inconsistency, and its own N = 156 differs from MEGC's 145. On (c), it establishes what these components do from scratch on one corpus, which is a different question from --- not an answer to --- how they behave under composite training; no cross-dataset claim is made anywhere. On (d), the thesis inherits parameter values rather than sweeping them: $\alpha = 10$, T = 32, $\lambda = 10^{-4}$ and $\gamma = 2.0$ are fixed throughout, so the ablation measures architecture at one point in preprocessing space, not across it.

\subsection{A consolidated ledger of declared divergences}
\label{sec:divergence-ledger}

Reviewing the literature closely enough to write this chapter surfaced twelve points at which the implemented system departs from the work it draws on. Several were discovered only through the review. They are collected in \autoref{tab:divergence-ledger} so that no result rests on an unstated assumption.

\begin{table}[htbp]
	\centering
	\footnotesize
	\begin{tabular}{@{}p{0.85cm} p{6.0cm} p{6.6cm}@{}}
		\hline
		\textbf{\S} & \textbf{Divergence} & \textbf{Consequence} \\
		\hline
		3.1.3 & Working set is N = 156, including sadness and fear in Negative; MEGC's CASME II subset is N = 145 & Comparisons against challenge figures are not exactly like-for-like \\
		3.1.2 & The corpus's own registered release is not used; the pipeline reads the original recorded frames, unregistered and uncropped & No landmarking or alignment of any kind precedes the network; tolerance to pose and scale must be learned from 156 clips \\
		3.2.6 & Temporal band is the unnarrowed 5--25 Hz duration rule, where Bai et al. narrowed it to 15--25 Hz as noisy & The pass-band admits the low-frequency range they rejected \\
		3.4.3 & Strain magnitude counts shear once, not twice as in the published Frobenius norm & Shear weighted lower by $\sqrt{2}$ relative to normal strain; untested \\
		3.4.7 & No dedicated strain filtering, where the literature applies Wiener and Gaussian filters & Spatial differentiation can amplify flow-estimation noise, and nothing downstream suppresses it \\
		3.5.7 & Temporal normalisation is uniform index sampling, \textbf{not} manifold-based TIM; no frames are synthesised & Avoids fabricating motion; forfeits the ability to lengthen short clips \\
		3.5.7 & Magnification is applied \textbf{after} subsampling, at a nominal 200 fps & Realised pass-band is clip-dependent and below the intended 5--25 Hz \\
		3.6.6 & Backbone kernels are (1, 3, 3) --- \textbf{no temporal mixing} & The ablation tests a learned \textit{spatial} stem, not spatio-temporal convolution \\
		3.6.6 & Backbone input is 224 $\times$ 224, against Xia et al.'s guidance of $\leq$ 100 $\times$ 100 --- though Bai et al. and Li, Huang and Zhao both use 224 $\times$ 224 & Whatever the stem is measured to contribute may be a resolution effect as much as an architecture effect \\
		3.7.7 & SimAM statistics pooled over the whole clip rather than per frame as in the source; $\lambda$ not re-tuned & Neurons judged against the clip, not the frame; defensible but untested \\
		3.8.6 & Six divergences from published SLSTT (temporal not spatial attention, mean not LSTM aggregation, from scratch, smaller, sinusoidal encoding, short-term flow) & Variable D is a temporal transformer encoder, \textbf{not} a reproduction of SLSTT \\
		3.9.7 & Label smoothing of 0.05, named in the corpus only by Dosovitskiy et al. and never examined on this task & A choice the micro-expression corpus supports neither way, held constant and unmeasured \\
		\hline
	\end{tabular}
	\caption[Every declared departure from the reviewed literature]{Every declared departure from the reviewed literature, with its consequence.}
	\label{tab:divergence-ledger}
\end{table}

Three of these --- the strain formula, the magnification ordering, and the input resolution --- were identified by the review itself rather than being known design decisions, and each is declared as proposed further work. None invalidates the ablation, since every configuration shares the same choices. What they constrain is the \textit{wording} of the conclusions: what is measured is the contribution of a spatial stem at 224 $\times$ 224 and of temporal self-attention in this pipeline, not a verdict on 3D convolution or on SLSTT.

\subsection{What this chapter contributes to the thesis}
\label{sec:ch3-contribution}

Three things carry forward. First, the \textbf{methodology that follows is determined by this review rather than chosen freely}: the corpus dictates the class grouping, representation, fold structure and metric, and each component's parameters come from published sweeps (\autoref{sec:corpus-implications}). Second, the \textbf{interpretation of the central result is prepared here rather than improvised later}: flow and strain already perform spatio-temporal feature extraction analytically (\autoref{sec:flow-implications}, \autoref{sec:strain-implications}), so a negative contribution from a spatial-only backbone at an oversized input resolution is expected rather than anomalous (\autoref{sec:cnn-implications}). Third, the \textbf{limits of what can be claimed are fixed before any number is reported}: a small, single-seed, demographically narrow corpus measured against richer published regimes, with the twelve divergences of \autoref{tab:divergence-ledger} standing throughout --- consequences of this review, not concessions made after the fact.
 in place of contents/chapter3.tex.
% Do not include both: they share every label.
%
% Assembled from chapter3_short.tex (3.1-3.2) and chapter3_short_part3..10.tex.
% Those part files are retained unchanged as the editing units; regenerate this
% file from them rather than editing it directly, or edit here and retire them.
%
% Structure that other chapters depend on, and that must not be disturbed:
%   * Chapter 2 cites Sections 3.1.2, 3.1.6, 3.1.8, 3.3.1, 3.3.3, 3.5.7,
%     3.9, 3.9.5 and 3.9.7 BY NUMBER, so those subsection positions are fixed.
%   * The divergence ledger (Table, sec:divergence-ledger) cites twelve
%     subsections by number. scratchpad/ledgercheck.py re-verifies all of them
%     plus every external numeric reference after any renumbering.
%   * Chapters 4, 5, 6 and the introduction reference sec:research-gap,
%     sec:four-limitations, sec:addressing-gap, sec:divergence-ledger,
%     sec:synthesis-and-gap, sec:imbalance-review, sec:imbalance-gap,
%     sec:imbalance-double-correction, sec:cnn-review, sec:cnn-implications,
%     sec:tim-implications, sec:evm-implications, sec:evm-evidence,
%     sec:evm-limitations, sec:evm-motion-hypothesis, sec:corpus,
%     sec:corpus-threats, sec:corpus-implications, sec:label-taxonomy,
%     sec:megc-metrics, sec:sibling-corpora, sec:flow-review, sec:flow-implications,
%     sec:strain-review, sec:strain-implications, sec:strain-divergence,
%     sec:simam-gap, sec:simam-implications, sec:slstt-results,
%     sec:transformer-implications, sec:addressing-gap, tab:megc-cde-results
%     and tab:spontaneous-corpora.
%
% Conventions held throughout: no sentence points at another chapter; every
% float carries a short \caption[...] for the List of Figures / List of Tables;
% figures are fig3_1_dataset.png and fig3_2_fold_composition.png only.
% =============================================================================

\chapter{Literature Review}
\label{ch:literature}

This chapter reviews the literature bearing on each component of the recognition pipeline evaluated in this thesis: the corpus on which every experiment is run, the four techniques placed under ablation, the three input transformations they operate on, and the training regime held constant across all of them.

It is organised as a single continuous argument rather than as ten independent surveys: a fact is established once, in the section that owns it, and cross-referenced thereafter. Each section runs from what the technique addresses, through how it works and what the published evidence shows, to its limitations and the implications for this thesis, closing with the gap and how this work differs. Two claims recur --- that on a corpus of 156 clips several decisions usually presented as modelling choices are in fact forced by properties of the data, and that several results usually attributed to architecture are attributable to the evaluation protocol. \autoref{sec:synthesis-and-gap} consolidates the nine closing statements into a single research gap and collects every declared departure from the reviewed methods.


% =============================================================================
\section{The Evaluation Corpus: CASME II and the Spontaneous Micro-Expression Benchmark Family}
\label{sec:corpus}
% =============================================================================

This section covers the corpus literature only --- the papers that introduce, specify or standardise evaluation on the database used throughout this thesis. It is treated first because on a database of this size several design decisions usually presented as modelling choices are \textit{forced} by properties of the data.

\subsection{Why a dedicated spontaneous corpus was necessary}
\label{sec:why-spontaneous}

Automatic facial expression recognition matured on large corpora of \textit{posed} expressions --- CK+, MUG, MMI, JAFFE and Multi-PIE --- reporting accuracies above 90\,\% on the six basic posed emotions by the early 2010s~\cite{yan2014}. Micro-expression research could not follow that path: as \autoref{sec:phenomenon} establishes, a micro-expression is involuntary and cannot be \textit{performed} on instruction. Yan et al.~\cite{yan2014} review a pre-2013 landscape of only four micro-expression databases, of which only two were spontaneous (\autoref{tab:spontaneous-corpora}). Two problems in it motivated CASME II directly.

The first is \textbf{construct validity}. Yan et al.~\cite{yan2014} and Li et al.~\cite{liX2013} both note that posed ``micro-expressions'', produced deliberately on instruction, differ systematically from spontaneous ones --- Li et al.~\cite{liX2013} further observing that USF-HD adopted a 2/3-second duration criterion longer than the definitional upper bound of \autoref{sec:phenomenon}. A system trained on posed data is therefore not demonstrably a micro-expression recogniser; it may be a recogniser of fast deliberate movement.

The second is \textbf{acquisition quality}. Defined jointly by very short duration and very low intensity, the phenomenon sits close to the noise floor of ordinary video: Yan et al.~\cite{yan2014} argue that CASME's 60~fps coarsely samples a movement lasting a fraction of a second --- while marking the benefit of higher rates as unresolved --- and that flicker under alternating-current lighting contaminated earlier recordings. Elicitation was a third difficulty: high-stakes deception paradigms produced samples heavily contaminated by conversational and other emotion-irrelevant movement~\cite{yan2014}, and the paradigm that proved workable, which CASME II adopts, is to have participants watch high-arousal video episodes while attempting to keep the face neutral~\cite{liX2013}.

\begin{table}[htbp]
	\centering
	\footnotesize
	\setlength{\tabcolsep}{4pt}
	\begin{tabular}{l c c c c c c c}
		\hline
		\textbf{Corpus} & \textbf{Subj.} & \textbf{Micro samples} & \textbf{fps} & \textbf{Facial res.} & \textbf{Spont.} & \textbf{FACS} & \textbf{Classes} \\
		\hline
		\multicolumn{8}{l}{\textit{The pre-CASME II landscape, reproduced from Yan et al.~\cite{yan2014} (Table 1)}} \\
		USF-HD~\cite{shreve2011}                       & ---      & 100 & 30  & ---              & No  & No  & 6 \\
		Polikovsky et al.\textsuperscript{\textdagger} & 10       & 42  & 200 & ---              & No  & No  & 6 \\
		SMIC~\cite{liX2013}                            & 16 valid & 164 & 100 & 190 $\times$ 230 & Yes & No  & 3 \\
		CASME (2013)\textsuperscript{\textdagger}      & 19 valid & 195 & 60  & 150 $\times$ 190 & Yes & Yes & 7 \\
		\hline
		\multicolumn{8}{l}{\textit{The spontaneous corpora that define the current context}} \\
		\textbf{CASME II}~\cite{yan2014} & \textbf{26 valid} & \textbf{247 / 255} & \textbf{200} & \textbf{280 $\times$ 340} & \textbf{Yes} & \textbf{Yes} & \textbf{5 (7 raw)} \\
		SAMM~\cite{davison2018}          & 32 & 159             & 200 & 400 $\times$ 400 & Yes & Yes & 7 \\
		CAS(ME)$^2$~\cite{qu2016}        & 22 & 53 (+250 macro) & 30  & ---              & Yes & Yes & 4 \\
		\hline
		\multicolumn{8}{l}{\textit{Ethnicities / mean age (SD):} SMIC 3 / 26.7 \quad CASME 1 / 22.03 (1.60) \quad \textbf{CASME II 1 / 22.03 (1.60)}} \\
		\multicolumn{8}{l}{\quad SAMM 13 / 33.24 (11.32) \quad CAS(ME)$^2$ 1\textsuperscript{*} / 22.59 (2.20)} \\
		\hline
	\end{tabular}
	\caption[Posed and spontaneous micro-expression corpora]{Posed and spontaneous micro-expression corpora, compiled from Yan et al.~\cite{yan2014} (Table 1), Davison et al.~\cite{davison2018} (Table 1), Li et al.~\cite{liX2013} and Qu et al.~\cite{qu2016}. \textdagger{}~Described in Yan et al.; the primary papers lie outside the review corpus and are not cited separately. \textsuperscript{*}~Inferred from the recruitment description; Qu et al. do not state it. SMIC's demographic figures are those of all 20 recruited participants; the 16 who produced usable clips have a mean age of 28.1 and span two ethnicities.}
	\label{tab:spontaneous-corpora}
\end{table}

\subsection{CASME II: construction, apparatus, and annotation}
\label{sec:casme2-construction}

CASME II~\cite{yan2014} is the corpus on which every experiment in this thesis is performed, and several of its specifics propagate into the data pipeline used here.

\textbf{Participants and apparatus.} Thirty-five participants were recruited, mean age 22.03 years (SD = 1.60), and watched high-valence video episodes pre-rated by an independent panel of twenty raters, under instruction to suppress all facial movement. Recording sustained \textbf{200 fps} at 640~$\times$~480 with no inter-frame compression, under four LED lamps with umbrella reflectors giving steady flicker-free illumination --- a deliberate correction of the artefacts observed in SMIC. The facial region is approximately \textbf{280 $\times$ 340 pixels}, over three times CASME's facial area and roughly double SMIC's.

\textbf{Selection and annotation.} From nearly 3,000 elicited movements a small subset was retained, very subtle movements being discarded because their onset and offset frames could not be reliably located. Each retained sample carries onset and offset frame indices, Facial Action Coding System (FACS) action-unit labels from two independent coders with disagreements arbitrated by discussion, and an emotion label. Yan et al.~\cite{yan2014} argue against forcing spontaneous micro-expressions into the six basic categories, since a single action unit is frequently ambiguous --- AU4 (brow lowerer) may indicate disgust, anger, attention or tension --- and assign the emotion label from a \textit{combination} of the coded AUs, the participant's self-report, and the emotional content of the eliciting episode.

\textbf{Preprocessing in the original release.} The baseline applies a three-step normalisation from which this thesis departs: a frontal neutral model face is chosen; 68 Active Shape Model landmarks are detected on the \textit{first frame only} of each clip and registered to the model face by a Local Weighted Mean transform; that transform is applied unchanged to every frame, and the face is cropped by eye position. Yan et al.~\cite{yan2014} justify landmarking only the first frame on two grounds --- that rigid head motion within so short a window is negligible, and that per-frame ASM landmarks are unstable enough to introduce jitter where the face did not in fact move.

\textbf{A note on sample count.} The published paper reports \textbf{247} micro-expressions from 26 valid participants; the coding file distributed with the released database contains \textbf{255}. This discrepancy, a consequence of post-publication revision of the coding, is rarely acknowledged and is one reason sample counts reported for ``CASME II'' vary between papers. The present work uses the released 255-sample coding.

\subsection{The label taxonomy, and the sparsity problem it creates}
\label{sec:label-taxonomy}

Yan et al.~\cite{yan2014} provide five main emotion categories with AU-based criteria, and are candid about the resulting imbalance: ``some types of facial expressions are difficult to elicit in laboratory situations, thus the samples in different categories distributed unequally.'' The released coding is finer-grained still, exposing seven raw labels; counting directly from the label table used in this thesis (\texttt{Processed\_Data/master\_thesis\_labels.csv}, 255 rows) gives the middle columns of \autoref{tab:casme2-labels}.

\begin{table}[htbp]
	\centering
	\footnotesize
	\begin{tabular}{l l c c l c}
		\hline
		\multicolumn{3}{c}{\textbf{Five categories, as published}} & \multicolumn{2}{c}{\textbf{Released coding (255)}} & \textbf{Grouped} \\
		\textbf{Emotion} & \textbf{AU criterion} & \textbf{N} & \textbf{N} & \textbf{Raw label} & \textbf{class} \\
		\hline
		Happiness  & AU6 or AU12                           & 33  & 32 & happiness  & Positive \\
		Disgust    & AU9, AU10, or AU4+AU7                 & 60  & 63 & disgust    & Negative \\
		Surprise   & AU1+2, AU25, or AU2                   & 25  & 25 & surprise   & Surprise \\
		Repression & AU15 or AU17, alone or combined       & 27  & 27 & repression & Negative \\
		Others     & other emotion-related facial movement & 102 & 99 & others     & \textit{discarded} \\
		---        & ---                                   & --- & 7  & sadness    & Negative \\
		---        & ---                                   & --- & 2  & fear       & Negative \\
		\hline
		\multicolumn{6}{l}{\textbf{Working set:} Negative $63+27+7+2 = $ \textbf{99}; Positive \textbf{32}; Surprise \textbf{25}; \textbf{N = 156}} \\
		\hline
	\end{tabular}
	\caption[CASME II label taxonomy and the three-class working set]{The CASME II label taxonomy at three granularities: the five published categories with their AU criteria and frequencies, after Yan et al.~\cite{yan2014} (Table 3); the seven raw labels of the released coding as parsed for this thesis; and the three-class valence grouping that yields the 156-clip working set used throughout.}
	\label{tab:casme2-labels}
\end{table}

\textbf{Seven-class classification is not viable on this corpus.} Fear is represented by two clips in the entire database and sadness by seven. Under any subject-disjoint protocol --- and \autoref{sec:original-baseline} shows subject-disjointness has been the norm for this benchmark since its introduction --- a two-clip class cannot be meaningfully trained or tested: a fold either contains the class or does not, and per-class metrics degenerate.

The field's response has been to group by affective valence. The convention originates with SMIC, released with exactly three classes --- positive, negative, surprise~\cite{liX2013} --- and was carried into cross-corpus work by the Second Facial Micro-Expression Grand Challenge (MEGC 2019)~\cite{see2019}. MEGC states the scheme twice and not identically: as a challenge-wide rule in prose, with Negative covering Repression, Anger, Contempt, Disgust, Fear and Sadness, and again as per-corpus footnotes under its sample table, where CASME II's Negative narrows to Disgust and Repression alone. \autoref{tab:casme2-labels} follows the challenge-wide rule.

\begin{figure}[htbp]
	\centering
	\includegraphics[width=\textwidth]{fig3_1_dataset.png}
	\caption[Composition of CASME II]{Composition of CASME II. Left: all 255 released clips by raw label, with the excluded ``others'' category in grey. Right: the three-class pool of 156 clips used throughout this thesis, in the ratio 99 : 32 : 25; the grey ``others'' bar is shown again for scale and is not part of the pool.}
	\label{fig:dataset-composition}
\end{figure}

\textbf{A declared divergence from the MEGC convention.} MEGC 2019 reports the CASME II portion of its composite database as \textbf{145} samples and footnotes that ``Negative class of CASME II: Disgust and Repression''~\cite{see2019} --- building its subset from the paper's \textit{five-category} release, in which sadness and fear are absorbed into ``others''. This thesis uses the released seven-label coding and assigns sadness and fear to Negative, following the mapping MEGC gives for \textbf{SAMM} rather than the one it gives for CASME II. The consequence is an \textbf{eleven-clip} difference: N = 156 against N = 145, with a Negative class of 99 rather than 88. Nine of the eleven are the sadness (7) and fear (2) clips the seven-label coding exposes; the remaining two are unaccounted for, since disgust and repression together number 90 in the released coding against the 88 the challenge reports --- a residual noted rather than resolved. The direction is mild: the nine added clips are all in the majority class and drawn from the categories the coders found hardest to assign, so if anything they make the present task marginally harder than the challenge's.

The residual imbalance is severe in either case --- 99 : 32 : 25, roughly 4 : 1.3 : 1 --- and everything in \autoref{sec:corpus-threats} concerning metric choice, sampling and loss follows from it.

\subsection{The original baseline, and what it does and does not establish}
\label{sec:original-baseline}

Yan et al.~\cite{yan2014} provide a baseline using Local Binary Patterns from Three Orthogonal Planes (LBP-TOP)~\cite{zhaoG2007} with a Support Vector Machine. LBP-TOP extends the static LBP operator by extracting codes on the XY, XT and YT planes of the video volume and concatenating the three histograms, so that the XT and YT planes encode dynamic rather than static texture. The baseline divides the face into 5~$\times$~5 blocks, sets $P = 4$ neighbouring points, and sweeps the spatial radii $R_X, R_Y \in \{1,2,3,4\}$ and temporal radius $R_T \in \{2,3,4\}$; $R_T = 1$ is deliberately excluded because at 200~fps the change between adjacent frames is negligible --- an early acknowledgement that a high sampling rate is not automatically an advantage. Five-class accuracies under \textbf{leave-one-subject-out cross-validation} cluster in the high fifties to low sixties, the best configuration ($R_X = R_Y = 1$, $R_T = 4$) reaching \textbf{63.41\,\%}.

Three observations matter here.

\textbf{First, leave-one-subject-out is native to this benchmark.} It is not a stricter protocol adopted later by cautious authors; it is the protocol of the paper that released the corpus. Any evaluation reporting a single random train/test split departs from the benchmark's own convention.

\textbf{Second, the baseline is reported as accuracy alone}, with no per-class or macro-averaged score. On a five-class problem in which ``others'' accounts for 102 of 247 samples, accuracy is a weak instrument: the majority-class-constant classifier scores 41\,\% without modelling anything. In the three-class working set the floor is higher still --- always predicting Negative yields \textbf{0.6346 accuracy} at a macro F1 of only \textbf{0.2588}. A headline accuracy near 63.41\,\% is therefore not, on its own, evidence of a working system.

\textbf{A caution on this figure.} Yan et al.~\cite{yan2014} state in both their abstract and their results prose that leave-one-subject-out was applied to obtain 63.41\,\% --- but the caption of the very table reporting it reads ``leave-one-out cross-validation'', without the subject qualifier. Li et al.~\cite{liX2018}, reproducing the same figure in their comparison table, mark it with an asterisk denoting ``results achieved using leave-one-sample-out cross validation'' --- a materially easier protocol. The sources therefore disagree on the protocol behind the most widely quoted number on this benchmark, and the originating paper disagrees with itself. This thesis follows that paper's prose, but records the discrepancy rather than resolving it, because it affects how the figure should be read in any comparison.

\textbf{Third, the baseline's sensitivity to a purely descriptive hyper-parameter is itself informative.} Accuracy moves by several points across the radius sweep with no change to the data, the split or the classifier. Where effect sizes of that magnitude arise from descriptor settings, differences of similar magnitude between architectures cannot be interpreted without a variance estimate --- a limitation this thesis inherits and declares.

\subsection{Sibling corpora, and why this study remains single-corpus}
\label{sec:sibling-corpora}

Three further databases define the context in which CASME II is used; \autoref{tab:spontaneous-corpora} gives their properties, and each isolates a specific limitation of CASME II.

\textbf{SMIC~\cite{liX2013}}, the first spontaneous micro-expression database, is the origin of the three-class taxonomy. Its high-speed subset of 164 clips is distributed 51 positive, 70 negative, 43 surprise, and its baseline --- LBP-TOP with a temporal interpolation model (TIM) normalising each clip to ten frames, under leave-one-subject-out --- reaches \textbf{48.78\,\%}. Two elements of its design are inherited directly here: the three-class valence grouping, and temporal interpolation to normalise clips of unequal length. It is not used as a corpus because it is neither FACS-coded nor apex-annotated.

\textbf{SAMM~\cite{davison2018}} addresses the demographic narrowness of the CASME family, contributing 13 ethnicities and a mean age of 33.24 years against CASME II's single ethnicity and 22.03. Davison et al.~\cite{davison2018} observe that a corpus drawn from one ethnicity and one narrow age band restricts analysis to ``similar looking participants'', and tailor stimulus selection to each participant in advance rather than relying on post-hoc self-report. The design is stronger than CASME II's on both counts, but the corpus is smaller in usable samples and spreads them across seven categories with the same sparsity difficulty.

\textbf{CAS(ME)$^2$~\cite{qu2016}} provides 250 macro-expression and 53 micro-expression samples from the same participants under identical conditions, which makes it the natural resource for studying macro-to-micro domain transfer. That strategy produced the strongest entry in MEGC 2019 --- Liu et al.'s EMR~\cite{liu2019}, the only reported submission to use macro-expression data, though it drew that data from CK+ rather than from CAS(ME)$^2$~\cite{see2019}. Two properties limit its use here: 30~fps is an order of magnitude below CASME II's and poorly matched to sub-200-ms movements, and only 53 micro-expression samples exist. Its baseline is also evaluated leave-\textit{one-video}-out, which does not enforce subject-disjointness and is therefore not comparable to the protocol used here.

\textbf{Why this thesis uses CASME II alone.} It has the highest temporal resolution of any spontaneous corpus of comparable size, a precondition for magnification in a defined frequency band and for dense frame-to-frame flow; it is FACS-coded with onset, apex and offset annotation, which makes temporal normalisation to a fixed length well-defined rather than heuristic; and it remains the most widely reported single-corpus benchmark, which makes comparison against published results meaningful. The cost --- that no claim here extends to other populations, capture conditions or corpora --- is stated as a declared limitation rather than argued away.

\subsection{How CASME II is evaluated in the field: MEGC 2019 and the metric convention}
\label{sec:megc-metrics}

The most consequential development in how CASME II is \textit{used}, as opposed to how it was built, is the Composite Database Evaluation (CDE) protocol of the Second Facial Micro-Expression Grand Challenge~\cite{see2019}. This thesis adopts the challenge's metric definitions while declining its composite training regime, so the protocol warrants precise description.

CDE merges SMIC, CASME II and SAMM into one composite database under the three-class grouping of \autoref{sec:label-taxonomy}, giving \textbf{442 samples} --- Negative 250, Positive 109, Surprise 83 --- of which CASME II contributes 145. Leave-one-subject-out is run across the composite, which the organisers note ``ensures subject-independent evaluation''. Plain accuracy is rejected explicitly on grounds of class imbalance, and two balanced metrics mandated: \textbf{Unweighted F1 (UF1)}, the macro-averaged F1-score, which \textit{accumulates} true positives, false positives and false negatives per class \textbf{over all $k$ folds}, computes a per-class F1 from those pooled counts and averages without weighting; and \textbf{Unweighted Average Recall (UAR)}, the balanced accuracy, obtained by averaging per-class recall.

\[
	F1_c = \frac{2 \cdot TP_c}{2 \cdot TP_c + FP_c + FN_c}, \qquad \text{UF1} = \frac{1}{C}\sum_c F1_c
\]

UF1 is a \textbf{pooled} quantity: counts accumulate across folds \textit{before} any F1 is computed, which \autoref{sec:pooled-vs-fold} establishes is a different estimator from the average of per-fold macro F1 scores. The distinction is not cosmetic here. Because clips are distributed very unequally across subjects, and many single-subject folds contain only one of the three classes, a per-fold macro F1 is bounded well below 1 by fold composition alone and its average across folds bounded correspondingly (\autoref{fig:fold-composition}). The primary metric of this thesis, \textit{pooled macro F1}, is therefore identical in construction to MEGC's UF1; the per-fold average appearing in the raw experimental logs is not a metric the field uses at all.

\begin{figure}[htbp]
	\centering
	\includegraphics[width=\textwidth]{fig3_2_fold_composition.png}
	\caption[Class composition of the 25 leave-one-subject-out folds]{Class composition of the 25 leave-one-subject-out folds, and the per-fold macro-F1 ceiling it implies.}
	\label{fig:fold-composition}
\end{figure}

\begin{table}[htbp]
	\centering
	\begin{tabular}{l c c}
		\hline
		\textbf{Method} & \textbf{Full composite UF1 / UAR} & \textbf{CASME II subset UF1 / UAR} \\
		\hline
		LBP-TOP~\cite{zhaoG2007}                  & 0.5882 / 0.5785 & 0.7026 / 0.7429 \\
		Bi-WOOF~\cite{liong2018}                  & 0.6296 / 0.6227 & 0.7805 / 0.8026 \\
		OFF-ApexNet~\cite{liong2019a}             & 0.7196 / 0.7096 & 0.8764 / 0.8681 \\
		Quang et al.\textsuperscript{\textdagger} & 0.6520 / 0.6506 & 0.7068 / 0.7018 \\
		Zhou et al.\textsuperscript{\textdagger}  & 0.7322 / 0.7278 & 0.8621 / 0.8560 \\
		Liong et al.~\cite{liong2019b}, STSTNet   & 0.7353 / 0.7605 & 0.8382 / 0.8686 \\
		Liu et al.~\cite{liu2019}, EMR            & \textbf{0.7885 / 0.7824} & 0.8293 / 0.8209 \\
		\hline
	\end{tabular}
	\caption[MEGC 2019 Composite Database Evaluation results]{MEGC 2019 CDE results, reproduced from See et al.~\cite{see2019} (Table IV). \textdagger{} outside the review corpus; reported here only as it appears in that table.}
	\label{tab:megc-cde-results}
\end{table}

\textbf{One structural caveat governs these numbers.} Every row of \autoref{tab:megc-cde-results} comes from a model trained on the \textit{composite} database --- roughly 442 samples across three corpora and 68 subjects --- then scored on the CASME II portion of the pooled predictions. The CASME II column is therefore not a CASME II-only result: it reports how a cross-corpus model performs on CASME II clips, with approximately \textbf{2.8 times the training data} available in any single-corpus fold, where a single-corpus LOSO experiment of the kind conducted here trains each fold on roughly 150 clips. The two regimes are not interchangeable, and single-corpus results should not be compared directly against \autoref{tab:megc-cde-results} without stating the difference.

Two further observations from the challenge bear on this thesis's design. The top three submissions all chose \textbf{optical flow} as their feature representation, a convergence that supports the motion-based input used here; and one submission substituted mid-position frames for annotated apex frames with little loss, suggesting that precise apex localisation may be less critical than assumed and that methods consuming the full onset-to-offset sequence are not obviously disadvantaged.

\subsection{Threats to validity carried by the corpus}
\label{sec:corpus-threats}

The properties in \autoref{tab:corpus-threats} are not criticisms of CASME II, which remains the best available resource for this task. They are constraints any study built on it inherits, and each is answered by a specific design decision --- which is the sense in which, on a corpus of this size, several choices usually presented as modelling decisions are determined by the data. Fold and duration statistics are computed directly from the 156-clip label table used in this thesis.

\begin{table}[htbp]
	\centering
	\scriptsize
	\begin{tabular}{@{}p{0.3cm} p{2.0cm} p{5.3cm} p{5.6cm}@{}}
		\hline
		& \textbf{Threat} & \textbf{What the corpus contains} & \textbf{Design decision it forces} \\
		\hline
		(a) & Scale
		& 156 usable clips across three classes --- small enough that parameter count is itself a risk factor.
		& \textbf{Parameter-free or low-parameter components} wherever a choice exists, motivating SimAM~\cite{yang2021}, which adds zero learnable parameters; shallow rather than deep spatial backbones, extending the model-shrinking argument of Xia et al.~\cite{xia2020b}, which they establish for the composite-database setting rather than a single small corpus. \\
		(b) & Class imbalance
		& 99 : 32 : 25. The majority-class-constant classifier attains 0.6346 accuracy at a macro F1 of 0.2588.
		& Minority oversampling via a balanced sampler with focal loss, inverse-frequency class weighting \textbf{disabled}, since both corrections together induced single-class collapse. Accuracy is never reported alone, and the always-Negative floor is quoted alongside every result. \\
		(c) & Subject-level non-uniformity
		& \textbf{Subject 17 alone contributes 33 clips, 21\,\% of the working set}; subjects 8, 10 and 21 contribute one each, so fold sizes vary by a factor of 33. \textbf{Ten of the 25 folds contain a single class}, eight contain two, only seven all three.
		& \textbf{Full 25-fold leave-one-subject-out}, every subject held out exactly once, no clip untested, subject-disjointness guaranteed by construction. Primary metric is \textbf{pooled macro F1}, identical in construction to MEGC's UF1~\cite{see2019}; per-fold averaged macro F1 is rejected, since fold composition bounds it below the study's own target. \\
		(d) & A subject that disappears
		& Of the 26 participants, subject 18's three clips are all ``others'' and are removed by the three-class filter.
		& The LOSO run has \textbf{25 folds, not 26} --- stated explicitly for reproducibility, since fold counts differ between papers using this corpus. \\
		(e) & Demographic homogeneity
		& All participants Chinese, mean age 22.03 (SD 1.60) --- named by Davison et al.~\cite{davison2018} as a principal motivation for SAMM. The MEGC organisers saw the effect in reverse: adaptation from CK+ helped on SMIC and SAMM but ``not so for CASME II, which contain predominantly Chinese subjects''~\cite{see2019}.
		& Input is a \textbf{motion representation} --- dense optical flow $(u,v)$ plus optical strain --- rather than raw pixels, since motion is invariant to identity and illumination whereas appearance is exactly what a 26-subject corpus would let a network memorise. \\
		(f) & Laboratory ecology
		& Frontal recording, fixed flicker-free lamps, participants told not to move and not speaking; Yan et al.~\cite{yan2014} flag ``natural conversation and interaction'' as future work. Almost none of the head motion, illumination variation, occlusion or speech-related movement a deployed system would face is present.
		& No cross-dataset or in-the-wild claim is made anywhere in this thesis; generalisation is declared as further work. \\
		(g) & Annotation dependence
		& Labels derive from coded AUs, self-report and stimulus content --- a construct produced by trained coders under a rubric, with disagreements arbitrated by discussion rather than reported as residual uncertainty.
		& The coding is taken as given; its construct status is recorded as a declared limitation rather than treated as measurement error. \\
		(h) & Duration heterogeneity
		& Onset-to-offset lengths range from \textbf{31 to 126 frames}, median 66 --- at 200~fps, \textbf{0.155\,s to 0.63\,s}, median 0.33\,s. The upper end exceeds the 1/2-second criterion of \autoref{sec:phenomenon}.
		& Every clip temporally resampled to a fixed \textbf{T = 32} frames, following the TIM precedent of SMIC's baseline~\cite{liX2013}. \\
		(i) & Version ambiguity
		& Three sample counts circulate under the name ``CASME II'' --- 247 (the paper), 255 (the released coding), 145 (the MEGC subset) --- and both five- and seven-class label sets are in use, rarely stated.
		& The released 255-clip coding is used, N is reported explicitly alongside every figure, and the eleven-clip divergence from the MEGC subset is declared in \autoref{sec:label-taxonomy}. \\
		\hline
	\end{tabular}
	\caption[Threats to validity carried by CASME II and the decisions they force]{Threats to validity carried by CASME II, and the design decision each forces.}
	\label{tab:corpus-threats}
\end{table}

\subsection{Implications for this research}
\label{sec:corpus-implications}

Three further corpus properties are not threats but still determine settings. The \textbf{200 fps} capture rate fixes the Eulerian magnification pass-band at \textbf{5--25 Hz}, matched to the recording rate. The \textbf{seven raw labels}, with fear at two clips and sadness at seven, force the three-class valence grouping of \autoref{sec:label-taxonomy}. And the baseline's \textbf{sensitivity to a descriptive hyper-parameter} (\autoref{sec:original-baseline}) sets the standard of evidence: effect sizes below the resolution of a 156-clip corpus are reported as unresolved rather than as findings, and no confidence interval is available, since the single-seed design and unretained per-clip predictions preclude one.

Three consequences then shape the interpretation of the results rather than merely the construction of the pipeline. \textbf{First, the corpus determines the metric, and the metric determines the conclusion:} the fold composition of \autoref{tab:corpus-threats}(c) caps a macro F1 averaged over folds by arithmetic alone, before any classifier is considered, so reporting the pooled quantity is not a stylistic preference but the correction MEGC 2019 already standardised. \textbf{Second, a single held-out split is not a reliable measuring instrument here:} with 25 subjects distributed as unevenly as \autoref{tab:corpus-threats}(c) describes, which subjects fall in the test set can move the result by more than any architectural change under study, so only a complete leave-one-subject-out run is reported and the incomplete evaluations that preceded it are set out with their exclusion reasons rather than suppressed. \textbf{Third, the small-sample regime predicts which components should help:} 156 clips cannot supply the data needed to learn a general spatio-temporal feature extractor from scratch, particularly once the input has been reduced to flow and strain --- that is, once much of that extraction has been performed analytically before the network sees the data --- whereas a component modelling the temporal arc of an expression exploits exactly the structure the corpus's onset/apex/offset annotation guarantees is present and correctly delimited.

\subsection{Limitations of the existing corpus literature, and how this work differs}
\label{sec:corpus-lit-limitations}

The corpus literature establishes three things well and one thing poorly. It establishes the phenomenon's definitional boundaries, an elicitation methodology producing genuinely spontaneous samples, and a capture standard --- 200 fps, flicker-free illumination, FACS coding with onset/apex/offset annotation --- that makes fine-grained motion analysis feasible. What it does not establish is a stable basis for comparing methods. The original baseline reports accuracy only, on a five-class problem whose majority class is 41\,\% of the data, and its own hyper-parameter sweep moves the headline figure by several points with no change to the model. Sample counts, label granularity and class groupings differ across papers using the same database --- 247, 255 and 145 all appear under the name ``CASME II'' --- and are frequently not stated. MEGC 2019's CDE protocol fixed the metric problem convincingly, but within a composite training regime that makes its per-database columns non-comparable to single-corpus work, a distinction subsequent papers routinely elide. Above all, only one study in this group quantifies how much of a reported difference between methods is attributable to the evaluation protocol rather than the methods themselves --- Li et al.~\cite{liX2018} add a single matched row of their own method under both protocols --- and none does so across methods, despite a subject distribution that makes the question unavoidable.

This thesis differs in three respects. It runs the \textbf{complete} 25-fold leave-one-subject-out protocol over \textbf{every} valid cell of a four-component ablation matrix rather than reporting a single proposed configuration, so each component's contribution is measured from matched pairs differing in exactly one factor. It adopts MEGC's pooled macro F1 as its primary metric while explicitly declining the composite training regime, so its numbers describe single-corpus performance and are labelled as such. And it reports one complete protocol rather than the strongest of several, setting out which earlier evaluations were run and why each is excluded rather than suppressing that history. The intended contribution is therefore not a recognition result on CASME II, but a measurement of how much of such a result is attributable to each stage of the pipeline that produces it.

The 200 fps capture rate established here is what lets Eulerian Video Magnification, reviewed next, target a matched temporal band.


% =============================================================================
\section{Motion Magnification: Eulerian Video Magnification as a Pre-Processing Stage}
\label{sec:evm-review}
% =============================================================================

This section reviews Eulerian Video Magnification (EVM) through the four review-corpus papers that apply it: Bai et al.~\cite{bai2021}, Li et al.~\cite{liX2018}, and Li, Huang and Zhao~\cite{liY2018,liY2021}. The original formulations they build on --- Wu et al.'s amplitude-based method and Wadhwa et al.'s phase-based method --- lie outside the corpus and are described only as attributed through those four papers.

\subsection{The problem magnification addresses}
\label{sec:evm-problem}

The corpus design of \autoref{sec:why-spontaneous} responds to the \textit{short-duration} half of the micro-expression definition but not to the \textit{low-intensity} half. CASME II's 200~fps sampling spreads a movement completing in under 200~ms across dozens of frames~\cite{yan2014}, yet that movement may still occupy only a few grey levels of change across a small patch of skin --- and Yan et al.~\cite{yan2014} report discarding samples outright because the movement was too faint for coders to place an onset or offset. Li et al.~\cite{liX2018} state the difficulty directly: ``the intensity levels of facial movements are too low to be distinguishable.'' Their remedy is to amplify the motion itself before any descriptor is computed, so that the \textit{difference between expression categories is enlarged} rather than the classifier being made more sensitive. This is a data-level intervention: the model is untouched and the modification is entirely upstream, which is what makes EVM cleanly separable as an ablation factor in this thesis.

\subsection{How Eulerian magnification works, and the two variants}
\label{sec:evm-how}

\textbf{Amplitude-based magnification (AMM)} is the formulation set out in \autoref{sec:motion-magnification} --- Laplacian pyramid, temporal band-pass, multiplication by $\alpha$, reconstruction --- realising its band-pass step with a \textbf{Butterworth} filter. It is the variant implemented in this thesis. \textbf{Phase-based magnification (PMM)}, of which Bai et al.~\cite{bai2021} give the clearest account in the corpus, decomposes each frame with \textbf{octave complex steerable pyramids} over four orientations and applies the temporal filter to the \textit{local phase} rather than the amplitude. Bai et al.~\cite{bai2021} explain the advantage in signal terms: the steerable pyramid ``has impulse response with spatial support,'' which makes it easier to isolate the intended temporal frequencies while suppressing the remainder. The practical consequences they report are that the phase-based method \textbf{supports larger amplification factors} and is \textbf{notably less sensitive to noise}.

Three parameters are therefore free in any EVM application, and all three recur below: the \textbf{temporal pass-band}, which decides \textit{which} motions are amplified; the \textbf{magnification factor $\alpha$}, which decides \textit{how much}; and the \textbf{pyramid depth}, which decides at which spatial scales amplification occurs.

\textbf{Setting the temporal band.} Bai et al.~\cite{bai2021} report the rule given by the original EVM authors: the pass-band should correspond to the duration of the phenomenon being amplified --- for micro-expressions, 1/5\,s to 1/25\,s, equivalently \textbf{5 Hz to 25 Hz}. They also report that applying this full band produced ``noticeable noise'' in their setting, and that they narrowed it to 15--25 Hz accordingly. This is the only explicit guidance on band selection anywhere in the review corpus, and it is the basis for the band used in this thesis.

\subsection{Evidence that magnification improves recognition}
\label{sec:evm-evidence}

The strongest empirical case in the review corpus is Li et al.~\cite{liX2018}, who sweep the magnification factor systematically rather than fixing it. Clips are magnified at what the paper describes as ten levels while listing nine values, $\alpha \in \{1, 2, 4, 8, 12, 16, 20, 24, 30\}$ --- an inconsistency in the source rather than in this account --- where $\alpha = 1$ denotes no magnification. Each magnified clip is temporally interpolated to ten frames (TIM10), and LBP, HOG and HIGO descriptors are extracted and classified, on both SMIC and CASME II under leave-one-subject-out. Two findings emerge.

\textbf{Magnification helps, consistently}, ``for all three kinds of features on all four testing datasets''. On CASME II the effect is substantial (\autoref{tab:evm-effect}), and the best configuration of their full framework --- magnification, TIM10 and HIGO --- reaches 67.21\,\% on CASME II and 68.29\,\% on SMIC-HS under LOSO, competitive with or better than the contemporaneous state of the art.

\begin{table}[htbp]
	\centering
	\begin{tabular}{l c c c}
		\hline
		\textbf{Descriptor} & \textbf{No magnification} & \textbf{With magnification} & \textbf{$\Delta$} \\
		\hline
		LBP  & 55.87\,\% & 60.73\,\% & \textbf{+4.86} \\
		HOG  & 57.49\,\% & 63.97\,\% & \textbf{+6.48} \\
		HIGO & 57.09\,\% & 67.21\,\% & \textbf{+10.12} \\
		\hline
	\end{tabular}
	\caption[Effect of Eulerian motion magnification on CASME II]{Effect of Eulerian motion magnification on CASME II under leave-one-subject-out, after Li et al.~\cite{liX2018} (Table 6).}
	\label{tab:evm-effect}
\end{table}

\textbf{The relationship between $\alpha$ and accuracy is non-monotonic.} Li et al.~\cite{liX2018} describe the resulting curves as ``rainbow-shaped,'' with best performance ``generally achieved when the motion is magnified in the range of $[8, 16]$,'' and give a mechanistic account of both failure modes: ``magnification at lower levels might not be enough to reveal the ME motion progress; on the other hand, magnification at higher levels degrades the performance because too many artifacts are induced.'' Magnification is therefore not a monotone improvement to be maximised but a parameter with an interior optimum that must be selected.

Bai et al.~\cite{bai2021} provide the second quantitative demonstration using a fundamentally different downstream model --- a pre-trained VGGFace2 (ResNet-50) spatial encoder feeding a bidirectional LSTM. On SMIC they reach \textbf{75.76\,\% with phase-based magnification against a 60.60\,\% baseline} drawn from an earlier VGGFace2 + Bi-LSTM publication rather than from a magnification-off arm of their own, and the phase-based variant outperforms the amplitude-based one. Because the pipeline shares nothing with Li et al.'s hand-crafted descriptors, the two results together establish that the benefit is not an artefact of a particular feature family --- at least when the downstream representation is appearance-based.

\subsection{Magnification in apex-frame pipelines, and the $\alpha$ discrepancy}
\label{sec:evm-apex}

Two further papers use EVM in a different regime, and the contrast is informative. Li, Huang and Zhao~\cite{liY2018} address whether a micro-expression can be recognised from a \textit{single} apex frame, arguing for magnification specifically in that setting: the apex frame carries the most emotional information, but ``the intensity level of apex frame is not distinguished due to subtle change, especially for deep model that describes high-layer abstract information.'' They magnify the spotted apex frame with \textbf{$\alpha = 30$} and fine-tune a VGG-Face descriptor on the result; because a single frame per clip is too little data to train on, they exploit the high frame rate to augment, retaining the apex frame plus its two preceding and two following frames and so expanding the training set fivefold. They restate the trade-off Li et al.~\cite{liX2018} quantified --- ``bigger values of motion amplification level lead to larger scale of motion amplification, but also can cause bigger displacement and artifacts.'' Li, Huang and Zhao~\cite{liY2021} extend this into a joint local--global learning framework with frequency-domain apex detection, carrying the same setting forward --- magnification level ``set as 30 in our framework'' --- with the same $\pm$2-frame augmentation producing what they name the Extended Magnified ME (EMME) database.

\textbf{The $\alpha$ discrepancy is easy to misread as a contradiction.} Li et al.~\cite{liX2018} find the optimum at $\alpha \in [8, 16]$ and observe degradation beyond it; Li, Huang and Zhao~\cite{liY2018,liY2021} adopt $\alpha = 30$, at the extreme end of that same sweep. The regimes differ in a way that plausibly explains the gap. The $[8, 16]$ optimum is measured on \textbf{full magnified sequences} described by spatio-temporal descriptors, where inter-frame artefacts accumulate along the temporal axis and corrupt exactly the dynamic texture the descriptor encodes. The $\alpha = 30$ setting is applied to a \textbf{narrow five-frame window around the apex}, consumed by a deep face encoder as largely static appearance, where temporal artefacts have far less opportunity to express themselves --- and where the encoder needs a strong appearance signal to distinguish categories at all. The implication for any new system is that $\alpha$ cannot be transferred between pipelines: it is a parameter of the \textit{combination} of magnification and downstream representation, not of the corpus.

\subsection{Known limitations of magnification}
\label{sec:evm-limitations}

The review corpus is unusually candid about EVM's failure modes, and four are documented.

\textbf{It amplifies everything in band, not only expressions.} This is the most important limitation and the most easily overlooked. Li et al.~\cite{liX2018} decline to use magnification for micro-expression \textit{spotting} at all, ``because it magnifies unwanted motions (e.g. head movements) at the same time.'' EVM has no notion of which motion is of interest: a rigid head translation whose temporal frequency falls inside the pass-band is amplified exactly as a brow raise is. In a recognition pipeline this is mitigated by prior face alignment and cropping --- Li et al.~\cite{liX2018} magnify only after normalising and cropping to the eye-defined rectangle --- but mitigated, not eliminated.

\textbf{Artefacts grow with $\alpha$.} All four papers report this: Li et al.~\cite{liX2018} as the descending limb of the rainbow curve, Li, Huang and Zhao~\cite{liY2018,liY2021} as a design caution, and Bai et al.~\cite{bai2021} directly, as noise that forced them to narrow their temporal band.

\textbf{There is no principled procedure for choosing $\alpha$.} One value in the corpus is obtained by exhaustive sweep~\cite{liX2018}; the apex-frame papers assert theirs --- Li et al.~\cite{liY2018} state $\alpha = 30$ without support, and Li et al.~\cite{liY2021} carry the value forward by citing them. No paper offers a way to select $\alpha$ from properties of the data.

\textbf{Evaluation is confined to appearance-based downstream representations:} LBP, HOG and HIGO in Li et al.~\cite{liX2018}, VGGFace2 features in Bai et al.~\cite{bai2021}, VGG-Face in both Li, Huang and Zhao papers. Every one of these descriptors is computed on \textbf{magnified pixel intensities}. No study in the review corpus measures what magnification contributes when the representation handed to the classifier is itself an explicit motion field rather than an appearance signal.

\subsection{Implications for this research}
\label{sec:evm-implications}
% sec:evm-motion-hypothesis is referenced from Chapter 5 (chapter5.tex:236). The
% hypothesis it names now opens this subsection instead of having one of its own,
% so the label is kept here and resolves to this subsection.
\label{sec:evm-motion-hypothesis}

\autoref{tab:evm-settings} sets out the settings the review determines. Behind them lies a hypothesis the review corpus does not test, stated here because it frames the interpretation of the measurement. EVM amplifies the \textit{amplitude} of intensity variation in a temporal band, and an appearance descriptor computed on magnified frames sees a directly stronger signal. An optical flow field is a different object: flow estimates \textit{displacement}, and amplifying the underlying intensity change by $\alpha$ does not simply scale the recovered displacement by $\alpha$ --- it changes the conditioning of the estimation problem, potentially making faint motion recoverable where it previously was not, while introducing displacement artefacts of the kind all four reviewed papers warn about. Any residual global scaling of flow magnitude is further attenuated by the per-channel normalisation a learned pipeline applies before training. The plausible outcome is that magnification in front of a flow-and-strain representation retains the part of its benefit that consists of \textit{making otherwise-unrecoverable motion recoverable} while losing the part that consists of \textit{increasing contrast for an appearance descriptor} --- so its measured contribution should be smaller on average than the roughly +5 to +10 percentage points of \autoref{tab:evm-effect}, and concentrated in configurations whose downstream model can exploit finer spatial deformation. EVM is measured across six matched pairs of configurations differing in nothing but the tensor directory, which is what allows that prediction to be tested.

\begin{table}[htbp]
	\centering
	\footnotesize
	\begin{tabular}{@{}p{6.4cm} p{7.4cm}@{}}
		\hline
		\textbf{Finding from the review} & \textbf{Decision taken in this thesis} \\
		\hline
		Pass-band should match the phenomenon's duration, 1/5--1/25\,s $\equiv$ 5--25 Hz~\cite{bai2021}, though Bai et al. found that band noisy and narrowed it to 15--25 Hz & Temporal band-pass set to \textbf{5--25 Hz}, the unnarrowed duration rule; Bai et al.'s narrowing was not followed, and this is a declared divergence \\
		Accuracy against $\alpha$ is rainbow-shaped, optimum in \textbf{$[8, 16]$} for full-sequence pipelines~\cite{liX2018} & \textbf{$\alpha = 10$}, inside the reported optimum; the $\alpha = 30$ of the apex-frame papers is explicitly \textit{not} adopted, since this thesis consumes full sequences \\
		Amplitude-based magnification uses a Laplacian pyramid plus temporal band-pass~\cite{bai2021} & Amplitude-based EVM over a \textbf{4-level Laplacian pyramid}, temporal filtering by FFT band-pass \\
		Magnification amplifies head motion as well as expression~\cite{liX2018} & Magnification applied only after the clip has been trimmed onset-to-offset; the risk is acknowledged as a residual threat rather than claimed to be eliminated \\
		Magnification is a purely upstream, data-level operation & Stage 1 is run \textbf{twice}, producing a magnified and a non-magnified tensor set; the EVM switch selects a directory and leaves the network unchanged, making it a clean ablation factor \\
		Magnification is part of the proposed data pipeline, not of the network & EVM is treated as the \textbf{baseline condition} (\texttt{config\_4}), not as a component to be added; the non-magnified arm (\texttt{config\_1}) is the control that isolates it \\
		\hline
	\end{tabular}
	\caption[Review findings and the magnification settings they determine]{Review findings and the magnification settings they determine.}
	\label{tab:evm-settings}
\end{table}

\textbf{A defect found and repaired, which the review made detectable.} In this project's earliest evaluations, each magnification-enabled configuration produced results \textit{bit-identical} to its magnification-disabled counterpart --- identical accuracies, macro F1 scores and confusion matrices. Given that all four reviewed papers report clear and repeatable effects from magnification, an effect of exactly zero across six independent pairs is not a plausible null result but a symptom: both arms were reading the same tensor directory, so half of the ablation matrix was duplicated and the magnification hypothesis had never in fact been exercised. The defect was repaired before the reported run, and wiring the existing verification helper in as a start-up assertion is a declared item of further work.

\subsection{Limitations of the existing work, and how this study differs}
\label{sec:evm-gap}

What the reviewed literature does not establish is threefold. First, every evaluation in the corpus places magnification in front of an \textbf{appearance-based} representation, so the reported gains conflate two distinct mechanisms --- making faint motion recoverable, and increasing contrast for a texture or gradient descriptor --- and no experiment separates them. Second, the magnification factor is selected empirically in every case and transferred between studies without re-validation, despite the $\alpha$ discrepancy of \autoref{sec:evm-apex}. Third, and most consequentially, magnification is always evaluated as an \textit{addition to a fixed pipeline} --- a two-cell comparison, with and without --- so its interaction with the other components of a system is never measured.

This thesis differs on all three counts. It places magnification in front of an explicit \textbf{motion} representation --- dense optical flow and optical strain --- rather than an appearance representation, isolating the recoverable-motion mechanism from the appearance-contrast one. It selects $\alpha$ from the full-sequence optimum rather than the apex-frame value, with the choice justified rather than inherited. And it evaluates magnification as one factor in a \textbf{full factorial ablation}, so its effect is reported not as a single delta but as six matched-pair measurements whose spread across architectures is itself the finding.

What magnification is placed in front of therefore matters as much as its own parameters, and \autoref{sec:flow-review} turns to that motion representation.


\section{Motion Representation: Dense Optical Flow}
\label{sec:flow-review}

Dense optical flow supplies the horizontal and vertical displacement channels used in this thesis. The relevant literature concerns three choices: the estimator, the frame pairs it processes, and whether flow serves as a recognition feature or an alignment tool. These choices determine what information remains available to the learned model.

\subsection{Why motion rather than appearance}
\label{sec:why-flow}

Xu et al.~\cite{xu2017} motivate facial-dynamics descriptors by the difficulty of interpreting micro-expressions through texture alone. Flow describes apparent movement directly, whereas image intensity also records appearance and illumination. This makes displacement a useful representation of subtle facial activity, although it does not guarantee independence from identity or recording conditions: facial geometry, alignment and brightness changes still affect its estimation.

The benchmark evidence supports this choice. The MEGC 2019 organisers report that their three highest-ranked submissions used optical flow~\cite{see2019}. OFF-ApexNet~\cite{liong2019a} uses horizontal and vertical flow as neural-network inputs; STSTNet~\cite{liong2019b} adds strain; Zhao et al.~\cite{zhaoS2021} retain a sequence of flow fields. Their convergence motivates a motion-based input, but is not a controlled demonstration that it always outperforms appearance-based representations.

\subsection{Estimation assumptions}
\label{sec:brightness-constancy}

OFF-ApexNet~\cite{liong2019a} formulates flow using brightness constancy and small-motion assumptions. A local brightness constraint supplies one equation for two displacement components, so estimation also needs spatial regularisation or a local motion model. The output is a pair of fields, $u(x,y)$ and $v(x,y)$; Zhao et al.~\cite{zhaoS2021} additionally describe their magnitude, $\sqrt{u^2+v^2}$.

These assumptions matter more than repeating the derivation. Illumination changes can resemble movement, while smoothing can suppress a weak local contraction together with noise. Magnification can make subtle motion easier to estimate but can also amplify unwanted movement. Flow therefore transforms the measurement problem; it does not remove uncertainty in the source frames.

\subsection{Choice of estimator}
\label{sec:flow-estimator}

The reviewed methods use three estimators. Zhao et al.~\cite{zhaoS2021} adopt Farneb\"ack through its OpenCV implementation, citing its established availability. It estimates displacement through local polynomial approximations. OFF-ApexNet~\cite{liong2019a} instead uses total-variation regularisation with an L1 data term (TV-L1), citing preservation of motion discontinuities and robustness. Xu et al.~\cite{xu2017} use the method of Sun, Roth and Black. These algorithm attributions are taken from the reviewed papers; the original estimator publications are outside the supplied reference corpus.

STSTNet~\cite{liong2019b} specifies onset--apex optical flow but does not name its estimator. Its similarity to OFF-ApexNet is therefore insufficient to attribute TV-L1 to it. Likewise, the advantages given for TV-L1 should not be read as a measured recognition gain over Farneb\"ack on CASME II. The cited studies do not provide a controlled three-estimator comparison supporting such a ranking. This thesis adopts Farneb\"ack as a fixed preprocessing choice and leaves that comparison open.

\subsection{Onset--apex flow or a temporal sequence}
\label{sec:flow-what-between}

OFF-ApexNet~\cite{liong2019a} represents each clip by flow from onset to apex, where expression intensity is greatest, and resizes the two component maps to $28\times28$. STSTNet~\cite{liong2019b} uses the same frame-pair construction. This is an economical representation of the displacement from neutral to peak expression, suitable for a small training set. It omits the intermediate trajectory and the relaxation after the apex.

Zhao et al.~\cite{zhaoS2021} instead construct an eleven-frame sequence using onset, apex and offset annotations, then calculate ten adjacent-pair flow fields. This retains changes in motion across the expression. The distinction is between one displacement summary and a sequence whose order a temporal model can exploit; it is not simply a difference in input size.

The present study follows the sequence approach but uses uniform sampling of existing frames rather than Zhao et al.'s key-frame construction. This preserves a temporal input for the transformer ablation. It does not, by itself, measure the recognition cost of replacing a sequence with an onset--apex pair: that representation comparison is not an experimental arm.

\subsection{Flow for fine alignment}
\label{sec:flow-alignment}

Xu et al.~\cite{xu2017} also use flow to refine alignment. Their pipeline first locates landmarks and coarsely aligns the face, then performs fine in-sequence alignment using pixel-level movement. The motivation is specific to micro-expressions: a small registration error can be comparable to the facial displacement being recognised.

This use differs from supplying flow to a classifier. An uncorrected displacement field contains residual head movement as well as expression. Magnifying the frames beforehand can increase both. The present pipeline performs no flow-based fine alignment, so corpus preprocessing must not be treated as proof that this source of error has been eliminated.

\subsection{Limitations of the representation}
\label{sec:flow-limitations}

Flow remains sensitive to brightness-constancy violations, alignment errors and estimator settings. Apex-based systems additionally depend on selecting an informative frame. OFF-ApexNet~\cite{liong2019a} attributes weaker performance on SMIC to imprecise apex selection and cites an earlier apex-spotting study reporting an average discrepancy of thirteen frames. That reported average is neither a maximum error nor a measurement made by this thesis.

Using an onset--offset sequence avoids dependence on one selected apex, but still relies on clip boundaries and temporal sampling. Repeated sampled frames can yield zero-flow pairs, and large sampling intervals can make the small-motion approximation less appropriate. These limitations constrain the representation before any network component is compared.

\subsection{Implications for this research}
\label{sec:flow-implications}

\begin{table}[htbp]
    \centering
    \small
    \begin{tabular}{@{}p{0.29\linewidth} p{0.65\linewidth}@{}}
        \hline
        \textbf{Review finding} & \textbf{Consequence for this study} \\
        \hline
        Flow is a recognition input~\cite{liong2019a,see2019}
        & Horizontal and vertical displacement form the first two channels; the classifier receives motion tensors rather than raw frames. \\
        Sequence flow preserves temporal variation~\cite{zhaoS2021}
        & Adjacent-pair fields are retained across the sampled clip, allowing the temporal encoder to operate on an ordered sequence. \\
        Estimator choices differ~\cite{liong2019a,zhaoS2021,xu2017}
        & Farneb\"ack remains fixed; a matched comparison with TV-L1 is an untested extension. \\
        Fine alignment matters~\cite{xu2017}
        & No flow-based fine alignment is performed. It remains a priority for improving preprocessing. \\
        \hline
    \end{tabular}
    \caption[Optical-flow evidence and study decisions]{Optical-flow evidence and the decisions it motivates. Fixed preprocessing choices constrain the interpretation of the architectural ablation.}
    \label{tab:flow-decisions}
\end{table}

The central interpretive point is that optical flow is an \textit{analytical spatio-temporal feature extractor}: it estimates motion from paired frames before the network is trained. A learned spatial stem therefore acts on a representation that already makes displacement explicit. With only 156 clips, learning additional spatial filters may offer limited benefit or duplicate structure already exposed by preprocessing. This is a hypothesis about the stem's marginal value, not proof that learned spatial processing is unnecessary.

A temporal encoder addresses a different task: modelling relations among the successive displacement fields. Pairwise flow estimation does not itself learn those sequence-level relations. Keeping these roles distinct is necessary when attributing performance to the stem and transformer.

\subsection{Research gap and scope}
\label{sec:flow-gap}

The relevant unresolved question is what learned spatial and temporal processing contributes \textit{after} motion has been extracted analytically. The thesis addresses this through matched architectural comparisons under complete leave-one-subject-out (LOSO) evaluation, retaining the same flow construction while switching the learned components. Magnification is varied separately through the paired preprocessing conditions.

The design neither ranks optical-flow estimators nor isolates alignment quality or onset--apex versus sequence representations. Those remain limitations rather than gaps claimed as solved. The resulting architectural conclusions apply to the implemented motion representation and evaluation setting; they do not establish a universal ordering of network families.


\section{Deformation Representation: The Optical Strain Tensor}
\label{sec:strain-review}

Optical strain supplies the third input channel. Its literature establishes a deformation descriptor derived from flow, two ways to use that descriptor, and limitations relevant to interpreting the architectural ablation. Shreve et al.~\cite{shreve2011}, the Liong strain series~\cite{liong2014a,liong2014b,liong2016}, and STSTNet~\cite{liong2019b} provide the principal evidence.

\subsection{What strain contributes}
\label{sec:strain-what-adds}

Flow describes displacement; strain describes its local spatial variation. A spatially uniform translation has zero displacement gradient, whereas neighbouring skin regions moving differently produce deformation. Shreve et al.~\cite{shreve2011} use this distinction to detect facial activity and describe robustness to moderate head translation. It does not imply immunity to arbitrary head movement or registration error.

Their illumination- and make-up-robustness statements are attributed to earlier studies rather than measured anew. They should therefore be read as reported motivation, not as evidence that the present pipeline is robust under those conditions. Strain can suppress a uniform displacement while still inheriting errors from the estimated flow.

\subsection{Tensor components and magnitude}
\label{sec:strain-tensor}

The relevant representation is the symmetric displacement-gradient tensor, with normal components $\varepsilon_{xx}$ and $\varepsilon_{yy}$ and shear components $\varepsilon_{xy}=\varepsilon_{yx}$. Shreve et al.~\cite{shreve2011} approximate spatial derivatives by central differences with offsets of approximately two to three pixels; Liong et al.~\cite{liong2014a} use one-pixel offsets. This discretisation is a modelling choice because it sets the scale at which local variation is measured.

Liong et al.~\cite{liong2014a,liong2014b,liong2016} use the four-component magnitude
\[
    \varepsilon_{\mathrm{F}}
    =\sqrt{\varepsilon_{xx}^{2}+\varepsilon_{yy}^{2}
        +\varepsilon_{xy}^{2}+\varepsilon_{yx}^{2}}
    =\sqrt{\varepsilon_{xx}^{2}+\varepsilon_{yy}^{2}
        +2\varepsilon_{xy}^{2}}.
\]
The scalar describes deformation intensity and discards its sign and direction. Supplying it together with $u$ and $v$ retains the displacement fields from which those spatial relations were calculated.

\subsection{The implemented magnitude and its divergence}
\label{sec:strain-divergence}

The implementation instead computes
\[
    \varepsilon_{\mathrm{os}}
    =\sqrt{\varepsilon_{xx}^{2}+\varepsilon_{yy}^{2}
        +\varepsilon_{xy}^{2}},
    \qquad
    \varepsilon_{xy}=\tfrac{1}{2}
        \left(\frac{\partial u}{\partial y}
             +\frac{\partial v}{\partial x}\right).
\]
It counts shear once rather than twice. NumPy's \texttt{gradient} supplies spatial differences at one-pixel spacing, using central differences in the array interior. The spacing follows Liong et al.~\cite{liong2014a}, but the scalar reduction differs from their convention.

STSTNet~\cite{liong2019b} cannot establish a clear precedent for the three-term form. Its printed magnitude equation repeats $\partial u$ in the mixed derivative where the symmetric tensor calls for $\partial v/\partial x$ alongside $\partial u/\partial y$. It also places the factor $\tfrac{1}{2}$ \emph{outside} the squared derivative sum. Correcting the repeated derivative while retaining that placement gives
\[
    \tfrac{1}{2}
       \left(\frac{\partial u}{\partial y}
            +\frac{\partial v}{\partial x}\right)^2
       =2\varepsilon_{xy}^{2},
\]
which agrees with the four-term convention. Moving the factor inside the square would instead give the implemented three-term form, but that is a further change to the printed expression. The intended implementation cannot be settled from this equation alone; the thesis therefore declares a departure from the unambiguous four-term convention rather than claiming to reproduce STSTNet's magnitude.

The four-term form weights pure shear by $\sqrt{2}$ relative to the implemented form. Independent min--max normalisation removes a uniform scale factor, but does not generally remove this difference when normal and shear contributions vary across pixels and time. Neither convention was compared experimentally here, so their practical equivalence is not assumed.

\subsection{Strain as a feature, a weight, or an input channel}
\label{sec:strain-feature-weight}

\textbf{Optical strain features (OSF)} use a temporally pooled strain map as a descriptor~\cite{liong2014a,liong2014b,liong2016}. Averaging accounts for sequence length: the early method already divides its accumulated map by the number of frame pairs~\cite{liong2014a}. The maps are filtered and normalised before classification. The early work explicitly applies Wiener and Gaussian filters to suppress noise~\cite{liong2014a}.

\textbf{Optical strain weighted features (OSW)} use pooled strain to weight the regional XY-plane histograms of local binary patterns on three orthogonal planes (LBP-TOP), emphasising active facial regions~\cite{liong2014b}. Liong et al.~\cite{liong2016} perform spatial mean pooling followed by temporal mean pooling and concatenate OSF and OSW for a combined descriptor. Strain then serves both as a feature and as a spatial weighting signal.

\textbf{STSTNet} supplies strain directly alongside horizontal and vertical onset--apex flow~\cite{liong2019b}. This is the precedent for the thesis's three-channel input structure, independently of the magnitude ambiguity. The present study retains all three channels across successive frame pairs rather than reducing each clip to one map.

\subsection{Recognition evidence and protocol limits}
\label{sec:strain-evidence}

\begin{table}[htbp]
    \centering
    \small
    \begin{tabular}{@{}lccc@{}}
        \hline
        \textbf{Method} & \multicolumn{2}{c}{\textbf{SMIC accuracy (\%)}} & \textbf{CASME II (\%)} \\
        & Micro-average & Macro-average & Accuracy \\
        \hline
        Baseline & 45.73 & 47.76 & 61.94 \\
        OSF & 41.46 & 46.00 & 51.01 \\
        OSW & 49.39 & 50.71 & 61.94 \\
        OSF + OSW & 52.44 & 58.15 & 63.16 \\
        OFF + OFW & 40.24 & 41.94 & 55.87 \\
        \hline
    \end{tabular}
    \caption[Recognition results for strain-based descriptors]{Recognition results reproduced from Liong et al.~\cite{liong2016}, Table 8. SMIC uses LOSO; CASME II uses leave-one-video-out. The summary combines block settings, so equal entries are not a matched-setting comparison. Macro-average here denotes subject-averaged accuracy, not macro F1.}
    \label{tab:strain-performance}
\end{table}

The CASME II evidence depends on how strain is used and on the comparator. OSF alone falls below the texture baseline. The equal baseline and OSW entries in \autoref{tab:strain-performance} reflect different block settings: under matched settings OSW improves accuracy by 0.81 percentage points. The combined OSF + OSW descriptor improves on the texture baseline by 0.41--1.22 points across block settings~\cite{liong2016}. These accuracy gains are not directly comparable with pooled macro F1, and leave-one-video-out does not enforce subject-disjoint testing.

Crucially, Liong et al.~\cite{liong2016} also test optical-flow features and weights (OFF + OFW) in the same descriptor framework. Their CASME II accuracy is 55.87\,\%, against 63.16\,\% for OSF + OSW, a 7.29-point difference. This is a matched flow-based control, so it would be incorrect to claim that the literature never isolates strain against flow. It concerns hand-crafted descriptors; it does not establish the gain from adding a strain channel to the learned sequence model used here.

\subsection{Limitations}
\label{sec:strain-limitations}

Spatial differentiation can amplify flow-estimation noise. The filtering used by Liong et al.~\cite{liong2014a,liong2016} is therefore relevant to a pipeline without a dedicated strain post-filter. A deterministic transformation of $u$ and $v$ adds no independent observation: its potential value lies in making deformation directly accessible to the classifier. Whether this representation helps must be measured rather than inferred from its construction.

The reviewed recognition methods also simplify time. OSF and OSW pool their sequences; STSTNet uses one onset--apex pair. Shreve et al.~\cite{shreve2011} retain temporal strain activity for spotting, which localises expressions rather than classifying their emotion. The temporal-collapse limitation thus concerns these recognition representations, not every use of strain in the literature.

\subsection{Implications for this research}
\label{sec:strain-implications}

\begin{table}[htbp]
    \centering
    \small
    \begin{tabular}{@{}p{0.29\linewidth} p{0.65\linewidth}@{}}
        \hline
        \textbf{Design choice} & \textbf{Interpretation or unresolved test} \\
        \hline
        Strain accompanies $u$ and $v$
        & Follows STSTNet's channel structure~\cite{liong2019b}; does not establish that its ambiguous magnitude equation was reproduced. \\
        Three-term magnitude
        & Differs from the four-term convention~\cite{liong2014a,liong2014b,liong2016}. Relative shear weighting remains an untested choice. \\
        Per-pair strain sequence
        & Preserves deformation over time instead of pooling it into one descriptor. Each channel is independently min--max normalised across the sequence. \\
        No dedicated strain post-filter
        & Departs from the explicit filtering of Liong et al.~\cite{liong2014a}. Farneb\"ack's internal smoothing is not a separately evaluated strain filter. \\
        Strain held constant
        & Present in all twelve evaluated configurations; its marginal contribution is not measured. \\
        \hline
    \end{tabular}
    \caption[Strain representation and unresolved choices]{The implemented strain representation and the limits it places on interpretation.}
    \label{tab:strain-decisions}
\end{table}

A strain-on/strain-off comparison is the natural next ablation. Neither published gains from a different descriptor nor the performance of the complete system substitutes for that measurement. Comparing magnitude conventions and adding a dedicated post-filter are separate follow-ups.

Strain also extends the analytical-feature argument: the input contains a \textit{hand-specified differential operator}, exposing local deformation without learning that operator from the training clips. A spatial stem may still learn useful combinations of those features, but it is operating after both motion estimation and deformation extraction. This motivates testing its incremental contribution; it does not establish the cause of any observed architectural effect.

\subsection{Research gap and scope}
\label{sec:strain-gap}

The study preserves the temporal evolution of strain for learned sequence processing. Its contribution is an architectural comparison on that representation, not an isolated demonstration of strain's benefit. The independent value of the third channel, the choice of shear weighting and the omitted post-filter remain unresolved. These boundaries preserve the distinction between evidence for a component, a reason to include it, and a result that the experimental design can actually establish.


\section{Temporal Normalisation: The Temporal Interpolation Model}
\label{sec:tim-review}

Flow and strain fix \textit{what} is computed from each frame pair; temporal normalisation fixes \textit{which} frame pairs, and how many. The Temporal Interpolation Model (TIM) is credited in this corpus to Zhou et al. by Li et al.~\cite{liX2013,liX2018} and to Pfister et al. by Lu et al.~\cite{lu2015}, who applied it to micro-expressions outside the review corpus. It is described here only through the reviewed papers that implement and evaluate it.

\subsection{The problem: sequences of unequal length}
\label{sec:tim-problem}

Micro-expressions vary in duration, and the three-class CASME II working set spans roughly a four-fold range of sequence length (\autoref{sec:corpus-threats}h). Three difficulties follow. Any batched network requires a fixed input shape, and the corpus supplies no constant. A descriptor with a temporal radius cannot be applied to a sequence shorter than that radius, so the shortest clips constrain the descriptor for the whole dataset: Li et al.~\cite{liX2018} note that at 25~fps some micro-expressions last four to five frames, restricting LBP-TOP to radius $r = 1$. And duration is itself a nuisance variable --- Lu et al.~\cite{lu2015} normalise to remove ``temporal fluctuations of micro-expressions, which may introduce noise irrelevant to classification''.

\subsection{What TIM is}
\label{sec:tim-what}

Lu et al.~\cite{lu2015} give the corpus's most compact description: ``TIM is a manifold-based interpolation method, which builds a low-dimensional embedding of an image sequence, and then interpolates a curve in the low-dimensional space. The interpolated frames are mapped back to the original high-dimensional space.''

The property that matters downstream is that TIM \textbf{synthesises frames that were never recorded}: Lu et al.~\cite{lu2015} note that Pfister et al. proposed it ``to generate sufficient frames from an image sequence'', that is, as an upsampler rather than only a length-standardiser.

\subsection{TIM in practice: the convergence on ten frames}
\label{sec:tim-ten-frames}

TIM appears throughout the corpus as an unremarked preprocessing step, almost always at ten frames: SMIC's baseline~\cite{liX2013} normalises to ten before LBP-TOP; Li et al.~\cite{liX2018} apply TIM10 after magnification; and Lu et al.~\cite{lu2015} apply TIM before their Delaunay-based coding model, following Pfister et al.'s protocol so as to compare against TIM10 baselines. CAS(ME)$^2$'s baseline~\cite{qu2016} is the exception, normalising to \textbf{120 frames by linear interpolation} and evaluating leave-one-video-out.

One study tests the setting on its own terms. Li et al.~\cite{liX2018} sweep interpolation length from 10 to 80 frames on all three SMIC variants under leave-one-subject-out, holding the descriptor fixed and skipping magnification so only length varies, and report that ``the best performance for all three datasets is achieved with interpolation to 10 frames'', because ``in longer sequences, changes along the time-dimension are diluted''.

\subsection{How many frames? The remaining systematic comparisons}
\label{sec:tim-how-many}

Ben et al.~\cite{ben2021} run a second controlled experiment and reach the opposite conclusion. Holding alignment, descriptor and classifier fixed --- JCFDA alignment, LBP-TOP, and an SVM with an RBF kernel --- they interpolate sequences to 10, 20, 30 and so on up to 110 frames by two algorithms, \textbf{Newton interpolation} and \textbf{TIM}, recording the best recognition rate at each length on MMEW.

Newton peaks at \textbf{33.3\,\%} at 60 frames; TIM peaks at \textbf{38.9\,\%} ``when the micro-expression sequence is respectively interpolated to 30 or 60 frames'', and Ben et al. adopt TIM at 60 accordingly. Accuracy is non-monotonic in length: rates ``increase initially and then decrease towards the end''.

Two caveats limit the transfer. The experiment uses MMEW with LBP-TOP, not CASME II and not a learned model; and it runs largely in the \textit{upsampling} regime, since Ben et al.~\cite{ben2021} interpolate ``to the maximal value of 110 frames'', so the longer settings add frames for most clips rather than discarding them. This thesis operates in the opposite regime (\autoref{sec:tim-implications}), and an optimum measured while synthesising frames need not hold when discarding them.

A third comparison is the only multi-corpus one. Xu et al.~\cite{xu2017} interpolate every sequence to 10, 15 and 20 frames for both their facial dynamics map and LBP-TOP, ``on all four datasets'' --- CASME, CASME II and SMIC's high-speed and visual subsets. They also depart from the convention in a way the others do not, using \textbf{linear} interpolation ``instead of Temporal Interpolation Model because in a high frame-rate scenario, linear interpolation suffices to characterize the motion pattern with less error''.

\subsection{Alternatives to uniform interpolation}
\label{sec:tim-keyframe}

Zhao et al.~\cite{zhaoS2021} do not interpolate uniformly. They build an eleven-frame key-frame sequence from the annotated onset, apex and offset by generating eight transition frames allocated between the onset--apex and apex--offset intervals \textbf{in proportion to each sub-interval's duration}, so sampling density follows the expression's own phase structure and the apex is emphasised. Optical flow between adjacent key-frames then gives a ten-frame flow sequence. OFF-ApexNet~\cite{liong2019a} and STSTNet~\cite{liong2019b} take the limit, keeping the onset--apex pair alone (\autoref{sec:flow-what-between}).

The corpus therefore contains temporal representations of length 2, 10, 11, 15, 20, 30, 60, 110 and 120 frames, chosen variously by convention, by sweep and by architecture. There is no consensus, and only Xu et al.'s~\cite{xu2017} sweep spans more than one corpus.

\subsection{Limitations}
\label{sec:tim-limitations}

\textbf{Ten frames is a convention, not a finding.} The most-used setting traces to one baseline paper's choice, and the three studies that test it do not agree: ten frames for Li et al.~\cite{liX2018}, 30 to 60 for Ben et al.~\cite{ben2021}, and a 10--20 frame range for Xu et al.~\cite{xu2017} (\autoref{sec:tim-how-many}). Interpolation length interacts with the temporal receptive field of whatever consumes it, and in all three that consumer is a hand-crafted descriptor.

\textbf{Frame synthesis and motion analysis are in tension.} TIM generates frames that were not recorded. Where the downstream representation is appearance-based this is benign; where it is \textit{motion}, synthesised frames yield synthesised displacements. One reviewed paper recognises this: Xu et al.~\cite{xu2017} choose linear interpolation over TIM precisely on the grounds that it ``suffices to characterize the motion pattern with less error'' at a high frame rate. But they assert the point rather than measure it --- TIM appears once in their paper, in that sentence --- and no reviewed study compares the two interpolators against a motion representation. The rest of the corpus does not raise the question, because those using TIM extract appearance descriptors, while those extracting flow either use two frames or, as Zhao et al.~\cite{zhaoS2021} do, generate transition frames without discussing the consequence.

\textbf{Downsampling forfeits part of the corpus's advantage.} \autoref{sec:corpus} establishes why CASME II was recorded at 200~fps rather than its predecessor's 60~fps; reducing a clip to a few tens of frames gives part of that back. Only Xu et al.~\cite{xu2017} tie the interpolation choice to the acquisition rate at all, and then only to justify the simpler interpolator. The rest apply TIM10 to 100~fps SMIC data as readily as to 200~fps CASME II data, without weighing what the higher rate bought.

\subsection{Implications for this research}
\label{sec:tim-implications}

\begin{table}[htbp]
	\centering
	\footnotesize
	\begin{tabular}{@{}p{5.6cm} p{8.0cm}@{}}
		\hline
		\textbf{Finding from the review} & \textbf{Decision taken in this thesis} \\
		\hline
		Fixed-length input is required, and unequal duration is nuisance variance~\cite{lu2015} & Every clip is normalised to a fixed temporal length before training \\
		TIM beats Newton interpolation by 5.6 points, with a joint optimum at 30 or 60 frames~\cite{ben2021} & \textbf{L = 33 frames sampled, giving T = 32 flow pairs} --- the lower of the two reported optima rather than the field-conventional 10 \\
		Sequence length has an interior optimum found only by sweep~\cite{ben2021} & T = 32 is fixed throughout and is \textbf{not} swept; a sweep is declared as further work \\
		Onset and offset are annotated in CASME II (\autoref{sec:casme2-construction}) & Sampling is bounded by the annotated onset and offset, so the window spans the expression rather than surrounding neutral frames. This holds on the image-folder path used for CASME II; the loader's \texttt{.avi} branch samples across a whole file and is not exercised here \\
		TIM synthesises frames, and flow between synthesised frames is synthesised motion (\autoref{sec:tim-limitations}) & \textbf{No frames are synthesised} --- see the declaration below \\
		\hline
	\end{tabular}
	\caption[Review findings and the temporal-normalisation decisions they determine]{Review findings and the temporal-normalisation decisions they determine.}
	\label{tab:tim-decisions}
\end{table}

\textbf{A declaration about what this pipeline implements.} The module performing this step is named for the Temporal Interpolation Model, but what it computes is \textbf{uniform index sampling} of existing frames (\autoref{sec:three-channel-tensor}). No embedding is built, no curve is interpolated, no frame is synthesised. That is a materially different operation from the TIM of \autoref{sec:tim-what}, and it is named accurately here rather than in the pipeline's own terms. It is also defensible: because the downstream representation is flow between adjacent sampled frames, interpolation would manufacture displacement fields for motion that never occurred, whereas subsampling discards temporal resolution without fabricating it. What is given up is TIM's ability to \textit{lengthen} short clips.

Three quantitative consequences follow.

\textit{Effective temporal resolution falls by a clip-dependent factor.} Sampling 33 frames from a clip of $N$ frames recorded at 200~fps gives an effective rate of $6600/N$ fps: \textbf{100 fps} at the median ($N = 66$), which is SMIC's recording rate and half what CASME II was built to provide, and roughly \textbf{52 fps} for the longest ($N = 126$), below the original CASME database's 60~fps~\cite{yan2014}. Eighty-one of the 156 clips are downsampled by a factor of two or more.

\textit{Magnification is applied after subsampling, which shifts the realised pass-band.} Magnification runs on the 33 sampled frames at the nominal 200~fps (\autoref{sec:evm-departures}), but those frames span the clip's full duration, so the true inter-frame interval is $N/(200 \times 32)$ seconds rather than $1/200$. At the median a band specified as 5--25 Hz therefore corresponds to roughly 2.5--12.5 Hz physically, and the discrepancy grows with clip length. The realised band is thus neither the 5--25 Hz \autoref{sec:evm-implications} derives from the duration rule of Bai et al.~\cite{bai2021} nor constant across clips. This is a declared divergence; magnifying before subsampling is declared as further work.

\textit{Three clips are shorter than the sampling window.} \textbf{Three of the 156 clips} contain 31 frames against a target of 33, so sampling repeats indices and yields near-zero flow at those positions.

\subsection{Research gap and scope}
\label{sec:tim-gap}

The unresolved question is what temporal normalisation costs a representation made of motion rather than appearance. Three studies sweep interpolation length and give three answers, none on CASME II and none with a learned temporal model; and although Xu et al.~\cite{xu2017} raise both the synthesis-versus-motion tension and the acquisition-rate trade-off, they settle neither by measurement (\autoref{sec:tim-limitations}). This thesis takes the position that follows from its own representation --- sampling only recorded frames --- and quantifies what that costs in effective frame rate (\autoref{sec:tim-implications}), which no reviewed paper reports.

The design does not sweep T, nor compare interpolators, nor separate the effect of sequence length from the effect of the sampling regime. Those remain limitations rather than gaps claimed as solved, and the resulting conclusions apply to the implemented sampling scheme at T = 32.


\section{The Learned Spatial Backbone: Shallow 3D Convolutional Networks}
\label{sec:cnn-review}

This section covers the convolutional backbone, switched on and off as Variable C of this thesis's ablation, through four review-corpus papers.

\subsection{The case for, and against, a learned spatial stem}
\label{sec:cnn-case}

By 2019 the field had converged on optical flow as the input representation (\autoref{sec:why-flow}); the open question was what should consume it. A learned backbone promises what hand-crafted descriptors cannot give~\cite{liong2019b}, but runs into the small-data constraint (\autoref{sec:capacity}, \autoref{sec:corpus-threats}). Xia et al.~\cite{xia2020b} describe a related failure mode from a different cause --- ``the important subtle dynamics are prone to disappearing in the domain shift \ldots especially for deep models'' --- where the mechanism is the domain shift induced by mixing corpora rather than sample count alone, and where they report deeper, higher-resolution models doing better on the individual-database task. Applying their shrinking argument to a single small corpus therefore extends it beyond the setting they establish.

\subsection{What ``3D CNN'' means in this literature}
\label{sec:cnn-meaning}

The term covers three materially different architectures.

\textbf{A convolution over spatial and modality axes, with no temporal axis.} STSTNet~\cite{liong2019b}, the canonical shallow model in this field, calls itself a ``Shallow Triple Stream Three-dimensional CNN''. Its input is a single \textbf{28 $\times$ 28 $\times$ 3} tensor from the onset and apex frames, whose third dimension holds horizontal flow, vertical flow and optical strain --- not time. Three parallel streams, each a single convolutional layer with 3, 5 and 8 kernels respectively, supplement ``the small scale input data \ldots to avoid the problem of underfitting''. \textbf{There is no time dimension in this network}; the clip was collapsed to a single frame pair upstream (\autoref{sec:flow-what-between}).

\textbf{A convolution genuinely extended across time.} STRCN~\cite{xia2020a} extends convolutional connectivity ``across temporal domain'' with recurrent convolutional layers, modelling appearance and geometry separately.

\textbf{A convolution over a short flow sequence.} Zhao et al.~\cite{zhaoS2021} apply a C3D-style network to the ten-frame flow sequence built from their eleven key-frames (\autoref{sec:tim-keyframe}), so the temporal axis is present but short.

\subsection{The central finding: shrink the model \textit{and} the input}
\label{sec:cnn-shrink}

Xia et al.~\cite{xia2020b} treat depth and input resolution as controlled variables. They build four models of increasing depth --- from a single convolutional layer (Model 1) to one convolutional plus three recurrent convolutional layers (Model 4) --- and evaluate each at nine input resolutions, \textbf{20, 40, 60, 80, 100, 150, 200, 250 and 300 pixels}, on optical flow with UAR as the metric. Two results follow.

\textbf{Deeper is worse.} Models 3 and 4 ``have worse recognition performances compared to shallower \ldots models'', even though all four are far shallower than ResNet-101 or DenseNet-169.

\textbf{Higher resolution is worse, with a threshold.} Performance ``will be degraded with larger input resolutions (Model 2, Model 3 and Model 4)'' --- the two-layer model they go on to adopt included --- and for the deeper models ``begins to decrease dramatically when the resolution is bigger than \textbf{100 $\times$ 100}''. Shallow models are nonetheless ``basically robust to the change of input resolution'', improving toward lower resolution above a floor of about 40 $\times$ 40.

STSTNet takes the same argument to its limit and remains competitive on the CASME II subset (\autoref{sec:megc-metrics}): Liong et al.~\cite{liong2019b} report \textbf{0.00167 M parameters across 2 layers at 28 $\times$ 28 $\times$ 3 input}, against OFF-ApexNet's 2.77 M across 5 layers and GoogLeNet's 7 M across 22.

\subsection{Coping with limited and imbalanced training data}
\label{sec:cnn-limited-data}

Two strategies in the corpus have counterparts in this thesis's training configuration. Xia et al.~\cite{xia2020a} pair temporal augmentation with a balanced loss (\autoref{sec:imbalance-balanced-loss}, \autoref{sec:imbalance-double-correction}); a two-stage few-shot scheme~\cite{zhaoS2021} is the one route not taken. Xia et al.~\cite{xia2020b} take a third --- adding capability without capacity --- by developing ``three parameter-free modules \ldots integrate with RCN without increasing any learnable parameters'', so representational power grows while the parameter count, and therefore the overfitting risk, does not.

\subsection{Limitations}
\label{sec:cnn-limitations}

\textbf{The terminology obscures what is being compared.} Papers titled as three-dimensional CNNs include networks with no temporal axis at all, so ``3D-CNN'' results across the corpus are not like with like.

\textbf{Input resolution is almost never controlled.} Xia et al.~\cite{xia2020b} are the exception; elsewhere resolution is an implementation detail --- 28 $\times$ 28 in STSTNet and OFF-ApexNet --- with no indication of sensitivity, despite the substantial sensitivity Xia et al. measure.

\textbf{The evidence base is composite-database.} Xia et al.'s shrinking study~\cite{xia2020b} and STSTNet~\cite{liong2019b} are developed for the MEGC composite setting, which supplies substantially more training data per fold than a single-corpus study (\autoref{sec:megc-metrics}); only the earlier recurrent-convolutional work~\cite{xia2020a} is evaluated per database. Conclusions about affordable capacity therefore understate the constraint here.

\textbf{No study isolates the backbone itself.} The corpus is not short of matched comparisons: it sets descriptor-only methods against learned ones on the same folds, and Xia et al.~\cite{xia2020b} vary individual architectural modules against a common baseline with every other parameter held constant. But those comparisons all \textit{add} to a backbone that stays in place. None removes the learned spatial stem entirely and passes the motion representation straight to a classifier or temporal model, so what the stem itself contributes, set against what it costs, is nowhere measured.

\subsection{Implications for this research}
\label{sec:cnn-implications}

The component ablated as Variable C is a \textbf{three-stream shallow convolutional backbone}: one unshared two-layer stream per input channel --- flow-$u$, flow-$v$, optical strain.

\begin{table}[htbp]
	\centering
	\footnotesize
	\begin{tabular}{@{}p{5.6cm} p{8.0cm}@{}}
		\hline
		\textbf{Finding from the review} & \textbf{Decision taken in this thesis} \\
		\hline
		Shallow beats deep on small micro-expression data~\cite{xia2020b}; STSTNet uses 2 layers and $\sim$1,670 parameters~\cite{liong2019b} & Two convolutional layers per stream, \textbf{14,544 parameters} in the backbone (15,027 including the classifier head) --- shallow by the corpus's standards, though roughly nine times STSTNet's \\
		Triple-stream design supplements small-scale input~\cite{liong2019b} & Three unshared streams, one per input modality. STSTNet's streams differ in \textit{filter count} over the same input; these differ in \textit{input channel} \\
		Parameter-free modules add capability without capacity~\cite{xia2020b} & Motivates the parameter-free attention evaluated as Variable B (\autoref{sec:simam-review}) \\
		Limited and imbalanced data require augmentation plus loss correction~\cite{xia2020a} & Balanced sampling with focal loss; loss class-weighting deliberately disabled after single-class collapse (\autoref{sec:corpus-implications}) \\
		\textbf{Performance degrades above 100 $\times$ 100 input, dramatically for deeper models~\cite{xia2020b}} & \textbf{Not followed.} The backbone is fed at \textbf{224 $\times$ 224}, more than twice that threshold, though the same resolution as Bai et al.~\cite{bai2021} and Li, Huang and Zhao~\cite{liY2021} use --- see below \\
		\hline
	\end{tabular}
	\caption[Review findings and the backbone decisions they determine]{Review findings and the backbone decisions they determine.}
	\label{tab:cnn-decisions}
\end{table}

\textbf{A declaration about what this component is.} Every kernel in this backbone is spatial only: no filter spans two time steps, so despite the \texttt{Conv3d} implementation it is a \textbf{purely spatial} feature extractor (\autoref{sec:conv3d}). The ablation therefore does not test spatio-temporal convolution; it tests a \textbf{learned spatial stem}, and whatever it measures for Variable C is evidence about \textit{this} stem --- spatial-only, at 224 $\times$ 224, on flow and strain input --- and \textbf{not} evidence that 3D convolution over time is unhelpful for micro-expression recognition, which this study never evaluated.

Two observations mitigate that, and one aggravates it. In the field's own terms the choice is defensible: STSTNet, one of the two strongest shallow baselines on the CASME II subset, likewise performs no temporal convolution (\autoref{sec:cnn-meaning}). The stem also competes against a representation that already performs spatio-temporal extraction analytically (\autoref{sec:flow-implications}, \autoref{sec:strain-implications}).

Against both, \textbf{input resolution is the clearest departure from Xia et al.'s guidance}. It is not a departure from the corpus as a whole: Bai et al.~\cite{bai2021} resize frames to 224 $\times$ 224, and without alignment or cropping, as here; Li, Huang and Zhao~\cite{liY2021} resize faces to the same; and Xia et al.'s own earlier work~\cite{xia2020a} runs one stream at 300 $\times$ 245. What 224 $\times$ 224 departs from is the specific threshold of \autoref{sec:cnn-shrink}, which is a composite-database result and which bites hardest on models deeper than this one. Since this stem is shallow, the reviewed evidence predicts a smaller resolution effect than the deep-model figures suggest, so re-running the backbone arm at a lower resolution is worth doing because it is cheap and would settle the question, not because the corpus implies a large gain. The choice is also the expensive one: a 224 $\times$ 224 map carries sixty-four times the spatial elements of STSTNet's 28 $\times$ 28, and the backbone convolves it at every one of the 32 time steps. Whatever this study measures for the stem is therefore a resolution result as much as an architecture result, and the re-run is the first item of declared further work.

\subsection{Limitations of the existing work, and how this study differs}
\label{sec:cnn-gap}

The four gaps of \autoref{sec:cnn-limitations} define what the reviewed literature does not establish.

This thesis differs on the third, which is what its research question requires: it removes the convolutional backbone entirely in four of its twelve configurations, and measures the difference under complete leave-one-subject-out with every other factor held constant. On the first gap it answers the field's terminological ambiguity by declaring precisely what its own component is (\autoref{sec:cnn-implications}); on the second it knowingly departs from the literature's guidance, at the 224 $\times$ 224 resolution flagged above.


\section{Parameter-Free Attention: SimAM}
\label{sec:simam-review}

SimAM re-weights the convolutional feature maps without adding learned parameters. Yang et al.~\cite{yang2021} provide the mechanism and its image-recognition evidence; Xia et al.~\cite{xia2020b} motivate parameter-free modules in composite-database micro-expression recognition. The latter support the design principle, rather than validating SimAM itself on this task.

\subsection{Why parameter-free attention}
\label{sec:simam-why}

Local facial activity occupies a small part of a motion tensor. Attention offers a way to emphasise informative activations, but its parameter and computational costs must be justified on a 156-clip corpus. Yang et al.'s~\cite{yang2021} comparison illustrates the range of parameter costs.

\begin{table}[htbp]
    \centering
    \small
    \begin{tabular}{@{}lc@{}}
        \hline
        \textbf{Module} & \textbf{Additional learnable parameters} \\
        \hline
        SE & $2C^2/r$ \\
        CBAM & $2C^2/r+2k^2$ \\
        GC & $2C^2/r+C$ \\
        ECA & $k$ \\
        SRM & $6C$ \\
        SimAM & $0$ \\
        \hline
    \end{tabular}
    \caption[Parameter costs of attention modules]{Parameter costs reported by Yang et al.~\cite{yang2021}, Table 1. $C$ denotes channels, $r$ the reduction ratio and $k$ the relevant kernel size. The competing modules are attributed through this source.}
    \label{tab:attention-modules}
\end{table}

Zero additional parameters makes SimAM a plausible candidate, not the only defensible one: alternatives such as ECA are also small. Xia et al.~\cite{xia2020b} demonstrate a related strategy of adding parameter-free modules to a recurrent convolutional model. Neither argument establishes that a parameter-free module has zero runtime cost or cannot worsen generalisation.

\subsection{The distinctiveness premise}
\label{sec:simam-premise}

Yang et al.~\cite{yang2021} motivate attention through spatial suppression: neurons distinguishable from their surroundings are assigned greater importance. They express this as the linear separability of one target activation from the other activations in its channel. The resulting weights scale the feature map rather than being added to it.

This is a prior about activation statistics, not a detector of facial action units or minority classes. A distinctive response may encode useful deformation, noise or an artefact. Its suitability for flow-derived features therefore requires evaluation rather than following automatically from the neuroscience motivation.

\subsection{Energy and implemented refinement}
\label{sec:simam-energy}

The paper minimises a regularised separation energy and obtains a closed-form importance score. Its reference implementation can be written compactly as
\[
    d_i=(x_i-\mu)^2,\qquad
    s^2=\frac{1}{M-1}\sum_{i=1}^{M}d_i,\qquad
    E_i^{-1}=\frac{d_i}{4(s^2+\lambda)}+\frac{1}{2},
\]
\[
    \widetilde{x}_i=x_i\,\operatorname{sigmoid}(E_i^{-1}),
    \qquad \mu=\frac{1}{M}\sum_{i=1}^{M}x_i.
\]
Here $M$ is the number of positions in one channel's statistical context and $\lambda$ is a positive regularisation constant. The $+\tfrac12$ follows from the reciprocal-energy expression rather than an arbitrary offset.

An implementation distinction matters: the shared variance printed with Yang et al.'s equation (5) divides by $M$, whereas their Figure 3 code divides by $M-1$~\cite{yang2021}. The equations above follow the reference \emph{code}, as does this thesis. The sigmoid preserves the ordering of the scores while bounding the multiplicative weights; it does not create a hard mask of inactive regions.

\subsection{What three-dimensional weights mean}
\label{sec:simam-3d-weights}

The original module produces a separate weight at each channel--height--width position~\cite{yang2021}. Its three-dimensional weights therefore describe $C\times H\times W$, not a temporal volume. Statistics are computed separately within each channel's spatial plane, with $M=HW$. Extending the context across video frames is an additional design decision, not part of the original image formulation.

\subsection{Evidence and hyper-parameter selection}
\label{sec:simam-evidence}

Yang et al.~\cite{yang2021} evaluate SimAM with several network families on CIFAR and ImageNet, reporting competitive classification performance against learned attention modules. Their timing results also show that unchanged parameter totals do not imply unchanged inference speed. The evidence supports testing SimAM, but does not establish its value on micro-expression motion tensors.

For CIFAR, they sweep $\lambda$ from $10^{-1}$ to $10^{-6}$ and select $10^{-4}$ as a balance between accuracy and variability. They repeat selection for ImageNet and choose $0.1$~\cite{yang2021}. Thus $10^{-4}$ is a dataset-selected operating point, not a universal setting justified by the formula.

\subsection{Limits of transfer to this task}
\label{sec:simam-limitations}

The source classification experiments concern natural images, different training-set sizes and learned feature distributions. In this thesis the input is flow and strain, so distinctiveness need not have the same relationship to useful evidence. The statistics are calculated from each feature map, not estimated by pooling the entire training corpus; the small dataset instead constrains how the surrounding network is learned.

The image experiments also do not decide whether video statistics should be computed per frame or over a clip. Nor does the lack of learned attention weights establish that runtime, memory or overfitting concerns disappear. These are empirical limits on transferring the source result, rather than reasons to assume either success or failure.

\subsection{Implications for this research}
\label{sec:simam-implications}

\begin{table}[htbp]
    \centering
    \small
    \begin{tabular}{@{}p{0.29\linewidth} p{0.65\linewidth}@{}}
        \hline
        \textbf{Choice} & \textbf{Implementation and interpretation} \\
        \hline
        Placement & Applied separately to each convolutional stream before concatenation; no trainable weights or biases are introduced. \\
        Refinement & Uses the reference-code reciprocal energy and sigmoid gate~\cite{yang2021}. \\
        Statistical context & Mean and variance span $(T,H,W)$ within each channel, with divisor $THW-1$, rather than one $(H,W)$ plane. \\
        Regularisation & $\lambda=10^{-4}$ is retained without re-tuning on CASME II. \\
        Architectural dependency & SimAM is enabled only with the CNN stem in this model. Four of sixteen candidate configurations are excluded, leaving twelve evaluated configurations and four SimAM matched pairs. \\
        \hline
    \end{tabular}
    \caption[SimAM implementation choices]{SimAM implementation choices and their relation to the reviewed image formulation. The dependency on the CNN is a rule of this architecture.}
    \label{tab:simam-decisions}
\end{table}

The five-dimensional input includes batch, channel, time, height and width. Pooling statistics across time makes an activation distinctive relative to its whole clip rather than its own frame. This could emphasise a brief contraction, but does not guarantee low weights at onset or high weights at the apex: the score depends on squared deviation, not expression labels. The two context choices are not compared experimentally.

The CNN dependency similarly describes this implementation, not a mathematical prohibition on applying SimAM to other tensors. Its measured effect is conditional on the stem being present. Retention should be judged against both recognition performance and computational cost; parameter count alone cannot settle that decision.

\subsection{Research gap and scope}
\label{sec:simam-gap}

The study evaluates SimAM in a motion-input, small-corpus setting through four matched comparisons under complete leave-one-subject-out evaluation. It tests the implemented module as a whole, rather than separately establishing the benefit of each design choice.

Three follow-ups remain: testing whether the distinctiveness prior suits flow-derived features, re-tuning $\lambda$, and comparing per-frame with whole-clip statistics. Each requires new evaluations despite introducing no learned attention weights. Per-class changes may also be examined, but a rare class benefiting most is a hypothesis of this study, not a prediction established by Yang et al.~\cite{yang2021}.


\section{The Temporal Encoder: Transformers for Micro-Expression Recognition}
\label{sec:transformer-review}

The temporal encoder is motivated by Zhang et al.'s Short and Long Range Relation Based Spatio-Temporal Transformer (SLSTT)~\cite{zhang2022} and its Vision Transformer (ViT) foundation~\cite{dosovitskiy2021}. The component evaluated here differs substantially from published SLSTT, particularly in the axis over which self-attention operates.

\subsection{Why temporal self-attention}
\label{sec:transformer-why}

A micro-expression evolves through onset, peak and relaxation. Retaining a motion sequence permits a model to relate separated stages rather than compressing them immediately into one descriptor. Zhang et al.~\cite{zhang2022} motivate their design through the need to capture both local and global patterns.

This thesis compares a temporal transformer with mean pooling over the same input sequence. Mean pooling discards order, whereas positional information and self-attention allow the encoder to model relations among frame-pair features. This motivates the comparison without establishing in advance which representation will generalise better.

\subsection{The Vision Transformer foundation}
\label{sec:transformer-vit}

ViT splits an image into patches, projects them into token vectors, adds learned positional embeddings and processes them with transformer encoder layers~\cite{dosovitskiy2021}. Attention then relates spatial patches. Dosovitskiy et al. also show that data scale and training regime matter: without the locality and translation-equivariance biases of convolution, image transformers can perform poorly when trained with insufficient data, while large-scale pre-training improves their performance.

That evidence is a caution for a 156-clip study, not a universal requirement that every transformer must be pre-trained. A short sequence of motion features presents a different learning problem from image patches, but remains subject to limited data and overfitting.

\subsection{Published SLSTT}
\label{sec:slstt-architecture}

Zhang et al.~\cite{zhang2022} combine three stages. First, \textbf{onset-referenced optical flow} describes sampled frames relative to the onset rather than to their immediate predecessors. They argue that this gives a more structured trajectory towards and away from the apex. Their experiments use this representation throughout; they do not provide a matched short-term-flow ablation establishing its independent gain.

Second, a \textbf{spatial ViT encoder} processes each flow field as patches. The implementation uses ImageNet-pre-trained ViT-B/16, twelve encoder layers and $384\times384$ input. Third, an \textbf{LSTM temporal aggregator} combines the per-frame features. A mean aggregator supplies the comparison variant. Self-attention is spatial in this architecture; recurrence handles time.

\subsection{Published results and the aggregation comparison}
\label{sec:slstt-results}

\begin{table}[htbp]
    \centering
    \small
    \begin{tabular}{@{}lccc@{}}
        \hline
        \textbf{Method} & \textbf{Composite UF1} & \textbf{CASME II UF1} & \textbf{CASME II UAR} \\
        \hline
        LBP-TOP & 0.588 & 0.703 & 0.743 \\
        Bi-WOOF & 0.630 & 0.781 & 0.803 \\
        OFF-ApexNet~\cite{liong2019a} & 0.720 & 0.876 & 0.868 \\
        STSTNet~\cite{liong2019b} & 0.735 & 0.838 & 0.869 \\
        EMR~\cite{liu2019} & 0.789 & 0.829 & 0.821 \\
        RCN~\cite{xia2020b} & 0.705 & 0.809 & 0.856 \\
        SLSTT-Mean & 0.788 & 0.844 & 0.830 \\
        SLSTT-LSTM & 0.816 & 0.901 & 0.885 \\
        \hline
    \end{tabular}
    \caption[SLSTT and comparison results]{Selected results from Zhang et al.~\cite{zhang2022}, Table 2, reported under their three-class Composite Database Evaluation heading. Columns retain the source's dataset labels. The listed methods are not all arms of one controlled architectural comparison.}
    \label{tab:slstt-results}
\end{table}

The SLSTT variants compare temporal aggregation with the spatial transformer retained in both. CASME II UF1 increases from 0.844 to 0.901, a difference of 0.057, or 5.7 UF1 percentage points~\cite{zhang2022}. This supports evaluating aggregation choices within that framework. It is not an estimate of what an LSTM would add to this thesis, where mean pooling follows \emph{temporal} self-attention rather than independent spatial encoding.

Protocol labels need care. The source describes three-class LOSO evaluations on the databases separately and together under its Composite Database Evaluation heading. That heading alone therefore does not establish that every CASME II column entry uses pooled cross-database training. It also reports a distinct Sole Database Evaluation using original labels, including five CASME II classes~\cite{zhang2022}. Neither setting should be equated automatically with this thesis's 156-clip three-class pool.

\subsection{Limitations of the evidence}
\label{sec:transformer-limitations}

Published SLSTT retains a pre-trained spatial encoder. Its results do not isolate the benefit of a transformer trained from scratch on this study's motion representation. The Mean--LSTM comparison isolates an aggregator, not the presence of a transformer, since both variants contain one.

Label interpretation also limits comparison. Zhang et al.~\cite{zhang2022} note that facial action units do not map uniquely to emotion labels, so visual evidence alone need not determine the annotation. This is a limitation of the recognition task, not a numerical upper bound measured by the present experiments.

\subsection{Implementation and six declared divergences}
\label{sec:transformer-implications}

The evaluated component receives 32 tokens of width 96. They come either from the convolutional stem or from $4\times4$ spatial pooling followed by a \emph{learned linear projection}. The encoder uses four pre-norm layers, eight heads, feed-forward width 256 and dropout 0.1. Fixed sinusoidal positions precede attention; mean pooling follows it. With the encoder disabled, the input tokens are averaged directly.

\begin{table}[htbp]
    \centering
    \small
    \begin{tabular}{@{}p{0.23\linewidth} p{0.34\linewidth} p{0.35\linewidth}@{}}
        \hline
        \textbf{Property} & \textbf{Published SLSTT~\cite{zhang2022}} & \textbf{This thesis} \\
        \hline
        Attention axis & Spatial patches within a frame & Temporal tokens within a clip \\
        Aggregation & LSTM over spatially encoded frames & Mean after temporal self-attention \\
        Initialisation & ImageNet-pre-trained ViT-B/16 & Trained from scratch \\
        Capacity & Twelve layers, width 768; $384\times384$ inputs & Four layers, width 96; 32 tokens \\
        Positions & Learned spatial embeddings & Fixed sinusoidal temporal encoding \\
        Input flow & Relative to onset & Relative to the preceding sampled frame \\
        \hline
    \end{tabular}
    \caption[Divergences from published SLSTT]{Six material divergences from published SLSTT. The component evaluated here is a small temporal transformer encoder, not a reproduction of the published architecture.}
    \label{tab:slstt-divergences}
\end{table}

Six matched pairs evaluate the temporal encoder with the other component choices fixed within each pair. They test the whole encoder, including positional encoding, feed-forward layers and normalisation, against direct mean pooling. They do not isolate self-attention from those accompanying operations or compare it with another learned temporal model.

Three follow-ups follow from the divergences: testing an LSTM aggregator, comparing onset-referenced with consecutive-frame flow, and replacing fixed positions with learned embeddings. The last would add $32\times96=3{,}072$ embedding parameters for this sequence length. All three remain untested; changing the input representation also requires re-training and evaluation before any recognition benefit can be claimed.

\subsection{Research gap and scope}
\label{sec:transformer-gap}

The study asks what a small temporal encoder contributes after motion extraction, using a single corpus without pre-training. This differs from SLSTT's spatial transformer and recurrent aggregation while retaining its motivation to model relations beyond isolated frames.

The contribution is a matched comparison of the implemented encoder and an order-blind baseline within this pipeline. It neither validates published SLSTT nor establishes a general ranking of attention, recurrence and convolution. The six declared divergences and the limited training regime define the scope of the architectural conclusions.


\section{Training Under Severe Class Imbalance}
\label{sec:imbalance-review}

This section covers the training configuration held constant across all twelve ablated configurations, and which \autoref{sec:corpus-implications} identified as forced by the corpus rather than chosen freely: the handling of a 4 : 1.3 : 1 class imbalance.

\textbf{A deliberate exclusion.} Two topics the project's reading list groups here --- Grad-CAM++ visual auditing and adversarial identity disentanglement --- are \textbf{not reviewed}, because neither is part of the implemented system. The loss module records that ``the Cross-Batch Memory (XBM) queue, and the adversarial identity loss --- are intentionally removed'', and no gradient-based visualisation exists anywhere in the codebase. Both are declared as proposed further work.

\subsection{The imbalance is a property of the phenomenon, not of the sampling}
\label{sec:imbalance-property}

\autoref{sec:label-taxonomy} and \autoref{sec:original-baseline} established the class counts and the always-Negative floor; what matters here is why the skew exists. It is not an artefact of careless collection. Yan et al.~\cite{yan2014} are explicit that it follows from what can be elicited in a laboratory: ``some types of facial expressions are difficult to elicit in laboratory situations, thus the samples in different categories distributed unequally, e.g., there are 60 disgust samples but only 7 sadness samples.'' Nor is it specific to CASME II --- the MEGC composite shows the same pattern (\autoref{sec:megc-metrics}). Skew must be coped with; it cannot be collected away.

\subsection{Three places to intervene}
\label{sec:imbalance-three-places}

The reviewed corpus corrects for imbalance at three distinct points, worth separating because this thesis intervenes at all three and the interaction between two of them turned out to matter.

\begin{enumerate}
	\item \textbf{At the data}, by changing how often each class is presented during training --- resampling or augmentation.
	\item \textbf{At the loss}, by changing how much each example contributes to the gradient --- class weighting or focal re-weighting.
	\item \textbf{At the metric}, by changing how performance is scored --- macro-averaged rather than sample-averaged measures.
\end{enumerate}

\subsection{Loss-level correction: focal loss}
\label{sec:imbalance-focal}

Zhao et al.~\cite{zhaoS2021} bring focal loss into micro-expression recognition ``to alleviate the inefficient model training caused by the class imbalance in the MEs Datasets'', since cross-entropy ``could be biased toward particular emotions that constitute a larger portion of the training set''. They adopt the formulation developed for one-stage object detection, where the imbalance is between foreground and background regions:

\[
	\text{FL}(p_t) = -\alpha_t (1 - p_t)^{\gamma}\log p_t
\]

where $p_t$ is the model's probability for the true class, $\gamma$ ``the balance factor for loss'' and $\alpha_t$ ``the weight balance factor for samples'' (\autoref{sec:focal-loss}); under skew, the examples it up-weights are disproportionately the minority classes. Adapting it to multi-class classification, Zhao et al. report that ``$\gamma$ is set as 2 in practice, and $\alpha$ is treated as a hyper-parameter to set by cross validation.''

The division of labour there matters: $\gamma$ re-weights by \textit{difficulty} and is treated as a constant, $\alpha$ by \textit{class frequency} and is treated as dataset-specific. Conflating them is where this project went wrong (\autoref{sec:imbalance-double-correction}).

\subsection{Data-level correction: resampling and augmentation}
\label{sec:imbalance-balanced-loss}

Xia et al.~\cite{xia2020a} address the same problem at the data end and pair the two interventions, using ``temporal data augmentation strategies as well as a balanced loss \ldots jointly'' to overcome ``limited and imbalanced training samples''. The pairing is not left unexamined: their subsection \textit{``The Impact of Data Augmentation and Balanced Loss''} removes each correction in turn from the joint system and reports the balanced loss improving performance ``slightly'', with one leave-one-subject-out cell on SMIC where removing it scores higher. This is the corpus's only controlled look at composing two corrections, and \autoref{sec:imbalance-double-correction} records a far larger effect in the same direction.

\subsection{Metric-level correction}
\label{sec:imbalance-metrics}

\autoref{sec:megc-metrics} defines UF1 and the case for macro-averaged metrics. Liong et al.~\cite{liong2016} add the argument from the subject side rather than the class side, reporting F1, precision and recall because these ``provide a more meaningful perspective than accuracy rates when the datasets used are naturally imbalanced since each subject has a different number of video samples.'' The two arguments compound on CASME II, where classes are skewed \textbf{and} per-subject clip counts range from 1 to 33 (\autoref{tab:corpus-threats}c).

\subsection{Limitations}
\label{sec:imbalance-limitations}

\textbf{The interaction between corrections is examined once, and only at small magnitude.} Xia et al.~\cite{xia2020a} ablate augmentation and balanced loss out of the joint system and report the balanced loss helping only ``slightly'' --- at most about two points across their twelve cells, and in one cell costing accuracy; Zhao et al.~\cite{zhaoS2021} apply focal loss alongside a large per-sample augmentation, but without class-proportional resampling and without ablating the pairing. What no reviewed paper reports is the pairing failing outright.

\textbf{$\gamma$ is inherited rather than tuned.} Zhao et al.~\cite{zhaoS2021} take the value from the object-detection literature ``in practice''. They do compare focal loss against cross-entropy across a sweep of $\alpha$, but that experiment holds $\gamma$ fixed, so the focusing exponent remains unvalidated on this task.

\textbf{Label smoothing is named in the corpus but never examined on this task.} Dosovitskiy et al.~\cite{dosovitskiy2021} list it among the regularisation parameters they tune, but no reviewed micro-expression paper uses or discusses it, or reports what it does under class skew. It is used here and declared as a choice this corpus supports neither way.

\subsection{Implications for this research}
\label{sec:imbalance-double-correction}

\begin{table}[htbp]
	\centering
	\footnotesize
	\begin{tabular}{@{}p{5.6cm} p{8.0cm}@{}}
		\hline
		\textbf{Finding from the review} & \textbf{Decision taken in this thesis} \\
		\hline
		Focal loss addresses class-frequency bias in micro-expression recognition~\cite{zhaoS2021} & Focal loss adopted as the training objective, with \textbf{$\gamma = 2.0$}, matching Zhao et al.'s reported setting \\
		$\alpha$ is a dataset-specific class-weight term~\cite{zhaoS2021} & Implemented as an optional per-class $\alpha$ vector, \textbf{switched off at run time} --- see the declaration below \\
		Correct at the data end as well as the loss end~\cite{xia2020a} & Minority oversampling via a weighted sampler, active in every configuration --- a sampler, where Xia et al. use uniform augmentation plus a balanced loss \\
		Imbalance-aware metrics are mandatory~\cite{see2019,liong2016} & Pooled macro F1 --- identical in construction to MEGC's UF1 --- as the primary metric, with accuracy quoted only alongside the always-Negative floor \\
		\textit{(no micro-expression source)} & Label smoothing of 0.05 applied inside the focal loss; named among the regularisers tuned by Dosovitskiy et al.~\cite{dosovitskiy2021}, but declared as unsupported either way by this corpus \\
		\hline
	\end{tabular}
	\caption[Review findings and the imbalance-handling decisions they determine]{Review findings and the imbalance-handling decisions they determine.}
	\label{tab:imbalance-decisions}
\end{table}

\textbf{The finding this project contributes: correct at exactly one point, and know which.} The configuration is resolved at run time by a single rule --- class weighting in the loss is enabled only when the balanced sampler is \textit{not} active --- and the run log records the decision each fold: \textit{``Class weights in loss disabled (balanced sampler already active).''}

That rule exists because of two failures in opposite directions, both recoverable from this project's own run history. In the earliest code era the loss carried inverse-frequency class weights and no balanced sampler existed; under that single correction the proposed configuration abandoned the \textbf{Positive} class outright, predicting none of its clips --- a failure initially misdiagnosed as over-capacity overfitting. A later era added the sampler while leaving the loss weights unconditional, so both corrections ran at once; the collapse then inverted, and configurations abandoned the \textbf{Negative} majority class instead, one of them predicting a single class for every clip in the corpus. Correcting once at the loss starved the minority; correcting twice over-served it. With the guard in place and exactly one correction active, no configuration in the reported study abandons any class.

This is a contribution rather than an incident, because of how little guidance the corpus offers (\autoref{sec:imbalance-limitations}): no reviewed paper warns that combining corrections can invert the intended effect rather than merely blunt it. On a 156-clip corpus with a 4 : 1 skew, correcting twice is as damaging as correcting once in the wrong place --- the two failures differ only in which class is abandoned.

\textbf{A second declaration.} The severity of the imbalance also constrains what the ablation can conclude. Because every configuration shares the same sampler, loss and $\gamma$, this thesis measures \textbf{no} effect attributable to the imbalance treatment itself: it is a constant, not a factor. The claim supported is that the four architectural components were compared under a training regime that does not abandon minority classes, not that this regime is optimal.

\subsection{Research gap and scope}
\label{sec:imbalance-gap}

The corrections themselves are well established; what the corpus does not settle are the three gaps of \autoref{sec:imbalance-limitations} --- how two corrections compose, what $\gamma$ should be on this task, and what over-correction looks like when it fails.

This thesis contributes to the first and third, documenting two opposite class-collapse failures and the run-time rule that resolved them (\autoref{sec:imbalance-double-correction}): a negative result the corpus does not contain and one directly actionable for anyone reproducing this class of pipeline. On the second it offers no improvement, since $\gamma = 2.0$ and label smoothing of 0.05 are held constant across all twelve configurations. Sweeping either, and the Grad-CAM++ auditing and identity-disentanglement work this project's earlier stages removed, are declared as further work.


\section{Synthesis and Research Gap}
\label{sec:synthesis-and-gap}

This section consolidates the gap statements closing \autoref{sec:corpus}--\autoref{sec:imbalance-review}: what the corpus collectively establishes, the four limitations recurring across every topic, the resulting research gap, and which parts of it this thesis closes. It also collects in one place every declared divergence between the system built here and the literature it draws on.

\subsection{What the corpus establishes}
\label{sec:corpus-establishes}

Read as a whole, the reviewed literature settles a good deal. \textbf{On data:} the corpora are small, demographically narrow, laboratory-bound and irreducibly imbalanced by what can be elicited, and CASME II remains the primary single-corpus benchmark under its native leave-one-subject-out protocol (\autoref{sec:corpus}, \autoref{sec:imbalance-review}). \textbf{On representation:} motion rather than appearance is the appropriate input; a strain field adds a further cue, though on CASME II only weakly and under a protocol that is not subject-disjoint; and both the magnification factor and the interpolated sequence length have interior optima located by sweep (\autoref{sec:evm-review}--\autoref{sec:tim-review}). \textbf{On architecture:} at this data size capability bought without parameters outperforms capability bought with them, and self-attention is the strongest-performing mechanism reported in the corpus (\autoref{sec:cnn-review}--\autoref{sec:transformer-review}). \textbf{On evaluation:} accuracy alone is inadequate under skew and must give way to macro-averaged measures accumulated across folds before averaging (\autoref{sec:megc-metrics}, \autoref{sec:imbalance-metrics}).

\subsection{Four limitations that recur across every topic}
\label{sec:four-limitations}

The per-section gap statements are not nine independent observations. Four structural problems generate most of them.

\textbf{(a) Components are isolated one at a time, inside one proposed design, and never jointly.} Matched-pair measurement is not absent from the corpus. On the input side \autoref{tab:evm-effect} isolates magnification and \autoref{tab:strain-performance} isolates strain; on the architectural side Xia et al. compare a basic recurrent convolutional network against three variants adding a wide expansion, a shortcut connection and an attention unit, holding every other parameter constant across all four (\autoref{sec:cnn-limited-data}), and separately sweep depth against input resolution (\autoref{sec:cnn-shrink}); SLSTT's Mean-versus-LSTM comparison isolates its aggregator (\autoref{sec:slstt-results}).

What no reviewed study does is vary \textit{several} components together, so that their interactions can be read. Every ablation in the corpus is internal to one proposed design and moves one factor at a time from that design's own defaults; across papers, systems are compared whole --- STSTNet against OFF-ApexNet, SLSTT against STSTNet --- so a difference in the reported number cannot be attributed to any one stage. There is therefore no basis for deciding which of the learned stages are load-bearing, nor whether one component's value depends on another's presence.

\textbf{(b) The evaluation protocol is unstable, and under-reported.} Sample counts, label-set size and the choice among leave-one-subject-out, leave-one-video-out and leave-one-sample-out all vary inconsistently, and the field's most-quoted baseline is reported under two different protocols by its two primary sources and under two more within the originating paper itself (\autoref{sec:original-baseline}). Given this corpus's fold composition (\autoref{sec:corpus-threats}), these are not bookkeeping differences: they change what the numbers mean.

\textbf{(c) The strongest results depend on regimes a single-corpus study cannot access.} The best CASME II figures in the corpus come from composite-database training (\autoref{sec:megc-metrics}) and, in SLSTT's case, ImageNet pre-training (\autoref{sec:transformer-limitations}) --- neither a property of the architecture being credited. Almost no evidence exists about what these components do when trained from scratch on one small corpus, which is the situation most practitioners are in.

\textbf{(d) Preprocessing parameters have interior optima, and are transferred without re-validation.} The magnification factor, the interpolated sequence length and the input resolution all rise then fall, all have their optima located only by sweep, and none has a principled selection rule (\autoref{sec:evm-evidence}, \autoref{sec:tim-how-many}, \autoref{sec:cnn-shrink}). Values are carried between pipelines whose downstream models differ entirely: $\alpha = 30$ moves from a single-frame appearance encoder to other settings, TIM10 propagates from one baseline paper to the whole field, and $\lambda = 10^{-4}$ transfers from CIFAR though its own authors re-tuned it for ImageNet (\autoref{sec:simam-evidence}).

\subsection{The research gap}
\label{sec:research-gap}

Taking (a)--(d) together:

\textbf{The field has established a standard pipeline for micro-expression recognition --- magnify, compute motion, normalise duration, extract spatially, model temporally --- but has never measured which of its stages are responsible for the result. Because components are only ever varied one at a time inside a single proposed system, under protocols that vary in unstated ways, and predominantly in training regimes richer than a single small corpus provides, it is not currently possible to say which parts of that pipeline earn their place, which are inherited convention, and which depend on the presence of another.}

This matters practically as well as scientifically. A practitioner assembling this pipeline today has no published basis for allocating effort: for judging whether the most computationally expensive stage is the most valuable, whether a component with a negligible average effect may still matter in a specific failure, or whether a ranking read off an incomplete evaluation of a corpus this size is stable at all.

\subsection{How this thesis addresses the gap --- and what it does not}
\label{sec:addressing-gap}

\textbf{What it does.} This thesis treats the pipeline as a factorial experiment rather than a proposal: four components varied independently across every architecturally valid cell of a 2\textsuperscript{4} matrix, twelve configurations, each evaluated under \textbf{complete 25-fold leave-one-subject-out on CASME II} with every non-varied factor held identical. That gives the joint, multi-component measurement (a) says the corpus lacks; reporting the protocol in full addresses (b) for this study at least; and training from scratch is the regime (c) identifies as unevidenced. Each component's marginal contribution is read from pairs of runs differing in one factor alone, and the distribution of those four effects --- how much of the pipeline's performance each stage accounts for, whether any subtracts, and what each costs --- is the substantive answer to the gap.

Three further contributions follow from the review rather than from the experiment:

\begin{itemize}
	\item \textbf{A caution about partial evaluations.} This project's earlier evaluations, on smaller held-out sets and incomplete fold counts, ranked the twelve configurations differently from the complete run reported here. They are not evidence of protocol sensitivity, since they differ in sample count, training budget and class coverage too, and the earliest carry a data-routing defect --- but the episode illustrates limitation (b) in practice, and only the complete run is reported.
	\item \textbf{A repaired measurement.} That defect left the magnification switch inert, producing bit-identical results across every pair those runs contained. What made it detectable as a defect rather than a null result was the review's evidence that magnification has a clear, repeatable effect (\autoref{sec:evm-implications}).
	\item \textbf{Two opposite class-collapse failures.} A single loss-level correction suppressed the minority class; adding a data-level correction on top of it suppressed the majority class instead. No reviewed paper reports either collapse (\autoref{sec:imbalance-double-correction}).
\end{itemize}

\textbf{What it does not do.} Three of the four limitations are only partially addressed. On (b), this study reports its own protocol completely but cannot repair the field's inconsistency, and its own N = 156 differs from MEGC's 145. On (c), it establishes what these components do from scratch on one corpus, which is a different question from --- not an answer to --- how they behave under composite training; no cross-dataset claim is made anywhere. On (d), the thesis inherits parameter values rather than sweeping them: $\alpha = 10$, T = 32, $\lambda = 10^{-4}$ and $\gamma = 2.0$ are fixed throughout, so the ablation measures architecture at one point in preprocessing space, not across it.

\subsection{A consolidated ledger of declared divergences}
\label{sec:divergence-ledger}

Reviewing the literature closely enough to write this chapter surfaced twelve points at which the implemented system departs from the work it draws on. Several were discovered only through the review. They are collected in \autoref{tab:divergence-ledger} so that no result rests on an unstated assumption.

\begin{table}[htbp]
	\centering
	\footnotesize
	\begin{tabular}{@{}p{0.85cm} p{6.0cm} p{6.6cm}@{}}
		\hline
		\textbf{\S} & \textbf{Divergence} & \textbf{Consequence} \\
		\hline
		3.1.3 & Working set is N = 156, including sadness and fear in Negative; MEGC's CASME II subset is N = 145 & Comparisons against challenge figures are not exactly like-for-like \\
		3.1.2 & The corpus's own registered release is not used; the pipeline reads the original recorded frames, unregistered and uncropped & No landmarking or alignment of any kind precedes the network; tolerance to pose and scale must be learned from 156 clips \\
		3.2.6 & Temporal band is the unnarrowed 5--25 Hz duration rule, where Bai et al. narrowed it to 15--25 Hz as noisy & The pass-band admits the low-frequency range they rejected \\
		3.4.3 & Strain magnitude counts shear once, not twice as in the published Frobenius norm & Shear weighted lower by $\sqrt{2}$ relative to normal strain; untested \\
		3.4.7 & No dedicated strain filtering, where the literature applies Wiener and Gaussian filters & Spatial differentiation can amplify flow-estimation noise, and nothing downstream suppresses it \\
		3.5.7 & Temporal normalisation is uniform index sampling, \textbf{not} manifold-based TIM; no frames are synthesised & Avoids fabricating motion; forfeits the ability to lengthen short clips \\
		3.5.7 & Magnification is applied \textbf{after} subsampling, at a nominal 200 fps & Realised pass-band is clip-dependent and below the intended 5--25 Hz \\
		3.6.6 & Backbone kernels are (1, 3, 3) --- \textbf{no temporal mixing} & The ablation tests a learned \textit{spatial} stem, not spatio-temporal convolution \\
		3.6.6 & Backbone input is 224 $\times$ 224, against Xia et al.'s guidance of $\leq$ 100 $\times$ 100 --- though Bai et al. and Li, Huang and Zhao both use 224 $\times$ 224 & Whatever the stem is measured to contribute may be a resolution effect as much as an architecture effect \\
		3.7.7 & SimAM statistics pooled over the whole clip rather than per frame as in the source; $\lambda$ not re-tuned & Neurons judged against the clip, not the frame; defensible but untested \\
		3.8.6 & Six divergences from published SLSTT (temporal not spatial attention, mean not LSTM aggregation, from scratch, smaller, sinusoidal encoding, short-term flow) & Variable D is a temporal transformer encoder, \textbf{not} a reproduction of SLSTT \\
		3.9.7 & Label smoothing of 0.05, named in the corpus only by Dosovitskiy et al. and never examined on this task & A choice the micro-expression corpus supports neither way, held constant and unmeasured \\
		\hline
	\end{tabular}
	\caption[Every declared departure from the reviewed literature]{Every declared departure from the reviewed literature, with its consequence.}
	\label{tab:divergence-ledger}
\end{table}

Three of these --- the strain formula, the magnification ordering, and the input resolution --- were identified by the review itself rather than being known design decisions, and each is declared as proposed further work. None invalidates the ablation, since every configuration shares the same choices. What they constrain is the \textit{wording} of the conclusions: what is measured is the contribution of a spatial stem at 224 $\times$ 224 and of temporal self-attention in this pipeline, not a verdict on 3D convolution or on SLSTT.

\subsection{What this chapter contributes to the thesis}
\label{sec:ch3-contribution}

Three things carry forward. First, the \textbf{methodology that follows is determined by this review rather than chosen freely}: the corpus dictates the class grouping, representation, fold structure and metric, and each component's parameters come from published sweeps (\autoref{sec:corpus-implications}). Second, the \textbf{interpretation of the central result is prepared here rather than improvised later}: flow and strain already perform spatio-temporal feature extraction analytically (\autoref{sec:flow-implications}, \autoref{sec:strain-implications}), so a negative contribution from a spatial-only backbone at an oversized input resolution is expected rather than anomalous (\autoref{sec:cnn-implications}). Third, the \textbf{limits of what can be claimed are fixed before any number is reported}: a small, single-seed, demographically narrow corpus measured against richer published regimes, with the twelve divergences of \autoref{tab:divergence-ledger} standing throughout --- consequences of this review, not concessions made after the fact.
 in place of contents/chapter3.tex.
% Do not include both: they share every label.
%
% Assembled from chapter3_short.tex (3.1-3.2) and chapter3_short_part3..10.tex.
% Those part files are retained unchanged as the editing units; regenerate this
% file from them rather than editing it directly, or edit here and retire them.
%
% Structure that other chapters depend on, and that must not be disturbed:
%   * Chapter 2 cites Sections 3.1.2, 3.1.6, 3.1.8, 3.3.1, 3.3.3, 3.5.7,
%     3.9, 3.9.5 and 3.9.7 BY NUMBER, so those subsection positions are fixed.
%   * The divergence ledger (Table, sec:divergence-ledger) cites twelve
%     subsections by number. scratchpad/ledgercheck.py re-verifies all of them
%     plus every external numeric reference after any renumbering.
%   * Chapters 4, 5, 6 and the introduction reference sec:research-gap,
%     sec:four-limitations, sec:addressing-gap, sec:divergence-ledger,
%     sec:synthesis-and-gap, sec:imbalance-review, sec:imbalance-gap,
%     sec:imbalance-double-correction, sec:cnn-review, sec:cnn-implications,
%     sec:tim-implications, sec:evm-implications, sec:evm-evidence,
%     sec:evm-limitations, sec:evm-motion-hypothesis, sec:corpus,
%     sec:corpus-threats, sec:corpus-implications, sec:label-taxonomy,
%     sec:megc-metrics, sec:sibling-corpora, sec:flow-review, sec:flow-implications,
%     sec:strain-review, sec:strain-implications, sec:strain-divergence,
%     sec:simam-gap, sec:simam-implications, sec:slstt-results,
%     sec:transformer-implications, sec:addressing-gap, tab:megc-cde-results
%     and tab:spontaneous-corpora.
%
% Conventions held throughout: no sentence points at another chapter; every
% float carries a short \caption[...] for the List of Figures / List of Tables;
% figures are fig3_1_dataset.png and fig3_2_fold_composition.png only.
% =============================================================================

\chapter{Literature Review}
\label{ch:literature}

This chapter reviews the literature bearing on each component of the recognition pipeline evaluated in this thesis: the corpus on which every experiment is run, the four techniques placed under ablation, the three input transformations they operate on, and the training regime held constant across all of them.

It is organised as a single continuous argument rather than as ten independent surveys: a fact is established once, in the section that owns it, and cross-referenced thereafter. Each section runs from what the technique addresses, through how it works and what the published evidence shows, to its limitations and the implications for this thesis, closing with the gap and how this work differs. Two claims recur --- that on a corpus of 156 clips several decisions usually presented as modelling choices are in fact forced by properties of the data, and that several results usually attributed to architecture are attributable to the evaluation protocol. \autoref{sec:synthesis-and-gap} consolidates the nine closing statements into a single research gap and collects every declared departure from the reviewed methods.


% =============================================================================
\section{The Evaluation Corpus: CASME II and the Spontaneous Micro-Expression Benchmark Family}
\label{sec:corpus}
% =============================================================================

This section covers the corpus literature only --- the papers that introduce, specify or standardise evaluation on the database used throughout this thesis. It is treated first because on a database of this size several design decisions usually presented as modelling choices are \textit{forced} by properties of the data.

\subsection{Why a dedicated spontaneous corpus was necessary}
\label{sec:why-spontaneous}

Automatic facial expression recognition matured on large corpora of \textit{posed} expressions --- CK+, MUG, MMI, JAFFE and Multi-PIE --- reporting accuracies above 90\,\% on the six basic posed emotions by the early 2010s~\cite{yan2014}. Micro-expression research could not follow that path: as \autoref{sec:phenomenon} establishes, a micro-expression is involuntary and cannot be \textit{performed} on instruction. Yan et al.~\cite{yan2014} review a pre-2013 landscape of only four micro-expression databases, of which only two were spontaneous (\autoref{tab:spontaneous-corpora}). Two problems in it motivated CASME II directly.

The first is \textbf{construct validity}. Yan et al.~\cite{yan2014} and Li et al.~\cite{liX2013} both note that posed ``micro-expressions'', produced deliberately on instruction, differ systematically from spontaneous ones --- Li et al.~\cite{liX2013} further observing that USF-HD adopted a 2/3-second duration criterion longer than the definitional upper bound of \autoref{sec:phenomenon}. A system trained on posed data is therefore not demonstrably a micro-expression recogniser; it may be a recogniser of fast deliberate movement.

The second is \textbf{acquisition quality}. Defined jointly by very short duration and very low intensity, the phenomenon sits close to the noise floor of ordinary video: Yan et al.~\cite{yan2014} argue that CASME's 60~fps coarsely samples a movement lasting a fraction of a second --- while marking the benefit of higher rates as unresolved --- and that flicker under alternating-current lighting contaminated earlier recordings. Elicitation was a third difficulty: high-stakes deception paradigms produced samples heavily contaminated by conversational and other emotion-irrelevant movement~\cite{yan2014}, and the paradigm that proved workable, which CASME II adopts, is to have participants watch high-arousal video episodes while attempting to keep the face neutral~\cite{liX2013}.

\begin{table}[htbp]
	\centering
	\footnotesize
	\setlength{\tabcolsep}{4pt}
	\begin{tabular}{l c c c c c c c}
		\hline
		\textbf{Corpus} & \textbf{Subj.} & \textbf{Micro samples} & \textbf{fps} & \textbf{Facial res.} & \textbf{Spont.} & \textbf{FACS} & \textbf{Classes} \\
		\hline
		\multicolumn{8}{l}{\textit{The pre-CASME II landscape, reproduced from Yan et al.~\cite{yan2014} (Table 1)}} \\
		USF-HD~\cite{shreve2011}                       & ---      & 100 & 30  & ---              & No  & No  & 6 \\
		Polikovsky et al.\textsuperscript{\textdagger} & 10       & 42  & 200 & ---              & No  & No  & 6 \\
		SMIC~\cite{liX2013}                            & 16 valid & 164 & 100 & 190 $\times$ 230 & Yes & No  & 3 \\
		CASME (2013)\textsuperscript{\textdagger}      & 19 valid & 195 & 60  & 150 $\times$ 190 & Yes & Yes & 7 \\
		\hline
		\multicolumn{8}{l}{\textit{The spontaneous corpora that define the current context}} \\
		\textbf{CASME II}~\cite{yan2014} & \textbf{26 valid} & \textbf{247 / 255} & \textbf{200} & \textbf{280 $\times$ 340} & \textbf{Yes} & \textbf{Yes} & \textbf{5 (7 raw)} \\
		SAMM~\cite{davison2018}          & 32 & 159             & 200 & 400 $\times$ 400 & Yes & Yes & 7 \\
		CAS(ME)$^2$~\cite{qu2016}        & 22 & 53 (+250 macro) & 30  & ---              & Yes & Yes & 4 \\
		\hline
		\multicolumn{8}{l}{\textit{Ethnicities / mean age (SD):} SMIC 3 / 26.7 \quad CASME 1 / 22.03 (1.60) \quad \textbf{CASME II 1 / 22.03 (1.60)}} \\
		\multicolumn{8}{l}{\quad SAMM 13 / 33.24 (11.32) \quad CAS(ME)$^2$ 1\textsuperscript{*} / 22.59 (2.20)} \\
		\hline
	\end{tabular}
	\caption[Posed and spontaneous micro-expression corpora]{Posed and spontaneous micro-expression corpora, compiled from Yan et al.~\cite{yan2014} (Table 1), Davison et al.~\cite{davison2018} (Table 1), Li et al.~\cite{liX2013} and Qu et al.~\cite{qu2016}. \textdagger{}~Described in Yan et al.; the primary papers lie outside the review corpus and are not cited separately. \textsuperscript{*}~Inferred from the recruitment description; Qu et al. do not state it. SMIC's demographic figures are those of all 20 recruited participants; the 16 who produced usable clips have a mean age of 28.1 and span two ethnicities.}
	\label{tab:spontaneous-corpora}
\end{table}

\subsection{CASME II: construction, apparatus, and annotation}
\label{sec:casme2-construction}

CASME II~\cite{yan2014} is the corpus on which every experiment in this thesis is performed, and several of its specifics propagate into the data pipeline used here.

\textbf{Participants and apparatus.} Thirty-five participants were recruited, mean age 22.03 years (SD = 1.60), and watched high-valence video episodes pre-rated by an independent panel of twenty raters, under instruction to suppress all facial movement. Recording sustained \textbf{200 fps} at 640~$\times$~480 with no inter-frame compression, under four LED lamps with umbrella reflectors giving steady flicker-free illumination --- a deliberate correction of the artefacts observed in SMIC. The facial region is approximately \textbf{280 $\times$ 340 pixels}, over three times CASME's facial area and roughly double SMIC's.

\textbf{Selection and annotation.} From nearly 3,000 elicited movements a small subset was retained, very subtle movements being discarded because their onset and offset frames could not be reliably located. Each retained sample carries onset and offset frame indices, Facial Action Coding System (FACS) action-unit labels from two independent coders with disagreements arbitrated by discussion, and an emotion label. Yan et al.~\cite{yan2014} argue against forcing spontaneous micro-expressions into the six basic categories, since a single action unit is frequently ambiguous --- AU4 (brow lowerer) may indicate disgust, anger, attention or tension --- and assign the emotion label from a \textit{combination} of the coded AUs, the participant's self-report, and the emotional content of the eliciting episode.

\textbf{Preprocessing in the original release.} The baseline applies a three-step normalisation from which this thesis departs: a frontal neutral model face is chosen; 68 Active Shape Model landmarks are detected on the \textit{first frame only} of each clip and registered to the model face by a Local Weighted Mean transform; that transform is applied unchanged to every frame, and the face is cropped by eye position. Yan et al.~\cite{yan2014} justify landmarking only the first frame on two grounds --- that rigid head motion within so short a window is negligible, and that per-frame ASM landmarks are unstable enough to introduce jitter where the face did not in fact move.

\textbf{A note on sample count.} The published paper reports \textbf{247} micro-expressions from 26 valid participants; the coding file distributed with the released database contains \textbf{255}. This discrepancy, a consequence of post-publication revision of the coding, is rarely acknowledged and is one reason sample counts reported for ``CASME II'' vary between papers. The present work uses the released 255-sample coding.

\subsection{The label taxonomy, and the sparsity problem it creates}
\label{sec:label-taxonomy}

Yan et al.~\cite{yan2014} provide five main emotion categories with AU-based criteria, and are candid about the resulting imbalance: ``some types of facial expressions are difficult to elicit in laboratory situations, thus the samples in different categories distributed unequally.'' The released coding is finer-grained still, exposing seven raw labels; counting directly from the label table used in this thesis (\texttt{Processed\_Data/master\_thesis\_labels.csv}, 255 rows) gives the middle columns of \autoref{tab:casme2-labels}.

\begin{table}[htbp]
	\centering
	\footnotesize
	\begin{tabular}{l l c c l c}
		\hline
		\multicolumn{3}{c}{\textbf{Five categories, as published}} & \multicolumn{2}{c}{\textbf{Released coding (255)}} & \textbf{Grouped} \\
		\textbf{Emotion} & \textbf{AU criterion} & \textbf{N} & \textbf{N} & \textbf{Raw label} & \textbf{class} \\
		\hline
		Happiness  & AU6 or AU12                           & 33  & 32 & happiness  & Positive \\
		Disgust    & AU9, AU10, or AU4+AU7                 & 60  & 63 & disgust    & Negative \\
		Surprise   & AU1+2, AU25, or AU2                   & 25  & 25 & surprise   & Surprise \\
		Repression & AU15 or AU17, alone or combined       & 27  & 27 & repression & Negative \\
		Others     & other emotion-related facial movement & 102 & 99 & others     & \textit{discarded} \\
		---        & ---                                   & --- & 7  & sadness    & Negative \\
		---        & ---                                   & --- & 2  & fear       & Negative \\
		\hline
		\multicolumn{6}{l}{\textbf{Working set:} Negative $63+27+7+2 = $ \textbf{99}; Positive \textbf{32}; Surprise \textbf{25}; \textbf{N = 156}} \\
		\hline
	\end{tabular}
	\caption[CASME II label taxonomy and the three-class working set]{The CASME II label taxonomy at three granularities: the five published categories with their AU criteria and frequencies, after Yan et al.~\cite{yan2014} (Table 3); the seven raw labels of the released coding as parsed for this thesis; and the three-class valence grouping that yields the 156-clip working set used throughout.}
	\label{tab:casme2-labels}
\end{table}

\textbf{Seven-class classification is not viable on this corpus.} Fear is represented by two clips in the entire database and sadness by seven. Under any subject-disjoint protocol --- and \autoref{sec:original-baseline} shows subject-disjointness has been the norm for this benchmark since its introduction --- a two-clip class cannot be meaningfully trained or tested: a fold either contains the class or does not, and per-class metrics degenerate.

The field's response has been to group by affective valence. The convention originates with SMIC, released with exactly three classes --- positive, negative, surprise~\cite{liX2013} --- and was carried into cross-corpus work by the Second Facial Micro-Expression Grand Challenge (MEGC 2019)~\cite{see2019}. MEGC states the scheme twice and not identically: as a challenge-wide rule in prose, with Negative covering Repression, Anger, Contempt, Disgust, Fear and Sadness, and again as per-corpus footnotes under its sample table, where CASME II's Negative narrows to Disgust and Repression alone. \autoref{tab:casme2-labels} follows the challenge-wide rule.

\begin{figure}[htbp]
	\centering
	\includegraphics[width=\textwidth]{fig3_1_dataset.png}
	\caption[Composition of CASME II]{Composition of CASME II. Left: all 255 released clips by raw label, with the excluded ``others'' category in grey. Right: the three-class pool of 156 clips used throughout this thesis, in the ratio 99 : 32 : 25; the grey ``others'' bar is shown again for scale and is not part of the pool.}
	\label{fig:dataset-composition}
\end{figure}

\textbf{A declared divergence from the MEGC convention.} MEGC 2019 reports the CASME II portion of its composite database as \textbf{145} samples and footnotes that ``Negative class of CASME II: Disgust and Repression''~\cite{see2019} --- building its subset from the paper's \textit{five-category} release, in which sadness and fear are absorbed into ``others''. This thesis uses the released seven-label coding and assigns sadness and fear to Negative, following the mapping MEGC gives for \textbf{SAMM} rather than the one it gives for CASME II. The consequence is an \textbf{eleven-clip} difference: N = 156 against N = 145, with a Negative class of 99 rather than 88. Nine of the eleven are the sadness (7) and fear (2) clips the seven-label coding exposes; the remaining two are unaccounted for, since disgust and repression together number 90 in the released coding against the 88 the challenge reports --- a residual noted rather than resolved. The direction is mild: the nine added clips are all in the majority class and drawn from the categories the coders found hardest to assign, so if anything they make the present task marginally harder than the challenge's.

The residual imbalance is severe in either case --- 99 : 32 : 25, roughly 4 : 1.3 : 1 --- and everything in \autoref{sec:corpus-threats} concerning metric choice, sampling and loss follows from it.

\subsection{The original baseline, and what it does and does not establish}
\label{sec:original-baseline}

Yan et al.~\cite{yan2014} provide a baseline using Local Binary Patterns from Three Orthogonal Planes (LBP-TOP)~\cite{zhaoG2007} with a Support Vector Machine. LBP-TOP extends the static LBP operator by extracting codes on the XY, XT and YT planes of the video volume and concatenating the three histograms, so that the XT and YT planes encode dynamic rather than static texture. The baseline divides the face into 5~$\times$~5 blocks, sets $P = 4$ neighbouring points, and sweeps the spatial radii $R_X, R_Y \in \{1,2,3,4\}$ and temporal radius $R_T \in \{2,3,4\}$; $R_T = 1$ is deliberately excluded because at 200~fps the change between adjacent frames is negligible --- an early acknowledgement that a high sampling rate is not automatically an advantage. Five-class accuracies under \textbf{leave-one-subject-out cross-validation} cluster in the high fifties to low sixties, the best configuration ($R_X = R_Y = 1$, $R_T = 4$) reaching \textbf{63.41\,\%}.

Three observations matter here.

\textbf{First, leave-one-subject-out is native to this benchmark.} It is not a stricter protocol adopted later by cautious authors; it is the protocol of the paper that released the corpus. Any evaluation reporting a single random train/test split departs from the benchmark's own convention.

\textbf{Second, the baseline is reported as accuracy alone}, with no per-class or macro-averaged score. On a five-class problem in which ``others'' accounts for 102 of 247 samples, accuracy is a weak instrument: the majority-class-constant classifier scores 41\,\% without modelling anything. In the three-class working set the floor is higher still --- always predicting Negative yields \textbf{0.6346 accuracy} at a macro F1 of only \textbf{0.2588}. A headline accuracy near 63.41\,\% is therefore not, on its own, evidence of a working system.

\textbf{A caution on this figure.} Yan et al.~\cite{yan2014} state in both their abstract and their results prose that leave-one-subject-out was applied to obtain 63.41\,\% --- but the caption of the very table reporting it reads ``leave-one-out cross-validation'', without the subject qualifier. Li et al.~\cite{liX2018}, reproducing the same figure in their comparison table, mark it with an asterisk denoting ``results achieved using leave-one-sample-out cross validation'' --- a materially easier protocol. The sources therefore disagree on the protocol behind the most widely quoted number on this benchmark, and the originating paper disagrees with itself. This thesis follows that paper's prose, but records the discrepancy rather than resolving it, because it affects how the figure should be read in any comparison.

\textbf{Third, the baseline's sensitivity to a purely descriptive hyper-parameter is itself informative.} Accuracy moves by several points across the radius sweep with no change to the data, the split or the classifier. Where effect sizes of that magnitude arise from descriptor settings, differences of similar magnitude between architectures cannot be interpreted without a variance estimate --- a limitation this thesis inherits and declares.

\subsection{Sibling corpora, and why this study remains single-corpus}
\label{sec:sibling-corpora}

Three further databases define the context in which CASME II is used; \autoref{tab:spontaneous-corpora} gives their properties, and each isolates a specific limitation of CASME II.

\textbf{SMIC~\cite{liX2013}}, the first spontaneous micro-expression database, is the origin of the three-class taxonomy. Its high-speed subset of 164 clips is distributed 51 positive, 70 negative, 43 surprise, and its baseline --- LBP-TOP with a temporal interpolation model (TIM) normalising each clip to ten frames, under leave-one-subject-out --- reaches \textbf{48.78\,\%}. Two elements of its design are inherited directly here: the three-class valence grouping, and temporal interpolation to normalise clips of unequal length. It is not used as a corpus because it is neither FACS-coded nor apex-annotated.

\textbf{SAMM~\cite{davison2018}} addresses the demographic narrowness of the CASME family, contributing 13 ethnicities and a mean age of 33.24 years against CASME II's single ethnicity and 22.03. Davison et al.~\cite{davison2018} observe that a corpus drawn from one ethnicity and one narrow age band restricts analysis to ``similar looking participants'', and tailor stimulus selection to each participant in advance rather than relying on post-hoc self-report. The design is stronger than CASME II's on both counts, but the corpus is smaller in usable samples and spreads them across seven categories with the same sparsity difficulty.

\textbf{CAS(ME)$^2$~\cite{qu2016}} provides 250 macro-expression and 53 micro-expression samples from the same participants under identical conditions, which makes it the natural resource for studying macro-to-micro domain transfer. That strategy produced the strongest entry in MEGC 2019 --- Liu et al.'s EMR~\cite{liu2019}, the only reported submission to use macro-expression data, though it drew that data from CK+ rather than from CAS(ME)$^2$~\cite{see2019}. Two properties limit its use here: 30~fps is an order of magnitude below CASME II's and poorly matched to sub-200-ms movements, and only 53 micro-expression samples exist. Its baseline is also evaluated leave-\textit{one-video}-out, which does not enforce subject-disjointness and is therefore not comparable to the protocol used here.

\textbf{Why this thesis uses CASME II alone.} It has the highest temporal resolution of any spontaneous corpus of comparable size, a precondition for magnification in a defined frequency band and for dense frame-to-frame flow; it is FACS-coded with onset, apex and offset annotation, which makes temporal normalisation to a fixed length well-defined rather than heuristic; and it remains the most widely reported single-corpus benchmark, which makes comparison against published results meaningful. The cost --- that no claim here extends to other populations, capture conditions or corpora --- is stated as a declared limitation rather than argued away.

\subsection{How CASME II is evaluated in the field: MEGC 2019 and the metric convention}
\label{sec:megc-metrics}

The most consequential development in how CASME II is \textit{used}, as opposed to how it was built, is the Composite Database Evaluation (CDE) protocol of the Second Facial Micro-Expression Grand Challenge~\cite{see2019}. This thesis adopts the challenge's metric definitions while declining its composite training regime, so the protocol warrants precise description.

CDE merges SMIC, CASME II and SAMM into one composite database under the three-class grouping of \autoref{sec:label-taxonomy}, giving \textbf{442 samples} --- Negative 250, Positive 109, Surprise 83 --- of which CASME II contributes 145. Leave-one-subject-out is run across the composite, which the organisers note ``ensures subject-independent evaluation''. Plain accuracy is rejected explicitly on grounds of class imbalance, and two balanced metrics mandated: \textbf{Unweighted F1 (UF1)}, the macro-averaged F1-score, which \textit{accumulates} true positives, false positives and false negatives per class \textbf{over all $k$ folds}, computes a per-class F1 from those pooled counts and averages without weighting; and \textbf{Unweighted Average Recall (UAR)}, the balanced accuracy, obtained by averaging per-class recall.

\[
	F1_c = \frac{2 \cdot TP_c}{2 \cdot TP_c + FP_c + FN_c}, \qquad \text{UF1} = \frac{1}{C}\sum_c F1_c
\]

UF1 is a \textbf{pooled} quantity: counts accumulate across folds \textit{before} any F1 is computed, which \autoref{sec:pooled-vs-fold} establishes is a different estimator from the average of per-fold macro F1 scores. The distinction is not cosmetic here. Because clips are distributed very unequally across subjects, and many single-subject folds contain only one of the three classes, a per-fold macro F1 is bounded well below 1 by fold composition alone and its average across folds bounded correspondingly (\autoref{fig:fold-composition}). The primary metric of this thesis, \textit{pooled macro F1}, is therefore identical in construction to MEGC's UF1; the per-fold average appearing in the raw experimental logs is not a metric the field uses at all.

\begin{figure}[htbp]
	\centering
	\includegraphics[width=\textwidth]{fig3_2_fold_composition.png}
	\caption[Class composition of the 25 leave-one-subject-out folds]{Class composition of the 25 leave-one-subject-out folds, and the per-fold macro-F1 ceiling it implies.}
	\label{fig:fold-composition}
\end{figure}

\begin{table}[htbp]
	\centering
	\begin{tabular}{l c c}
		\hline
		\textbf{Method} & \textbf{Full composite UF1 / UAR} & \textbf{CASME II subset UF1 / UAR} \\
		\hline
		LBP-TOP~\cite{zhaoG2007}                  & 0.5882 / 0.5785 & 0.7026 / 0.7429 \\
		Bi-WOOF~\cite{liong2018}                  & 0.6296 / 0.6227 & 0.7805 / 0.8026 \\
		OFF-ApexNet~\cite{liong2019a}             & 0.7196 / 0.7096 & 0.8764 / 0.8681 \\
		Quang et al.\textsuperscript{\textdagger} & 0.6520 / 0.6506 & 0.7068 / 0.7018 \\
		Zhou et al.\textsuperscript{\textdagger}  & 0.7322 / 0.7278 & 0.8621 / 0.8560 \\
		Liong et al.~\cite{liong2019b}, STSTNet   & 0.7353 / 0.7605 & 0.8382 / 0.8686 \\
		Liu et al.~\cite{liu2019}, EMR            & \textbf{0.7885 / 0.7824} & 0.8293 / 0.8209 \\
		\hline
	\end{tabular}
	\caption[MEGC 2019 Composite Database Evaluation results]{MEGC 2019 CDE results, reproduced from See et al.~\cite{see2019} (Table IV). \textdagger{} outside the review corpus; reported here only as it appears in that table.}
	\label{tab:megc-cde-results}
\end{table}

\textbf{One structural caveat governs these numbers.} Every row of \autoref{tab:megc-cde-results} comes from a model trained on the \textit{composite} database --- roughly 442 samples across three corpora and 68 subjects --- then scored on the CASME II portion of the pooled predictions. The CASME II column is therefore not a CASME II-only result: it reports how a cross-corpus model performs on CASME II clips, with approximately \textbf{2.8 times the training data} available in any single-corpus fold, where a single-corpus LOSO experiment of the kind conducted here trains each fold on roughly 150 clips. The two regimes are not interchangeable, and single-corpus results should not be compared directly against \autoref{tab:megc-cde-results} without stating the difference.

Two further observations from the challenge bear on this thesis's design. The top three submissions all chose \textbf{optical flow} as their feature representation, a convergence that supports the motion-based input used here; and one submission substituted mid-position frames for annotated apex frames with little loss, suggesting that precise apex localisation may be less critical than assumed and that methods consuming the full onset-to-offset sequence are not obviously disadvantaged.

\subsection{Threats to validity carried by the corpus}
\label{sec:corpus-threats}

The properties in \autoref{tab:corpus-threats} are not criticisms of CASME II, which remains the best available resource for this task. They are constraints any study built on it inherits, and each is answered by a specific design decision --- which is the sense in which, on a corpus of this size, several choices usually presented as modelling decisions are determined by the data. Fold and duration statistics are computed directly from the 156-clip label table used in this thesis.

\begin{table}[htbp]
	\centering
	\scriptsize
	\begin{tabular}{@{}p{0.3cm} p{2.0cm} p{5.3cm} p{5.6cm}@{}}
		\hline
		& \textbf{Threat} & \textbf{What the corpus contains} & \textbf{Design decision it forces} \\
		\hline
		(a) & Scale
		& 156 usable clips across three classes --- small enough that parameter count is itself a risk factor.
		& \textbf{Parameter-free or low-parameter components} wherever a choice exists, motivating SimAM~\cite{yang2021}, which adds zero learnable parameters; shallow rather than deep spatial backbones, extending the model-shrinking argument of Xia et al.~\cite{xia2020b}, which they establish for the composite-database setting rather than a single small corpus. \\
		(b) & Class imbalance
		& 99 : 32 : 25. The majority-class-constant classifier attains 0.6346 accuracy at a macro F1 of 0.2588.
		& Minority oversampling via a balanced sampler with focal loss, inverse-frequency class weighting \textbf{disabled}, since both corrections together induced single-class collapse. Accuracy is never reported alone, and the always-Negative floor is quoted alongside every result. \\
		(c) & Subject-level non-uniformity
		& \textbf{Subject 17 alone contributes 33 clips, 21\,\% of the working set}; subjects 8, 10 and 21 contribute one each, so fold sizes vary by a factor of 33. \textbf{Ten of the 25 folds contain a single class}, eight contain two, only seven all three.
		& \textbf{Full 25-fold leave-one-subject-out}, every subject held out exactly once, no clip untested, subject-disjointness guaranteed by construction. Primary metric is \textbf{pooled macro F1}, identical in construction to MEGC's UF1~\cite{see2019}; per-fold averaged macro F1 is rejected, since fold composition bounds it below the study's own target. \\
		(d) & A subject that disappears
		& Of the 26 participants, subject 18's three clips are all ``others'' and are removed by the three-class filter.
		& The LOSO run has \textbf{25 folds, not 26} --- stated explicitly for reproducibility, since fold counts differ between papers using this corpus. \\
		(e) & Demographic homogeneity
		& All participants Chinese, mean age 22.03 (SD 1.60) --- named by Davison et al.~\cite{davison2018} as a principal motivation for SAMM. The MEGC organisers saw the effect in reverse: adaptation from CK+ helped on SMIC and SAMM but ``not so for CASME II, which contain predominantly Chinese subjects''~\cite{see2019}.
		& Input is a \textbf{motion representation} --- dense optical flow $(u,v)$ plus optical strain --- rather than raw pixels, since motion is invariant to identity and illumination whereas appearance is exactly what a 26-subject corpus would let a network memorise. \\
		(f) & Laboratory ecology
		& Frontal recording, fixed flicker-free lamps, participants told not to move and not speaking; Yan et al.~\cite{yan2014} flag ``natural conversation and interaction'' as future work. Almost none of the head motion, illumination variation, occlusion or speech-related movement a deployed system would face is present.
		& No cross-dataset or in-the-wild claim is made anywhere in this thesis; generalisation is declared as further work. \\
		(g) & Annotation dependence
		& Labels derive from coded AUs, self-report and stimulus content --- a construct produced by trained coders under a rubric, with disagreements arbitrated by discussion rather than reported as residual uncertainty.
		& The coding is taken as given; its construct status is recorded as a declared limitation rather than treated as measurement error. \\
		(h) & Duration heterogeneity
		& Onset-to-offset lengths range from \textbf{31 to 126 frames}, median 66 --- at 200~fps, \textbf{0.155\,s to 0.63\,s}, median 0.33\,s. The upper end exceeds the 1/2-second criterion of \autoref{sec:phenomenon}.
		& Every clip temporally resampled to a fixed \textbf{T = 32} frames, following the TIM precedent of SMIC's baseline~\cite{liX2013}. \\
		(i) & Version ambiguity
		& Three sample counts circulate under the name ``CASME II'' --- 247 (the paper), 255 (the released coding), 145 (the MEGC subset) --- and both five- and seven-class label sets are in use, rarely stated.
		& The released 255-clip coding is used, N is reported explicitly alongside every figure, and the eleven-clip divergence from the MEGC subset is declared in \autoref{sec:label-taxonomy}. \\
		\hline
	\end{tabular}
	\caption[Threats to validity carried by CASME II and the decisions they force]{Threats to validity carried by CASME II, and the design decision each forces.}
	\label{tab:corpus-threats}
\end{table}

\subsection{Implications for this research}
\label{sec:corpus-implications}

Three further corpus properties are not threats but still determine settings. The \textbf{200 fps} capture rate fixes the Eulerian magnification pass-band at \textbf{5--25 Hz}, matched to the recording rate. The \textbf{seven raw labels}, with fear at two clips and sadness at seven, force the three-class valence grouping of \autoref{sec:label-taxonomy}. And the baseline's \textbf{sensitivity to a descriptive hyper-parameter} (\autoref{sec:original-baseline}) sets the standard of evidence: effect sizes below the resolution of a 156-clip corpus are reported as unresolved rather than as findings, and no confidence interval is available, since the single-seed design and unretained per-clip predictions preclude one.

Three consequences then shape the interpretation of the results rather than merely the construction of the pipeline. \textbf{First, the corpus determines the metric, and the metric determines the conclusion:} the fold composition of \autoref{tab:corpus-threats}(c) caps a macro F1 averaged over folds by arithmetic alone, before any classifier is considered, so reporting the pooled quantity is not a stylistic preference but the correction MEGC 2019 already standardised. \textbf{Second, a single held-out split is not a reliable measuring instrument here:} with 25 subjects distributed as unevenly as \autoref{tab:corpus-threats}(c) describes, which subjects fall in the test set can move the result by more than any architectural change under study, so only a complete leave-one-subject-out run is reported and the incomplete evaluations that preceded it are set out with their exclusion reasons rather than suppressed. \textbf{Third, the small-sample regime predicts which components should help:} 156 clips cannot supply the data needed to learn a general spatio-temporal feature extractor from scratch, particularly once the input has been reduced to flow and strain --- that is, once much of that extraction has been performed analytically before the network sees the data --- whereas a component modelling the temporal arc of an expression exploits exactly the structure the corpus's onset/apex/offset annotation guarantees is present and correctly delimited.

\subsection{Limitations of the existing corpus literature, and how this work differs}
\label{sec:corpus-lit-limitations}

The corpus literature establishes three things well and one thing poorly. It establishes the phenomenon's definitional boundaries, an elicitation methodology producing genuinely spontaneous samples, and a capture standard --- 200 fps, flicker-free illumination, FACS coding with onset/apex/offset annotation --- that makes fine-grained motion analysis feasible. What it does not establish is a stable basis for comparing methods. The original baseline reports accuracy only, on a five-class problem whose majority class is 41\,\% of the data, and its own hyper-parameter sweep moves the headline figure by several points with no change to the model. Sample counts, label granularity and class groupings differ across papers using the same database --- 247, 255 and 145 all appear under the name ``CASME II'' --- and are frequently not stated. MEGC 2019's CDE protocol fixed the metric problem convincingly, but within a composite training regime that makes its per-database columns non-comparable to single-corpus work, a distinction subsequent papers routinely elide. Above all, only one study in this group quantifies how much of a reported difference between methods is attributable to the evaluation protocol rather than the methods themselves --- Li et al.~\cite{liX2018} add a single matched row of their own method under both protocols --- and none does so across methods, despite a subject distribution that makes the question unavoidable.

This thesis differs in three respects. It runs the \textbf{complete} 25-fold leave-one-subject-out protocol over \textbf{every} valid cell of a four-component ablation matrix rather than reporting a single proposed configuration, so each component's contribution is measured from matched pairs differing in exactly one factor. It adopts MEGC's pooled macro F1 as its primary metric while explicitly declining the composite training regime, so its numbers describe single-corpus performance and are labelled as such. And it reports one complete protocol rather than the strongest of several, setting out which earlier evaluations were run and why each is excluded rather than suppressing that history. The intended contribution is therefore not a recognition result on CASME II, but a measurement of how much of such a result is attributable to each stage of the pipeline that produces it.

The 200 fps capture rate established here is what lets Eulerian Video Magnification, reviewed next, target a matched temporal band.


% =============================================================================
\section{Motion Magnification: Eulerian Video Magnification as a Pre-Processing Stage}
\label{sec:evm-review}
% =============================================================================

This section reviews Eulerian Video Magnification (EVM) through the four review-corpus papers that apply it: Bai et al.~\cite{bai2021}, Li et al.~\cite{liX2018}, and Li, Huang and Zhao~\cite{liY2018,liY2021}. The original formulations they build on --- Wu et al.'s amplitude-based method and Wadhwa et al.'s phase-based method --- lie outside the corpus and are described only as attributed through those four papers.

\subsection{The problem magnification addresses}
\label{sec:evm-problem}

The corpus design of \autoref{sec:why-spontaneous} responds to the \textit{short-duration} half of the micro-expression definition but not to the \textit{low-intensity} half. CASME II's 200~fps sampling spreads a movement completing in under 200~ms across dozens of frames~\cite{yan2014}, yet that movement may still occupy only a few grey levels of change across a small patch of skin --- and Yan et al.~\cite{yan2014} report discarding samples outright because the movement was too faint for coders to place an onset or offset. Li et al.~\cite{liX2018} state the difficulty directly: ``the intensity levels of facial movements are too low to be distinguishable.'' Their remedy is to amplify the motion itself before any descriptor is computed, so that the \textit{difference between expression categories is enlarged} rather than the classifier being made more sensitive. This is a data-level intervention: the model is untouched and the modification is entirely upstream, which is what makes EVM cleanly separable as an ablation factor in this thesis.

\subsection{How Eulerian magnification works, and the two variants}
\label{sec:evm-how}

\textbf{Amplitude-based magnification (AMM)} is the formulation set out in \autoref{sec:motion-magnification} --- Laplacian pyramid, temporal band-pass, multiplication by $\alpha$, reconstruction --- realising its band-pass step with a \textbf{Butterworth} filter. It is the variant implemented in this thesis. \textbf{Phase-based magnification (PMM)}, of which Bai et al.~\cite{bai2021} give the clearest account in the corpus, decomposes each frame with \textbf{octave complex steerable pyramids} over four orientations and applies the temporal filter to the \textit{local phase} rather than the amplitude. Bai et al.~\cite{bai2021} explain the advantage in signal terms: the steerable pyramid ``has impulse response with spatial support,'' which makes it easier to isolate the intended temporal frequencies while suppressing the remainder. The practical consequences they report are that the phase-based method \textbf{supports larger amplification factors} and is \textbf{notably less sensitive to noise}.

Three parameters are therefore free in any EVM application, and all three recur below: the \textbf{temporal pass-band}, which decides \textit{which} motions are amplified; the \textbf{magnification factor $\alpha$}, which decides \textit{how much}; and the \textbf{pyramid depth}, which decides at which spatial scales amplification occurs.

\textbf{Setting the temporal band.} Bai et al.~\cite{bai2021} report the rule given by the original EVM authors: the pass-band should correspond to the duration of the phenomenon being amplified --- for micro-expressions, 1/5\,s to 1/25\,s, equivalently \textbf{5 Hz to 25 Hz}. They also report that applying this full band produced ``noticeable noise'' in their setting, and that they narrowed it to 15--25 Hz accordingly. This is the only explicit guidance on band selection anywhere in the review corpus, and it is the basis for the band used in this thesis.

\subsection{Evidence that magnification improves recognition}
\label{sec:evm-evidence}

The strongest empirical case in the review corpus is Li et al.~\cite{liX2018}, who sweep the magnification factor systematically rather than fixing it. Clips are magnified at what the paper describes as ten levels while listing nine values, $\alpha \in \{1, 2, 4, 8, 12, 16, 20, 24, 30\}$ --- an inconsistency in the source rather than in this account --- where $\alpha = 1$ denotes no magnification. Each magnified clip is temporally interpolated to ten frames (TIM10), and LBP, HOG and HIGO descriptors are extracted and classified, on both SMIC and CASME II under leave-one-subject-out. Two findings emerge.

\textbf{Magnification helps, consistently}, ``for all three kinds of features on all four testing datasets''. On CASME II the effect is substantial (\autoref{tab:evm-effect}), and the best configuration of their full framework --- magnification, TIM10 and HIGO --- reaches 67.21\,\% on CASME II and 68.29\,\% on SMIC-HS under LOSO, competitive with or better than the contemporaneous state of the art.

\begin{table}[htbp]
	\centering
	\begin{tabular}{l c c c}
		\hline
		\textbf{Descriptor} & \textbf{No magnification} & \textbf{With magnification} & \textbf{$\Delta$} \\
		\hline
		LBP  & 55.87\,\% & 60.73\,\% & \textbf{+4.86} \\
		HOG  & 57.49\,\% & 63.97\,\% & \textbf{+6.48} \\
		HIGO & 57.09\,\% & 67.21\,\% & \textbf{+10.12} \\
		\hline
	\end{tabular}
	\caption[Effect of Eulerian motion magnification on CASME II]{Effect of Eulerian motion magnification on CASME II under leave-one-subject-out, after Li et al.~\cite{liX2018} (Table 6).}
	\label{tab:evm-effect}
\end{table}

\textbf{The relationship between $\alpha$ and accuracy is non-monotonic.} Li et al.~\cite{liX2018} describe the resulting curves as ``rainbow-shaped,'' with best performance ``generally achieved when the motion is magnified in the range of $[8, 16]$,'' and give a mechanistic account of both failure modes: ``magnification at lower levels might not be enough to reveal the ME motion progress; on the other hand, magnification at higher levels degrades the performance because too many artifacts are induced.'' Magnification is therefore not a monotone improvement to be maximised but a parameter with an interior optimum that must be selected.

Bai et al.~\cite{bai2021} provide the second quantitative demonstration using a fundamentally different downstream model --- a pre-trained VGGFace2 (ResNet-50) spatial encoder feeding a bidirectional LSTM. On SMIC they reach \textbf{75.76\,\% with phase-based magnification against a 60.60\,\% baseline} drawn from an earlier VGGFace2 + Bi-LSTM publication rather than from a magnification-off arm of their own, and the phase-based variant outperforms the amplitude-based one. Because the pipeline shares nothing with Li et al.'s hand-crafted descriptors, the two results together establish that the benefit is not an artefact of a particular feature family --- at least when the downstream representation is appearance-based.

\subsection{Magnification in apex-frame pipelines, and the $\alpha$ discrepancy}
\label{sec:evm-apex}

Two further papers use EVM in a different regime, and the contrast is informative. Li, Huang and Zhao~\cite{liY2018} address whether a micro-expression can be recognised from a \textit{single} apex frame, arguing for magnification specifically in that setting: the apex frame carries the most emotional information, but ``the intensity level of apex frame is not distinguished due to subtle change, especially for deep model that describes high-layer abstract information.'' They magnify the spotted apex frame with \textbf{$\alpha = 30$} and fine-tune a VGG-Face descriptor on the result; because a single frame per clip is too little data to train on, they exploit the high frame rate to augment, retaining the apex frame plus its two preceding and two following frames and so expanding the training set fivefold. They restate the trade-off Li et al.~\cite{liX2018} quantified --- ``bigger values of motion amplification level lead to larger scale of motion amplification, but also can cause bigger displacement and artifacts.'' Li, Huang and Zhao~\cite{liY2021} extend this into a joint local--global learning framework with frequency-domain apex detection, carrying the same setting forward --- magnification level ``set as 30 in our framework'' --- with the same $\pm$2-frame augmentation producing what they name the Extended Magnified ME (EMME) database.

\textbf{The $\alpha$ discrepancy is easy to misread as a contradiction.} Li et al.~\cite{liX2018} find the optimum at $\alpha \in [8, 16]$ and observe degradation beyond it; Li, Huang and Zhao~\cite{liY2018,liY2021} adopt $\alpha = 30$, at the extreme end of that same sweep. The regimes differ in a way that plausibly explains the gap. The $[8, 16]$ optimum is measured on \textbf{full magnified sequences} described by spatio-temporal descriptors, where inter-frame artefacts accumulate along the temporal axis and corrupt exactly the dynamic texture the descriptor encodes. The $\alpha = 30$ setting is applied to a \textbf{narrow five-frame window around the apex}, consumed by a deep face encoder as largely static appearance, where temporal artefacts have far less opportunity to express themselves --- and where the encoder needs a strong appearance signal to distinguish categories at all. The implication for any new system is that $\alpha$ cannot be transferred between pipelines: it is a parameter of the \textit{combination} of magnification and downstream representation, not of the corpus.

\subsection{Known limitations of magnification}
\label{sec:evm-limitations}

The review corpus is unusually candid about EVM's failure modes, and four are documented.

\textbf{It amplifies everything in band, not only expressions.} This is the most important limitation and the most easily overlooked. Li et al.~\cite{liX2018} decline to use magnification for micro-expression \textit{spotting} at all, ``because it magnifies unwanted motions (e.g. head movements) at the same time.'' EVM has no notion of which motion is of interest: a rigid head translation whose temporal frequency falls inside the pass-band is amplified exactly as a brow raise is. In a recognition pipeline this is mitigated by prior face alignment and cropping --- Li et al.~\cite{liX2018} magnify only after normalising and cropping to the eye-defined rectangle --- but mitigated, not eliminated.

\textbf{Artefacts grow with $\alpha$.} All four papers report this: Li et al.~\cite{liX2018} as the descending limb of the rainbow curve, Li, Huang and Zhao~\cite{liY2018,liY2021} as a design caution, and Bai et al.~\cite{bai2021} directly, as noise that forced them to narrow their temporal band.

\textbf{There is no principled procedure for choosing $\alpha$.} One value in the corpus is obtained by exhaustive sweep~\cite{liX2018}; the apex-frame papers assert theirs --- Li et al.~\cite{liY2018} state $\alpha = 30$ without support, and Li et al.~\cite{liY2021} carry the value forward by citing them. No paper offers a way to select $\alpha$ from properties of the data.

\textbf{Evaluation is confined to appearance-based downstream representations:} LBP, HOG and HIGO in Li et al.~\cite{liX2018}, VGGFace2 features in Bai et al.~\cite{bai2021}, VGG-Face in both Li, Huang and Zhao papers. Every one of these descriptors is computed on \textbf{magnified pixel intensities}. No study in the review corpus measures what magnification contributes when the representation handed to the classifier is itself an explicit motion field rather than an appearance signal.

\subsection{Implications for this research}
\label{sec:evm-implications}
% sec:evm-motion-hypothesis is referenced from Chapter 5 (chapter5.tex:236). The
% hypothesis it names now opens this subsection instead of having one of its own,
% so the label is kept here and resolves to this subsection.
\label{sec:evm-motion-hypothesis}

\autoref{tab:evm-settings} sets out the settings the review determines. Behind them lies a hypothesis the review corpus does not test, stated here because it frames the interpretation of the measurement. EVM amplifies the \textit{amplitude} of intensity variation in a temporal band, and an appearance descriptor computed on magnified frames sees a directly stronger signal. An optical flow field is a different object: flow estimates \textit{displacement}, and amplifying the underlying intensity change by $\alpha$ does not simply scale the recovered displacement by $\alpha$ --- it changes the conditioning of the estimation problem, potentially making faint motion recoverable where it previously was not, while introducing displacement artefacts of the kind all four reviewed papers warn about. Any residual global scaling of flow magnitude is further attenuated by the per-channel normalisation a learned pipeline applies before training. The plausible outcome is that magnification in front of a flow-and-strain representation retains the part of its benefit that consists of \textit{making otherwise-unrecoverable motion recoverable} while losing the part that consists of \textit{increasing contrast for an appearance descriptor} --- so its measured contribution should be smaller on average than the roughly +5 to +10 percentage points of \autoref{tab:evm-effect}, and concentrated in configurations whose downstream model can exploit finer spatial deformation. EVM is measured across six matched pairs of configurations differing in nothing but the tensor directory, which is what allows that prediction to be tested.

\begin{table}[htbp]
	\centering
	\footnotesize
	\begin{tabular}{@{}p{6.4cm} p{7.4cm}@{}}
		\hline
		\textbf{Finding from the review} & \textbf{Decision taken in this thesis} \\
		\hline
		Pass-band should match the phenomenon's duration, 1/5--1/25\,s $\equiv$ 5--25 Hz~\cite{bai2021}, though Bai et al. found that band noisy and narrowed it to 15--25 Hz & Temporal band-pass set to \textbf{5--25 Hz}, the unnarrowed duration rule; Bai et al.'s narrowing was not followed, and this is a declared divergence \\
		Accuracy against $\alpha$ is rainbow-shaped, optimum in \textbf{$[8, 16]$} for full-sequence pipelines~\cite{liX2018} & \textbf{$\alpha = 10$}, inside the reported optimum; the $\alpha = 30$ of the apex-frame papers is explicitly \textit{not} adopted, since this thesis consumes full sequences \\
		Amplitude-based magnification uses a Laplacian pyramid plus temporal band-pass~\cite{bai2021} & Amplitude-based EVM over a \textbf{4-level Laplacian pyramid}, temporal filtering by FFT band-pass \\
		Magnification amplifies head motion as well as expression~\cite{liX2018} & Magnification applied only after the clip has been trimmed onset-to-offset; the risk is acknowledged as a residual threat rather than claimed to be eliminated \\
		Magnification is a purely upstream, data-level operation & Stage 1 is run \textbf{twice}, producing a magnified and a non-magnified tensor set; the EVM switch selects a directory and leaves the network unchanged, making it a clean ablation factor \\
		Magnification is part of the proposed data pipeline, not of the network & EVM is treated as the \textbf{baseline condition} (\texttt{config\_4}), not as a component to be added; the non-magnified arm (\texttt{config\_1}) is the control that isolates it \\
		\hline
	\end{tabular}
	\caption[Review findings and the magnification settings they determine]{Review findings and the magnification settings they determine.}
	\label{tab:evm-settings}
\end{table}

\textbf{A defect found and repaired, which the review made detectable.} In this project's earliest evaluations, each magnification-enabled configuration produced results \textit{bit-identical} to its magnification-disabled counterpart --- identical accuracies, macro F1 scores and confusion matrices. Given that all four reviewed papers report clear and repeatable effects from magnification, an effect of exactly zero across six independent pairs is not a plausible null result but a symptom: both arms were reading the same tensor directory, so half of the ablation matrix was duplicated and the magnification hypothesis had never in fact been exercised. The defect was repaired before the reported run, and wiring the existing verification helper in as a start-up assertion is a declared item of further work.

\subsection{Limitations of the existing work, and how this study differs}
\label{sec:evm-gap}

What the reviewed literature does not establish is threefold. First, every evaluation in the corpus places magnification in front of an \textbf{appearance-based} representation, so the reported gains conflate two distinct mechanisms --- making faint motion recoverable, and increasing contrast for a texture or gradient descriptor --- and no experiment separates them. Second, the magnification factor is selected empirically in every case and transferred between studies without re-validation, despite the $\alpha$ discrepancy of \autoref{sec:evm-apex}. Third, and most consequentially, magnification is always evaluated as an \textit{addition to a fixed pipeline} --- a two-cell comparison, with and without --- so its interaction with the other components of a system is never measured.

This thesis differs on all three counts. It places magnification in front of an explicit \textbf{motion} representation --- dense optical flow and optical strain --- rather than an appearance representation, isolating the recoverable-motion mechanism from the appearance-contrast one. It selects $\alpha$ from the full-sequence optimum rather than the apex-frame value, with the choice justified rather than inherited. And it evaluates magnification as one factor in a \textbf{full factorial ablation}, so its effect is reported not as a single delta but as six matched-pair measurements whose spread across architectures is itself the finding.

What magnification is placed in front of therefore matters as much as its own parameters, and \autoref{sec:flow-review} turns to that motion representation.


\section{Motion Representation: Dense Optical Flow}
\label{sec:flow-review}

Dense optical flow supplies the horizontal and vertical displacement channels used in this thesis. The relevant literature concerns three choices: the estimator, the frame pairs it processes, and whether flow serves as a recognition feature or an alignment tool. These choices determine what information remains available to the learned model.

\subsection{Why motion rather than appearance}
\label{sec:why-flow}

Xu et al.~\cite{xu2017} motivate facial-dynamics descriptors by the difficulty of interpreting micro-expressions through texture alone. Flow describes apparent movement directly, whereas image intensity also records appearance and illumination. This makes displacement a useful representation of subtle facial activity, although it does not guarantee independence from identity or recording conditions: facial geometry, alignment and brightness changes still affect its estimation.

The benchmark evidence supports this choice. The MEGC 2019 organisers report that their three highest-ranked submissions used optical flow~\cite{see2019}. OFF-ApexNet~\cite{liong2019a} uses horizontal and vertical flow as neural-network inputs; STSTNet~\cite{liong2019b} adds strain; Zhao et al.~\cite{zhaoS2021} retain a sequence of flow fields. Their convergence motivates a motion-based input, but is not a controlled demonstration that it always outperforms appearance-based representations.

\subsection{Estimation assumptions}
\label{sec:brightness-constancy}

OFF-ApexNet~\cite{liong2019a} formulates flow using brightness constancy and small-motion assumptions. A local brightness constraint supplies one equation for two displacement components, so estimation also needs spatial regularisation or a local motion model. The output is a pair of fields, $u(x,y)$ and $v(x,y)$; Zhao et al.~\cite{zhaoS2021} additionally describe their magnitude, $\sqrt{u^2+v^2}$.

These assumptions matter more than repeating the derivation. Illumination changes can resemble movement, while smoothing can suppress a weak local contraction together with noise. Magnification can make subtle motion easier to estimate but can also amplify unwanted movement. Flow therefore transforms the measurement problem; it does not remove uncertainty in the source frames.

\subsection{Choice of estimator}
\label{sec:flow-estimator}

The reviewed methods use three estimators. Zhao et al.~\cite{zhaoS2021} adopt Farneb\"ack through its OpenCV implementation, citing its established availability. It estimates displacement through local polynomial approximations. OFF-ApexNet~\cite{liong2019a} instead uses total-variation regularisation with an L1 data term (TV-L1), citing preservation of motion discontinuities and robustness. Xu et al.~\cite{xu2017} use the method of Sun, Roth and Black. These algorithm attributions are taken from the reviewed papers; the original estimator publications are outside the supplied reference corpus.

STSTNet~\cite{liong2019b} specifies onset--apex optical flow but does not name its estimator. Its similarity to OFF-ApexNet is therefore insufficient to attribute TV-L1 to it. Likewise, the advantages given for TV-L1 should not be read as a measured recognition gain over Farneb\"ack on CASME II. The cited studies do not provide a controlled three-estimator comparison supporting such a ranking. This thesis adopts Farneb\"ack as a fixed preprocessing choice and leaves that comparison open.

\subsection{Onset--apex flow or a temporal sequence}
\label{sec:flow-what-between}

OFF-ApexNet~\cite{liong2019a} represents each clip by flow from onset to apex, where expression intensity is greatest, and resizes the two component maps to $28\times28$. STSTNet~\cite{liong2019b} uses the same frame-pair construction. This is an economical representation of the displacement from neutral to peak expression, suitable for a small training set. It omits the intermediate trajectory and the relaxation after the apex.

Zhao et al.~\cite{zhaoS2021} instead construct an eleven-frame sequence using onset, apex and offset annotations, then calculate ten adjacent-pair flow fields. This retains changes in motion across the expression. The distinction is between one displacement summary and a sequence whose order a temporal model can exploit; it is not simply a difference in input size.

The present study follows the sequence approach but uses uniform sampling of existing frames rather than Zhao et al.'s key-frame construction. This preserves a temporal input for the transformer ablation. It does not, by itself, measure the recognition cost of replacing a sequence with an onset--apex pair: that representation comparison is not an experimental arm.

\subsection{Flow for fine alignment}
\label{sec:flow-alignment}

Xu et al.~\cite{xu2017} also use flow to refine alignment. Their pipeline first locates landmarks and coarsely aligns the face, then performs fine in-sequence alignment using pixel-level movement. The motivation is specific to micro-expressions: a small registration error can be comparable to the facial displacement being recognised.

This use differs from supplying flow to a classifier. An uncorrected displacement field contains residual head movement as well as expression. Magnifying the frames beforehand can increase both. The present pipeline performs no flow-based fine alignment, so corpus preprocessing must not be treated as proof that this source of error has been eliminated.

\subsection{Limitations of the representation}
\label{sec:flow-limitations}

Flow remains sensitive to brightness-constancy violations, alignment errors and estimator settings. Apex-based systems additionally depend on selecting an informative frame. OFF-ApexNet~\cite{liong2019a} attributes weaker performance on SMIC to imprecise apex selection and cites an earlier apex-spotting study reporting an average discrepancy of thirteen frames. That reported average is neither a maximum error nor a measurement made by this thesis.

Using an onset--offset sequence avoids dependence on one selected apex, but still relies on clip boundaries and temporal sampling. Repeated sampled frames can yield zero-flow pairs, and large sampling intervals can make the small-motion approximation less appropriate. These limitations constrain the representation before any network component is compared.

\subsection{Implications for this research}
\label{sec:flow-implications}

\begin{table}[htbp]
    \centering
    \small
    \begin{tabular}{@{}p{0.29\linewidth} p{0.65\linewidth}@{}}
        \hline
        \textbf{Review finding} & \textbf{Consequence for this study} \\
        \hline
        Flow is a recognition input~\cite{liong2019a,see2019}
        & Horizontal and vertical displacement form the first two channels; the classifier receives motion tensors rather than raw frames. \\
        Sequence flow preserves temporal variation~\cite{zhaoS2021}
        & Adjacent-pair fields are retained across the sampled clip, allowing the temporal encoder to operate on an ordered sequence. \\
        Estimator choices differ~\cite{liong2019a,zhaoS2021,xu2017}
        & Farneb\"ack remains fixed; a matched comparison with TV-L1 is an untested extension. \\
        Fine alignment matters~\cite{xu2017}
        & No flow-based fine alignment is performed. It remains a priority for improving preprocessing. \\
        \hline
    \end{tabular}
    \caption[Optical-flow evidence and study decisions]{Optical-flow evidence and the decisions it motivates. Fixed preprocessing choices constrain the interpretation of the architectural ablation.}
    \label{tab:flow-decisions}
\end{table}

The central interpretive point is that optical flow is an \textit{analytical spatio-temporal feature extractor}: it estimates motion from paired frames before the network is trained. A learned spatial stem therefore acts on a representation that already makes displacement explicit. With only 156 clips, learning additional spatial filters may offer limited benefit or duplicate structure already exposed by preprocessing. This is a hypothesis about the stem's marginal value, not proof that learned spatial processing is unnecessary.

A temporal encoder addresses a different task: modelling relations among the successive displacement fields. Pairwise flow estimation does not itself learn those sequence-level relations. Keeping these roles distinct is necessary when attributing performance to the stem and transformer.

\subsection{Research gap and scope}
\label{sec:flow-gap}

The relevant unresolved question is what learned spatial and temporal processing contributes \textit{after} motion has been extracted analytically. The thesis addresses this through matched architectural comparisons under complete leave-one-subject-out (LOSO) evaluation, retaining the same flow construction while switching the learned components. Magnification is varied separately through the paired preprocessing conditions.

The design neither ranks optical-flow estimators nor isolates alignment quality or onset--apex versus sequence representations. Those remain limitations rather than gaps claimed as solved. The resulting architectural conclusions apply to the implemented motion representation and evaluation setting; they do not establish a universal ordering of network families.


\section{Deformation Representation: The Optical Strain Tensor}
\label{sec:strain-review}

Optical strain supplies the third input channel. Its literature establishes a deformation descriptor derived from flow, two ways to use that descriptor, and limitations relevant to interpreting the architectural ablation. Shreve et al.~\cite{shreve2011}, the Liong strain series~\cite{liong2014a,liong2014b,liong2016}, and STSTNet~\cite{liong2019b} provide the principal evidence.

\subsection{What strain contributes}
\label{sec:strain-what-adds}

Flow describes displacement; strain describes its local spatial variation. A spatially uniform translation has zero displacement gradient, whereas neighbouring skin regions moving differently produce deformation. Shreve et al.~\cite{shreve2011} use this distinction to detect facial activity and describe robustness to moderate head translation. It does not imply immunity to arbitrary head movement or registration error.

Their illumination- and make-up-robustness statements are attributed to earlier studies rather than measured anew. They should therefore be read as reported motivation, not as evidence that the present pipeline is robust under those conditions. Strain can suppress a uniform displacement while still inheriting errors from the estimated flow.

\subsection{Tensor components and magnitude}
\label{sec:strain-tensor}

The relevant representation is the symmetric displacement-gradient tensor, with normal components $\varepsilon_{xx}$ and $\varepsilon_{yy}$ and shear components $\varepsilon_{xy}=\varepsilon_{yx}$. Shreve et al.~\cite{shreve2011} approximate spatial derivatives by central differences with offsets of approximately two to three pixels; Liong et al.~\cite{liong2014a} use one-pixel offsets. This discretisation is a modelling choice because it sets the scale at which local variation is measured.

Liong et al.~\cite{liong2014a,liong2014b,liong2016} use the four-component magnitude
\[
    \varepsilon_{\mathrm{F}}
    =\sqrt{\varepsilon_{xx}^{2}+\varepsilon_{yy}^{2}
        +\varepsilon_{xy}^{2}+\varepsilon_{yx}^{2}}
    =\sqrt{\varepsilon_{xx}^{2}+\varepsilon_{yy}^{2}
        +2\varepsilon_{xy}^{2}}.
\]
The scalar describes deformation intensity and discards its sign and direction. Supplying it together with $u$ and $v$ retains the displacement fields from which those spatial relations were calculated.

\subsection{The implemented magnitude and its divergence}
\label{sec:strain-divergence}

The implementation instead computes
\[
    \varepsilon_{\mathrm{os}}
    =\sqrt{\varepsilon_{xx}^{2}+\varepsilon_{yy}^{2}
        +\varepsilon_{xy}^{2}},
    \qquad
    \varepsilon_{xy}=\tfrac{1}{2}
        \left(\frac{\partial u}{\partial y}
             +\frac{\partial v}{\partial x}\right).
\]
It counts shear once rather than twice. NumPy's \texttt{gradient} supplies spatial differences at one-pixel spacing, using central differences in the array interior. The spacing follows Liong et al.~\cite{liong2014a}, but the scalar reduction differs from their convention.

STSTNet~\cite{liong2019b} cannot establish a clear precedent for the three-term form. Its printed magnitude equation repeats $\partial u$ in the mixed derivative where the symmetric tensor calls for $\partial v/\partial x$ alongside $\partial u/\partial y$. It also places the factor $\tfrac{1}{2}$ \emph{outside} the squared derivative sum. Correcting the repeated derivative while retaining that placement gives
\[
    \tfrac{1}{2}
       \left(\frac{\partial u}{\partial y}
            +\frac{\partial v}{\partial x}\right)^2
       =2\varepsilon_{xy}^{2},
\]
which agrees with the four-term convention. Moving the factor inside the square would instead give the implemented three-term form, but that is a further change to the printed expression. The intended implementation cannot be settled from this equation alone; the thesis therefore declares a departure from the unambiguous four-term convention rather than claiming to reproduce STSTNet's magnitude.

The four-term form weights pure shear by $\sqrt{2}$ relative to the implemented form. Independent min--max normalisation removes a uniform scale factor, but does not generally remove this difference when normal and shear contributions vary across pixels and time. Neither convention was compared experimentally here, so their practical equivalence is not assumed.

\subsection{Strain as a feature, a weight, or an input channel}
\label{sec:strain-feature-weight}

\textbf{Optical strain features (OSF)} use a temporally pooled strain map as a descriptor~\cite{liong2014a,liong2014b,liong2016}. Averaging accounts for sequence length: the early method already divides its accumulated map by the number of frame pairs~\cite{liong2014a}. The maps are filtered and normalised before classification. The early work explicitly applies Wiener and Gaussian filters to suppress noise~\cite{liong2014a}.

\textbf{Optical strain weighted features (OSW)} use pooled strain to weight the regional XY-plane histograms of local binary patterns on three orthogonal planes (LBP-TOP), emphasising active facial regions~\cite{liong2014b}. Liong et al.~\cite{liong2016} perform spatial mean pooling followed by temporal mean pooling and concatenate OSF and OSW for a combined descriptor. Strain then serves both as a feature and as a spatial weighting signal.

\textbf{STSTNet} supplies strain directly alongside horizontal and vertical onset--apex flow~\cite{liong2019b}. This is the precedent for the thesis's three-channel input structure, independently of the magnitude ambiguity. The present study retains all three channels across successive frame pairs rather than reducing each clip to one map.

\subsection{Recognition evidence and protocol limits}
\label{sec:strain-evidence}

\begin{table}[htbp]
    \centering
    \small
    \begin{tabular}{@{}lccc@{}}
        \hline
        \textbf{Method} & \multicolumn{2}{c}{\textbf{SMIC accuracy (\%)}} & \textbf{CASME II (\%)} \\
        & Micro-average & Macro-average & Accuracy \\
        \hline
        Baseline & 45.73 & 47.76 & 61.94 \\
        OSF & 41.46 & 46.00 & 51.01 \\
        OSW & 49.39 & 50.71 & 61.94 \\
        OSF + OSW & 52.44 & 58.15 & 63.16 \\
        OFF + OFW & 40.24 & 41.94 & 55.87 \\
        \hline
    \end{tabular}
    \caption[Recognition results for strain-based descriptors]{Recognition results reproduced from Liong et al.~\cite{liong2016}, Table 8. SMIC uses LOSO; CASME II uses leave-one-video-out. The summary combines block settings, so equal entries are not a matched-setting comparison. Macro-average here denotes subject-averaged accuracy, not macro F1.}
    \label{tab:strain-performance}
\end{table}

The CASME II evidence depends on how strain is used and on the comparator. OSF alone falls below the texture baseline. The equal baseline and OSW entries in \autoref{tab:strain-performance} reflect different block settings: under matched settings OSW improves accuracy by 0.81 percentage points. The combined OSF + OSW descriptor improves on the texture baseline by 0.41--1.22 points across block settings~\cite{liong2016}. These accuracy gains are not directly comparable with pooled macro F1, and leave-one-video-out does not enforce subject-disjoint testing.

Crucially, Liong et al.~\cite{liong2016} also test optical-flow features and weights (OFF + OFW) in the same descriptor framework. Their CASME II accuracy is 55.87\,\%, against 63.16\,\% for OSF + OSW, a 7.29-point difference. This is a matched flow-based control, so it would be incorrect to claim that the literature never isolates strain against flow. It concerns hand-crafted descriptors; it does not establish the gain from adding a strain channel to the learned sequence model used here.

\subsection{Limitations}
\label{sec:strain-limitations}

Spatial differentiation can amplify flow-estimation noise. The filtering used by Liong et al.~\cite{liong2014a,liong2016} is therefore relevant to a pipeline without a dedicated strain post-filter. A deterministic transformation of $u$ and $v$ adds no independent observation: its potential value lies in making deformation directly accessible to the classifier. Whether this representation helps must be measured rather than inferred from its construction.

The reviewed recognition methods also simplify time. OSF and OSW pool their sequences; STSTNet uses one onset--apex pair. Shreve et al.~\cite{shreve2011} retain temporal strain activity for spotting, which localises expressions rather than classifying their emotion. The temporal-collapse limitation thus concerns these recognition representations, not every use of strain in the literature.

\subsection{Implications for this research}
\label{sec:strain-implications}

\begin{table}[htbp]
    \centering
    \small
    \begin{tabular}{@{}p{0.29\linewidth} p{0.65\linewidth}@{}}
        \hline
        \textbf{Design choice} & \textbf{Interpretation or unresolved test} \\
        \hline
        Strain accompanies $u$ and $v$
        & Follows STSTNet's channel structure~\cite{liong2019b}; does not establish that its ambiguous magnitude equation was reproduced. \\
        Three-term magnitude
        & Differs from the four-term convention~\cite{liong2014a,liong2014b,liong2016}. Relative shear weighting remains an untested choice. \\
        Per-pair strain sequence
        & Preserves deformation over time instead of pooling it into one descriptor. Each channel is independently min--max normalised across the sequence. \\
        No dedicated strain post-filter
        & Departs from the explicit filtering of Liong et al.~\cite{liong2014a}. Farneb\"ack's internal smoothing is not a separately evaluated strain filter. \\
        Strain held constant
        & Present in all twelve evaluated configurations; its marginal contribution is not measured. \\
        \hline
    \end{tabular}
    \caption[Strain representation and unresolved choices]{The implemented strain representation and the limits it places on interpretation.}
    \label{tab:strain-decisions}
\end{table}

A strain-on/strain-off comparison is the natural next ablation. Neither published gains from a different descriptor nor the performance of the complete system substitutes for that measurement. Comparing magnitude conventions and adding a dedicated post-filter are separate follow-ups.

Strain also extends the analytical-feature argument: the input contains a \textit{hand-specified differential operator}, exposing local deformation without learning that operator from the training clips. A spatial stem may still learn useful combinations of those features, but it is operating after both motion estimation and deformation extraction. This motivates testing its incremental contribution; it does not establish the cause of any observed architectural effect.

\subsection{Research gap and scope}
\label{sec:strain-gap}

The study preserves the temporal evolution of strain for learned sequence processing. Its contribution is an architectural comparison on that representation, not an isolated demonstration of strain's benefit. The independent value of the third channel, the choice of shear weighting and the omitted post-filter remain unresolved. These boundaries preserve the distinction between evidence for a component, a reason to include it, and a result that the experimental design can actually establish.


\section{Temporal Normalisation: The Temporal Interpolation Model}
\label{sec:tim-review}

Flow and strain fix \textit{what} is computed from each frame pair; temporal normalisation fixes \textit{which} frame pairs, and how many. The Temporal Interpolation Model (TIM) is credited in this corpus to Zhou et al. by Li et al.~\cite{liX2013,liX2018} and to Pfister et al. by Lu et al.~\cite{lu2015}, who applied it to micro-expressions outside the review corpus. It is described here only through the reviewed papers that implement and evaluate it.

\subsection{The problem: sequences of unequal length}
\label{sec:tim-problem}

Micro-expressions vary in duration, and the three-class CASME II working set spans roughly a four-fold range of sequence length (\autoref{sec:corpus-threats}h). Three difficulties follow. Any batched network requires a fixed input shape, and the corpus supplies no constant. A descriptor with a temporal radius cannot be applied to a sequence shorter than that radius, so the shortest clips constrain the descriptor for the whole dataset: Li et al.~\cite{liX2018} note that at 25~fps some micro-expressions last four to five frames, restricting LBP-TOP to radius $r = 1$. And duration is itself a nuisance variable --- Lu et al.~\cite{lu2015} normalise to remove ``temporal fluctuations of micro-expressions, which may introduce noise irrelevant to classification''.

\subsection{What TIM is}
\label{sec:tim-what}

Lu et al.~\cite{lu2015} give the corpus's most compact description: ``TIM is a manifold-based interpolation method, which builds a low-dimensional embedding of an image sequence, and then interpolates a curve in the low-dimensional space. The interpolated frames are mapped back to the original high-dimensional space.''

The property that matters downstream is that TIM \textbf{synthesises frames that were never recorded}: Lu et al.~\cite{lu2015} note that Pfister et al. proposed it ``to generate sufficient frames from an image sequence'', that is, as an upsampler rather than only a length-standardiser.

\subsection{TIM in practice: the convergence on ten frames}
\label{sec:tim-ten-frames}

TIM appears throughout the corpus as an unremarked preprocessing step, almost always at ten frames: SMIC's baseline~\cite{liX2013} normalises to ten before LBP-TOP; Li et al.~\cite{liX2018} apply TIM10 after magnification; and Lu et al.~\cite{lu2015} apply TIM before their Delaunay-based coding model, following Pfister et al.'s protocol so as to compare against TIM10 baselines. CAS(ME)$^2$'s baseline~\cite{qu2016} is the exception, normalising to \textbf{120 frames by linear interpolation} and evaluating leave-one-video-out.

One study tests the setting on its own terms. Li et al.~\cite{liX2018} sweep interpolation length from 10 to 80 frames on all three SMIC variants under leave-one-subject-out, holding the descriptor fixed and skipping magnification so only length varies, and report that ``the best performance for all three datasets is achieved with interpolation to 10 frames'', because ``in longer sequences, changes along the time-dimension are diluted''.

\subsection{How many frames? The remaining systematic comparisons}
\label{sec:tim-how-many}

Ben et al.~\cite{ben2021} run a second controlled experiment and reach the opposite conclusion. Holding alignment, descriptor and classifier fixed --- JCFDA alignment, LBP-TOP, and an SVM with an RBF kernel --- they interpolate sequences to 10, 20, 30 and so on up to 110 frames by two algorithms, \textbf{Newton interpolation} and \textbf{TIM}, recording the best recognition rate at each length on MMEW.

Newton peaks at \textbf{33.3\,\%} at 60 frames; TIM peaks at \textbf{38.9\,\%} ``when the micro-expression sequence is respectively interpolated to 30 or 60 frames'', and Ben et al. adopt TIM at 60 accordingly. Accuracy is non-monotonic in length: rates ``increase initially and then decrease towards the end''.

Two caveats limit the transfer. The experiment uses MMEW with LBP-TOP, not CASME II and not a learned model; and it runs largely in the \textit{upsampling} regime, since Ben et al.~\cite{ben2021} interpolate ``to the maximal value of 110 frames'', so the longer settings add frames for most clips rather than discarding them. This thesis operates in the opposite regime (\autoref{sec:tim-implications}), and an optimum measured while synthesising frames need not hold when discarding them.

A third comparison is the only multi-corpus one. Xu et al.~\cite{xu2017} interpolate every sequence to 10, 15 and 20 frames for both their facial dynamics map and LBP-TOP, ``on all four datasets'' --- CASME, CASME II and SMIC's high-speed and visual subsets. They also depart from the convention in a way the others do not, using \textbf{linear} interpolation ``instead of Temporal Interpolation Model because in a high frame-rate scenario, linear interpolation suffices to characterize the motion pattern with less error''.

\subsection{Alternatives to uniform interpolation}
\label{sec:tim-keyframe}

Zhao et al.~\cite{zhaoS2021} do not interpolate uniformly. They build an eleven-frame key-frame sequence from the annotated onset, apex and offset by generating eight transition frames allocated between the onset--apex and apex--offset intervals \textbf{in proportion to each sub-interval's duration}, so sampling density follows the expression's own phase structure and the apex is emphasised. Optical flow between adjacent key-frames then gives a ten-frame flow sequence. OFF-ApexNet~\cite{liong2019a} and STSTNet~\cite{liong2019b} take the limit, keeping the onset--apex pair alone (\autoref{sec:flow-what-between}).

The corpus therefore contains temporal representations of length 2, 10, 11, 15, 20, 30, 60, 110 and 120 frames, chosen variously by convention, by sweep and by architecture. There is no consensus, and only Xu et al.'s~\cite{xu2017} sweep spans more than one corpus.

\subsection{Limitations}
\label{sec:tim-limitations}

\textbf{Ten frames is a convention, not a finding.} The most-used setting traces to one baseline paper's choice, and the three studies that test it do not agree: ten frames for Li et al.~\cite{liX2018}, 30 to 60 for Ben et al.~\cite{ben2021}, and a 10--20 frame range for Xu et al.~\cite{xu2017} (\autoref{sec:tim-how-many}). Interpolation length interacts with the temporal receptive field of whatever consumes it, and in all three that consumer is a hand-crafted descriptor.

\textbf{Frame synthesis and motion analysis are in tension.} TIM generates frames that were not recorded. Where the downstream representation is appearance-based this is benign; where it is \textit{motion}, synthesised frames yield synthesised displacements. One reviewed paper recognises this: Xu et al.~\cite{xu2017} choose linear interpolation over TIM precisely on the grounds that it ``suffices to characterize the motion pattern with less error'' at a high frame rate. But they assert the point rather than measure it --- TIM appears once in their paper, in that sentence --- and no reviewed study compares the two interpolators against a motion representation. The rest of the corpus does not raise the question, because those using TIM extract appearance descriptors, while those extracting flow either use two frames or, as Zhao et al.~\cite{zhaoS2021} do, generate transition frames without discussing the consequence.

\textbf{Downsampling forfeits part of the corpus's advantage.} \autoref{sec:corpus} establishes why CASME II was recorded at 200~fps rather than its predecessor's 60~fps; reducing a clip to a few tens of frames gives part of that back. Only Xu et al.~\cite{xu2017} tie the interpolation choice to the acquisition rate at all, and then only to justify the simpler interpolator. The rest apply TIM10 to 100~fps SMIC data as readily as to 200~fps CASME II data, without weighing what the higher rate bought.

\subsection{Implications for this research}
\label{sec:tim-implications}

\begin{table}[htbp]
	\centering
	\footnotesize
	\begin{tabular}{@{}p{5.6cm} p{8.0cm}@{}}
		\hline
		\textbf{Finding from the review} & \textbf{Decision taken in this thesis} \\
		\hline
		Fixed-length input is required, and unequal duration is nuisance variance~\cite{lu2015} & Every clip is normalised to a fixed temporal length before training \\
		TIM beats Newton interpolation by 5.6 points, with a joint optimum at 30 or 60 frames~\cite{ben2021} & \textbf{L = 33 frames sampled, giving T = 32 flow pairs} --- the lower of the two reported optima rather than the field-conventional 10 \\
		Sequence length has an interior optimum found only by sweep~\cite{ben2021} & T = 32 is fixed throughout and is \textbf{not} swept; a sweep is declared as further work \\
		Onset and offset are annotated in CASME II (\autoref{sec:casme2-construction}) & Sampling is bounded by the annotated onset and offset, so the window spans the expression rather than surrounding neutral frames. This holds on the image-folder path used for CASME II; the loader's \texttt{.avi} branch samples across a whole file and is not exercised here \\
		TIM synthesises frames, and flow between synthesised frames is synthesised motion (\autoref{sec:tim-limitations}) & \textbf{No frames are synthesised} --- see the declaration below \\
		\hline
	\end{tabular}
	\caption[Review findings and the temporal-normalisation decisions they determine]{Review findings and the temporal-normalisation decisions they determine.}
	\label{tab:tim-decisions}
\end{table}

\textbf{A declaration about what this pipeline implements.} The module performing this step is named for the Temporal Interpolation Model, but what it computes is \textbf{uniform index sampling} of existing frames (\autoref{sec:three-channel-tensor}). No embedding is built, no curve is interpolated, no frame is synthesised. That is a materially different operation from the TIM of \autoref{sec:tim-what}, and it is named accurately here rather than in the pipeline's own terms. It is also defensible: because the downstream representation is flow between adjacent sampled frames, interpolation would manufacture displacement fields for motion that never occurred, whereas subsampling discards temporal resolution without fabricating it. What is given up is TIM's ability to \textit{lengthen} short clips.

Three quantitative consequences follow.

\textit{Effective temporal resolution falls by a clip-dependent factor.} Sampling 33 frames from a clip of $N$ frames recorded at 200~fps gives an effective rate of $6600/N$ fps: \textbf{100 fps} at the median ($N = 66$), which is SMIC's recording rate and half what CASME II was built to provide, and roughly \textbf{52 fps} for the longest ($N = 126$), below the original CASME database's 60~fps~\cite{yan2014}. Eighty-one of the 156 clips are downsampled by a factor of two or more.

\textit{Magnification is applied after subsampling, which shifts the realised pass-band.} Magnification runs on the 33 sampled frames at the nominal 200~fps (\autoref{sec:evm-departures}), but those frames span the clip's full duration, so the true inter-frame interval is $N/(200 \times 32)$ seconds rather than $1/200$. At the median a band specified as 5--25 Hz therefore corresponds to roughly 2.5--12.5 Hz physically, and the discrepancy grows with clip length. The realised band is thus neither the 5--25 Hz \autoref{sec:evm-implications} derives from the duration rule of Bai et al.~\cite{bai2021} nor constant across clips. This is a declared divergence; magnifying before subsampling is declared as further work.

\textit{Three clips are shorter than the sampling window.} \textbf{Three of the 156 clips} contain 31 frames against a target of 33, so sampling repeats indices and yields near-zero flow at those positions.

\subsection{Research gap and scope}
\label{sec:tim-gap}

The unresolved question is what temporal normalisation costs a representation made of motion rather than appearance. Three studies sweep interpolation length and give three answers, none on CASME II and none with a learned temporal model; and although Xu et al.~\cite{xu2017} raise both the synthesis-versus-motion tension and the acquisition-rate trade-off, they settle neither by measurement (\autoref{sec:tim-limitations}). This thesis takes the position that follows from its own representation --- sampling only recorded frames --- and quantifies what that costs in effective frame rate (\autoref{sec:tim-implications}), which no reviewed paper reports.

The design does not sweep T, nor compare interpolators, nor separate the effect of sequence length from the effect of the sampling regime. Those remain limitations rather than gaps claimed as solved, and the resulting conclusions apply to the implemented sampling scheme at T = 32.


\section{The Learned Spatial Backbone: Shallow 3D Convolutional Networks}
\label{sec:cnn-review}

This section covers the convolutional backbone, switched on and off as Variable C of this thesis's ablation, through four review-corpus papers.

\subsection{The case for, and against, a learned spatial stem}
\label{sec:cnn-case}

By 2019 the field had converged on optical flow as the input representation (\autoref{sec:why-flow}); the open question was what should consume it. A learned backbone promises what hand-crafted descriptors cannot give~\cite{liong2019b}, but runs into the small-data constraint (\autoref{sec:capacity}, \autoref{sec:corpus-threats}). Xia et al.~\cite{xia2020b} describe a related failure mode from a different cause --- ``the important subtle dynamics are prone to disappearing in the domain shift \ldots especially for deep models'' --- where the mechanism is the domain shift induced by mixing corpora rather than sample count alone, and where they report deeper, higher-resolution models doing better on the individual-database task. Applying their shrinking argument to a single small corpus therefore extends it beyond the setting they establish.

\subsection{What ``3D CNN'' means in this literature}
\label{sec:cnn-meaning}

The term covers three materially different architectures.

\textbf{A convolution over spatial and modality axes, with no temporal axis.} STSTNet~\cite{liong2019b}, the canonical shallow model in this field, calls itself a ``Shallow Triple Stream Three-dimensional CNN''. Its input is a single \textbf{28 $\times$ 28 $\times$ 3} tensor from the onset and apex frames, whose third dimension holds horizontal flow, vertical flow and optical strain --- not time. Three parallel streams, each a single convolutional layer with 3, 5 and 8 kernels respectively, supplement ``the small scale input data \ldots to avoid the problem of underfitting''. \textbf{There is no time dimension in this network}; the clip was collapsed to a single frame pair upstream (\autoref{sec:flow-what-between}).

\textbf{A convolution genuinely extended across time.} STRCN~\cite{xia2020a} extends convolutional connectivity ``across temporal domain'' with recurrent convolutional layers, modelling appearance and geometry separately.

\textbf{A convolution over a short flow sequence.} Zhao et al.~\cite{zhaoS2021} apply a C3D-style network to the ten-frame flow sequence built from their eleven key-frames (\autoref{sec:tim-keyframe}), so the temporal axis is present but short.

\subsection{The central finding: shrink the model \textit{and} the input}
\label{sec:cnn-shrink}

Xia et al.~\cite{xia2020b} treat depth and input resolution as controlled variables. They build four models of increasing depth --- from a single convolutional layer (Model 1) to one convolutional plus three recurrent convolutional layers (Model 4) --- and evaluate each at nine input resolutions, \textbf{20, 40, 60, 80, 100, 150, 200, 250 and 300 pixels}, on optical flow with UAR as the metric. Two results follow.

\textbf{Deeper is worse.} Models 3 and 4 ``have worse recognition performances compared to shallower \ldots models'', even though all four are far shallower than ResNet-101 or DenseNet-169.

\textbf{Higher resolution is worse, with a threshold.} Performance ``will be degraded with larger input resolutions (Model 2, Model 3 and Model 4)'' --- the two-layer model they go on to adopt included --- and for the deeper models ``begins to decrease dramatically when the resolution is bigger than \textbf{100 $\times$ 100}''. Shallow models are nonetheless ``basically robust to the change of input resolution'', improving toward lower resolution above a floor of about 40 $\times$ 40.

STSTNet takes the same argument to its limit and remains competitive on the CASME II subset (\autoref{sec:megc-metrics}): Liong et al.~\cite{liong2019b} report \textbf{0.00167 M parameters across 2 layers at 28 $\times$ 28 $\times$ 3 input}, against OFF-ApexNet's 2.77 M across 5 layers and GoogLeNet's 7 M across 22.

\subsection{Coping with limited and imbalanced training data}
\label{sec:cnn-limited-data}

Two strategies in the corpus have counterparts in this thesis's training configuration. Xia et al.~\cite{xia2020a} pair temporal augmentation with a balanced loss (\autoref{sec:imbalance-balanced-loss}, \autoref{sec:imbalance-double-correction}); a two-stage few-shot scheme~\cite{zhaoS2021} is the one route not taken. Xia et al.~\cite{xia2020b} take a third --- adding capability without capacity --- by developing ``three parameter-free modules \ldots integrate with RCN without increasing any learnable parameters'', so representational power grows while the parameter count, and therefore the overfitting risk, does not.

\subsection{Limitations}
\label{sec:cnn-limitations}

\textbf{The terminology obscures what is being compared.} Papers titled as three-dimensional CNNs include networks with no temporal axis at all, so ``3D-CNN'' results across the corpus are not like with like.

\textbf{Input resolution is almost never controlled.} Xia et al.~\cite{xia2020b} are the exception; elsewhere resolution is an implementation detail --- 28 $\times$ 28 in STSTNet and OFF-ApexNet --- with no indication of sensitivity, despite the substantial sensitivity Xia et al. measure.

\textbf{The evidence base is composite-database.} Xia et al.'s shrinking study~\cite{xia2020b} and STSTNet~\cite{liong2019b} are developed for the MEGC composite setting, which supplies substantially more training data per fold than a single-corpus study (\autoref{sec:megc-metrics}); only the earlier recurrent-convolutional work~\cite{xia2020a} is evaluated per database. Conclusions about affordable capacity therefore understate the constraint here.

\textbf{No study isolates the backbone itself.} The corpus is not short of matched comparisons: it sets descriptor-only methods against learned ones on the same folds, and Xia et al.~\cite{xia2020b} vary individual architectural modules against a common baseline with every other parameter held constant. But those comparisons all \textit{add} to a backbone that stays in place. None removes the learned spatial stem entirely and passes the motion representation straight to a classifier or temporal model, so what the stem itself contributes, set against what it costs, is nowhere measured.

\subsection{Implications for this research}
\label{sec:cnn-implications}

The component ablated as Variable C is a \textbf{three-stream shallow convolutional backbone}: one unshared two-layer stream per input channel --- flow-$u$, flow-$v$, optical strain.

\begin{table}[htbp]
	\centering
	\footnotesize
	\begin{tabular}{@{}p{5.6cm} p{8.0cm}@{}}
		\hline
		\textbf{Finding from the review} & \textbf{Decision taken in this thesis} \\
		\hline
		Shallow beats deep on small micro-expression data~\cite{xia2020b}; STSTNet uses 2 layers and $\sim$1,670 parameters~\cite{liong2019b} & Two convolutional layers per stream, \textbf{14,544 parameters} in the backbone (15,027 including the classifier head) --- shallow by the corpus's standards, though roughly nine times STSTNet's \\
		Triple-stream design supplements small-scale input~\cite{liong2019b} & Three unshared streams, one per input modality. STSTNet's streams differ in \textit{filter count} over the same input; these differ in \textit{input channel} \\
		Parameter-free modules add capability without capacity~\cite{xia2020b} & Motivates the parameter-free attention evaluated as Variable B (\autoref{sec:simam-review}) \\
		Limited and imbalanced data require augmentation plus loss correction~\cite{xia2020a} & Balanced sampling with focal loss; loss class-weighting deliberately disabled after single-class collapse (\autoref{sec:corpus-implications}) \\
		\textbf{Performance degrades above 100 $\times$ 100 input, dramatically for deeper models~\cite{xia2020b}} & \textbf{Not followed.} The backbone is fed at \textbf{224 $\times$ 224}, more than twice that threshold, though the same resolution as Bai et al.~\cite{bai2021} and Li, Huang and Zhao~\cite{liY2021} use --- see below \\
		\hline
	\end{tabular}
	\caption[Review findings and the backbone decisions they determine]{Review findings and the backbone decisions they determine.}
	\label{tab:cnn-decisions}
\end{table}

\textbf{A declaration about what this component is.} Every kernel in this backbone is spatial only: no filter spans two time steps, so despite the \texttt{Conv3d} implementation it is a \textbf{purely spatial} feature extractor (\autoref{sec:conv3d}). The ablation therefore does not test spatio-temporal convolution; it tests a \textbf{learned spatial stem}, and whatever it measures for Variable C is evidence about \textit{this} stem --- spatial-only, at 224 $\times$ 224, on flow and strain input --- and \textbf{not} evidence that 3D convolution over time is unhelpful for micro-expression recognition, which this study never evaluated.

Two observations mitigate that, and one aggravates it. In the field's own terms the choice is defensible: STSTNet, one of the two strongest shallow baselines on the CASME II subset, likewise performs no temporal convolution (\autoref{sec:cnn-meaning}). The stem also competes against a representation that already performs spatio-temporal extraction analytically (\autoref{sec:flow-implications}, \autoref{sec:strain-implications}).

Against both, \textbf{input resolution is the clearest departure from Xia et al.'s guidance}. It is not a departure from the corpus as a whole: Bai et al.~\cite{bai2021} resize frames to 224 $\times$ 224, and without alignment or cropping, as here; Li, Huang and Zhao~\cite{liY2021} resize faces to the same; and Xia et al.'s own earlier work~\cite{xia2020a} runs one stream at 300 $\times$ 245. What 224 $\times$ 224 departs from is the specific threshold of \autoref{sec:cnn-shrink}, which is a composite-database result and which bites hardest on models deeper than this one. Since this stem is shallow, the reviewed evidence predicts a smaller resolution effect than the deep-model figures suggest, so re-running the backbone arm at a lower resolution is worth doing because it is cheap and would settle the question, not because the corpus implies a large gain. The choice is also the expensive one: a 224 $\times$ 224 map carries sixty-four times the spatial elements of STSTNet's 28 $\times$ 28, and the backbone convolves it at every one of the 32 time steps. Whatever this study measures for the stem is therefore a resolution result as much as an architecture result, and the re-run is the first item of declared further work.

\subsection{Limitations of the existing work, and how this study differs}
\label{sec:cnn-gap}

The four gaps of \autoref{sec:cnn-limitations} define what the reviewed literature does not establish.

This thesis differs on the third, which is what its research question requires: it removes the convolutional backbone entirely in four of its twelve configurations, and measures the difference under complete leave-one-subject-out with every other factor held constant. On the first gap it answers the field's terminological ambiguity by declaring precisely what its own component is (\autoref{sec:cnn-implications}); on the second it knowingly departs from the literature's guidance, at the 224 $\times$ 224 resolution flagged above.


\section{Parameter-Free Attention: SimAM}
\label{sec:simam-review}

SimAM re-weights the convolutional feature maps without adding learned parameters. Yang et al.~\cite{yang2021} provide the mechanism and its image-recognition evidence; Xia et al.~\cite{xia2020b} motivate parameter-free modules in composite-database micro-expression recognition. The latter support the design principle, rather than validating SimAM itself on this task.

\subsection{Why parameter-free attention}
\label{sec:simam-why}

Local facial activity occupies a small part of a motion tensor. Attention offers a way to emphasise informative activations, but its parameter and computational costs must be justified on a 156-clip corpus. Yang et al.'s~\cite{yang2021} comparison illustrates the range of parameter costs.

\begin{table}[htbp]
    \centering
    \small
    \begin{tabular}{@{}lc@{}}
        \hline
        \textbf{Module} & \textbf{Additional learnable parameters} \\
        \hline
        SE & $2C^2/r$ \\
        CBAM & $2C^2/r+2k^2$ \\
        GC & $2C^2/r+C$ \\
        ECA & $k$ \\
        SRM & $6C$ \\
        SimAM & $0$ \\
        \hline
    \end{tabular}
    \caption[Parameter costs of attention modules]{Parameter costs reported by Yang et al.~\cite{yang2021}, Table 1. $C$ denotes channels, $r$ the reduction ratio and $k$ the relevant kernel size. The competing modules are attributed through this source.}
    \label{tab:attention-modules}
\end{table}

Zero additional parameters makes SimAM a plausible candidate, not the only defensible one: alternatives such as ECA are also small. Xia et al.~\cite{xia2020b} demonstrate a related strategy of adding parameter-free modules to a recurrent convolutional model. Neither argument establishes that a parameter-free module has zero runtime cost or cannot worsen generalisation.

\subsection{The distinctiveness premise}
\label{sec:simam-premise}

Yang et al.~\cite{yang2021} motivate attention through spatial suppression: neurons distinguishable from their surroundings are assigned greater importance. They express this as the linear separability of one target activation from the other activations in its channel. The resulting weights scale the feature map rather than being added to it.

This is a prior about activation statistics, not a detector of facial action units or minority classes. A distinctive response may encode useful deformation, noise or an artefact. Its suitability for flow-derived features therefore requires evaluation rather than following automatically from the neuroscience motivation.

\subsection{Energy and implemented refinement}
\label{sec:simam-energy}

The paper minimises a regularised separation energy and obtains a closed-form importance score. Its reference implementation can be written compactly as
\[
    d_i=(x_i-\mu)^2,\qquad
    s^2=\frac{1}{M-1}\sum_{i=1}^{M}d_i,\qquad
    E_i^{-1}=\frac{d_i}{4(s^2+\lambda)}+\frac{1}{2},
\]
\[
    \widetilde{x}_i=x_i\,\operatorname{sigmoid}(E_i^{-1}),
    \qquad \mu=\frac{1}{M}\sum_{i=1}^{M}x_i.
\]
Here $M$ is the number of positions in one channel's statistical context and $\lambda$ is a positive regularisation constant. The $+\tfrac12$ follows from the reciprocal-energy expression rather than an arbitrary offset.

An implementation distinction matters: the shared variance printed with Yang et al.'s equation (5) divides by $M$, whereas their Figure 3 code divides by $M-1$~\cite{yang2021}. The equations above follow the reference \emph{code}, as does this thesis. The sigmoid preserves the ordering of the scores while bounding the multiplicative weights; it does not create a hard mask of inactive regions.

\subsection{What three-dimensional weights mean}
\label{sec:simam-3d-weights}

The original module produces a separate weight at each channel--height--width position~\cite{yang2021}. Its three-dimensional weights therefore describe $C\times H\times W$, not a temporal volume. Statistics are computed separately within each channel's spatial plane, with $M=HW$. Extending the context across video frames is an additional design decision, not part of the original image formulation.

\subsection{Evidence and hyper-parameter selection}
\label{sec:simam-evidence}

Yang et al.~\cite{yang2021} evaluate SimAM with several network families on CIFAR and ImageNet, reporting competitive classification performance against learned attention modules. Their timing results also show that unchanged parameter totals do not imply unchanged inference speed. The evidence supports testing SimAM, but does not establish its value on micro-expression motion tensors.

For CIFAR, they sweep $\lambda$ from $10^{-1}$ to $10^{-6}$ and select $10^{-4}$ as a balance between accuracy and variability. They repeat selection for ImageNet and choose $0.1$~\cite{yang2021}. Thus $10^{-4}$ is a dataset-selected operating point, not a universal setting justified by the formula.

\subsection{Limits of transfer to this task}
\label{sec:simam-limitations}

The source classification experiments concern natural images, different training-set sizes and learned feature distributions. In this thesis the input is flow and strain, so distinctiveness need not have the same relationship to useful evidence. The statistics are calculated from each feature map, not estimated by pooling the entire training corpus; the small dataset instead constrains how the surrounding network is learned.

The image experiments also do not decide whether video statistics should be computed per frame or over a clip. Nor does the lack of learned attention weights establish that runtime, memory or overfitting concerns disappear. These are empirical limits on transferring the source result, rather than reasons to assume either success or failure.

\subsection{Implications for this research}
\label{sec:simam-implications}

\begin{table}[htbp]
    \centering
    \small
    \begin{tabular}{@{}p{0.29\linewidth} p{0.65\linewidth}@{}}
        \hline
        \textbf{Choice} & \textbf{Implementation and interpretation} \\
        \hline
        Placement & Applied separately to each convolutional stream before concatenation; no trainable weights or biases are introduced. \\
        Refinement & Uses the reference-code reciprocal energy and sigmoid gate~\cite{yang2021}. \\
        Statistical context & Mean and variance span $(T,H,W)$ within each channel, with divisor $THW-1$, rather than one $(H,W)$ plane. \\
        Regularisation & $\lambda=10^{-4}$ is retained without re-tuning on CASME II. \\
        Architectural dependency & SimAM is enabled only with the CNN stem in this model. Four of sixteen candidate configurations are excluded, leaving twelve evaluated configurations and four SimAM matched pairs. \\
        \hline
    \end{tabular}
    \caption[SimAM implementation choices]{SimAM implementation choices and their relation to the reviewed image formulation. The dependency on the CNN is a rule of this architecture.}
    \label{tab:simam-decisions}
\end{table}

The five-dimensional input includes batch, channel, time, height and width. Pooling statistics across time makes an activation distinctive relative to its whole clip rather than its own frame. This could emphasise a brief contraction, but does not guarantee low weights at onset or high weights at the apex: the score depends on squared deviation, not expression labels. The two context choices are not compared experimentally.

The CNN dependency similarly describes this implementation, not a mathematical prohibition on applying SimAM to other tensors. Its measured effect is conditional on the stem being present. Retention should be judged against both recognition performance and computational cost; parameter count alone cannot settle that decision.

\subsection{Research gap and scope}
\label{sec:simam-gap}

The study evaluates SimAM in a motion-input, small-corpus setting through four matched comparisons under complete leave-one-subject-out evaluation. It tests the implemented module as a whole, rather than separately establishing the benefit of each design choice.

Three follow-ups remain: testing whether the distinctiveness prior suits flow-derived features, re-tuning $\lambda$, and comparing per-frame with whole-clip statistics. Each requires new evaluations despite introducing no learned attention weights. Per-class changes may also be examined, but a rare class benefiting most is a hypothesis of this study, not a prediction established by Yang et al.~\cite{yang2021}.


\section{The Temporal Encoder: Transformers for Micro-Expression Recognition}
\label{sec:transformer-review}

The temporal encoder is motivated by Zhang et al.'s Short and Long Range Relation Based Spatio-Temporal Transformer (SLSTT)~\cite{zhang2022} and its Vision Transformer (ViT) foundation~\cite{dosovitskiy2021}. The component evaluated here differs substantially from published SLSTT, particularly in the axis over which self-attention operates.

\subsection{Why temporal self-attention}
\label{sec:transformer-why}

A micro-expression evolves through onset, peak and relaxation. Retaining a motion sequence permits a model to relate separated stages rather than compressing them immediately into one descriptor. Zhang et al.~\cite{zhang2022} motivate their design through the need to capture both local and global patterns.

This thesis compares a temporal transformer with mean pooling over the same input sequence. Mean pooling discards order, whereas positional information and self-attention allow the encoder to model relations among frame-pair features. This motivates the comparison without establishing in advance which representation will generalise better.

\subsection{The Vision Transformer foundation}
\label{sec:transformer-vit}

ViT splits an image into patches, projects them into token vectors, adds learned positional embeddings and processes them with transformer encoder layers~\cite{dosovitskiy2021}. Attention then relates spatial patches. Dosovitskiy et al. also show that data scale and training regime matter: without the locality and translation-equivariance biases of convolution, image transformers can perform poorly when trained with insufficient data, while large-scale pre-training improves their performance.

That evidence is a caution for a 156-clip study, not a universal requirement that every transformer must be pre-trained. A short sequence of motion features presents a different learning problem from image patches, but remains subject to limited data and overfitting.

\subsection{Published SLSTT}
\label{sec:slstt-architecture}

Zhang et al.~\cite{zhang2022} combine three stages. First, \textbf{onset-referenced optical flow} describes sampled frames relative to the onset rather than to their immediate predecessors. They argue that this gives a more structured trajectory towards and away from the apex. Their experiments use this representation throughout; they do not provide a matched short-term-flow ablation establishing its independent gain.

Second, a \textbf{spatial ViT encoder} processes each flow field as patches. The implementation uses ImageNet-pre-trained ViT-B/16, twelve encoder layers and $384\times384$ input. Third, an \textbf{LSTM temporal aggregator} combines the per-frame features. A mean aggregator supplies the comparison variant. Self-attention is spatial in this architecture; recurrence handles time.

\subsection{Published results and the aggregation comparison}
\label{sec:slstt-results}

\begin{table}[htbp]
    \centering
    \small
    \begin{tabular}{@{}lccc@{}}
        \hline
        \textbf{Method} & \textbf{Composite UF1} & \textbf{CASME II UF1} & \textbf{CASME II UAR} \\
        \hline
        LBP-TOP & 0.588 & 0.703 & 0.743 \\
        Bi-WOOF & 0.630 & 0.781 & 0.803 \\
        OFF-ApexNet~\cite{liong2019a} & 0.720 & 0.876 & 0.868 \\
        STSTNet~\cite{liong2019b} & 0.735 & 0.838 & 0.869 \\
        EMR~\cite{liu2019} & 0.789 & 0.829 & 0.821 \\
        RCN~\cite{xia2020b} & 0.705 & 0.809 & 0.856 \\
        SLSTT-Mean & 0.788 & 0.844 & 0.830 \\
        SLSTT-LSTM & 0.816 & 0.901 & 0.885 \\
        \hline
    \end{tabular}
    \caption[SLSTT and comparison results]{Selected results from Zhang et al.~\cite{zhang2022}, Table 2, reported under their three-class Composite Database Evaluation heading. Columns retain the source's dataset labels. The listed methods are not all arms of one controlled architectural comparison.}
    \label{tab:slstt-results}
\end{table}

The SLSTT variants compare temporal aggregation with the spatial transformer retained in both. CASME II UF1 increases from 0.844 to 0.901, a difference of 0.057, or 5.7 UF1 percentage points~\cite{zhang2022}. This supports evaluating aggregation choices within that framework. It is not an estimate of what an LSTM would add to this thesis, where mean pooling follows \emph{temporal} self-attention rather than independent spatial encoding.

Protocol labels need care. The source describes three-class LOSO evaluations on the databases separately and together under its Composite Database Evaluation heading. That heading alone therefore does not establish that every CASME II column entry uses pooled cross-database training. It also reports a distinct Sole Database Evaluation using original labels, including five CASME II classes~\cite{zhang2022}. Neither setting should be equated automatically with this thesis's 156-clip three-class pool.

\subsection{Limitations of the evidence}
\label{sec:transformer-limitations}

Published SLSTT retains a pre-trained spatial encoder. Its results do not isolate the benefit of a transformer trained from scratch on this study's motion representation. The Mean--LSTM comparison isolates an aggregator, not the presence of a transformer, since both variants contain one.

Label interpretation also limits comparison. Zhang et al.~\cite{zhang2022} note that facial action units do not map uniquely to emotion labels, so visual evidence alone need not determine the annotation. This is a limitation of the recognition task, not a numerical upper bound measured by the present experiments.

\subsection{Implementation and six declared divergences}
\label{sec:transformer-implications}

The evaluated component receives 32 tokens of width 96. They come either from the convolutional stem or from $4\times4$ spatial pooling followed by a \emph{learned linear projection}. The encoder uses four pre-norm layers, eight heads, feed-forward width 256 and dropout 0.1. Fixed sinusoidal positions precede attention; mean pooling follows it. With the encoder disabled, the input tokens are averaged directly.

\begin{table}[htbp]
    \centering
    \small
    \begin{tabular}{@{}p{0.23\linewidth} p{0.34\linewidth} p{0.35\linewidth}@{}}
        \hline
        \textbf{Property} & \textbf{Published SLSTT~\cite{zhang2022}} & \textbf{This thesis} \\
        \hline
        Attention axis & Spatial patches within a frame & Temporal tokens within a clip \\
        Aggregation & LSTM over spatially encoded frames & Mean after temporal self-attention \\
        Initialisation & ImageNet-pre-trained ViT-B/16 & Trained from scratch \\
        Capacity & Twelve layers, width 768; $384\times384$ inputs & Four layers, width 96; 32 tokens \\
        Positions & Learned spatial embeddings & Fixed sinusoidal temporal encoding \\
        Input flow & Relative to onset & Relative to the preceding sampled frame \\
        \hline
    \end{tabular}
    \caption[Divergences from published SLSTT]{Six material divergences from published SLSTT. The component evaluated here is a small temporal transformer encoder, not a reproduction of the published architecture.}
    \label{tab:slstt-divergences}
\end{table}

Six matched pairs evaluate the temporal encoder with the other component choices fixed within each pair. They test the whole encoder, including positional encoding, feed-forward layers and normalisation, against direct mean pooling. They do not isolate self-attention from those accompanying operations or compare it with another learned temporal model.

Three follow-ups follow from the divergences: testing an LSTM aggregator, comparing onset-referenced with consecutive-frame flow, and replacing fixed positions with learned embeddings. The last would add $32\times96=3{,}072$ embedding parameters for this sequence length. All three remain untested; changing the input representation also requires re-training and evaluation before any recognition benefit can be claimed.

\subsection{Research gap and scope}
\label{sec:transformer-gap}

The study asks what a small temporal encoder contributes after motion extraction, using a single corpus without pre-training. This differs from SLSTT's spatial transformer and recurrent aggregation while retaining its motivation to model relations beyond isolated frames.

The contribution is a matched comparison of the implemented encoder and an order-blind baseline within this pipeline. It neither validates published SLSTT nor establishes a general ranking of attention, recurrence and convolution. The six declared divergences and the limited training regime define the scope of the architectural conclusions.


\section{Training Under Severe Class Imbalance}
\label{sec:imbalance-review}

This section covers the training configuration held constant across all twelve ablated configurations, and which \autoref{sec:corpus-implications} identified as forced by the corpus rather than chosen freely: the handling of a 4 : 1.3 : 1 class imbalance.

\textbf{A deliberate exclusion.} Two topics the project's reading list groups here --- Grad-CAM++ visual auditing and adversarial identity disentanglement --- are \textbf{not reviewed}, because neither is part of the implemented system. The loss module records that ``the Cross-Batch Memory (XBM) queue, and the adversarial identity loss --- are intentionally removed'', and no gradient-based visualisation exists anywhere in the codebase. Both are declared as proposed further work.

\subsection{The imbalance is a property of the phenomenon, not of the sampling}
\label{sec:imbalance-property}

\autoref{sec:label-taxonomy} and \autoref{sec:original-baseline} established the class counts and the always-Negative floor; what matters here is why the skew exists. It is not an artefact of careless collection. Yan et al.~\cite{yan2014} are explicit that it follows from what can be elicited in a laboratory: ``some types of facial expressions are difficult to elicit in laboratory situations, thus the samples in different categories distributed unequally, e.g., there are 60 disgust samples but only 7 sadness samples.'' Nor is it specific to CASME II --- the MEGC composite shows the same pattern (\autoref{sec:megc-metrics}). Skew must be coped with; it cannot be collected away.

\subsection{Three places to intervene}
\label{sec:imbalance-three-places}

The reviewed corpus corrects for imbalance at three distinct points, worth separating because this thesis intervenes at all three and the interaction between two of them turned out to matter.

\begin{enumerate}
	\item \textbf{At the data}, by changing how often each class is presented during training --- resampling or augmentation.
	\item \textbf{At the loss}, by changing how much each example contributes to the gradient --- class weighting or focal re-weighting.
	\item \textbf{At the metric}, by changing how performance is scored --- macro-averaged rather than sample-averaged measures.
\end{enumerate}

\subsection{Loss-level correction: focal loss}
\label{sec:imbalance-focal}

Zhao et al.~\cite{zhaoS2021} bring focal loss into micro-expression recognition ``to alleviate the inefficient model training caused by the class imbalance in the MEs Datasets'', since cross-entropy ``could be biased toward particular emotions that constitute a larger portion of the training set''. They adopt the formulation developed for one-stage object detection, where the imbalance is between foreground and background regions:

\[
	\text{FL}(p_t) = -\alpha_t (1 - p_t)^{\gamma}\log p_t
\]

where $p_t$ is the model's probability for the true class, $\gamma$ ``the balance factor for loss'' and $\alpha_t$ ``the weight balance factor for samples'' (\autoref{sec:focal-loss}); under skew, the examples it up-weights are disproportionately the minority classes. Adapting it to multi-class classification, Zhao et al. report that ``$\gamma$ is set as 2 in practice, and $\alpha$ is treated as a hyper-parameter to set by cross validation.''

The division of labour there matters: $\gamma$ re-weights by \textit{difficulty} and is treated as a constant, $\alpha$ by \textit{class frequency} and is treated as dataset-specific. Conflating them is where this project went wrong (\autoref{sec:imbalance-double-correction}).

\subsection{Data-level correction: resampling and augmentation}
\label{sec:imbalance-balanced-loss}

Xia et al.~\cite{xia2020a} address the same problem at the data end and pair the two interventions, using ``temporal data augmentation strategies as well as a balanced loss \ldots jointly'' to overcome ``limited and imbalanced training samples''. The pairing is not left unexamined: their subsection \textit{``The Impact of Data Augmentation and Balanced Loss''} removes each correction in turn from the joint system and reports the balanced loss improving performance ``slightly'', with one leave-one-subject-out cell on SMIC where removing it scores higher. This is the corpus's only controlled look at composing two corrections, and \autoref{sec:imbalance-double-correction} records a far larger effect in the same direction.

\subsection{Metric-level correction}
\label{sec:imbalance-metrics}

\autoref{sec:megc-metrics} defines UF1 and the case for macro-averaged metrics. Liong et al.~\cite{liong2016} add the argument from the subject side rather than the class side, reporting F1, precision and recall because these ``provide a more meaningful perspective than accuracy rates when the datasets used are naturally imbalanced since each subject has a different number of video samples.'' The two arguments compound on CASME II, where classes are skewed \textbf{and} per-subject clip counts range from 1 to 33 (\autoref{tab:corpus-threats}c).

\subsection{Limitations}
\label{sec:imbalance-limitations}

\textbf{The interaction between corrections is examined once, and only at small magnitude.} Xia et al.~\cite{xia2020a} ablate augmentation and balanced loss out of the joint system and report the balanced loss helping only ``slightly'' --- at most about two points across their twelve cells, and in one cell costing accuracy; Zhao et al.~\cite{zhaoS2021} apply focal loss alongside a large per-sample augmentation, but without class-proportional resampling and without ablating the pairing. What no reviewed paper reports is the pairing failing outright.

\textbf{$\gamma$ is inherited rather than tuned.} Zhao et al.~\cite{zhaoS2021} take the value from the object-detection literature ``in practice''. They do compare focal loss against cross-entropy across a sweep of $\alpha$, but that experiment holds $\gamma$ fixed, so the focusing exponent remains unvalidated on this task.

\textbf{Label smoothing is named in the corpus but never examined on this task.} Dosovitskiy et al.~\cite{dosovitskiy2021} list it among the regularisation parameters they tune, but no reviewed micro-expression paper uses or discusses it, or reports what it does under class skew. It is used here and declared as a choice this corpus supports neither way.

\subsection{Implications for this research}
\label{sec:imbalance-double-correction}

\begin{table}[htbp]
	\centering
	\footnotesize
	\begin{tabular}{@{}p{5.6cm} p{8.0cm}@{}}
		\hline
		\textbf{Finding from the review} & \textbf{Decision taken in this thesis} \\
		\hline
		Focal loss addresses class-frequency bias in micro-expression recognition~\cite{zhaoS2021} & Focal loss adopted as the training objective, with \textbf{$\gamma = 2.0$}, matching Zhao et al.'s reported setting \\
		$\alpha$ is a dataset-specific class-weight term~\cite{zhaoS2021} & Implemented as an optional per-class $\alpha$ vector, \textbf{switched off at run time} --- see the declaration below \\
		Correct at the data end as well as the loss end~\cite{xia2020a} & Minority oversampling via a weighted sampler, active in every configuration --- a sampler, where Xia et al. use uniform augmentation plus a balanced loss \\
		Imbalance-aware metrics are mandatory~\cite{see2019,liong2016} & Pooled macro F1 --- identical in construction to MEGC's UF1 --- as the primary metric, with accuracy quoted only alongside the always-Negative floor \\
		\textit{(no micro-expression source)} & Label smoothing of 0.05 applied inside the focal loss; named among the regularisers tuned by Dosovitskiy et al.~\cite{dosovitskiy2021}, but declared as unsupported either way by this corpus \\
		\hline
	\end{tabular}
	\caption[Review findings and the imbalance-handling decisions they determine]{Review findings and the imbalance-handling decisions they determine.}
	\label{tab:imbalance-decisions}
\end{table}

\textbf{The finding this project contributes: correct at exactly one point, and know which.} The configuration is resolved at run time by a single rule --- class weighting in the loss is enabled only when the balanced sampler is \textit{not} active --- and the run log records the decision each fold: \textit{``Class weights in loss disabled (balanced sampler already active).''}

That rule exists because of two failures in opposite directions, both recoverable from this project's own run history. In the earliest code era the loss carried inverse-frequency class weights and no balanced sampler existed; under that single correction the proposed configuration abandoned the \textbf{Positive} class outright, predicting none of its clips --- a failure initially misdiagnosed as over-capacity overfitting. A later era added the sampler while leaving the loss weights unconditional, so both corrections ran at once; the collapse then inverted, and configurations abandoned the \textbf{Negative} majority class instead, one of them predicting a single class for every clip in the corpus. Correcting once at the loss starved the minority; correcting twice over-served it. With the guard in place and exactly one correction active, no configuration in the reported study abandons any class.

This is a contribution rather than an incident, because of how little guidance the corpus offers (\autoref{sec:imbalance-limitations}): no reviewed paper warns that combining corrections can invert the intended effect rather than merely blunt it. On a 156-clip corpus with a 4 : 1 skew, correcting twice is as damaging as correcting once in the wrong place --- the two failures differ only in which class is abandoned.

\textbf{A second declaration.} The severity of the imbalance also constrains what the ablation can conclude. Because every configuration shares the same sampler, loss and $\gamma$, this thesis measures \textbf{no} effect attributable to the imbalance treatment itself: it is a constant, not a factor. The claim supported is that the four architectural components were compared under a training regime that does not abandon minority classes, not that this regime is optimal.

\subsection{Research gap and scope}
\label{sec:imbalance-gap}

The corrections themselves are well established; what the corpus does not settle are the three gaps of \autoref{sec:imbalance-limitations} --- how two corrections compose, what $\gamma$ should be on this task, and what over-correction looks like when it fails.

This thesis contributes to the first and third, documenting two opposite class-collapse failures and the run-time rule that resolved them (\autoref{sec:imbalance-double-correction}): a negative result the corpus does not contain and one directly actionable for anyone reproducing this class of pipeline. On the second it offers no improvement, since $\gamma = 2.0$ and label smoothing of 0.05 are held constant across all twelve configurations. Sweeping either, and the Grad-CAM++ auditing and identity-disentanglement work this project's earlier stages removed, are declared as further work.


\section{Synthesis and Research Gap}
\label{sec:synthesis-and-gap}

This section consolidates the gap statements closing \autoref{sec:corpus}--\autoref{sec:imbalance-review}: what the corpus collectively establishes, the four limitations recurring across every topic, the resulting research gap, and which parts of it this thesis closes. It also collects in one place every declared divergence between the system built here and the literature it draws on.

\subsection{What the corpus establishes}
\label{sec:corpus-establishes}

Read as a whole, the reviewed literature settles a good deal. \textbf{On data:} the corpora are small, demographically narrow, laboratory-bound and irreducibly imbalanced by what can be elicited, and CASME II remains the primary single-corpus benchmark under its native leave-one-subject-out protocol (\autoref{sec:corpus}, \autoref{sec:imbalance-review}). \textbf{On representation:} motion rather than appearance is the appropriate input; a strain field adds a further cue, though on CASME II only weakly and under a protocol that is not subject-disjoint; and both the magnification factor and the interpolated sequence length have interior optima located by sweep (\autoref{sec:evm-review}--\autoref{sec:tim-review}). \textbf{On architecture:} at this data size capability bought without parameters outperforms capability bought with them, and self-attention is the strongest-performing mechanism reported in the corpus (\autoref{sec:cnn-review}--\autoref{sec:transformer-review}). \textbf{On evaluation:} accuracy alone is inadequate under skew and must give way to macro-averaged measures accumulated across folds before averaging (\autoref{sec:megc-metrics}, \autoref{sec:imbalance-metrics}).

\subsection{Four limitations that recur across every topic}
\label{sec:four-limitations}

The per-section gap statements are not nine independent observations. Four structural problems generate most of them.

\textbf{(a) Components are isolated one at a time, inside one proposed design, and never jointly.} Matched-pair measurement is not absent from the corpus. On the input side \autoref{tab:evm-effect} isolates magnification and \autoref{tab:strain-performance} isolates strain; on the architectural side Xia et al. compare a basic recurrent convolutional network against three variants adding a wide expansion, a shortcut connection and an attention unit, holding every other parameter constant across all four (\autoref{sec:cnn-limited-data}), and separately sweep depth against input resolution (\autoref{sec:cnn-shrink}); SLSTT's Mean-versus-LSTM comparison isolates its aggregator (\autoref{sec:slstt-results}).

What no reviewed study does is vary \textit{several} components together, so that their interactions can be read. Every ablation in the corpus is internal to one proposed design and moves one factor at a time from that design's own defaults; across papers, systems are compared whole --- STSTNet against OFF-ApexNet, SLSTT against STSTNet --- so a difference in the reported number cannot be attributed to any one stage. There is therefore no basis for deciding which of the learned stages are load-bearing, nor whether one component's value depends on another's presence.

\textbf{(b) The evaluation protocol is unstable, and under-reported.} Sample counts, label-set size and the choice among leave-one-subject-out, leave-one-video-out and leave-one-sample-out all vary inconsistently, and the field's most-quoted baseline is reported under two different protocols by its two primary sources and under two more within the originating paper itself (\autoref{sec:original-baseline}). Given this corpus's fold composition (\autoref{sec:corpus-threats}), these are not bookkeeping differences: they change what the numbers mean.

\textbf{(c) The strongest results depend on regimes a single-corpus study cannot access.} The best CASME II figures in the corpus come from composite-database training (\autoref{sec:megc-metrics}) and, in SLSTT's case, ImageNet pre-training (\autoref{sec:transformer-limitations}) --- neither a property of the architecture being credited. Almost no evidence exists about what these components do when trained from scratch on one small corpus, which is the situation most practitioners are in.

\textbf{(d) Preprocessing parameters have interior optima, and are transferred without re-validation.} The magnification factor, the interpolated sequence length and the input resolution all rise then fall, all have their optima located only by sweep, and none has a principled selection rule (\autoref{sec:evm-evidence}, \autoref{sec:tim-how-many}, \autoref{sec:cnn-shrink}). Values are carried between pipelines whose downstream models differ entirely: $\alpha = 30$ moves from a single-frame appearance encoder to other settings, TIM10 propagates from one baseline paper to the whole field, and $\lambda = 10^{-4}$ transfers from CIFAR though its own authors re-tuned it for ImageNet (\autoref{sec:simam-evidence}).

\subsection{The research gap}
\label{sec:research-gap}

Taking (a)--(d) together:

\textbf{The field has established a standard pipeline for micro-expression recognition --- magnify, compute motion, normalise duration, extract spatially, model temporally --- but has never measured which of its stages are responsible for the result. Because components are only ever varied one at a time inside a single proposed system, under protocols that vary in unstated ways, and predominantly in training regimes richer than a single small corpus provides, it is not currently possible to say which parts of that pipeline earn their place, which are inherited convention, and which depend on the presence of another.}

This matters practically as well as scientifically. A practitioner assembling this pipeline today has no published basis for allocating effort: for judging whether the most computationally expensive stage is the most valuable, whether a component with a negligible average effect may still matter in a specific failure, or whether a ranking read off an incomplete evaluation of a corpus this size is stable at all.

\subsection{How this thesis addresses the gap --- and what it does not}
\label{sec:addressing-gap}

\textbf{What it does.} This thesis treats the pipeline as a factorial experiment rather than a proposal: four components varied independently across every architecturally valid cell of a 2\textsuperscript{4} matrix, twelve configurations, each evaluated under \textbf{complete 25-fold leave-one-subject-out on CASME II} with every non-varied factor held identical. That gives the joint, multi-component measurement (a) says the corpus lacks; reporting the protocol in full addresses (b) for this study at least; and training from scratch is the regime (c) identifies as unevidenced. Each component's marginal contribution is read from pairs of runs differing in one factor alone, and the distribution of those four effects --- how much of the pipeline's performance each stage accounts for, whether any subtracts, and what each costs --- is the substantive answer to the gap.

Three further contributions follow from the review rather than from the experiment:

\begin{itemize}
	\item \textbf{A caution about partial evaluations.} This project's earlier evaluations, on smaller held-out sets and incomplete fold counts, ranked the twelve configurations differently from the complete run reported here. They are not evidence of protocol sensitivity, since they differ in sample count, training budget and class coverage too, and the earliest carry a data-routing defect --- but the episode illustrates limitation (b) in practice, and only the complete run is reported.
	\item \textbf{A repaired measurement.} That defect left the magnification switch inert, producing bit-identical results across every pair those runs contained. What made it detectable as a defect rather than a null result was the review's evidence that magnification has a clear, repeatable effect (\autoref{sec:evm-implications}).
	\item \textbf{Two opposite class-collapse failures.} A single loss-level correction suppressed the minority class; adding a data-level correction on top of it suppressed the majority class instead. No reviewed paper reports either collapse (\autoref{sec:imbalance-double-correction}).
\end{itemize}

\textbf{What it does not do.} Three of the four limitations are only partially addressed. On (b), this study reports its own protocol completely but cannot repair the field's inconsistency, and its own N = 156 differs from MEGC's 145. On (c), it establishes what these components do from scratch on one corpus, which is a different question from --- not an answer to --- how they behave under composite training; no cross-dataset claim is made anywhere. On (d), the thesis inherits parameter values rather than sweeping them: $\alpha = 10$, T = 32, $\lambda = 10^{-4}$ and $\gamma = 2.0$ are fixed throughout, so the ablation measures architecture at one point in preprocessing space, not across it.

\subsection{A consolidated ledger of declared divergences}
\label{sec:divergence-ledger}

Reviewing the literature closely enough to write this chapter surfaced twelve points at which the implemented system departs from the work it draws on. Several were discovered only through the review. They are collected in \autoref{tab:divergence-ledger} so that no result rests on an unstated assumption.

\begin{table}[htbp]
	\centering
	\footnotesize
	\begin{tabular}{@{}p{0.85cm} p{6.0cm} p{6.6cm}@{}}
		\hline
		\textbf{\S} & \textbf{Divergence} & \textbf{Consequence} \\
		\hline
		3.1.3 & Working set is N = 156, including sadness and fear in Negative; MEGC's CASME II subset is N = 145 & Comparisons against challenge figures are not exactly like-for-like \\
		3.1.2 & The corpus's own registered release is not used; the pipeline reads the original recorded frames, unregistered and uncropped & No landmarking or alignment of any kind precedes the network; tolerance to pose and scale must be learned from 156 clips \\
		3.2.6 & Temporal band is the unnarrowed 5--25 Hz duration rule, where Bai et al. narrowed it to 15--25 Hz as noisy & The pass-band admits the low-frequency range they rejected \\
		3.4.3 & Strain magnitude counts shear once, not twice as in the published Frobenius norm & Shear weighted lower by $\sqrt{2}$ relative to normal strain; untested \\
		3.4.7 & No dedicated strain filtering, where the literature applies Wiener and Gaussian filters & Spatial differentiation can amplify flow-estimation noise, and nothing downstream suppresses it \\
		3.5.7 & Temporal normalisation is uniform index sampling, \textbf{not} manifold-based TIM; no frames are synthesised & Avoids fabricating motion; forfeits the ability to lengthen short clips \\
		3.5.7 & Magnification is applied \textbf{after} subsampling, at a nominal 200 fps & Realised pass-band is clip-dependent and below the intended 5--25 Hz \\
		3.6.6 & Backbone kernels are (1, 3, 3) --- \textbf{no temporal mixing} & The ablation tests a learned \textit{spatial} stem, not spatio-temporal convolution \\
		3.6.6 & Backbone input is 224 $\times$ 224, against Xia et al.'s guidance of $\leq$ 100 $\times$ 100 --- though Bai et al. and Li, Huang and Zhao both use 224 $\times$ 224 & Whatever the stem is measured to contribute may be a resolution effect as much as an architecture effect \\
		3.7.7 & SimAM statistics pooled over the whole clip rather than per frame as in the source; $\lambda$ not re-tuned & Neurons judged against the clip, not the frame; defensible but untested \\
		3.8.6 & Six divergences from published SLSTT (temporal not spatial attention, mean not LSTM aggregation, from scratch, smaller, sinusoidal encoding, short-term flow) & Variable D is a temporal transformer encoder, \textbf{not} a reproduction of SLSTT \\
		3.9.7 & Label smoothing of 0.05, named in the corpus only by Dosovitskiy et al. and never examined on this task & A choice the micro-expression corpus supports neither way, held constant and unmeasured \\
		\hline
	\end{tabular}
	\caption[Every declared departure from the reviewed literature]{Every declared departure from the reviewed literature, with its consequence.}
	\label{tab:divergence-ledger}
\end{table}

Three of these --- the strain formula, the magnification ordering, and the input resolution --- were identified by the review itself rather than being known design decisions, and each is declared as proposed further work. None invalidates the ablation, since every configuration shares the same choices. What they constrain is the \textit{wording} of the conclusions: what is measured is the contribution of a spatial stem at 224 $\times$ 224 and of temporal self-attention in this pipeline, not a verdict on 3D convolution or on SLSTT.

\subsection{What this chapter contributes to the thesis}
\label{sec:ch3-contribution}

Three things carry forward. First, the \textbf{methodology that follows is determined by this review rather than chosen freely}: the corpus dictates the class grouping, representation, fold structure and metric, and each component's parameters come from published sweeps (\autoref{sec:corpus-implications}). Second, the \textbf{interpretation of the central result is prepared here rather than improvised later}: flow and strain already perform spatio-temporal feature extraction analytically (\autoref{sec:flow-implications}, \autoref{sec:strain-implications}), so a negative contribution from a spatial-only backbone at an oversized input resolution is expected rather than anomalous (\autoref{sec:cnn-implications}). Third, the \textbf{limits of what can be claimed are fixed before any number is reported}: a small, single-seed, demographically narrow corpus measured against richer published regimes, with the twelve divergences of \autoref{tab:divergence-ledger} standing throughout --- consequences of this review, not concessions made after the fact.
 in place of contents/chapter3.tex.
% Do not include both: they share every label.
%
% Assembled from chapter3_short.tex (3.1-3.2) and chapter3_short_part3..10.tex.
% Those part files are retained unchanged as the editing units; regenerate this
% file from them rather than editing it directly, or edit here and retire them.
%
% Structure that other chapters depend on, and that must not be disturbed:
%   * Chapter 2 cites Sections 3.1.2, 3.1.6, 3.1.8, 3.3.1, 3.3.3, 3.5.7,
%     3.9, 3.9.5 and 3.9.7 BY NUMBER, so those subsection positions are fixed.
%   * The divergence ledger (Table, sec:divergence-ledger) cites twelve
%     subsections by number. scratchpad/ledgercheck.py re-verifies all of them
%     plus every external numeric reference after any renumbering.
%   * Chapters 4, 5, 6 and the introduction reference sec:research-gap,
%     sec:four-limitations, sec:addressing-gap, sec:divergence-ledger,
%     sec:synthesis-and-gap, sec:imbalance-review, sec:imbalance-gap,
%     sec:imbalance-double-correction, sec:cnn-review, sec:cnn-implications,
%     sec:tim-implications, sec:evm-implications, sec:evm-evidence,
%     sec:evm-limitations, sec:evm-motion-hypothesis, sec:corpus,
%     sec:corpus-threats, sec:corpus-implications, sec:label-taxonomy,
%     sec:megc-metrics, sec:sibling-corpora, sec:flow-review, sec:flow-implications,
%     sec:strain-review, sec:strain-implications, sec:strain-divergence,
%     sec:simam-gap, sec:simam-implications, sec:slstt-results,
%     sec:transformer-implications, sec:addressing-gap, tab:megc-cde-results
%     and tab:spontaneous-corpora.
%
% Conventions held throughout: no sentence points at another chapter; every
% float carries a short \caption[...] for the List of Figures / List of Tables;
% figures are fig3_1_dataset.png and fig3_2_fold_composition.png only.
% =============================================================================

\chapter{Literature Review}
\label{ch:literature}

This chapter reviews the literature bearing on each component of the recognition pipeline evaluated in this thesis: the corpus on which every experiment is run, the four techniques placed under ablation, the three input transformations they operate on, and the training regime held constant across all of them.

It is organised as a single continuous argument rather than as ten independent surveys: a fact is established once, in the section that owns it, and cross-referenced thereafter. Each section runs from what the technique addresses, through how it works and what the published evidence shows, to its limitations and the implications for this thesis, closing with the gap and how this work differs. Two claims recur --- that on a corpus of 156 clips several decisions usually presented as modelling choices are in fact forced by properties of the data, and that several results usually attributed to architecture are attributable to the evaluation protocol. \autoref{sec:synthesis-and-gap} consolidates the nine closing statements into a single research gap and collects every declared departure from the reviewed methods.


% =============================================================================
\section{The Evaluation Corpus: CASME II and the Spontaneous Micro-Expression Benchmark Family}
\label{sec:corpus}
% =============================================================================

This section covers the corpus literature only --- the papers that introduce, specify or standardise evaluation on the database used throughout this thesis. It is treated first because on a database of this size several design decisions usually presented as modelling choices are \textit{forced} by properties of the data.

\subsection{Why a dedicated spontaneous corpus was necessary}
\label{sec:why-spontaneous}

Automatic facial expression recognition matured on large corpora of \textit{posed} expressions --- CK+, MUG, MMI, JAFFE and Multi-PIE --- reporting accuracies above 90\,\% on the six basic posed emotions by the early 2010s~\cite{yan2014}. Micro-expression research could not follow that path: as \autoref{sec:phenomenon} establishes, a micro-expression is involuntary and cannot be \textit{performed} on instruction. Yan et al.~\cite{yan2014} review a pre-2013 landscape of only four micro-expression databases, of which only two were spontaneous (\autoref{tab:spontaneous-corpora}). Two problems in it motivated CASME II directly.

The first is \textbf{construct validity}. Yan et al.~\cite{yan2014} and Li et al.~\cite{liX2013} both note that posed ``micro-expressions'', produced deliberately on instruction, differ systematically from spontaneous ones --- Li et al.~\cite{liX2013} further observing that USF-HD adopted a 2/3-second duration criterion longer than the definitional upper bound of \autoref{sec:phenomenon}. A system trained on posed data is therefore not demonstrably a micro-expression recogniser; it may be a recogniser of fast deliberate movement.

The second is \textbf{acquisition quality}. Defined jointly by very short duration and very low intensity, the phenomenon sits close to the noise floor of ordinary video: Yan et al.~\cite{yan2014} argue that CASME's 60~fps coarsely samples a movement lasting a fraction of a second --- while marking the benefit of higher rates as unresolved --- and that flicker under alternating-current lighting contaminated earlier recordings. Elicitation was a third difficulty: high-stakes deception paradigms produced samples heavily contaminated by conversational and other emotion-irrelevant movement~\cite{yan2014}, and the paradigm that proved workable, which CASME II adopts, is to have participants watch high-arousal video episodes while attempting to keep the face neutral~\cite{liX2013}.

\begin{table}[htbp]
	\centering
	\footnotesize
	\setlength{\tabcolsep}{4pt}
	\begin{tabular}{l c c c c c c c}
		\hline
		\textbf{Corpus} & \textbf{Subj.} & \textbf{Micro samples} & \textbf{fps} & \textbf{Facial res.} & \textbf{Spont.} & \textbf{FACS} & \textbf{Classes} \\
		\hline
		\multicolumn{8}{l}{\textit{The pre-CASME II landscape, reproduced from Yan et al.~\cite{yan2014} (Table 1)}} \\
		USF-HD~\cite{shreve2011}                       & ---      & 100 & 30  & ---              & No  & No  & 6 \\
		Polikovsky et al.\textsuperscript{\textdagger} & 10       & 42  & 200 & ---              & No  & No  & 6 \\
		SMIC~\cite{liX2013}                            & 16 valid & 164 & 100 & 190 $\times$ 230 & Yes & No  & 3 \\
		CASME (2013)\textsuperscript{\textdagger}      & 19 valid & 195 & 60  & 150 $\times$ 190 & Yes & Yes & 7 \\
		\hline
		\multicolumn{8}{l}{\textit{The spontaneous corpora that define the current context}} \\
		\textbf{CASME II}~\cite{yan2014} & \textbf{26 valid} & \textbf{247 / 255} & \textbf{200} & \textbf{280 $\times$ 340} & \textbf{Yes} & \textbf{Yes} & \textbf{5 (7 raw)} \\
		SAMM~\cite{davison2018}          & 32 & 159             & 200 & 400 $\times$ 400 & Yes & Yes & 7 \\
		CAS(ME)$^2$~\cite{qu2016}        & 22 & 53 (+250 macro) & 30  & ---              & Yes & Yes & 4 \\
		\hline
		\multicolumn{8}{l}{\textit{Ethnicities / mean age (SD):} SMIC 3 / 26.7 \quad CASME 1 / 22.03 (1.60) \quad \textbf{CASME II 1 / 22.03 (1.60)}} \\
		\multicolumn{8}{l}{\quad SAMM 13 / 33.24 (11.32) \quad CAS(ME)$^2$ 1\textsuperscript{*} / 22.59 (2.20)} \\
		\hline
	\end{tabular}
	\caption[Posed and spontaneous micro-expression corpora]{Posed and spontaneous micro-expression corpora, compiled from Yan et al.~\cite{yan2014} (Table 1), Davison et al.~\cite{davison2018} (Table 1), Li et al.~\cite{liX2013} and Qu et al.~\cite{qu2016}. \textdagger{}~Described in Yan et al.; the primary papers lie outside the review corpus and are not cited separately. \textsuperscript{*}~Inferred from the recruitment description; Qu et al. do not state it. SMIC's demographic figures are those of all 20 recruited participants; the 16 who produced usable clips have a mean age of 28.1 and span two ethnicities.}
	\label{tab:spontaneous-corpora}
\end{table}

\subsection{CASME II: construction, apparatus, and annotation}
\label{sec:casme2-construction}

CASME II~\cite{yan2014} is the corpus on which every experiment in this thesis is performed, and several of its specifics propagate into the data pipeline used here.

\textbf{Participants and apparatus.} Thirty-five participants were recruited, mean age 22.03 years (SD = 1.60), and watched high-valence video episodes pre-rated by an independent panel of twenty raters, under instruction to suppress all facial movement. Recording sustained \textbf{200 fps} at 640~$\times$~480 with no inter-frame compression, under four LED lamps with umbrella reflectors giving steady flicker-free illumination --- a deliberate correction of the artefacts observed in SMIC. The facial region is approximately \textbf{280 $\times$ 340 pixels}, over three times CASME's facial area and roughly double SMIC's.

\textbf{Selection and annotation.} From nearly 3,000 elicited movements a small subset was retained, very subtle movements being discarded because their onset and offset frames could not be reliably located. Each retained sample carries onset and offset frame indices, Facial Action Coding System (FACS) action-unit labels from two independent coders with disagreements arbitrated by discussion, and an emotion label. Yan et al.~\cite{yan2014} argue against forcing spontaneous micro-expressions into the six basic categories, since a single action unit is frequently ambiguous --- AU4 (brow lowerer) may indicate disgust, anger, attention or tension --- and assign the emotion label from a \textit{combination} of the coded AUs, the participant's self-report, and the emotional content of the eliciting episode.

\textbf{Preprocessing in the original release.} The baseline applies a three-step normalisation from which this thesis departs: a frontal neutral model face is chosen; 68 Active Shape Model landmarks are detected on the \textit{first frame only} of each clip and registered to the model face by a Local Weighted Mean transform; that transform is applied unchanged to every frame, and the face is cropped by eye position. Yan et al.~\cite{yan2014} justify landmarking only the first frame on two grounds --- that rigid head motion within so short a window is negligible, and that per-frame ASM landmarks are unstable enough to introduce jitter where the face did not in fact move.

\textbf{A note on sample count.} The published paper reports \textbf{247} micro-expressions from 26 valid participants; the coding file distributed with the released database contains \textbf{255}. This discrepancy, a consequence of post-publication revision of the coding, is rarely acknowledged and is one reason sample counts reported for ``CASME II'' vary between papers. The present work uses the released 255-sample coding.

\subsection{The label taxonomy, and the sparsity problem it creates}
\label{sec:label-taxonomy}

Yan et al.~\cite{yan2014} provide five main emotion categories with AU-based criteria, and are candid about the resulting imbalance: ``some types of facial expressions are difficult to elicit in laboratory situations, thus the samples in different categories distributed unequally.'' The released coding is finer-grained still, exposing seven raw labels; counting directly from the label table used in this thesis (\texttt{Processed\_Data/master\_thesis\_labels.csv}, 255 rows) gives the middle columns of \autoref{tab:casme2-labels}.

\begin{table}[htbp]
	\centering
	\footnotesize
	\begin{tabular}{l l c c l c}
		\hline
		\multicolumn{3}{c}{\textbf{Five categories, as published}} & \multicolumn{2}{c}{\textbf{Released coding (255)}} & \textbf{Grouped} \\
		\textbf{Emotion} & \textbf{AU criterion} & \textbf{N} & \textbf{N} & \textbf{Raw label} & \textbf{class} \\
		\hline
		Happiness  & AU6 or AU12                           & 33  & 32 & happiness  & Positive \\
		Disgust    & AU9, AU10, or AU4+AU7                 & 60  & 63 & disgust    & Negative \\
		Surprise   & AU1+2, AU25, or AU2                   & 25  & 25 & surprise   & Surprise \\
		Repression & AU15 or AU17, alone or combined       & 27  & 27 & repression & Negative \\
		Others     & other emotion-related facial movement & 102 & 99 & others     & \textit{discarded} \\
		---        & ---                                   & --- & 7  & sadness    & Negative \\
		---        & ---                                   & --- & 2  & fear       & Negative \\
		\hline
		\multicolumn{6}{l}{\textbf{Working set:} Negative $63+27+7+2 = $ \textbf{99}; Positive \textbf{32}; Surprise \textbf{25}; \textbf{N = 156}} \\
		\hline
	\end{tabular}
	\caption[CASME II label taxonomy and the three-class working set]{The CASME II label taxonomy at three granularities: the five published categories with their AU criteria and frequencies, after Yan et al.~\cite{yan2014} (Table 3); the seven raw labels of the released coding as parsed for this thesis; and the three-class valence grouping that yields the 156-clip working set used throughout.}
	\label{tab:casme2-labels}
\end{table}

\textbf{Seven-class classification is not viable on this corpus.} Fear is represented by two clips in the entire database and sadness by seven. Under any subject-disjoint protocol --- and \autoref{sec:original-baseline} shows subject-disjointness has been the norm for this benchmark since its introduction --- a two-clip class cannot be meaningfully trained or tested: a fold either contains the class or does not, and per-class metrics degenerate.

The field's response has been to group by affective valence. The convention originates with SMIC, released with exactly three classes --- positive, negative, surprise~\cite{liX2013} --- and was carried into cross-corpus work by the Second Facial Micro-Expression Grand Challenge (MEGC 2019)~\cite{see2019}. MEGC states the scheme twice and not identically: as a challenge-wide rule in prose, with Negative covering Repression, Anger, Contempt, Disgust, Fear and Sadness, and again as per-corpus footnotes under its sample table, where CASME II's Negative narrows to Disgust and Repression alone. \autoref{tab:casme2-labels} follows the challenge-wide rule.

\begin{figure}[htbp]
	\centering
	\includegraphics[width=\textwidth]{fig3_1_dataset.png}
	\caption[Composition of CASME II]{Composition of CASME II. Left: all 255 released clips by raw label, with the excluded ``others'' category in grey. Right: the three-class pool of 156 clips used throughout this thesis, in the ratio 99 : 32 : 25; the grey ``others'' bar is shown again for scale and is not part of the pool.}
	\label{fig:dataset-composition}
\end{figure}

\textbf{A declared divergence from the MEGC convention.} MEGC 2019 reports the CASME II portion of its composite database as \textbf{145} samples and footnotes that ``Negative class of CASME II: Disgust and Repression''~\cite{see2019} --- building its subset from the paper's \textit{five-category} release, in which sadness and fear are absorbed into ``others''. This thesis uses the released seven-label coding and assigns sadness and fear to Negative, following the mapping MEGC gives for \textbf{SAMM} rather than the one it gives for CASME II. The consequence is an \textbf{eleven-clip} difference: N = 156 against N = 145, with a Negative class of 99 rather than 88. Nine of the eleven are the sadness (7) and fear (2) clips the seven-label coding exposes; the remaining two are unaccounted for, since disgust and repression together number 90 in the released coding against the 88 the challenge reports --- a residual noted rather than resolved. The direction is mild: the nine added clips are all in the majority class and drawn from the categories the coders found hardest to assign, so if anything they make the present task marginally harder than the challenge's.

The residual imbalance is severe in either case --- 99 : 32 : 25, roughly 4 : 1.3 : 1 --- and everything in \autoref{sec:corpus-threats} concerning metric choice, sampling and loss follows from it.

\subsection{The original baseline, and what it does and does not establish}
\label{sec:original-baseline}

Yan et al.~\cite{yan2014} provide a baseline using Local Binary Patterns from Three Orthogonal Planes (LBP-TOP)~\cite{zhaoG2007} with a Support Vector Machine. LBP-TOP extends the static LBP operator by extracting codes on the XY, XT and YT planes of the video volume and concatenating the three histograms, so that the XT and YT planes encode dynamic rather than static texture. The baseline divides the face into 5~$\times$~5 blocks, sets $P = 4$ neighbouring points, and sweeps the spatial radii $R_X, R_Y \in \{1,2,3,4\}$ and temporal radius $R_T \in \{2,3,4\}$; $R_T = 1$ is deliberately excluded because at 200~fps the change between adjacent frames is negligible --- an early acknowledgement that a high sampling rate is not automatically an advantage. Five-class accuracies under \textbf{leave-one-subject-out cross-validation} cluster in the high fifties to low sixties, the best configuration ($R_X = R_Y = 1$, $R_T = 4$) reaching \textbf{63.41\,\%}.

Three observations matter here.

\textbf{First, leave-one-subject-out is native to this benchmark.} It is not a stricter protocol adopted later by cautious authors; it is the protocol of the paper that released the corpus. Any evaluation reporting a single random train/test split departs from the benchmark's own convention.

\textbf{Second, the baseline is reported as accuracy alone}, with no per-class or macro-averaged score. On a five-class problem in which ``others'' accounts for 102 of 247 samples, accuracy is a weak instrument: the majority-class-constant classifier scores 41\,\% without modelling anything. In the three-class working set the floor is higher still --- always predicting Negative yields \textbf{0.6346 accuracy} at a macro F1 of only \textbf{0.2588}. A headline accuracy near 63.41\,\% is therefore not, on its own, evidence of a working system.

\textbf{A caution on this figure.} Yan et al.~\cite{yan2014} state in both their abstract and their results prose that leave-one-subject-out was applied to obtain 63.41\,\% --- but the caption of the very table reporting it reads ``leave-one-out cross-validation'', without the subject qualifier. Li et al.~\cite{liX2018}, reproducing the same figure in their comparison table, mark it with an asterisk denoting ``results achieved using leave-one-sample-out cross validation'' --- a materially easier protocol. The sources therefore disagree on the protocol behind the most widely quoted number on this benchmark, and the originating paper disagrees with itself. This thesis follows that paper's prose, but records the discrepancy rather than resolving it, because it affects how the figure should be read in any comparison.

\textbf{Third, the baseline's sensitivity to a purely descriptive hyper-parameter is itself informative.} Accuracy moves by several points across the radius sweep with no change to the data, the split or the classifier. Where effect sizes of that magnitude arise from descriptor settings, differences of similar magnitude between architectures cannot be interpreted without a variance estimate --- a limitation this thesis inherits and declares.

\subsection{Sibling corpora, and why this study remains single-corpus}
\label{sec:sibling-corpora}

Three further databases define the context in which CASME II is used; \autoref{tab:spontaneous-corpora} gives their properties, and each isolates a specific limitation of CASME II.

\textbf{SMIC~\cite{liX2013}}, the first spontaneous micro-expression database, is the origin of the three-class taxonomy. Its high-speed subset of 164 clips is distributed 51 positive, 70 negative, 43 surprise, and its baseline --- LBP-TOP with a temporal interpolation model (TIM) normalising each clip to ten frames, under leave-one-subject-out --- reaches \textbf{48.78\,\%}. Two elements of its design are inherited directly here: the three-class valence grouping, and temporal interpolation to normalise clips of unequal length. It is not used as a corpus because it is neither FACS-coded nor apex-annotated.

\textbf{SAMM~\cite{davison2018}} addresses the demographic narrowness of the CASME family, contributing 13 ethnicities and a mean age of 33.24 years against CASME II's single ethnicity and 22.03. Davison et al.~\cite{davison2018} observe that a corpus drawn from one ethnicity and one narrow age band restricts analysis to ``similar looking participants'', and tailor stimulus selection to each participant in advance rather than relying on post-hoc self-report. The design is stronger than CASME II's on both counts, but the corpus is smaller in usable samples and spreads them across seven categories with the same sparsity difficulty.

\textbf{CAS(ME)$^2$~\cite{qu2016}} provides 250 macro-expression and 53 micro-expression samples from the same participants under identical conditions, which makes it the natural resource for studying macro-to-micro domain transfer. That strategy produced the strongest entry in MEGC 2019 --- Liu et al.'s EMR~\cite{liu2019}, the only reported submission to use macro-expression data, though it drew that data from CK+ rather than from CAS(ME)$^2$~\cite{see2019}. Two properties limit its use here: 30~fps is an order of magnitude below CASME II's and poorly matched to sub-200-ms movements, and only 53 micro-expression samples exist. Its baseline is also evaluated leave-\textit{one-video}-out, which does not enforce subject-disjointness and is therefore not comparable to the protocol used here.

\textbf{Why this thesis uses CASME II alone.} It has the highest temporal resolution of any spontaneous corpus of comparable size, a precondition for magnification in a defined frequency band and for dense frame-to-frame flow; it is FACS-coded with onset, apex and offset annotation, which makes temporal normalisation to a fixed length well-defined rather than heuristic; and it remains the most widely reported single-corpus benchmark, which makes comparison against published results meaningful. The cost --- that no claim here extends to other populations, capture conditions or corpora --- is stated as a declared limitation rather than argued away.

\subsection{How CASME II is evaluated in the field: MEGC 2019 and the metric convention}
\label{sec:megc-metrics}

The most consequential development in how CASME II is \textit{used}, as opposed to how it was built, is the Composite Database Evaluation (CDE) protocol of the Second Facial Micro-Expression Grand Challenge~\cite{see2019}. This thesis adopts the challenge's metric definitions while declining its composite training regime, so the protocol warrants precise description.

CDE merges SMIC, CASME II and SAMM into one composite database under the three-class grouping of \autoref{sec:label-taxonomy}, giving \textbf{442 samples} --- Negative 250, Positive 109, Surprise 83 --- of which CASME II contributes 145. Leave-one-subject-out is run across the composite, which the organisers note ``ensures subject-independent evaluation''. Plain accuracy is rejected explicitly on grounds of class imbalance, and two balanced metrics mandated: \textbf{Unweighted F1 (UF1)}, the macro-averaged F1-score, which \textit{accumulates} true positives, false positives and false negatives per class \textbf{over all $k$ folds}, computes a per-class F1 from those pooled counts and averages without weighting; and \textbf{Unweighted Average Recall (UAR)}, the balanced accuracy, obtained by averaging per-class recall.

\[
	F1_c = \frac{2 \cdot TP_c}{2 \cdot TP_c + FP_c + FN_c}, \qquad \text{UF1} = \frac{1}{C}\sum_c F1_c
\]

UF1 is a \textbf{pooled} quantity: counts accumulate across folds \textit{before} any F1 is computed, which \autoref{sec:pooled-vs-fold} establishes is a different estimator from the average of per-fold macro F1 scores. The distinction is not cosmetic here. Because clips are distributed very unequally across subjects, and many single-subject folds contain only one of the three classes, a per-fold macro F1 is bounded well below 1 by fold composition alone and its average across folds bounded correspondingly (\autoref{fig:fold-composition}). The primary metric of this thesis, \textit{pooled macro F1}, is therefore identical in construction to MEGC's UF1; the per-fold average appearing in the raw experimental logs is not a metric the field uses at all.

\begin{figure}[htbp]
	\centering
	\includegraphics[width=\textwidth]{fig3_2_fold_composition.png}
	\caption[Class composition of the 25 leave-one-subject-out folds]{Class composition of the 25 leave-one-subject-out folds, and the per-fold macro-F1 ceiling it implies.}
	\label{fig:fold-composition}
\end{figure}

\begin{table}[htbp]
	\centering
	\begin{tabular}{l c c}
		\hline
		\textbf{Method} & \textbf{Full composite UF1 / UAR} & \textbf{CASME II subset UF1 / UAR} \\
		\hline
		LBP-TOP~\cite{zhaoG2007}                  & 0.5882 / 0.5785 & 0.7026 / 0.7429 \\
		Bi-WOOF~\cite{liong2018}                  & 0.6296 / 0.6227 & 0.7805 / 0.8026 \\
		OFF-ApexNet~\cite{liong2019a}             & 0.7196 / 0.7096 & 0.8764 / 0.8681 \\
		Quang et al.\textsuperscript{\textdagger} & 0.6520 / 0.6506 & 0.7068 / 0.7018 \\
		Zhou et al.\textsuperscript{\textdagger}  & 0.7322 / 0.7278 & 0.8621 / 0.8560 \\
		Liong et al.~\cite{liong2019b}, STSTNet   & 0.7353 / 0.7605 & 0.8382 / 0.8686 \\
		Liu et al.~\cite{liu2019}, EMR            & \textbf{0.7885 / 0.7824} & 0.8293 / 0.8209 \\
		\hline
	\end{tabular}
	\caption[MEGC 2019 Composite Database Evaluation results]{MEGC 2019 CDE results, reproduced from See et al.~\cite{see2019} (Table IV). \textdagger{} outside the review corpus; reported here only as it appears in that table.}
	\label{tab:megc-cde-results}
\end{table}

\textbf{One structural caveat governs these numbers.} Every row of \autoref{tab:megc-cde-results} comes from a model trained on the \textit{composite} database --- roughly 442 samples across three corpora and 68 subjects --- then scored on the CASME II portion of the pooled predictions. The CASME II column is therefore not a CASME II-only result: it reports how a cross-corpus model performs on CASME II clips, with approximately \textbf{2.8 times the training data} available in any single-corpus fold, where a single-corpus LOSO experiment of the kind conducted here trains each fold on roughly 150 clips. The two regimes are not interchangeable, and single-corpus results should not be compared directly against \autoref{tab:megc-cde-results} without stating the difference.

Two further observations from the challenge bear on this thesis's design. The top three submissions all chose \textbf{optical flow} as their feature representation, a convergence that supports the motion-based input used here; and one submission substituted mid-position frames for annotated apex frames with little loss, suggesting that precise apex localisation may be less critical than assumed and that methods consuming the full onset-to-offset sequence are not obviously disadvantaged.

\subsection{Threats to validity carried by the corpus}
\label{sec:corpus-threats}

The properties in \autoref{tab:corpus-threats} are not criticisms of CASME II, which remains the best available resource for this task. They are constraints any study built on it inherits, and each is answered by a specific design decision --- which is the sense in which, on a corpus of this size, several choices usually presented as modelling decisions are determined by the data. Fold and duration statistics are computed directly from the 156-clip label table used in this thesis.

\begin{table}[htbp]
	\centering
	\scriptsize
	\begin{tabular}{@{}p{0.3cm} p{2.0cm} p{5.3cm} p{5.6cm}@{}}
		\hline
		& \textbf{Threat} & \textbf{What the corpus contains} & \textbf{Design decision it forces} \\
		\hline
		(a) & Scale
		& 156 usable clips across three classes --- small enough that parameter count is itself a risk factor.
		& \textbf{Parameter-free or low-parameter components} wherever a choice exists, motivating SimAM~\cite{yang2021}, which adds zero learnable parameters; shallow rather than deep spatial backbones, extending the model-shrinking argument of Xia et al.~\cite{xia2020b}, which they establish for the composite-database setting rather than a single small corpus. \\
		(b) & Class imbalance
		& 99 : 32 : 25. The majority-class-constant classifier attains 0.6346 accuracy at a macro F1 of 0.2588.
		& Minority oversampling via a balanced sampler with focal loss, inverse-frequency class weighting \textbf{disabled}, since both corrections together induced single-class collapse. Accuracy is never reported alone, and the always-Negative floor is quoted alongside every result. \\
		(c) & Subject-level non-uniformity
		& \textbf{Subject 17 alone contributes 33 clips, 21\,\% of the working set}; subjects 8, 10 and 21 contribute one each, so fold sizes vary by a factor of 33. \textbf{Ten of the 25 folds contain a single class}, eight contain two, only seven all three.
		& \textbf{Full 25-fold leave-one-subject-out}, every subject held out exactly once, no clip untested, subject-disjointness guaranteed by construction. Primary metric is \textbf{pooled macro F1}, identical in construction to MEGC's UF1~\cite{see2019}; per-fold averaged macro F1 is rejected, since fold composition bounds it below the study's own target. \\
		(d) & A subject that disappears
		& Of the 26 participants, subject 18's three clips are all ``others'' and are removed by the three-class filter.
		& The LOSO run has \textbf{25 folds, not 26} --- stated explicitly for reproducibility, since fold counts differ between papers using this corpus. \\
		(e) & Demographic homogeneity
		& All participants Chinese, mean age 22.03 (SD 1.60) --- named by Davison et al.~\cite{davison2018} as a principal motivation for SAMM. The MEGC organisers saw the effect in reverse: adaptation from CK+ helped on SMIC and SAMM but ``not so for CASME II, which contain predominantly Chinese subjects''~\cite{see2019}.
		& Input is a \textbf{motion representation} --- dense optical flow $(u,v)$ plus optical strain --- rather than raw pixels, since motion is invariant to identity and illumination whereas appearance is exactly what a 26-subject corpus would let a network memorise. \\
		(f) & Laboratory ecology
		& Frontal recording, fixed flicker-free lamps, participants told not to move and not speaking; Yan et al.~\cite{yan2014} flag ``natural conversation and interaction'' as future work. Almost none of the head motion, illumination variation, occlusion or speech-related movement a deployed system would face is present.
		& No cross-dataset or in-the-wild claim is made anywhere in this thesis; generalisation is declared as further work. \\
		(g) & Annotation dependence
		& Labels derive from coded AUs, self-report and stimulus content --- a construct produced by trained coders under a rubric, with disagreements arbitrated by discussion rather than reported as residual uncertainty.
		& The coding is taken as given; its construct status is recorded as a declared limitation rather than treated as measurement error. \\
		(h) & Duration heterogeneity
		& Onset-to-offset lengths range from \textbf{31 to 126 frames}, median 66 --- at 200~fps, \textbf{0.155\,s to 0.63\,s}, median 0.33\,s. The upper end exceeds the 1/2-second criterion of \autoref{sec:phenomenon}.
		& Every clip temporally resampled to a fixed \textbf{T = 32} frames, following the TIM precedent of SMIC's baseline~\cite{liX2013}. \\
		(i) & Version ambiguity
		& Three sample counts circulate under the name ``CASME II'' --- 247 (the paper), 255 (the released coding), 145 (the MEGC subset) --- and both five- and seven-class label sets are in use, rarely stated.
		& The released 255-clip coding is used, N is reported explicitly alongside every figure, and the eleven-clip divergence from the MEGC subset is declared in \autoref{sec:label-taxonomy}. \\
		\hline
	\end{tabular}
	\caption[Threats to validity carried by CASME II and the decisions they force]{Threats to validity carried by CASME II, and the design decision each forces.}
	\label{tab:corpus-threats}
\end{table}

\subsection{Implications for this research}
\label{sec:corpus-implications}

Three further corpus properties are not threats but still determine settings. The \textbf{200 fps} capture rate fixes the Eulerian magnification pass-band at \textbf{5--25 Hz}, matched to the recording rate. The \textbf{seven raw labels}, with fear at two clips and sadness at seven, force the three-class valence grouping of \autoref{sec:label-taxonomy}. And the baseline's \textbf{sensitivity to a descriptive hyper-parameter} (\autoref{sec:original-baseline}) sets the standard of evidence: effect sizes below the resolution of a 156-clip corpus are reported as unresolved rather than as findings, and no confidence interval is available, since the single-seed design and unretained per-clip predictions preclude one.

Three consequences then shape the interpretation of the results rather than merely the construction of the pipeline. \textbf{First, the corpus determines the metric, and the metric determines the conclusion:} the fold composition of \autoref{tab:corpus-threats}(c) caps a macro F1 averaged over folds by arithmetic alone, before any classifier is considered, so reporting the pooled quantity is not a stylistic preference but the correction MEGC 2019 already standardised. \textbf{Second, a single held-out split is not a reliable measuring instrument here:} with 25 subjects distributed as unevenly as \autoref{tab:corpus-threats}(c) describes, which subjects fall in the test set can move the result by more than any architectural change under study, so only a complete leave-one-subject-out run is reported and the incomplete evaluations that preceded it are set out with their exclusion reasons rather than suppressed. \textbf{Third, the small-sample regime predicts which components should help:} 156 clips cannot supply the data needed to learn a general spatio-temporal feature extractor from scratch, particularly once the input has been reduced to flow and strain --- that is, once much of that extraction has been performed analytically before the network sees the data --- whereas a component modelling the temporal arc of an expression exploits exactly the structure the corpus's onset/apex/offset annotation guarantees is present and correctly delimited.

\subsection{Limitations of the existing corpus literature, and how this work differs}
\label{sec:corpus-lit-limitations}

The corpus literature establishes three things well and one thing poorly. It establishes the phenomenon's definitional boundaries, an elicitation methodology producing genuinely spontaneous samples, and a capture standard --- 200 fps, flicker-free illumination, FACS coding with onset/apex/offset annotation --- that makes fine-grained motion analysis feasible. What it does not establish is a stable basis for comparing methods. The original baseline reports accuracy only, on a five-class problem whose majority class is 41\,\% of the data, and its own hyper-parameter sweep moves the headline figure by several points with no change to the model. Sample counts, label granularity and class groupings differ across papers using the same database --- 247, 255 and 145 all appear under the name ``CASME II'' --- and are frequently not stated. MEGC 2019's CDE protocol fixed the metric problem convincingly, but within a composite training regime that makes its per-database columns non-comparable to single-corpus work, a distinction subsequent papers routinely elide. Above all, only one study in this group quantifies how much of a reported difference between methods is attributable to the evaluation protocol rather than the methods themselves --- Li et al.~\cite{liX2018} add a single matched row of their own method under both protocols --- and none does so across methods, despite a subject distribution that makes the question unavoidable.

This thesis differs in three respects. It runs the \textbf{complete} 25-fold leave-one-subject-out protocol over \textbf{every} valid cell of a four-component ablation matrix rather than reporting a single proposed configuration, so each component's contribution is measured from matched pairs differing in exactly one factor. It adopts MEGC's pooled macro F1 as its primary metric while explicitly declining the composite training regime, so its numbers describe single-corpus performance and are labelled as such. And it reports one complete protocol rather than the strongest of several, setting out which earlier evaluations were run and why each is excluded rather than suppressing that history. The intended contribution is therefore not a recognition result on CASME II, but a measurement of how much of such a result is attributable to each stage of the pipeline that produces it.

The 200 fps capture rate established here is what lets Eulerian Video Magnification, reviewed next, target a matched temporal band.


% =============================================================================
\section{Motion Magnification: Eulerian Video Magnification as a Pre-Processing Stage}
\label{sec:evm-review}
% =============================================================================

This section reviews Eulerian Video Magnification (EVM) through the four review-corpus papers that apply it: Bai et al.~\cite{bai2021}, Li et al.~\cite{liX2018}, and Li, Huang and Zhao~\cite{liY2018,liY2021}. The original formulations they build on --- Wu et al.'s amplitude-based method and Wadhwa et al.'s phase-based method --- lie outside the corpus and are described only as attributed through those four papers.

\subsection{The problem magnification addresses}
\label{sec:evm-problem}

The corpus design of \autoref{sec:why-spontaneous} responds to the \textit{short-duration} half of the micro-expression definition but not to the \textit{low-intensity} half. CASME II's 200~fps sampling spreads a movement completing in under 200~ms across dozens of frames~\cite{yan2014}, yet that movement may still occupy only a few grey levels of change across a small patch of skin --- and Yan et al.~\cite{yan2014} report discarding samples outright because the movement was too faint for coders to place an onset or offset. Li et al.~\cite{liX2018} state the difficulty directly: ``the intensity levels of facial movements are too low to be distinguishable.'' Their remedy is to amplify the motion itself before any descriptor is computed, so that the \textit{difference between expression categories is enlarged} rather than the classifier being made more sensitive. This is a data-level intervention: the model is untouched and the modification is entirely upstream, which is what makes EVM cleanly separable as an ablation factor in this thesis.

\subsection{How Eulerian magnification works, and the two variants}
\label{sec:evm-how}

\textbf{Amplitude-based magnification (AMM)} is the formulation set out in \autoref{sec:motion-magnification} --- Laplacian pyramid, temporal band-pass, multiplication by $\alpha$, reconstruction --- realising its band-pass step with a \textbf{Butterworth} filter. It is the variant implemented in this thesis. \textbf{Phase-based magnification (PMM)}, of which Bai et al.~\cite{bai2021} give the clearest account in the corpus, decomposes each frame with \textbf{octave complex steerable pyramids} over four orientations and applies the temporal filter to the \textit{local phase} rather than the amplitude. Bai et al.~\cite{bai2021} explain the advantage in signal terms: the steerable pyramid ``has impulse response with spatial support,'' which makes it easier to isolate the intended temporal frequencies while suppressing the remainder. The practical consequences they report are that the phase-based method \textbf{supports larger amplification factors} and is \textbf{notably less sensitive to noise}.

Three parameters are therefore free in any EVM application, and all three recur below: the \textbf{temporal pass-band}, which decides \textit{which} motions are amplified; the \textbf{magnification factor $\alpha$}, which decides \textit{how much}; and the \textbf{pyramid depth}, which decides at which spatial scales amplification occurs.

\textbf{Setting the temporal band.} Bai et al.~\cite{bai2021} report the rule given by the original EVM authors: the pass-band should correspond to the duration of the phenomenon being amplified --- for micro-expressions, 1/5\,s to 1/25\,s, equivalently \textbf{5 Hz to 25 Hz}. They also report that applying this full band produced ``noticeable noise'' in their setting, and that they narrowed it to 15--25 Hz accordingly. This is the only explicit guidance on band selection anywhere in the review corpus, and it is the basis for the band used in this thesis.

\subsection{Evidence that magnification improves recognition}
\label{sec:evm-evidence}

The strongest empirical case in the review corpus is Li et al.~\cite{liX2018}, who sweep the magnification factor systematically rather than fixing it. Clips are magnified at what the paper describes as ten levels while listing nine values, $\alpha \in \{1, 2, 4, 8, 12, 16, 20, 24, 30\}$ --- an inconsistency in the source rather than in this account --- where $\alpha = 1$ denotes no magnification. Each magnified clip is temporally interpolated to ten frames (TIM10), and LBP, HOG and HIGO descriptors are extracted and classified, on both SMIC and CASME II under leave-one-subject-out. Two findings emerge.

\textbf{Magnification helps, consistently}, ``for all three kinds of features on all four testing datasets''. On CASME II the effect is substantial (\autoref{tab:evm-effect}), and the best configuration of their full framework --- magnification, TIM10 and HIGO --- reaches 67.21\,\% on CASME II and 68.29\,\% on SMIC-HS under LOSO, competitive with or better than the contemporaneous state of the art.

\begin{table}[htbp]
	\centering
	\begin{tabular}{l c c c}
		\hline
		\textbf{Descriptor} & \textbf{No magnification} & \textbf{With magnification} & \textbf{$\Delta$} \\
		\hline
		LBP  & 55.87\,\% & 60.73\,\% & \textbf{+4.86} \\
		HOG  & 57.49\,\% & 63.97\,\% & \textbf{+6.48} \\
		HIGO & 57.09\,\% & 67.21\,\% & \textbf{+10.12} \\
		\hline
	\end{tabular}
	\caption[Effect of Eulerian motion magnification on CASME II]{Effect of Eulerian motion magnification on CASME II under leave-one-subject-out, after Li et al.~\cite{liX2018} (Table 6).}
	\label{tab:evm-effect}
\end{table}

\textbf{The relationship between $\alpha$ and accuracy is non-monotonic.} Li et al.~\cite{liX2018} describe the resulting curves as ``rainbow-shaped,'' with best performance ``generally achieved when the motion is magnified in the range of $[8, 16]$,'' and give a mechanistic account of both failure modes: ``magnification at lower levels might not be enough to reveal the ME motion progress; on the other hand, magnification at higher levels degrades the performance because too many artifacts are induced.'' Magnification is therefore not a monotone improvement to be maximised but a parameter with an interior optimum that must be selected.

Bai et al.~\cite{bai2021} provide the second quantitative demonstration using a fundamentally different downstream model --- a pre-trained VGGFace2 (ResNet-50) spatial encoder feeding a bidirectional LSTM. On SMIC they reach \textbf{75.76\,\% with phase-based magnification against a 60.60\,\% baseline} drawn from an earlier VGGFace2 + Bi-LSTM publication rather than from a magnification-off arm of their own, and the phase-based variant outperforms the amplitude-based one. Because the pipeline shares nothing with Li et al.'s hand-crafted descriptors, the two results together establish that the benefit is not an artefact of a particular feature family --- at least when the downstream representation is appearance-based.

\subsection{Magnification in apex-frame pipelines, and the $\alpha$ discrepancy}
\label{sec:evm-apex}

Two further papers use EVM in a different regime, and the contrast is informative. Li, Huang and Zhao~\cite{liY2018} address whether a micro-expression can be recognised from a \textit{single} apex frame, arguing for magnification specifically in that setting: the apex frame carries the most emotional information, but ``the intensity level of apex frame is not distinguished due to subtle change, especially for deep model that describes high-layer abstract information.'' They magnify the spotted apex frame with \textbf{$\alpha = 30$} and fine-tune a VGG-Face descriptor on the result; because a single frame per clip is too little data to train on, they exploit the high frame rate to augment, retaining the apex frame plus its two preceding and two following frames and so expanding the training set fivefold. They restate the trade-off Li et al.~\cite{liX2018} quantified --- ``bigger values of motion amplification level lead to larger scale of motion amplification, but also can cause bigger displacement and artifacts.'' Li, Huang and Zhao~\cite{liY2021} extend this into a joint local--global learning framework with frequency-domain apex detection, carrying the same setting forward --- magnification level ``set as 30 in our framework'' --- with the same $\pm$2-frame augmentation producing what they name the Extended Magnified ME (EMME) database.

\textbf{The $\alpha$ discrepancy is easy to misread as a contradiction.} Li et al.~\cite{liX2018} find the optimum at $\alpha \in [8, 16]$ and observe degradation beyond it; Li, Huang and Zhao~\cite{liY2018,liY2021} adopt $\alpha = 30$, at the extreme end of that same sweep. The regimes differ in a way that plausibly explains the gap. The $[8, 16]$ optimum is measured on \textbf{full magnified sequences} described by spatio-temporal descriptors, where inter-frame artefacts accumulate along the temporal axis and corrupt exactly the dynamic texture the descriptor encodes. The $\alpha = 30$ setting is applied to a \textbf{narrow five-frame window around the apex}, consumed by a deep face encoder as largely static appearance, where temporal artefacts have far less opportunity to express themselves --- and where the encoder needs a strong appearance signal to distinguish categories at all. The implication for any new system is that $\alpha$ cannot be transferred between pipelines: it is a parameter of the \textit{combination} of magnification and downstream representation, not of the corpus.

\subsection{Known limitations of magnification}
\label{sec:evm-limitations}

The review corpus is unusually candid about EVM's failure modes, and four are documented.

\textbf{It amplifies everything in band, not only expressions.} This is the most important limitation and the most easily overlooked. Li et al.~\cite{liX2018} decline to use magnification for micro-expression \textit{spotting} at all, ``because it magnifies unwanted motions (e.g. head movements) at the same time.'' EVM has no notion of which motion is of interest: a rigid head translation whose temporal frequency falls inside the pass-band is amplified exactly as a brow raise is. In a recognition pipeline this is mitigated by prior face alignment and cropping --- Li et al.~\cite{liX2018} magnify only after normalising and cropping to the eye-defined rectangle --- but mitigated, not eliminated.

\textbf{Artefacts grow with $\alpha$.} All four papers report this: Li et al.~\cite{liX2018} as the descending limb of the rainbow curve, Li, Huang and Zhao~\cite{liY2018,liY2021} as a design caution, and Bai et al.~\cite{bai2021} directly, as noise that forced them to narrow their temporal band.

\textbf{There is no principled procedure for choosing $\alpha$.} One value in the corpus is obtained by exhaustive sweep~\cite{liX2018}; the apex-frame papers assert theirs --- Li et al.~\cite{liY2018} state $\alpha = 30$ without support, and Li et al.~\cite{liY2021} carry the value forward by citing them. No paper offers a way to select $\alpha$ from properties of the data.

\textbf{Evaluation is confined to appearance-based downstream representations:} LBP, HOG and HIGO in Li et al.~\cite{liX2018}, VGGFace2 features in Bai et al.~\cite{bai2021}, VGG-Face in both Li, Huang and Zhao papers. Every one of these descriptors is computed on \textbf{magnified pixel intensities}. No study in the review corpus measures what magnification contributes when the representation handed to the classifier is itself an explicit motion field rather than an appearance signal.

\subsection{Implications for this research}
\label{sec:evm-implications}
% sec:evm-motion-hypothesis is referenced from Chapter 5 (chapter5.tex:236). The
% hypothesis it names now opens this subsection instead of having one of its own,
% so the label is kept here and resolves to this subsection.
\label{sec:evm-motion-hypothesis}

\autoref{tab:evm-settings} sets out the settings the review determines. Behind them lies a hypothesis the review corpus does not test, stated here because it frames the interpretation of the measurement. EVM amplifies the \textit{amplitude} of intensity variation in a temporal band, and an appearance descriptor computed on magnified frames sees a directly stronger signal. An optical flow field is a different object: flow estimates \textit{displacement}, and amplifying the underlying intensity change by $\alpha$ does not simply scale the recovered displacement by $\alpha$ --- it changes the conditioning of the estimation problem, potentially making faint motion recoverable where it previously was not, while introducing displacement artefacts of the kind all four reviewed papers warn about. Any residual global scaling of flow magnitude is further attenuated by the per-channel normalisation a learned pipeline applies before training. The plausible outcome is that magnification in front of a flow-and-strain representation retains the part of its benefit that consists of \textit{making otherwise-unrecoverable motion recoverable} while losing the part that consists of \textit{increasing contrast for an appearance descriptor} --- so its measured contribution should be smaller on average than the roughly +5 to +10 percentage points of \autoref{tab:evm-effect}, and concentrated in configurations whose downstream model can exploit finer spatial deformation. EVM is measured across six matched pairs of configurations differing in nothing but the tensor directory, which is what allows that prediction to be tested.

\begin{table}[htbp]
	\centering
	\footnotesize
	\begin{tabular}{@{}p{6.4cm} p{7.4cm}@{}}
		\hline
		\textbf{Finding from the review} & \textbf{Decision taken in this thesis} \\
		\hline
		Pass-band should match the phenomenon's duration, 1/5--1/25\,s $\equiv$ 5--25 Hz~\cite{bai2021}, though Bai et al. found that band noisy and narrowed it to 15--25 Hz & Temporal band-pass set to \textbf{5--25 Hz}, the unnarrowed duration rule; Bai et al.'s narrowing was not followed, and this is a declared divergence \\
		Accuracy against $\alpha$ is rainbow-shaped, optimum in \textbf{$[8, 16]$} for full-sequence pipelines~\cite{liX2018} & \textbf{$\alpha = 10$}, inside the reported optimum; the $\alpha = 30$ of the apex-frame papers is explicitly \textit{not} adopted, since this thesis consumes full sequences \\
		Amplitude-based magnification uses a Laplacian pyramid plus temporal band-pass~\cite{bai2021} & Amplitude-based EVM over a \textbf{4-level Laplacian pyramid}, temporal filtering by FFT band-pass \\
		Magnification amplifies head motion as well as expression~\cite{liX2018} & Magnification applied only after the clip has been trimmed onset-to-offset; the risk is acknowledged as a residual threat rather than claimed to be eliminated \\
		Magnification is a purely upstream, data-level operation & Stage 1 is run \textbf{twice}, producing a magnified and a non-magnified tensor set; the EVM switch selects a directory and leaves the network unchanged, making it a clean ablation factor \\
		Magnification is part of the proposed data pipeline, not of the network & EVM is treated as the \textbf{baseline condition} (\texttt{config\_4}), not as a component to be added; the non-magnified arm (\texttt{config\_1}) is the control that isolates it \\
		\hline
	\end{tabular}
	\caption[Review findings and the magnification settings they determine]{Review findings and the magnification settings they determine.}
	\label{tab:evm-settings}
\end{table}

\textbf{A defect found and repaired, which the review made detectable.} In this project's earliest evaluations, each magnification-enabled configuration produced results \textit{bit-identical} to its magnification-disabled counterpart --- identical accuracies, macro F1 scores and confusion matrices. Given that all four reviewed papers report clear and repeatable effects from magnification, an effect of exactly zero across six independent pairs is not a plausible null result but a symptom: both arms were reading the same tensor directory, so half of the ablation matrix was duplicated and the magnification hypothesis had never in fact been exercised. The defect was repaired before the reported run, and wiring the existing verification helper in as a start-up assertion is a declared item of further work.

\subsection{Limitations of the existing work, and how this study differs}
\label{sec:evm-gap}

What the reviewed literature does not establish is threefold. First, every evaluation in the corpus places magnification in front of an \textbf{appearance-based} representation, so the reported gains conflate two distinct mechanisms --- making faint motion recoverable, and increasing contrast for a texture or gradient descriptor --- and no experiment separates them. Second, the magnification factor is selected empirically in every case and transferred between studies without re-validation, despite the $\alpha$ discrepancy of \autoref{sec:evm-apex}. Third, and most consequentially, magnification is always evaluated as an \textit{addition to a fixed pipeline} --- a two-cell comparison, with and without --- so its interaction with the other components of a system is never measured.

This thesis differs on all three counts. It places magnification in front of an explicit \textbf{motion} representation --- dense optical flow and optical strain --- rather than an appearance representation, isolating the recoverable-motion mechanism from the appearance-contrast one. It selects $\alpha$ from the full-sequence optimum rather than the apex-frame value, with the choice justified rather than inherited. And it evaluates magnification as one factor in a \textbf{full factorial ablation}, so its effect is reported not as a single delta but as six matched-pair measurements whose spread across architectures is itself the finding.

What magnification is placed in front of therefore matters as much as its own parameters, and \autoref{sec:flow-review} turns to that motion representation.


\section{Motion Representation: Dense Optical Flow}
\label{sec:flow-review}

Dense optical flow supplies the horizontal and vertical displacement channels used in this thesis. The relevant literature concerns three choices: the estimator, the frame pairs it processes, and whether flow serves as a recognition feature or an alignment tool. These choices determine what information remains available to the learned model.

\subsection{Why motion rather than appearance}
\label{sec:why-flow}

Xu et al.~\cite{xu2017} motivate facial-dynamics descriptors by the difficulty of interpreting micro-expressions through texture alone. Flow describes apparent movement directly, whereas image intensity also records appearance and illumination. This makes displacement a useful representation of subtle facial activity, although it does not guarantee independence from identity or recording conditions: facial geometry, alignment and brightness changes still affect its estimation.

The benchmark evidence supports this choice. The MEGC 2019 organisers report that their three highest-ranked submissions used optical flow~\cite{see2019}. OFF-ApexNet~\cite{liong2019a} uses horizontal and vertical flow as neural-network inputs; STSTNet~\cite{liong2019b} adds strain; Zhao et al.~\cite{zhaoS2021} retain a sequence of flow fields. Their convergence motivates a motion-based input, but is not a controlled demonstration that it always outperforms appearance-based representations.

\subsection{Estimation assumptions}
\label{sec:brightness-constancy}

OFF-ApexNet~\cite{liong2019a} formulates flow using brightness constancy and small-motion assumptions. A local brightness constraint supplies one equation for two displacement components, so estimation also needs spatial regularisation or a local motion model. The output is a pair of fields, $u(x,y)$ and $v(x,y)$; Zhao et al.~\cite{zhaoS2021} additionally describe their magnitude, $\sqrt{u^2+v^2}$.

These assumptions matter more than repeating the derivation. Illumination changes can resemble movement, while smoothing can suppress a weak local contraction together with noise. Magnification can make subtle motion easier to estimate but can also amplify unwanted movement. Flow therefore transforms the measurement problem; it does not remove uncertainty in the source frames.

\subsection{Choice of estimator}
\label{sec:flow-estimator}

The reviewed methods use three estimators. Zhao et al.~\cite{zhaoS2021} adopt Farneb\"ack through its OpenCV implementation, citing its established availability. It estimates displacement through local polynomial approximations. OFF-ApexNet~\cite{liong2019a} instead uses total-variation regularisation with an L1 data term (TV-L1), citing preservation of motion discontinuities and robustness. Xu et al.~\cite{xu2017} use the method of Sun, Roth and Black. These algorithm attributions are taken from the reviewed papers; the original estimator publications are outside the supplied reference corpus.

STSTNet~\cite{liong2019b} specifies onset--apex optical flow but does not name its estimator. Its similarity to OFF-ApexNet is therefore insufficient to attribute TV-L1 to it. Likewise, the advantages given for TV-L1 should not be read as a measured recognition gain over Farneb\"ack on CASME II. The cited studies do not provide a controlled three-estimator comparison supporting such a ranking. This thesis adopts Farneb\"ack as a fixed preprocessing choice and leaves that comparison open.

\subsection{Onset--apex flow or a temporal sequence}
\label{sec:flow-what-between}

OFF-ApexNet~\cite{liong2019a} represents each clip by flow from onset to apex, where expression intensity is greatest, and resizes the two component maps to $28\times28$. STSTNet~\cite{liong2019b} uses the same frame-pair construction. This is an economical representation of the displacement from neutral to peak expression, suitable for a small training set. It omits the intermediate trajectory and the relaxation after the apex.

Zhao et al.~\cite{zhaoS2021} instead construct an eleven-frame sequence using onset, apex and offset annotations, then calculate ten adjacent-pair flow fields. This retains changes in motion across the expression. The distinction is between one displacement summary and a sequence whose order a temporal model can exploit; it is not simply a difference in input size.

The present study follows the sequence approach but uses uniform sampling of existing frames rather than Zhao et al.'s key-frame construction. This preserves a temporal input for the transformer ablation. It does not, by itself, measure the recognition cost of replacing a sequence with an onset--apex pair: that representation comparison is not an experimental arm.

\subsection{Flow for fine alignment}
\label{sec:flow-alignment}

Xu et al.~\cite{xu2017} also use flow to refine alignment. Their pipeline first locates landmarks and coarsely aligns the face, then performs fine in-sequence alignment using pixel-level movement. The motivation is specific to micro-expressions: a small registration error can be comparable to the facial displacement being recognised.

This use differs from supplying flow to a classifier. An uncorrected displacement field contains residual head movement as well as expression. Magnifying the frames beforehand can increase both. The present pipeline performs no flow-based fine alignment, so corpus preprocessing must not be treated as proof that this source of error has been eliminated.

\subsection{Limitations of the representation}
\label{sec:flow-limitations}

Flow remains sensitive to brightness-constancy violations, alignment errors and estimator settings. Apex-based systems additionally depend on selecting an informative frame. OFF-ApexNet~\cite{liong2019a} attributes weaker performance on SMIC to imprecise apex selection and cites an earlier apex-spotting study reporting an average discrepancy of thirteen frames. That reported average is neither a maximum error nor a measurement made by this thesis.

Using an onset--offset sequence avoids dependence on one selected apex, but still relies on clip boundaries and temporal sampling. Repeated sampled frames can yield zero-flow pairs, and large sampling intervals can make the small-motion approximation less appropriate. These limitations constrain the representation before any network component is compared.

\subsection{Implications for this research}
\label{sec:flow-implications}

\begin{table}[htbp]
    \centering
    \small
    \begin{tabular}{@{}p{0.29\linewidth} p{0.65\linewidth}@{}}
        \hline
        \textbf{Review finding} & \textbf{Consequence for this study} \\
        \hline
        Flow is a recognition input~\cite{liong2019a,see2019}
        & Horizontal and vertical displacement form the first two channels; the classifier receives motion tensors rather than raw frames. \\
        Sequence flow preserves temporal variation~\cite{zhaoS2021}
        & Adjacent-pair fields are retained across the sampled clip, allowing the temporal encoder to operate on an ordered sequence. \\
        Estimator choices differ~\cite{liong2019a,zhaoS2021,xu2017}
        & Farneb\"ack remains fixed; a matched comparison with TV-L1 is an untested extension. \\
        Fine alignment matters~\cite{xu2017}
        & No flow-based fine alignment is performed. It remains a priority for improving preprocessing. \\
        \hline
    \end{tabular}
    \caption[Optical-flow evidence and study decisions]{Optical-flow evidence and the decisions it motivates. Fixed preprocessing choices constrain the interpretation of the architectural ablation.}
    \label{tab:flow-decisions}
\end{table}

The central interpretive point is that optical flow is an \textit{analytical spatio-temporal feature extractor}: it estimates motion from paired frames before the network is trained. A learned spatial stem therefore acts on a representation that already makes displacement explicit. With only 156 clips, learning additional spatial filters may offer limited benefit or duplicate structure already exposed by preprocessing. This is a hypothesis about the stem's marginal value, not proof that learned spatial processing is unnecessary.

A temporal encoder addresses a different task: modelling relations among the successive displacement fields. Pairwise flow estimation does not itself learn those sequence-level relations. Keeping these roles distinct is necessary when attributing performance to the stem and transformer.

\subsection{Research gap and scope}
\label{sec:flow-gap}

The relevant unresolved question is what learned spatial and temporal processing contributes \textit{after} motion has been extracted analytically. The thesis addresses this through matched architectural comparisons under complete leave-one-subject-out (LOSO) evaluation, retaining the same flow construction while switching the learned components. Magnification is varied separately through the paired preprocessing conditions.

The design neither ranks optical-flow estimators nor isolates alignment quality or onset--apex versus sequence representations. Those remain limitations rather than gaps claimed as solved. The resulting architectural conclusions apply to the implemented motion representation and evaluation setting; they do not establish a universal ordering of network families.


\section{Deformation Representation: The Optical Strain Tensor}
\label{sec:strain-review}

Optical strain supplies the third input channel. Its literature establishes a deformation descriptor derived from flow, two ways to use that descriptor, and limitations relevant to interpreting the architectural ablation. Shreve et al.~\cite{shreve2011}, the Liong strain series~\cite{liong2014a,liong2014b,liong2016}, and STSTNet~\cite{liong2019b} provide the principal evidence.

\subsection{What strain contributes}
\label{sec:strain-what-adds}

Flow describes displacement; strain describes its local spatial variation. A spatially uniform translation has zero displacement gradient, whereas neighbouring skin regions moving differently produce deformation. Shreve et al.~\cite{shreve2011} use this distinction to detect facial activity and describe robustness to moderate head translation. It does not imply immunity to arbitrary head movement or registration error.

Their illumination- and make-up-robustness statements are attributed to earlier studies rather than measured anew. They should therefore be read as reported motivation, not as evidence that the present pipeline is robust under those conditions. Strain can suppress a uniform displacement while still inheriting errors from the estimated flow.

\subsection{Tensor components and magnitude}
\label{sec:strain-tensor}

The relevant representation is the symmetric displacement-gradient tensor, with normal components $\varepsilon_{xx}$ and $\varepsilon_{yy}$ and shear components $\varepsilon_{xy}=\varepsilon_{yx}$. Shreve et al.~\cite{shreve2011} approximate spatial derivatives by central differences with offsets of approximately two to three pixels; Liong et al.~\cite{liong2014a} use one-pixel offsets. This discretisation is a modelling choice because it sets the scale at which local variation is measured.

Liong et al.~\cite{liong2014a,liong2014b,liong2016} use the four-component magnitude
\[
    \varepsilon_{\mathrm{F}}
    =\sqrt{\varepsilon_{xx}^{2}+\varepsilon_{yy}^{2}
        +\varepsilon_{xy}^{2}+\varepsilon_{yx}^{2}}
    =\sqrt{\varepsilon_{xx}^{2}+\varepsilon_{yy}^{2}
        +2\varepsilon_{xy}^{2}}.
\]
The scalar describes deformation intensity and discards its sign and direction. Supplying it together with $u$ and $v$ retains the displacement fields from which those spatial relations were calculated.

\subsection{The implemented magnitude and its divergence}
\label{sec:strain-divergence}

The implementation instead computes
\[
    \varepsilon_{\mathrm{os}}
    =\sqrt{\varepsilon_{xx}^{2}+\varepsilon_{yy}^{2}
        +\varepsilon_{xy}^{2}},
    \qquad
    \varepsilon_{xy}=\tfrac{1}{2}
        \left(\frac{\partial u}{\partial y}
             +\frac{\partial v}{\partial x}\right).
\]
It counts shear once rather than twice. NumPy's \texttt{gradient} supplies spatial differences at one-pixel spacing, using central differences in the array interior. The spacing follows Liong et al.~\cite{liong2014a}, but the scalar reduction differs from their convention.

STSTNet~\cite{liong2019b} cannot establish a clear precedent for the three-term form. Its printed magnitude equation repeats $\partial u$ in the mixed derivative where the symmetric tensor calls for $\partial v/\partial x$ alongside $\partial u/\partial y$. It also places the factor $\tfrac{1}{2}$ \emph{outside} the squared derivative sum. Correcting the repeated derivative while retaining that placement gives
\[
    \tfrac{1}{2}
       \left(\frac{\partial u}{\partial y}
            +\frac{\partial v}{\partial x}\right)^2
       =2\varepsilon_{xy}^{2},
\]
which agrees with the four-term convention. Moving the factor inside the square would instead give the implemented three-term form, but that is a further change to the printed expression. The intended implementation cannot be settled from this equation alone; the thesis therefore declares a departure from the unambiguous four-term convention rather than claiming to reproduce STSTNet's magnitude.

The four-term form weights pure shear by $\sqrt{2}$ relative to the implemented form. Independent min--max normalisation removes a uniform scale factor, but does not generally remove this difference when normal and shear contributions vary across pixels and time. Neither convention was compared experimentally here, so their practical equivalence is not assumed.

\subsection{Strain as a feature, a weight, or an input channel}
\label{sec:strain-feature-weight}

\textbf{Optical strain features (OSF)} use a temporally pooled strain map as a descriptor~\cite{liong2014a,liong2014b,liong2016}. Averaging accounts for sequence length: the early method already divides its accumulated map by the number of frame pairs~\cite{liong2014a}. The maps are filtered and normalised before classification. The early work explicitly applies Wiener and Gaussian filters to suppress noise~\cite{liong2014a}.

\textbf{Optical strain weighted features (OSW)} use pooled strain to weight the regional XY-plane histograms of local binary patterns on three orthogonal planes (LBP-TOP), emphasising active facial regions~\cite{liong2014b}. Liong et al.~\cite{liong2016} perform spatial mean pooling followed by temporal mean pooling and concatenate OSF and OSW for a combined descriptor. Strain then serves both as a feature and as a spatial weighting signal.

\textbf{STSTNet} supplies strain directly alongside horizontal and vertical onset--apex flow~\cite{liong2019b}. This is the precedent for the thesis's three-channel input structure, independently of the magnitude ambiguity. The present study retains all three channels across successive frame pairs rather than reducing each clip to one map.

\subsection{Recognition evidence and protocol limits}
\label{sec:strain-evidence}

\begin{table}[htbp]
    \centering
    \small
    \begin{tabular}{@{}lccc@{}}
        \hline
        \textbf{Method} & \multicolumn{2}{c}{\textbf{SMIC accuracy (\%)}} & \textbf{CASME II (\%)} \\
        & Micro-average & Macro-average & Accuracy \\
        \hline
        Baseline & 45.73 & 47.76 & 61.94 \\
        OSF & 41.46 & 46.00 & 51.01 \\
        OSW & 49.39 & 50.71 & 61.94 \\
        OSF + OSW & 52.44 & 58.15 & 63.16 \\
        OFF + OFW & 40.24 & 41.94 & 55.87 \\
        \hline
    \end{tabular}
    \caption[Recognition results for strain-based descriptors]{Recognition results reproduced from Liong et al.~\cite{liong2016}, Table 8. SMIC uses LOSO; CASME II uses leave-one-video-out. The summary combines block settings, so equal entries are not a matched-setting comparison. Macro-average here denotes subject-averaged accuracy, not macro F1.}
    \label{tab:strain-performance}
\end{table}

The CASME II evidence depends on how strain is used and on the comparator. OSF alone falls below the texture baseline. The equal baseline and OSW entries in \autoref{tab:strain-performance} reflect different block settings: under matched settings OSW improves accuracy by 0.81 percentage points. The combined OSF + OSW descriptor improves on the texture baseline by 0.41--1.22 points across block settings~\cite{liong2016}. These accuracy gains are not directly comparable with pooled macro F1, and leave-one-video-out does not enforce subject-disjoint testing.

Crucially, Liong et al.~\cite{liong2016} also test optical-flow features and weights (OFF + OFW) in the same descriptor framework. Their CASME II accuracy is 55.87\,\%, against 63.16\,\% for OSF + OSW, a 7.29-point difference. This is a matched flow-based control, so it would be incorrect to claim that the literature never isolates strain against flow. It concerns hand-crafted descriptors; it does not establish the gain from adding a strain channel to the learned sequence model used here.

\subsection{Limitations}
\label{sec:strain-limitations}

Spatial differentiation can amplify flow-estimation noise. The filtering used by Liong et al.~\cite{liong2014a,liong2016} is therefore relevant to a pipeline without a dedicated strain post-filter. A deterministic transformation of $u$ and $v$ adds no independent observation: its potential value lies in making deformation directly accessible to the classifier. Whether this representation helps must be measured rather than inferred from its construction.

The reviewed recognition methods also simplify time. OSF and OSW pool their sequences; STSTNet uses one onset--apex pair. Shreve et al.~\cite{shreve2011} retain temporal strain activity for spotting, which localises expressions rather than classifying their emotion. The temporal-collapse limitation thus concerns these recognition representations, not every use of strain in the literature.

\subsection{Implications for this research}
\label{sec:strain-implications}

\begin{table}[htbp]
    \centering
    \small
    \begin{tabular}{@{}p{0.29\linewidth} p{0.65\linewidth}@{}}
        \hline
        \textbf{Design choice} & \textbf{Interpretation or unresolved test} \\
        \hline
        Strain accompanies $u$ and $v$
        & Follows STSTNet's channel structure~\cite{liong2019b}; does not establish that its ambiguous magnitude equation was reproduced. \\
        Three-term magnitude
        & Differs from the four-term convention~\cite{liong2014a,liong2014b,liong2016}. Relative shear weighting remains an untested choice. \\
        Per-pair strain sequence
        & Preserves deformation over time instead of pooling it into one descriptor. Each channel is independently min--max normalised across the sequence. \\
        No dedicated strain post-filter
        & Departs from the explicit filtering of Liong et al.~\cite{liong2014a}. Farneb\"ack's internal smoothing is not a separately evaluated strain filter. \\
        Strain held constant
        & Present in all twelve evaluated configurations; its marginal contribution is not measured. \\
        \hline
    \end{tabular}
    \caption[Strain representation and unresolved choices]{The implemented strain representation and the limits it places on interpretation.}
    \label{tab:strain-decisions}
\end{table}

A strain-on/strain-off comparison is the natural next ablation. Neither published gains from a different descriptor nor the performance of the complete system substitutes for that measurement. Comparing magnitude conventions and adding a dedicated post-filter are separate follow-ups.

Strain also extends the analytical-feature argument: the input contains a \textit{hand-specified differential operator}, exposing local deformation without learning that operator from the training clips. A spatial stem may still learn useful combinations of those features, but it is operating after both motion estimation and deformation extraction. This motivates testing its incremental contribution; it does not establish the cause of any observed architectural effect.

\subsection{Research gap and scope}
\label{sec:strain-gap}

The study preserves the temporal evolution of strain for learned sequence processing. Its contribution is an architectural comparison on that representation, not an isolated demonstration of strain's benefit. The independent value of the third channel, the choice of shear weighting and the omitted post-filter remain unresolved. These boundaries preserve the distinction between evidence for a component, a reason to include it, and a result that the experimental design can actually establish.


\section{Temporal Normalisation: The Temporal Interpolation Model}
\label{sec:tim-review}

Flow and strain fix \textit{what} is computed from each frame pair; temporal normalisation fixes \textit{which} frame pairs, and how many. The Temporal Interpolation Model (TIM) is credited in this corpus to Zhou et al. by Li et al.~\cite{liX2013,liX2018} and to Pfister et al. by Lu et al.~\cite{lu2015}, who applied it to micro-expressions outside the review corpus. It is described here only through the reviewed papers that implement and evaluate it.

\subsection{The problem: sequences of unequal length}
\label{sec:tim-problem}

Micro-expressions vary in duration, and the three-class CASME II working set spans roughly a four-fold range of sequence length (\autoref{sec:corpus-threats}h). Three difficulties follow. Any batched network requires a fixed input shape, and the corpus supplies no constant. A descriptor with a temporal radius cannot be applied to a sequence shorter than that radius, so the shortest clips constrain the descriptor for the whole dataset: Li et al.~\cite{liX2018} note that at 25~fps some micro-expressions last four to five frames, restricting LBP-TOP to radius $r = 1$. And duration is itself a nuisance variable --- Lu et al.~\cite{lu2015} normalise to remove ``temporal fluctuations of micro-expressions, which may introduce noise irrelevant to classification''.

\subsection{What TIM is}
\label{sec:tim-what}

Lu et al.~\cite{lu2015} give the corpus's most compact description: ``TIM is a manifold-based interpolation method, which builds a low-dimensional embedding of an image sequence, and then interpolates a curve in the low-dimensional space. The interpolated frames are mapped back to the original high-dimensional space.''

The property that matters downstream is that TIM \textbf{synthesises frames that were never recorded}: Lu et al.~\cite{lu2015} note that Pfister et al. proposed it ``to generate sufficient frames from an image sequence'', that is, as an upsampler rather than only a length-standardiser.

\subsection{TIM in practice: the convergence on ten frames}
\label{sec:tim-ten-frames}

TIM appears throughout the corpus as an unremarked preprocessing step, almost always at ten frames: SMIC's baseline~\cite{liX2013} normalises to ten before LBP-TOP; Li et al.~\cite{liX2018} apply TIM10 after magnification; and Lu et al.~\cite{lu2015} apply TIM before their Delaunay-based coding model, following Pfister et al.'s protocol so as to compare against TIM10 baselines. CAS(ME)$^2$'s baseline~\cite{qu2016} is the exception, normalising to \textbf{120 frames by linear interpolation} and evaluating leave-one-video-out.

One study tests the setting on its own terms. Li et al.~\cite{liX2018} sweep interpolation length from 10 to 80 frames on all three SMIC variants under leave-one-subject-out, holding the descriptor fixed and skipping magnification so only length varies, and report that ``the best performance for all three datasets is achieved with interpolation to 10 frames'', because ``in longer sequences, changes along the time-dimension are diluted''.

\subsection{How many frames? The remaining systematic comparisons}
\label{sec:tim-how-many}

Ben et al.~\cite{ben2021} run a second controlled experiment and reach the opposite conclusion. Holding alignment, descriptor and classifier fixed --- JCFDA alignment, LBP-TOP, and an SVM with an RBF kernel --- they interpolate sequences to 10, 20, 30 and so on up to 110 frames by two algorithms, \textbf{Newton interpolation} and \textbf{TIM}, recording the best recognition rate at each length on MMEW.

Newton peaks at \textbf{33.3\,\%} at 60 frames; TIM peaks at \textbf{38.9\,\%} ``when the micro-expression sequence is respectively interpolated to 30 or 60 frames'', and Ben et al. adopt TIM at 60 accordingly. Accuracy is non-monotonic in length: rates ``increase initially and then decrease towards the end''.

Two caveats limit the transfer. The experiment uses MMEW with LBP-TOP, not CASME II and not a learned model; and it runs largely in the \textit{upsampling} regime, since Ben et al.~\cite{ben2021} interpolate ``to the maximal value of 110 frames'', so the longer settings add frames for most clips rather than discarding them. This thesis operates in the opposite regime (\autoref{sec:tim-implications}), and an optimum measured while synthesising frames need not hold when discarding them.

A third comparison is the only multi-corpus one. Xu et al.~\cite{xu2017} interpolate every sequence to 10, 15 and 20 frames for both their facial dynamics map and LBP-TOP, ``on all four datasets'' --- CASME, CASME II and SMIC's high-speed and visual subsets. They also depart from the convention in a way the others do not, using \textbf{linear} interpolation ``instead of Temporal Interpolation Model because in a high frame-rate scenario, linear interpolation suffices to characterize the motion pattern with less error''.

\subsection{Alternatives to uniform interpolation}
\label{sec:tim-keyframe}

Zhao et al.~\cite{zhaoS2021} do not interpolate uniformly. They build an eleven-frame key-frame sequence from the annotated onset, apex and offset by generating eight transition frames allocated between the onset--apex and apex--offset intervals \textbf{in proportion to each sub-interval's duration}, so sampling density follows the expression's own phase structure and the apex is emphasised. Optical flow between adjacent key-frames then gives a ten-frame flow sequence. OFF-ApexNet~\cite{liong2019a} and STSTNet~\cite{liong2019b} take the limit, keeping the onset--apex pair alone (\autoref{sec:flow-what-between}).

The corpus therefore contains temporal representations of length 2, 10, 11, 15, 20, 30, 60, 110 and 120 frames, chosen variously by convention, by sweep and by architecture. There is no consensus, and only Xu et al.'s~\cite{xu2017} sweep spans more than one corpus.

\subsection{Limitations}
\label{sec:tim-limitations}

\textbf{Ten frames is a convention, not a finding.} The most-used setting traces to one baseline paper's choice, and the three studies that test it do not agree: ten frames for Li et al.~\cite{liX2018}, 30 to 60 for Ben et al.~\cite{ben2021}, and a 10--20 frame range for Xu et al.~\cite{xu2017} (\autoref{sec:tim-how-many}). Interpolation length interacts with the temporal receptive field of whatever consumes it, and in all three that consumer is a hand-crafted descriptor.

\textbf{Frame synthesis and motion analysis are in tension.} TIM generates frames that were not recorded. Where the downstream representation is appearance-based this is benign; where it is \textit{motion}, synthesised frames yield synthesised displacements. One reviewed paper recognises this: Xu et al.~\cite{xu2017} choose linear interpolation over TIM precisely on the grounds that it ``suffices to characterize the motion pattern with less error'' at a high frame rate. But they assert the point rather than measure it --- TIM appears once in their paper, in that sentence --- and no reviewed study compares the two interpolators against a motion representation. The rest of the corpus does not raise the question, because those using TIM extract appearance descriptors, while those extracting flow either use two frames or, as Zhao et al.~\cite{zhaoS2021} do, generate transition frames without discussing the consequence.

\textbf{Downsampling forfeits part of the corpus's advantage.} \autoref{sec:corpus} establishes why CASME II was recorded at 200~fps rather than its predecessor's 60~fps; reducing a clip to a few tens of frames gives part of that back. Only Xu et al.~\cite{xu2017} tie the interpolation choice to the acquisition rate at all, and then only to justify the simpler interpolator. The rest apply TIM10 to 100~fps SMIC data as readily as to 200~fps CASME II data, without weighing what the higher rate bought.

\subsection{Implications for this research}
\label{sec:tim-implications}

\begin{table}[htbp]
	\centering
	\footnotesize
	\begin{tabular}{@{}p{5.6cm} p{8.0cm}@{}}
		\hline
		\textbf{Finding from the review} & \textbf{Decision taken in this thesis} \\
		\hline
		Fixed-length input is required, and unequal duration is nuisance variance~\cite{lu2015} & Every clip is normalised to a fixed temporal length before training \\
		TIM beats Newton interpolation by 5.6 points, with a joint optimum at 30 or 60 frames~\cite{ben2021} & \textbf{L = 33 frames sampled, giving T = 32 flow pairs} --- the lower of the two reported optima rather than the field-conventional 10 \\
		Sequence length has an interior optimum found only by sweep~\cite{ben2021} & T = 32 is fixed throughout and is \textbf{not} swept; a sweep is declared as further work \\
		Onset and offset are annotated in CASME II (\autoref{sec:casme2-construction}) & Sampling is bounded by the annotated onset and offset, so the window spans the expression rather than surrounding neutral frames. This holds on the image-folder path used for CASME II; the loader's \texttt{.avi} branch samples across a whole file and is not exercised here \\
		TIM synthesises frames, and flow between synthesised frames is synthesised motion (\autoref{sec:tim-limitations}) & \textbf{No frames are synthesised} --- see the declaration below \\
		\hline
	\end{tabular}
	\caption[Review findings and the temporal-normalisation decisions they determine]{Review findings and the temporal-normalisation decisions they determine.}
	\label{tab:tim-decisions}
\end{table}

\textbf{A declaration about what this pipeline implements.} The module performing this step is named for the Temporal Interpolation Model, but what it computes is \textbf{uniform index sampling} of existing frames (\autoref{sec:three-channel-tensor}). No embedding is built, no curve is interpolated, no frame is synthesised. That is a materially different operation from the TIM of \autoref{sec:tim-what}, and it is named accurately here rather than in the pipeline's own terms. It is also defensible: because the downstream representation is flow between adjacent sampled frames, interpolation would manufacture displacement fields for motion that never occurred, whereas subsampling discards temporal resolution without fabricating it. What is given up is TIM's ability to \textit{lengthen} short clips.

Three quantitative consequences follow.

\textit{Effective temporal resolution falls by a clip-dependent factor.} Sampling 33 frames from a clip of $N$ frames recorded at 200~fps gives an effective rate of $6600/N$ fps: \textbf{100 fps} at the median ($N = 66$), which is SMIC's recording rate and half what CASME II was built to provide, and roughly \textbf{52 fps} for the longest ($N = 126$), below the original CASME database's 60~fps~\cite{yan2014}. Eighty-one of the 156 clips are downsampled by a factor of two or more.

\textit{Magnification is applied after subsampling, which shifts the realised pass-band.} Magnification runs on the 33 sampled frames at the nominal 200~fps (\autoref{sec:evm-departures}), but those frames span the clip's full duration, so the true inter-frame interval is $N/(200 \times 32)$ seconds rather than $1/200$. At the median a band specified as 5--25 Hz therefore corresponds to roughly 2.5--12.5 Hz physically, and the discrepancy grows with clip length. The realised band is thus neither the 5--25 Hz \autoref{sec:evm-implications} derives from the duration rule of Bai et al.~\cite{bai2021} nor constant across clips. This is a declared divergence; magnifying before subsampling is declared as further work.

\textit{Three clips are shorter than the sampling window.} \textbf{Three of the 156 clips} contain 31 frames against a target of 33, so sampling repeats indices and yields near-zero flow at those positions.

\subsection{Research gap and scope}
\label{sec:tim-gap}

The unresolved question is what temporal normalisation costs a representation made of motion rather than appearance. Three studies sweep interpolation length and give three answers, none on CASME II and none with a learned temporal model; and although Xu et al.~\cite{xu2017} raise both the synthesis-versus-motion tension and the acquisition-rate trade-off, they settle neither by measurement (\autoref{sec:tim-limitations}). This thesis takes the position that follows from its own representation --- sampling only recorded frames --- and quantifies what that costs in effective frame rate (\autoref{sec:tim-implications}), which no reviewed paper reports.

The design does not sweep T, nor compare interpolators, nor separate the effect of sequence length from the effect of the sampling regime. Those remain limitations rather than gaps claimed as solved, and the resulting conclusions apply to the implemented sampling scheme at T = 32.


\section{The Learned Spatial Backbone: Shallow 3D Convolutional Networks}
\label{sec:cnn-review}

This section covers the convolutional backbone, switched on and off as Variable C of this thesis's ablation, through four review-corpus papers.

\subsection{The case for, and against, a learned spatial stem}
\label{sec:cnn-case}

By 2019 the field had converged on optical flow as the input representation (\autoref{sec:why-flow}); the open question was what should consume it. A learned backbone promises what hand-crafted descriptors cannot give~\cite{liong2019b}, but runs into the small-data constraint (\autoref{sec:capacity}, \autoref{sec:corpus-threats}). Xia et al.~\cite{xia2020b} describe a related failure mode from a different cause --- ``the important subtle dynamics are prone to disappearing in the domain shift \ldots especially for deep models'' --- where the mechanism is the domain shift induced by mixing corpora rather than sample count alone, and where they report deeper, higher-resolution models doing better on the individual-database task. Applying their shrinking argument to a single small corpus therefore extends it beyond the setting they establish.

\subsection{What ``3D CNN'' means in this literature}
\label{sec:cnn-meaning}

The term covers three materially different architectures.

\textbf{A convolution over spatial and modality axes, with no temporal axis.} STSTNet~\cite{liong2019b}, the canonical shallow model in this field, calls itself a ``Shallow Triple Stream Three-dimensional CNN''. Its input is a single \textbf{28 $\times$ 28 $\times$ 3} tensor from the onset and apex frames, whose third dimension holds horizontal flow, vertical flow and optical strain --- not time. Three parallel streams, each a single convolutional layer with 3, 5 and 8 kernels respectively, supplement ``the small scale input data \ldots to avoid the problem of underfitting''. \textbf{There is no time dimension in this network}; the clip was collapsed to a single frame pair upstream (\autoref{sec:flow-what-between}).

\textbf{A convolution genuinely extended across time.} STRCN~\cite{xia2020a} extends convolutional connectivity ``across temporal domain'' with recurrent convolutional layers, modelling appearance and geometry separately.

\textbf{A convolution over a short flow sequence.} Zhao et al.~\cite{zhaoS2021} apply a C3D-style network to the ten-frame flow sequence built from their eleven key-frames (\autoref{sec:tim-keyframe}), so the temporal axis is present but short.

\subsection{The central finding: shrink the model \textit{and} the input}
\label{sec:cnn-shrink}

Xia et al.~\cite{xia2020b} treat depth and input resolution as controlled variables. They build four models of increasing depth --- from a single convolutional layer (Model 1) to one convolutional plus three recurrent convolutional layers (Model 4) --- and evaluate each at nine input resolutions, \textbf{20, 40, 60, 80, 100, 150, 200, 250 and 300 pixels}, on optical flow with UAR as the metric. Two results follow.

\textbf{Deeper is worse.} Models 3 and 4 ``have worse recognition performances compared to shallower \ldots models'', even though all four are far shallower than ResNet-101 or DenseNet-169.

\textbf{Higher resolution is worse, with a threshold.} Performance ``will be degraded with larger input resolutions (Model 2, Model 3 and Model 4)'' --- the two-layer model they go on to adopt included --- and for the deeper models ``begins to decrease dramatically when the resolution is bigger than \textbf{100 $\times$ 100}''. Shallow models are nonetheless ``basically robust to the change of input resolution'', improving toward lower resolution above a floor of about 40 $\times$ 40.

STSTNet takes the same argument to its limit and remains competitive on the CASME II subset (\autoref{sec:megc-metrics}): Liong et al.~\cite{liong2019b} report \textbf{0.00167 M parameters across 2 layers at 28 $\times$ 28 $\times$ 3 input}, against OFF-ApexNet's 2.77 M across 5 layers and GoogLeNet's 7 M across 22.

\subsection{Coping with limited and imbalanced training data}
\label{sec:cnn-limited-data}

Two strategies in the corpus have counterparts in this thesis's training configuration. Xia et al.~\cite{xia2020a} pair temporal augmentation with a balanced loss (\autoref{sec:imbalance-balanced-loss}, \autoref{sec:imbalance-double-correction}); a two-stage few-shot scheme~\cite{zhaoS2021} is the one route not taken. Xia et al.~\cite{xia2020b} take a third --- adding capability without capacity --- by developing ``three parameter-free modules \ldots integrate with RCN without increasing any learnable parameters'', so representational power grows while the parameter count, and therefore the overfitting risk, does not.

\subsection{Limitations}
\label{sec:cnn-limitations}

\textbf{The terminology obscures what is being compared.} Papers titled as three-dimensional CNNs include networks with no temporal axis at all, so ``3D-CNN'' results across the corpus are not like with like.

\textbf{Input resolution is almost never controlled.} Xia et al.~\cite{xia2020b} are the exception; elsewhere resolution is an implementation detail --- 28 $\times$ 28 in STSTNet and OFF-ApexNet --- with no indication of sensitivity, despite the substantial sensitivity Xia et al. measure.

\textbf{The evidence base is composite-database.} Xia et al.'s shrinking study~\cite{xia2020b} and STSTNet~\cite{liong2019b} are developed for the MEGC composite setting, which supplies substantially more training data per fold than a single-corpus study (\autoref{sec:megc-metrics}); only the earlier recurrent-convolutional work~\cite{xia2020a} is evaluated per database. Conclusions about affordable capacity therefore understate the constraint here.

\textbf{No study isolates the backbone itself.} The corpus is not short of matched comparisons: it sets descriptor-only methods against learned ones on the same folds, and Xia et al.~\cite{xia2020b} vary individual architectural modules against a common baseline with every other parameter held constant. But those comparisons all \textit{add} to a backbone that stays in place. None removes the learned spatial stem entirely and passes the motion representation straight to a classifier or temporal model, so what the stem itself contributes, set against what it costs, is nowhere measured.

\subsection{Implications for this research}
\label{sec:cnn-implications}

The component ablated as Variable C is a \textbf{three-stream shallow convolutional backbone}: one unshared two-layer stream per input channel --- flow-$u$, flow-$v$, optical strain.

\begin{table}[htbp]
	\centering
	\footnotesize
	\begin{tabular}{@{}p{5.6cm} p{8.0cm}@{}}
		\hline
		\textbf{Finding from the review} & \textbf{Decision taken in this thesis} \\
		\hline
		Shallow beats deep on small micro-expression data~\cite{xia2020b}; STSTNet uses 2 layers and $\sim$1,670 parameters~\cite{liong2019b} & Two convolutional layers per stream, \textbf{14,544 parameters} in the backbone (15,027 including the classifier head) --- shallow by the corpus's standards, though roughly nine times STSTNet's \\
		Triple-stream design supplements small-scale input~\cite{liong2019b} & Three unshared streams, one per input modality. STSTNet's streams differ in \textit{filter count} over the same input; these differ in \textit{input channel} \\
		Parameter-free modules add capability without capacity~\cite{xia2020b} & Motivates the parameter-free attention evaluated as Variable B (\autoref{sec:simam-review}) \\
		Limited and imbalanced data require augmentation plus loss correction~\cite{xia2020a} & Balanced sampling with focal loss; loss class-weighting deliberately disabled after single-class collapse (\autoref{sec:corpus-implications}) \\
		\textbf{Performance degrades above 100 $\times$ 100 input, dramatically for deeper models~\cite{xia2020b}} & \textbf{Not followed.} The backbone is fed at \textbf{224 $\times$ 224}, more than twice that threshold, though the same resolution as Bai et al.~\cite{bai2021} and Li, Huang and Zhao~\cite{liY2021} use --- see below \\
		\hline
	\end{tabular}
	\caption[Review findings and the backbone decisions they determine]{Review findings and the backbone decisions they determine.}
	\label{tab:cnn-decisions}
\end{table}

\textbf{A declaration about what this component is.} Every kernel in this backbone is spatial only: no filter spans two time steps, so despite the \texttt{Conv3d} implementation it is a \textbf{purely spatial} feature extractor (\autoref{sec:conv3d}). The ablation therefore does not test spatio-temporal convolution; it tests a \textbf{learned spatial stem}, and whatever it measures for Variable C is evidence about \textit{this} stem --- spatial-only, at 224 $\times$ 224, on flow and strain input --- and \textbf{not} evidence that 3D convolution over time is unhelpful for micro-expression recognition, which this study never evaluated.

Two observations mitigate that, and one aggravates it. In the field's own terms the choice is defensible: STSTNet, one of the two strongest shallow baselines on the CASME II subset, likewise performs no temporal convolution (\autoref{sec:cnn-meaning}). The stem also competes against a representation that already performs spatio-temporal extraction analytically (\autoref{sec:flow-implications}, \autoref{sec:strain-implications}).

Against both, \textbf{input resolution is the clearest departure from Xia et al.'s guidance}. It is not a departure from the corpus as a whole: Bai et al.~\cite{bai2021} resize frames to 224 $\times$ 224, and without alignment or cropping, as here; Li, Huang and Zhao~\cite{liY2021} resize faces to the same; and Xia et al.'s own earlier work~\cite{xia2020a} runs one stream at 300 $\times$ 245. What 224 $\times$ 224 departs from is the specific threshold of \autoref{sec:cnn-shrink}, which is a composite-database result and which bites hardest on models deeper than this one. Since this stem is shallow, the reviewed evidence predicts a smaller resolution effect than the deep-model figures suggest, so re-running the backbone arm at a lower resolution is worth doing because it is cheap and would settle the question, not because the corpus implies a large gain. The choice is also the expensive one: a 224 $\times$ 224 map carries sixty-four times the spatial elements of STSTNet's 28 $\times$ 28, and the backbone convolves it at every one of the 32 time steps. Whatever this study measures for the stem is therefore a resolution result as much as an architecture result, and the re-run is the first item of declared further work.

\subsection{Limitations of the existing work, and how this study differs}
\label{sec:cnn-gap}

The four gaps of \autoref{sec:cnn-limitations} define what the reviewed literature does not establish.

This thesis differs on the third, which is what its research question requires: it removes the convolutional backbone entirely in four of its twelve configurations, and measures the difference under complete leave-one-subject-out with every other factor held constant. On the first gap it answers the field's terminological ambiguity by declaring precisely what its own component is (\autoref{sec:cnn-implications}); on the second it knowingly departs from the literature's guidance, at the 224 $\times$ 224 resolution flagged above.


\section{Parameter-Free Attention: SimAM}
\label{sec:simam-review}

SimAM re-weights the convolutional feature maps without adding learned parameters. Yang et al.~\cite{yang2021} provide the mechanism and its image-recognition evidence; Xia et al.~\cite{xia2020b} motivate parameter-free modules in composite-database micro-expression recognition. The latter support the design principle, rather than validating SimAM itself on this task.

\subsection{Why parameter-free attention}
\label{sec:simam-why}

Local facial activity occupies a small part of a motion tensor. Attention offers a way to emphasise informative activations, but its parameter and computational costs must be justified on a 156-clip corpus. Yang et al.'s~\cite{yang2021} comparison illustrates the range of parameter costs.

\begin{table}[htbp]
    \centering
    \small
    \begin{tabular}{@{}lc@{}}
        \hline
        \textbf{Module} & \textbf{Additional learnable parameters} \\
        \hline
        SE & $2C^2/r$ \\
        CBAM & $2C^2/r+2k^2$ \\
        GC & $2C^2/r+C$ \\
        ECA & $k$ \\
        SRM & $6C$ \\
        SimAM & $0$ \\
        \hline
    \end{tabular}
    \caption[Parameter costs of attention modules]{Parameter costs reported by Yang et al.~\cite{yang2021}, Table 1. $C$ denotes channels, $r$ the reduction ratio and $k$ the relevant kernel size. The competing modules are attributed through this source.}
    \label{tab:attention-modules}
\end{table}

Zero additional parameters makes SimAM a plausible candidate, not the only defensible one: alternatives such as ECA are also small. Xia et al.~\cite{xia2020b} demonstrate a related strategy of adding parameter-free modules to a recurrent convolutional model. Neither argument establishes that a parameter-free module has zero runtime cost or cannot worsen generalisation.

\subsection{The distinctiveness premise}
\label{sec:simam-premise}

Yang et al.~\cite{yang2021} motivate attention through spatial suppression: neurons distinguishable from their surroundings are assigned greater importance. They express this as the linear separability of one target activation from the other activations in its channel. The resulting weights scale the feature map rather than being added to it.

This is a prior about activation statistics, not a detector of facial action units or minority classes. A distinctive response may encode useful deformation, noise or an artefact. Its suitability for flow-derived features therefore requires evaluation rather than following automatically from the neuroscience motivation.

\subsection{Energy and implemented refinement}
\label{sec:simam-energy}

The paper minimises a regularised separation energy and obtains a closed-form importance score. Its reference implementation can be written compactly as
\[
    d_i=(x_i-\mu)^2,\qquad
    s^2=\frac{1}{M-1}\sum_{i=1}^{M}d_i,\qquad
    E_i^{-1}=\frac{d_i}{4(s^2+\lambda)}+\frac{1}{2},
\]
\[
    \widetilde{x}_i=x_i\,\operatorname{sigmoid}(E_i^{-1}),
    \qquad \mu=\frac{1}{M}\sum_{i=1}^{M}x_i.
\]
Here $M$ is the number of positions in one channel's statistical context and $\lambda$ is a positive regularisation constant. The $+\tfrac12$ follows from the reciprocal-energy expression rather than an arbitrary offset.

An implementation distinction matters: the shared variance printed with Yang et al.'s equation (5) divides by $M$, whereas their Figure 3 code divides by $M-1$~\cite{yang2021}. The equations above follow the reference \emph{code}, as does this thesis. The sigmoid preserves the ordering of the scores while bounding the multiplicative weights; it does not create a hard mask of inactive regions.

\subsection{What three-dimensional weights mean}
\label{sec:simam-3d-weights}

The original module produces a separate weight at each channel--height--width position~\cite{yang2021}. Its three-dimensional weights therefore describe $C\times H\times W$, not a temporal volume. Statistics are computed separately within each channel's spatial plane, with $M=HW$. Extending the context across video frames is an additional design decision, not part of the original image formulation.

\subsection{Evidence and hyper-parameter selection}
\label{sec:simam-evidence}

Yang et al.~\cite{yang2021} evaluate SimAM with several network families on CIFAR and ImageNet, reporting competitive classification performance against learned attention modules. Their timing results also show that unchanged parameter totals do not imply unchanged inference speed. The evidence supports testing SimAM, but does not establish its value on micro-expression motion tensors.

For CIFAR, they sweep $\lambda$ from $10^{-1}$ to $10^{-6}$ and select $10^{-4}$ as a balance between accuracy and variability. They repeat selection for ImageNet and choose $0.1$~\cite{yang2021}. Thus $10^{-4}$ is a dataset-selected operating point, not a universal setting justified by the formula.

\subsection{Limits of transfer to this task}
\label{sec:simam-limitations}

The source classification experiments concern natural images, different training-set sizes and learned feature distributions. In this thesis the input is flow and strain, so distinctiveness need not have the same relationship to useful evidence. The statistics are calculated from each feature map, not estimated by pooling the entire training corpus; the small dataset instead constrains how the surrounding network is learned.

The image experiments also do not decide whether video statistics should be computed per frame or over a clip. Nor does the lack of learned attention weights establish that runtime, memory or overfitting concerns disappear. These are empirical limits on transferring the source result, rather than reasons to assume either success or failure.

\subsection{Implications for this research}
\label{sec:simam-implications}

\begin{table}[htbp]
    \centering
    \small
    \begin{tabular}{@{}p{0.29\linewidth} p{0.65\linewidth}@{}}
        \hline
        \textbf{Choice} & \textbf{Implementation and interpretation} \\
        \hline
        Placement & Applied separately to each convolutional stream before concatenation; no trainable weights or biases are introduced. \\
        Refinement & Uses the reference-code reciprocal energy and sigmoid gate~\cite{yang2021}. \\
        Statistical context & Mean and variance span $(T,H,W)$ within each channel, with divisor $THW-1$, rather than one $(H,W)$ plane. \\
        Regularisation & $\lambda=10^{-4}$ is retained without re-tuning on CASME II. \\
        Architectural dependency & SimAM is enabled only with the CNN stem in this model. Four of sixteen candidate configurations are excluded, leaving twelve evaluated configurations and four SimAM matched pairs. \\
        \hline
    \end{tabular}
    \caption[SimAM implementation choices]{SimAM implementation choices and their relation to the reviewed image formulation. The dependency on the CNN is a rule of this architecture.}
    \label{tab:simam-decisions}
\end{table}

The five-dimensional input includes batch, channel, time, height and width. Pooling statistics across time makes an activation distinctive relative to its whole clip rather than its own frame. This could emphasise a brief contraction, but does not guarantee low weights at onset or high weights at the apex: the score depends on squared deviation, not expression labels. The two context choices are not compared experimentally.

The CNN dependency similarly describes this implementation, not a mathematical prohibition on applying SimAM to other tensors. Its measured effect is conditional on the stem being present. Retention should be judged against both recognition performance and computational cost; parameter count alone cannot settle that decision.

\subsection{Research gap and scope}
\label{sec:simam-gap}

The study evaluates SimAM in a motion-input, small-corpus setting through four matched comparisons under complete leave-one-subject-out evaluation. It tests the implemented module as a whole, rather than separately establishing the benefit of each design choice.

Three follow-ups remain: testing whether the distinctiveness prior suits flow-derived features, re-tuning $\lambda$, and comparing per-frame with whole-clip statistics. Each requires new evaluations despite introducing no learned attention weights. Per-class changes may also be examined, but a rare class benefiting most is a hypothesis of this study, not a prediction established by Yang et al.~\cite{yang2021}.


\section{The Temporal Encoder: Transformers for Micro-Expression Recognition}
\label{sec:transformer-review}

The temporal encoder is motivated by Zhang et al.'s Short and Long Range Relation Based Spatio-Temporal Transformer (SLSTT)~\cite{zhang2022} and its Vision Transformer (ViT) foundation~\cite{dosovitskiy2021}. The component evaluated here differs substantially from published SLSTT, particularly in the axis over which self-attention operates.

\subsection{Why temporal self-attention}
\label{sec:transformer-why}

A micro-expression evolves through onset, peak and relaxation. Retaining a motion sequence permits a model to relate separated stages rather than compressing them immediately into one descriptor. Zhang et al.~\cite{zhang2022} motivate their design through the need to capture both local and global patterns.

This thesis compares a temporal transformer with mean pooling over the same input sequence. Mean pooling discards order, whereas positional information and self-attention allow the encoder to model relations among frame-pair features. This motivates the comparison without establishing in advance which representation will generalise better.

\subsection{The Vision Transformer foundation}
\label{sec:transformer-vit}

ViT splits an image into patches, projects them into token vectors, adds learned positional embeddings and processes them with transformer encoder layers~\cite{dosovitskiy2021}. Attention then relates spatial patches. Dosovitskiy et al. also show that data scale and training regime matter: without the locality and translation-equivariance biases of convolution, image transformers can perform poorly when trained with insufficient data, while large-scale pre-training improves their performance.

That evidence is a caution for a 156-clip study, not a universal requirement that every transformer must be pre-trained. A short sequence of motion features presents a different learning problem from image patches, but remains subject to limited data and overfitting.

\subsection{Published SLSTT}
\label{sec:slstt-architecture}

Zhang et al.~\cite{zhang2022} combine three stages. First, \textbf{onset-referenced optical flow} describes sampled frames relative to the onset rather than to their immediate predecessors. They argue that this gives a more structured trajectory towards and away from the apex. Their experiments use this representation throughout; they do not provide a matched short-term-flow ablation establishing its independent gain.

Second, a \textbf{spatial ViT encoder} processes each flow field as patches. The implementation uses ImageNet-pre-trained ViT-B/16, twelve encoder layers and $384\times384$ input. Third, an \textbf{LSTM temporal aggregator} combines the per-frame features. A mean aggregator supplies the comparison variant. Self-attention is spatial in this architecture; recurrence handles time.

\subsection{Published results and the aggregation comparison}
\label{sec:slstt-results}

\begin{table}[htbp]
    \centering
    \small
    \begin{tabular}{@{}lccc@{}}
        \hline
        \textbf{Method} & \textbf{Composite UF1} & \textbf{CASME II UF1} & \textbf{CASME II UAR} \\
        \hline
        LBP-TOP & 0.588 & 0.703 & 0.743 \\
        Bi-WOOF & 0.630 & 0.781 & 0.803 \\
        OFF-ApexNet~\cite{liong2019a} & 0.720 & 0.876 & 0.868 \\
        STSTNet~\cite{liong2019b} & 0.735 & 0.838 & 0.869 \\
        EMR~\cite{liu2019} & 0.789 & 0.829 & 0.821 \\
        RCN~\cite{xia2020b} & 0.705 & 0.809 & 0.856 \\
        SLSTT-Mean & 0.788 & 0.844 & 0.830 \\
        SLSTT-LSTM & 0.816 & 0.901 & 0.885 \\
        \hline
    \end{tabular}
    \caption[SLSTT and comparison results]{Selected results from Zhang et al.~\cite{zhang2022}, Table 2, reported under their three-class Composite Database Evaluation heading. Columns retain the source's dataset labels. The listed methods are not all arms of one controlled architectural comparison.}
    \label{tab:slstt-results}
\end{table}

The SLSTT variants compare temporal aggregation with the spatial transformer retained in both. CASME II UF1 increases from 0.844 to 0.901, a difference of 0.057, or 5.7 UF1 percentage points~\cite{zhang2022}. This supports evaluating aggregation choices within that framework. It is not an estimate of what an LSTM would add to this thesis, where mean pooling follows \emph{temporal} self-attention rather than independent spatial encoding.

Protocol labels need care. The source describes three-class LOSO evaluations on the databases separately and together under its Composite Database Evaluation heading. That heading alone therefore does not establish that every CASME II column entry uses pooled cross-database training. It also reports a distinct Sole Database Evaluation using original labels, including five CASME II classes~\cite{zhang2022}. Neither setting should be equated automatically with this thesis's 156-clip three-class pool.

\subsection{Limitations of the evidence}
\label{sec:transformer-limitations}

Published SLSTT retains a pre-trained spatial encoder. Its results do not isolate the benefit of a transformer trained from scratch on this study's motion representation. The Mean--LSTM comparison isolates an aggregator, not the presence of a transformer, since both variants contain one.

Label interpretation also limits comparison. Zhang et al.~\cite{zhang2022} note that facial action units do not map uniquely to emotion labels, so visual evidence alone need not determine the annotation. This is a limitation of the recognition task, not a numerical upper bound measured by the present experiments.

\subsection{Implementation and six declared divergences}
\label{sec:transformer-implications}

The evaluated component receives 32 tokens of width 96. They come either from the convolutional stem or from $4\times4$ spatial pooling followed by a \emph{learned linear projection}. The encoder uses four pre-norm layers, eight heads, feed-forward width 256 and dropout 0.1. Fixed sinusoidal positions precede attention; mean pooling follows it. With the encoder disabled, the input tokens are averaged directly.

\begin{table}[htbp]
    \centering
    \small
    \begin{tabular}{@{}p{0.23\linewidth} p{0.34\linewidth} p{0.35\linewidth}@{}}
        \hline
        \textbf{Property} & \textbf{Published SLSTT~\cite{zhang2022}} & \textbf{This thesis} \\
        \hline
        Attention axis & Spatial patches within a frame & Temporal tokens within a clip \\
        Aggregation & LSTM over spatially encoded frames & Mean after temporal self-attention \\
        Initialisation & ImageNet-pre-trained ViT-B/16 & Trained from scratch \\
        Capacity & Twelve layers, width 768; $384\times384$ inputs & Four layers, width 96; 32 tokens \\
        Positions & Learned spatial embeddings & Fixed sinusoidal temporal encoding \\
        Input flow & Relative to onset & Relative to the preceding sampled frame \\
        \hline
    \end{tabular}
    \caption[Divergences from published SLSTT]{Six material divergences from published SLSTT. The component evaluated here is a small temporal transformer encoder, not a reproduction of the published architecture.}
    \label{tab:slstt-divergences}
\end{table}

Six matched pairs evaluate the temporal encoder with the other component choices fixed within each pair. They test the whole encoder, including positional encoding, feed-forward layers and normalisation, against direct mean pooling. They do not isolate self-attention from those accompanying operations or compare it with another learned temporal model.

Three follow-ups follow from the divergences: testing an LSTM aggregator, comparing onset-referenced with consecutive-frame flow, and replacing fixed positions with learned embeddings. The last would add $32\times96=3{,}072$ embedding parameters for this sequence length. All three remain untested; changing the input representation also requires re-training and evaluation before any recognition benefit can be claimed.

\subsection{Research gap and scope}
\label{sec:transformer-gap}

The study asks what a small temporal encoder contributes after motion extraction, using a single corpus without pre-training. This differs from SLSTT's spatial transformer and recurrent aggregation while retaining its motivation to model relations beyond isolated frames.

The contribution is a matched comparison of the implemented encoder and an order-blind baseline within this pipeline. It neither validates published SLSTT nor establishes a general ranking of attention, recurrence and convolution. The six declared divergences and the limited training regime define the scope of the architectural conclusions.


\section{Training Under Severe Class Imbalance}
\label{sec:imbalance-review}

This section covers the training configuration held constant across all twelve ablated configurations, and which \autoref{sec:corpus-implications} identified as forced by the corpus rather than chosen freely: the handling of a 4 : 1.3 : 1 class imbalance.

\textbf{A deliberate exclusion.} Two topics the project's reading list groups here --- Grad-CAM++ visual auditing and adversarial identity disentanglement --- are \textbf{not reviewed}, because neither is part of the implemented system. The loss module records that ``the Cross-Batch Memory (XBM) queue, and the adversarial identity loss --- are intentionally removed'', and no gradient-based visualisation exists anywhere in the codebase. Both are declared as proposed further work.

\subsection{The imbalance is a property of the phenomenon, not of the sampling}
\label{sec:imbalance-property}

\autoref{sec:label-taxonomy} and \autoref{sec:original-baseline} established the class counts and the always-Negative floor; what matters here is why the skew exists. It is not an artefact of careless collection. Yan et al.~\cite{yan2014} are explicit that it follows from what can be elicited in a laboratory: ``some types of facial expressions are difficult to elicit in laboratory situations, thus the samples in different categories distributed unequally, e.g., there are 60 disgust samples but only 7 sadness samples.'' Nor is it specific to CASME II --- the MEGC composite shows the same pattern (\autoref{sec:megc-metrics}). Skew must be coped with; it cannot be collected away.

\subsection{Three places to intervene}
\label{sec:imbalance-three-places}

The reviewed corpus corrects for imbalance at three distinct points, worth separating because this thesis intervenes at all three and the interaction between two of them turned out to matter.

\begin{enumerate}
	\item \textbf{At the data}, by changing how often each class is presented during training --- resampling or augmentation.
	\item \textbf{At the loss}, by changing how much each example contributes to the gradient --- class weighting or focal re-weighting.
	\item \textbf{At the metric}, by changing how performance is scored --- macro-averaged rather than sample-averaged measures.
\end{enumerate}

\subsection{Loss-level correction: focal loss}
\label{sec:imbalance-focal}

Zhao et al.~\cite{zhaoS2021} bring focal loss into micro-expression recognition ``to alleviate the inefficient model training caused by the class imbalance in the MEs Datasets'', since cross-entropy ``could be biased toward particular emotions that constitute a larger portion of the training set''. They adopt the formulation developed for one-stage object detection, where the imbalance is between foreground and background regions:

\[
	\text{FL}(p_t) = -\alpha_t (1 - p_t)^{\gamma}\log p_t
\]

where $p_t$ is the model's probability for the true class, $\gamma$ ``the balance factor for loss'' and $\alpha_t$ ``the weight balance factor for samples'' (\autoref{sec:focal-loss}); under skew, the examples it up-weights are disproportionately the minority classes. Adapting it to multi-class classification, Zhao et al. report that ``$\gamma$ is set as 2 in practice, and $\alpha$ is treated as a hyper-parameter to set by cross validation.''

The division of labour there matters: $\gamma$ re-weights by \textit{difficulty} and is treated as a constant, $\alpha$ by \textit{class frequency} and is treated as dataset-specific. Conflating them is where this project went wrong (\autoref{sec:imbalance-double-correction}).

\subsection{Data-level correction: resampling and augmentation}
\label{sec:imbalance-balanced-loss}

Xia et al.~\cite{xia2020a} address the same problem at the data end and pair the two interventions, using ``temporal data augmentation strategies as well as a balanced loss \ldots jointly'' to overcome ``limited and imbalanced training samples''. The pairing is not left unexamined: their subsection \textit{``The Impact of Data Augmentation and Balanced Loss''} removes each correction in turn from the joint system and reports the balanced loss improving performance ``slightly'', with one leave-one-subject-out cell on SMIC where removing it scores higher. This is the corpus's only controlled look at composing two corrections, and \autoref{sec:imbalance-double-correction} records a far larger effect in the same direction.

\subsection{Metric-level correction}
\label{sec:imbalance-metrics}

\autoref{sec:megc-metrics} defines UF1 and the case for macro-averaged metrics. Liong et al.~\cite{liong2016} add the argument from the subject side rather than the class side, reporting F1, precision and recall because these ``provide a more meaningful perspective than accuracy rates when the datasets used are naturally imbalanced since each subject has a different number of video samples.'' The two arguments compound on CASME II, where classes are skewed \textbf{and} per-subject clip counts range from 1 to 33 (\autoref{tab:corpus-threats}c).

\subsection{Limitations}
\label{sec:imbalance-limitations}

\textbf{The interaction between corrections is examined once, and only at small magnitude.} Xia et al.~\cite{xia2020a} ablate augmentation and balanced loss out of the joint system and report the balanced loss helping only ``slightly'' --- at most about two points across their twelve cells, and in one cell costing accuracy; Zhao et al.~\cite{zhaoS2021} apply focal loss alongside a large per-sample augmentation, but without class-proportional resampling and without ablating the pairing. What no reviewed paper reports is the pairing failing outright.

\textbf{$\gamma$ is inherited rather than tuned.} Zhao et al.~\cite{zhaoS2021} take the value from the object-detection literature ``in practice''. They do compare focal loss against cross-entropy across a sweep of $\alpha$, but that experiment holds $\gamma$ fixed, so the focusing exponent remains unvalidated on this task.

\textbf{Label smoothing is named in the corpus but never examined on this task.} Dosovitskiy et al.~\cite{dosovitskiy2021} list it among the regularisation parameters they tune, but no reviewed micro-expression paper uses or discusses it, or reports what it does under class skew. It is used here and declared as a choice this corpus supports neither way.

\subsection{Implications for this research}
\label{sec:imbalance-double-correction}

\begin{table}[htbp]
	\centering
	\footnotesize
	\begin{tabular}{@{}p{5.6cm} p{8.0cm}@{}}
		\hline
		\textbf{Finding from the review} & \textbf{Decision taken in this thesis} \\
		\hline
		Focal loss addresses class-frequency bias in micro-expression recognition~\cite{zhaoS2021} & Focal loss adopted as the training objective, with \textbf{$\gamma = 2.0$}, matching Zhao et al.'s reported setting \\
		$\alpha$ is a dataset-specific class-weight term~\cite{zhaoS2021} & Implemented as an optional per-class $\alpha$ vector, \textbf{switched off at run time} --- see the declaration below \\
		Correct at the data end as well as the loss end~\cite{xia2020a} & Minority oversampling via a weighted sampler, active in every configuration --- a sampler, where Xia et al. use uniform augmentation plus a balanced loss \\
		Imbalance-aware metrics are mandatory~\cite{see2019,liong2016} & Pooled macro F1 --- identical in construction to MEGC's UF1 --- as the primary metric, with accuracy quoted only alongside the always-Negative floor \\
		\textit{(no micro-expression source)} & Label smoothing of 0.05 applied inside the focal loss; named among the regularisers tuned by Dosovitskiy et al.~\cite{dosovitskiy2021}, but declared as unsupported either way by this corpus \\
		\hline
	\end{tabular}
	\caption[Review findings and the imbalance-handling decisions they determine]{Review findings and the imbalance-handling decisions they determine.}
	\label{tab:imbalance-decisions}
\end{table}

\textbf{The finding this project contributes: correct at exactly one point, and know which.} The configuration is resolved at run time by a single rule --- class weighting in the loss is enabled only when the balanced sampler is \textit{not} active --- and the run log records the decision each fold: \textit{``Class weights in loss disabled (balanced sampler already active).''}

That rule exists because of two failures in opposite directions, both recoverable from this project's own run history. In the earliest code era the loss carried inverse-frequency class weights and no balanced sampler existed; under that single correction the proposed configuration abandoned the \textbf{Positive} class outright, predicting none of its clips --- a failure initially misdiagnosed as over-capacity overfitting. A later era added the sampler while leaving the loss weights unconditional, so both corrections ran at once; the collapse then inverted, and configurations abandoned the \textbf{Negative} majority class instead, one of them predicting a single class for every clip in the corpus. Correcting once at the loss starved the minority; correcting twice over-served it. With the guard in place and exactly one correction active, no configuration in the reported study abandons any class.

This is a contribution rather than an incident, because of how little guidance the corpus offers (\autoref{sec:imbalance-limitations}): no reviewed paper warns that combining corrections can invert the intended effect rather than merely blunt it. On a 156-clip corpus with a 4 : 1 skew, correcting twice is as damaging as correcting once in the wrong place --- the two failures differ only in which class is abandoned.

\textbf{A second declaration.} The severity of the imbalance also constrains what the ablation can conclude. Because every configuration shares the same sampler, loss and $\gamma$, this thesis measures \textbf{no} effect attributable to the imbalance treatment itself: it is a constant, not a factor. The claim supported is that the four architectural components were compared under a training regime that does not abandon minority classes, not that this regime is optimal.

\subsection{Research gap and scope}
\label{sec:imbalance-gap}

The corrections themselves are well established; what the corpus does not settle are the three gaps of \autoref{sec:imbalance-limitations} --- how two corrections compose, what $\gamma$ should be on this task, and what over-correction looks like when it fails.

This thesis contributes to the first and third, documenting two opposite class-collapse failures and the run-time rule that resolved them (\autoref{sec:imbalance-double-correction}): a negative result the corpus does not contain and one directly actionable for anyone reproducing this class of pipeline. On the second it offers no improvement, since $\gamma = 2.0$ and label smoothing of 0.05 are held constant across all twelve configurations. Sweeping either, and the Grad-CAM++ auditing and identity-disentanglement work this project's earlier stages removed, are declared as further work.


\section{Synthesis and Research Gap}
\label{sec:synthesis-and-gap}

This section consolidates the gap statements closing \autoref{sec:corpus}--\autoref{sec:imbalance-review}: what the corpus collectively establishes, the four limitations recurring across every topic, the resulting research gap, and which parts of it this thesis closes. It also collects in one place every declared divergence between the system built here and the literature it draws on.

\subsection{What the corpus establishes}
\label{sec:corpus-establishes}

Read as a whole, the reviewed literature settles a good deal. \textbf{On data:} the corpora are small, demographically narrow, laboratory-bound and irreducibly imbalanced by what can be elicited, and CASME II remains the primary single-corpus benchmark under its native leave-one-subject-out protocol (\autoref{sec:corpus}, \autoref{sec:imbalance-review}). \textbf{On representation:} motion rather than appearance is the appropriate input; a strain field adds a further cue, though on CASME II only weakly and under a protocol that is not subject-disjoint; and both the magnification factor and the interpolated sequence length have interior optima located by sweep (\autoref{sec:evm-review}--\autoref{sec:tim-review}). \textbf{On architecture:} at this data size capability bought without parameters outperforms capability bought with them, and self-attention is the strongest-performing mechanism reported in the corpus (\autoref{sec:cnn-review}--\autoref{sec:transformer-review}). \textbf{On evaluation:} accuracy alone is inadequate under skew and must give way to macro-averaged measures accumulated across folds before averaging (\autoref{sec:megc-metrics}, \autoref{sec:imbalance-metrics}).

\subsection{Four limitations that recur across every topic}
\label{sec:four-limitations}

The per-section gap statements are not nine independent observations. Four structural problems generate most of them.

\textbf{(a) Components are isolated one at a time, inside one proposed design, and never jointly.} Matched-pair measurement is not absent from the corpus. On the input side \autoref{tab:evm-effect} isolates magnification and \autoref{tab:strain-performance} isolates strain; on the architectural side Xia et al. compare a basic recurrent convolutional network against three variants adding a wide expansion, a shortcut connection and an attention unit, holding every other parameter constant across all four (\autoref{sec:cnn-limited-data}), and separately sweep depth against input resolution (\autoref{sec:cnn-shrink}); SLSTT's Mean-versus-LSTM comparison isolates its aggregator (\autoref{sec:slstt-results}).

What no reviewed study does is vary \textit{several} components together, so that their interactions can be read. Every ablation in the corpus is internal to one proposed design and moves one factor at a time from that design's own defaults; across papers, systems are compared whole --- STSTNet against OFF-ApexNet, SLSTT against STSTNet --- so a difference in the reported number cannot be attributed to any one stage. There is therefore no basis for deciding which of the learned stages are load-bearing, nor whether one component's value depends on another's presence.

\textbf{(b) The evaluation protocol is unstable, and under-reported.} Sample counts, label-set size and the choice among leave-one-subject-out, leave-one-video-out and leave-one-sample-out all vary inconsistently, and the field's most-quoted baseline is reported under two different protocols by its two primary sources and under two more within the originating paper itself (\autoref{sec:original-baseline}). Given this corpus's fold composition (\autoref{sec:corpus-threats}), these are not bookkeeping differences: they change what the numbers mean.

\textbf{(c) The strongest results depend on regimes a single-corpus study cannot access.} The best CASME II figures in the corpus come from composite-database training (\autoref{sec:megc-metrics}) and, in SLSTT's case, ImageNet pre-training (\autoref{sec:transformer-limitations}) --- neither a property of the architecture being credited. Almost no evidence exists about what these components do when trained from scratch on one small corpus, which is the situation most practitioners are in.

\textbf{(d) Preprocessing parameters have interior optima, and are transferred without re-validation.} The magnification factor, the interpolated sequence length and the input resolution all rise then fall, all have their optima located only by sweep, and none has a principled selection rule (\autoref{sec:evm-evidence}, \autoref{sec:tim-how-many}, \autoref{sec:cnn-shrink}). Values are carried between pipelines whose downstream models differ entirely: $\alpha = 30$ moves from a single-frame appearance encoder to other settings, TIM10 propagates from one baseline paper to the whole field, and $\lambda = 10^{-4}$ transfers from CIFAR though its own authors re-tuned it for ImageNet (\autoref{sec:simam-evidence}).

\subsection{The research gap}
\label{sec:research-gap}

Taking (a)--(d) together:

\textbf{The field has established a standard pipeline for micro-expression recognition --- magnify, compute motion, normalise duration, extract spatially, model temporally --- but has never measured which of its stages are responsible for the result. Because components are only ever varied one at a time inside a single proposed system, under protocols that vary in unstated ways, and predominantly in training regimes richer than a single small corpus provides, it is not currently possible to say which parts of that pipeline earn their place, which are inherited convention, and which depend on the presence of another.}

This matters practically as well as scientifically. A practitioner assembling this pipeline today has no published basis for allocating effort: for judging whether the most computationally expensive stage is the most valuable, whether a component with a negligible average effect may still matter in a specific failure, or whether a ranking read off an incomplete evaluation of a corpus this size is stable at all.

\subsection{How this thesis addresses the gap --- and what it does not}
\label{sec:addressing-gap}

\textbf{What it does.} This thesis treats the pipeline as a factorial experiment rather than a proposal: four components varied independently across every architecturally valid cell of a 2\textsuperscript{4} matrix, twelve configurations, each evaluated under \textbf{complete 25-fold leave-one-subject-out on CASME II} with every non-varied factor held identical. That gives the joint, multi-component measurement (a) says the corpus lacks; reporting the protocol in full addresses (b) for this study at least; and training from scratch is the regime (c) identifies as unevidenced. Each component's marginal contribution is read from pairs of runs differing in one factor alone, and the distribution of those four effects --- how much of the pipeline's performance each stage accounts for, whether any subtracts, and what each costs --- is the substantive answer to the gap.

Three further contributions follow from the review rather than from the experiment:

\begin{itemize}
	\item \textbf{A caution about partial evaluations.} This project's earlier evaluations, on smaller held-out sets and incomplete fold counts, ranked the twelve configurations differently from the complete run reported here. They are not evidence of protocol sensitivity, since they differ in sample count, training budget and class coverage too, and the earliest carry a data-routing defect --- but the episode illustrates limitation (b) in practice, and only the complete run is reported.
	\item \textbf{A repaired measurement.} That defect left the magnification switch inert, producing bit-identical results across every pair those runs contained. What made it detectable as a defect rather than a null result was the review's evidence that magnification has a clear, repeatable effect (\autoref{sec:evm-implications}).
	\item \textbf{Two opposite class-collapse failures.} A single loss-level correction suppressed the minority class; adding a data-level correction on top of it suppressed the majority class instead. No reviewed paper reports either collapse (\autoref{sec:imbalance-double-correction}).
\end{itemize}

\textbf{What it does not do.} Three of the four limitations are only partially addressed. On (b), this study reports its own protocol completely but cannot repair the field's inconsistency, and its own N = 156 differs from MEGC's 145. On (c), it establishes what these components do from scratch on one corpus, which is a different question from --- not an answer to --- how they behave under composite training; no cross-dataset claim is made anywhere. On (d), the thesis inherits parameter values rather than sweeping them: $\alpha = 10$, T = 32, $\lambda = 10^{-4}$ and $\gamma = 2.0$ are fixed throughout, so the ablation measures architecture at one point in preprocessing space, not across it.

\subsection{A consolidated ledger of declared divergences}
\label{sec:divergence-ledger}

Reviewing the literature closely enough to write this chapter surfaced twelve points at which the implemented system departs from the work it draws on. Several were discovered only through the review. They are collected in \autoref{tab:divergence-ledger} so that no result rests on an unstated assumption.

\begin{table}[htbp]
	\centering
	\footnotesize
	\begin{tabular}{@{}p{0.85cm} p{6.0cm} p{6.6cm}@{}}
		\hline
		\textbf{\S} & \textbf{Divergence} & \textbf{Consequence} \\
		\hline
		3.1.3 & Working set is N = 156, including sadness and fear in Negative; MEGC's CASME II subset is N = 145 & Comparisons against challenge figures are not exactly like-for-like \\
		3.1.2 & The corpus's own registered release is not used; the pipeline reads the original recorded frames, unregistered and uncropped & No landmarking or alignment of any kind precedes the network; tolerance to pose and scale must be learned from 156 clips \\
		3.2.6 & Temporal band is the unnarrowed 5--25 Hz duration rule, where Bai et al. narrowed it to 15--25 Hz as noisy & The pass-band admits the low-frequency range they rejected \\
		3.4.3 & Strain magnitude counts shear once, not twice as in the published Frobenius norm & Shear weighted lower by $\sqrt{2}$ relative to normal strain; untested \\
		3.4.7 & No dedicated strain filtering, where the literature applies Wiener and Gaussian filters & Spatial differentiation can amplify flow-estimation noise, and nothing downstream suppresses it \\
		3.5.7 & Temporal normalisation is uniform index sampling, \textbf{not} manifold-based TIM; no frames are synthesised & Avoids fabricating motion; forfeits the ability to lengthen short clips \\
		3.5.7 & Magnification is applied \textbf{after} subsampling, at a nominal 200 fps & Realised pass-band is clip-dependent and below the intended 5--25 Hz \\
		3.6.6 & Backbone kernels are (1, 3, 3) --- \textbf{no temporal mixing} & The ablation tests a learned \textit{spatial} stem, not spatio-temporal convolution \\
		3.6.6 & Backbone input is 224 $\times$ 224, against Xia et al.'s guidance of $\leq$ 100 $\times$ 100 --- though Bai et al. and Li, Huang and Zhao both use 224 $\times$ 224 & Whatever the stem is measured to contribute may be a resolution effect as much as an architecture effect \\
		3.7.7 & SimAM statistics pooled over the whole clip rather than per frame as in the source; $\lambda$ not re-tuned & Neurons judged against the clip, not the frame; defensible but untested \\
		3.8.6 & Six divergences from published SLSTT (temporal not spatial attention, mean not LSTM aggregation, from scratch, smaller, sinusoidal encoding, short-term flow) & Variable D is a temporal transformer encoder, \textbf{not} a reproduction of SLSTT \\
		3.9.7 & Label smoothing of 0.05, named in the corpus only by Dosovitskiy et al. and never examined on this task & A choice the micro-expression corpus supports neither way, held constant and unmeasured \\
		\hline
	\end{tabular}
	\caption[Every declared departure from the reviewed literature]{Every declared departure from the reviewed literature, with its consequence.}
	\label{tab:divergence-ledger}
\end{table}

Three of these --- the strain formula, the magnification ordering, and the input resolution --- were identified by the review itself rather than being known design decisions, and each is declared as proposed further work. None invalidates the ablation, since every configuration shares the same choices. What they constrain is the \textit{wording} of the conclusions: what is measured is the contribution of a spatial stem at 224 $\times$ 224 and of temporal self-attention in this pipeline, not a verdict on 3D convolution or on SLSTT.

\subsection{What this chapter contributes to the thesis}
\label{sec:ch3-contribution}

Three things carry forward. First, the \textbf{methodology that follows is determined by this review rather than chosen freely}: the corpus dictates the class grouping, representation, fold structure and metric, and each component's parameters come from published sweeps (\autoref{sec:corpus-implications}). Second, the \textbf{interpretation of the central result is prepared here rather than improvised later}: flow and strain already perform spatio-temporal feature extraction analytically (\autoref{sec:flow-implications}, \autoref{sec:strain-implications}), so a negative contribution from a spatial-only backbone at an oversized input resolution is expected rather than anomalous (\autoref{sec:cnn-implications}). Third, the \textbf{limits of what can be claimed are fixed before any number is reported}: a small, single-seed, demographically narrow corpus measured against richer published regimes, with the twelve divergences of \autoref{tab:divergence-ledger} standing throughout --- consequences of this review, not concessions made after the fact.
